% Vietnamese translation for OI Public Library
% Source-PDF: 模板 TEMPLATES/kuangbin的ACM模板（新）.pdf
\documentclass[11pt]{article}
\usepackage{oipl-vn}
\usepackage{fvextra}
\setmonofont{Unifont}[Scale=MatchLowercase]
\DefineVerbatimEnvironment{Code}{Verbatim}{fontsize=\scriptsize,breaklines=true,breakanywhere=true}
\sloppy

\title{Mẫu thuật toán ACM thường dùng của kuangbin}
\author{kuangbin\\\url{http://www.cnblogs.com/kuangbin/}\\Email: \texttt{kuangbin2009@126.com}}
\date{Bản gốc biên dịch lần cuối: July 8, 2018}

\begin{document}
\maketitle
\tableofcontents
\clearpage

\section{Xử lý chuỗi}

\subsection{KMP}

\begin{Code}
 1   /*
 2    * next[] 的含义：x[i-next[i]...i-1]=x[0...next[i]-1]
 3    * next[i] 为满足 x[i-z...i-1]=x[0...z-1] 的最大 z 值（就是 x 的自身匹配）
 4    */
 5   void kmp_pre(char x[],int m,int next[]){
 6        int i,j;
 7        j=next[0]=-1;
 8        i=0;
 9        while(i<m){
10            while(-1!=j && x[i]!=x[j])j=next[j];
11            next[++i]=++j;
12        }
13   }
14   /*
15     * kmpNext[i] 的意思:next'[i]=next[next[...[next[i]]]](直到
          next'[i]<0 或者 x[next'[i]]!=x[i])
16     * 这样的预处理可以快一些
17     */
18   void preKMP(char x[],int m,int kmpNext[]){
19        int i,j;
20        j=kmpNext[0]=-1;
21        i=0;
22        while(i<m){
23            while(-1!=j && x[i]!=x[j])j=kmpNext[j];
24            if(x[++i]==x[++j])kmpNext[i]=kmpNext[j];
25            else kmpNext[i]=j;
26        }
27   }
28   /*
29     * 返回 x 在 y 中出现的次数，可以重叠
30     */
31   int next[10010];
32   int KMP_Count(char x[],int m,char y[],int n){//x 是模式串，y 是主串
33        int i,j;
34        int ans=0;
35        //preKMP(x,m,next);
36        kmp_pre(x,m,next);
37        i=j=0;
38        while(i<n){
39            while(-1!=j && y[i]!=x[j])j=next[j];
40            i++;j++;
41            if(j>=m){
42                 ans++;
43                 j=next[j];
44            }
45        }
46        return ans;


47   }
48   //经典题目：POJ 3167
49   /*
50     * POJ 3167 Cow Patterns
51     * 模式串可以浮动的模式匹配问题
52     * 给出模式串的相对大小，需要找出模式串匹配次数和位置
53     * 比如说模式串：1，4，4，2，3，1 而主串：5,6,2,10,10,7,3,2,9
54     * 那么 2,10,10,7,3,2 就是匹配的
55     *
56     * 统计比当前数小，和于当前数相等的，然后进行 kmp
57     */
58   const int MAXN=100010;
59   const int MAXM=25010;
60   int a[MAXN];
61   int b[MAXN];
62   int n,m,s;
63   int as[MAXN][30];
64   int bs[MAXM][30];
65   void init(){
66        for(int i=0;i<n;i++){
67            if(i==0){
68                 for(int j=1;j<=25;j++)as[i][j]=0;
69            }
70            else{
71                 for(int j=1;j<=25;j++)as[i][j]=as[i-1][j];
72            }
73            as[i][a[i]]++;
74        }
75        for(int i=0;i<m;i++){
76            if(i==0){
77                 for(int j=1;j<=25;j++)bs[i][j]=0;
78            }
79            else{
80                 for(int j=1;j<=25;j++)bs[i][j]=bs[i-1][j];
81            }
82            bs[i][b[i]]++;
83        }
84   }
85   int next[MAXM];
86   void kmp_pre(){
87        int i,j;
88        j=next[0]=-1;
89        i=0;
90        while(i<m){
91            int t11=0,t12=0,t21=0,t22=0;
92            for(int k=1;k<b[i];k++){
93                 if(i-j>0)t11+=bs[i][k]-bs[i-j-1][k];
94                 else t11+=bs[i][k];
95            }
96            if(i-j>0)t12=bs[i][b[i]]-bs[i-j-1][b[i]];
97            else t12=bs[i][b[i]];


 98
 99              for(int k=1;k<b[j];k++){
100                  t21+=bs[j][k];
101              }
102              t22=bs[j][b[j]];
103              if(j==-1 || (t11==t21&&t12==t22)){
104                  next[++i]=++j;
105              }
106              else j=next[j];
107       }
108   }
109   vector<int>ans;
110   void kmp(){
111       ans.clear();
112       int i,j;
113       kmp_pre();
114       i=j=0;
115       while(i<n){
116           int t11=0,t12=0,t21=0,t22=0;
117           for(int k=1;k<a[i];k++){
118                if(i-j>0)t11+=as[i][k]-as[i-j-1][k];
119                else t11+=as[i][k];
120           }
121           if(i-j>0)t12=as[i][a[i]]-as[i-j-1][a[i]];
122           else t12=as[i][a[i]];
123
124              for(int k=1;k<b[j];k++){
125                  t21+=bs[j][k];
126              }
127              t22=bs[j][b[j]];
128              if(j==-1 || (t11==t21&&t12==t22)){
129                  i++;j++;
130                  if(j>=m){
131                      ans.push_back(i-m+1);
132                      j=next[j];
133                  }
134              }
135              else j=next[j];
136       }
137   }
138   int main(){
139       while(scanf("%d%d%d",&n,&m,&s)==3){
140           for(int i=0;i<n;i++)scanf("%d",&a[i]);
141           for(int i=0;i<m;i++)scanf("%d",&b[i]);
142           init();
143           kmp();
144           printf("%d\n",ans.size());
145           for(int i=0;i<ans.size();i++)
146               printf("%d\n",ans[i]);
147       }
148       return 0;


149 }
\end{Code}

\subsection{e-KMP}

\begin{Code}
 1   /*
 2    * 扩展 KMP 算法
 3    */
 4   //next[i]:x[i...m-1] 与 x[0...m-1] 的最长公共前缀
 5   //extend[i]:y[i...n-1] 与 x[0...m-1] 的最长公共前缀
 6   void pre_EKMP(char x[],int m,int next[]){
 7       next[0] = m;
 8       int j = 0;
 9       while( j+1 < m && x[j] == x[j+1] )j++;
10       next[1] = j;
11       int k = 1;
12       for(int i = 2; i < m; i++){
13           int p = next[k]+k-1;
14           int L = next[i-k];
15           if( i+L < p+1 )next[i] = L;
16           else{
17               j = max(0,p-i+1);
18               while( i+j < m && x[i+j] == x[j])j++;
19               next[i] = j;
20               k = i;
21           }
22       }
23   }
24   void EKMP(char x[],int m,char y[],int n,int next[],int extend[]){
25       pre_EKMP(x,m,next);
26       int j = 0;
27       while(j < n && j < m && x[j] == y[j])j++;
28       extend[0] = j;
29       int k = 0;
30       for(int i = 1;i < n;i++){
31           int p = extend[k]+k-1;
32           int L = next[i-k];
33           if(i+L < p+1)extend[i] = L;
34           else{
35               j = max(0,p-i+1);
36               while( i+j < n && j < m && y[i+j] == x[j] )j++;
37               extend[i] = j;
38               k = i;
39           }
40       }
41   }
\end{Code}

\subsection{Manacher}

\begin{Code}
 1 /*
 2 * 求最长回文子串
 3 */
 4 const int MAXN=110010;


 5   char Ma[MAXN*2];
 6   int Mp[MAXN*2];
 7   void Manacher(char s[],int len){
 8        int l=0;
 9        Ma[l++]='$';
10        Ma[l++]='#';
11        for(int i=0;i<len;i++){
12            Ma[l++]=s[i];
13            Ma[l++]='#';
14        }
15        Ma[l]=0;
16        int mx=0,id=0;
17        for(int i=0;i<l;i++){
18            Mp[i]=mx>i?min(Mp[2*id-i],mx-i):1;
19            while(Ma[i+Mp[i]]==Ma[i-Mp[i]])Mp[i]++;
20            if(i+Mp[i]>mx){
21                 mx=i+Mp[i];
22                 id=i;
23            }
24        }
25   }
26   /*
27     * abaaba
28     * i:     0 1 2 3 4 5 6 7 8 9 10 11 12 13
29     * Ma[i]: $ # a # b # a # a # b # a #
30     * Mp[i]: 1 1 2 1 4 1 2 7 2 1 4 1 2 1
31     */
32   char s[MAXN];
33   int main(){
34        while(scanf("%s",s)==1){
35            int len=strlen(s);
36            Manacher(s,len);
37            int ans=0;
38            for(int i=0;i<2*len+2;i++)
39                 ans=max(ans,Mp[i]-1);
40            printf("%d\n",ans);
41        }
42        return 0;
43   }
\end{Code}

\subsection{Tự động cơ AC}

\begin{Code}
 1   //======================
 2   // HDU 2222
 3   // 求目标串中出现了几个模式串
 4   //====================
 5   struct Trie{
 6       int next[500010][26],fail[500010],end[500010];
 7       int root,L;
 8       int newnode(){
 9           for(int i = 0;i < 26;i++)
10                next[L][i] = -1;


11              end[L++] = 0;
12              return L-1;
13       }
14       void init(){
15           L = 0;
16           root = newnode();
17       }
18       void insert(char buf[]){
19           int len = strlen(buf);
20           int now = root;
21           for(int i = 0;i < len;i++){
22               if(next[now][buf[i]-'a'] == -1)
23                   next[now][buf[i]-'a'] = newnode();
24               now = next[now][buf[i]-'a'];
25           }
26           end[now]++;
27       }
28       void build(){
29           queue<int>Q;
30           fail[root] = root;
31           for(int i = 0;i < 26;i++)
32               if(next[root][i] == -1)
33                   next[root][i] = root;
34               else{
35                   fail[next[root][i]] = root;
36                   Q.push(next[root][i]);
37               }
38           while( !Q.empty() ){
39               int now = Q.front();
40               Q.pop();
41               for(int i = 0;i < 26;i++)
42                   if(next[now][i] == -1)
43                        next[now][i] = next[fail[now]][i];
44                   else
45                   {
46                        fail[next[now][i]]=next[fail[now]][i];
47                        Q.push(next[now][i]);
48                   }
49           }
50       }
51       int query(char buf[]){
52           int len = strlen(buf);
53           int now = root;
54           int res = 0;
55           for(int i = 0;i < len;i++){
56               now = next[now][buf[i]-'a'];
57               int temp = now;
58               while( temp != root ){
59                   res += end[temp];
60                   end[temp] = 0;
61                   temp = fail[temp];


62                  }
63              }
64              return res;
65         }
66         void debug(){
67             for(int i = 0;i < L;i++){
68                 printf("id = %3d,fail = %3d,end = %3d,chi = [",i,fail[i
                      ],end[i]);
69                 for(int j = 0;j < 26;j++)
70                     printf("%2d",next[i][j]);
71                 printf("]\n");
72             }
73         }
74   };
75   char buf[1000010];
76   Trie ac;
77   int main(){
78       int T;
79       int n;
80       scanf("%d",&T);
81       while( T-- ){
82            scanf("%d",&n);
83            ac.init();
84            for(int i = 0;i < n;i++){
85                scanf("%s",buf);
86                ac.insert(buf);
87            }
88            ac.build();
89            scanf("%s",buf);
90            printf("%d\n",ac.query(buf));
91       }
92       return 0;
93   }
\end{Code}

\subsection{Mảng hậu tố}

\subsubsection{DA}

\begin{Code}
1  /*
2  *suffix array
3  *倍增算法 O(n*logn)
4  *待排序数组长度为 n, 放在 0 n-1 中，在最后面补一个 0
5  *da(str ,sa,rank,height, n , );//注意是 n;
6  *例如：
7  *n = 8;
8  * num[]   = { 1, 1, 2, 1, 1, 1, 1, 2, $ }; 注意 num 最后一位为 0，其他
      大于 0
 9 *rank[] = 4, 6, 8, 1, 2, 3, 5, 7, 0 ;rank[0 n-1] 为有效值，rank[n]
      必定为 0 无效值
10 *sa[] = 8, 3, 4, 5, 0, 6, 1, 7, 2 ;sa[1 n] 为有效值，sa[0] 必定为 n 是
      无效值
11 *height[]= 0, 0, 3, 2, 3, 1, 2, 0, 1 ;height[2 n] 为有效值


12   *
13   */
14   const int MAXN=20010;
15   int t1[MAXN],t2[MAXN],c[MAXN];//求 SA 数组需要的中间变量，不需要赋值
16   //待排序的字符串放在 s 数组中，从 s[0] 到 s[n-1], 长度为 n, 且最大值小于 m,
17   //除 s[n-1] 外的所有 s[i] 都大于 0，r[n-1]=0
18   //函数结束以后结果放在 sa 数组中
19   bool cmp(int *r,int a,int b,int l){
20       return r[a] == r[b] && r[a+l] == r[b+l];
21   }
22   void da(int str[],int sa[],int rank[],int height[],int n,int m){
23       n++;
24       int i, j, p, *x = t1, *y = t2;
25       //第一轮基数排序，如果 s 的最大值很大，可改为快速排序
26       for(i = 0;i < m;i++)c[i] = 0;
27       for(i = 0;i < n;i++)c[x[i] = str[i]]++;
28       for(i = 1;i < m;i++)c[i] += c[i-1];
29       for(i = n-1;i >= 0;i--)sa[--c[x[i]]] = i;
30       for(j = 1;j <= n; j <<= 1){
31           p = 0;
32           //直接利用 sa 数组排序第二关键字
33           for(i = n-j; i < n; i++)y[p++] = i;//后面的 j 个数第二关键字为
                空的最小
34           for(i = 0; i < n; i++)if(sa[i] >= j)y[p++] = sa[i] - j;
35           //这样数组 y 保存的就是按照第二关键字排序的结果
36           //基数排序第一关键字
37           for(i = 0; i < m; i++)c[i] = 0;
38           for(i = 0; i < n; i++)c[x[y[i]]]++;
39           for(i = 1; i < m;i++)c[i] += c[i-1];
40           for(i = n-1; i >= 0;i--)sa[--c[x[y[i]]]] = y[i];
41           //根据 sa 和 x 数组计算新的 x 数组
42           swap(x,y);
43           p = 1; x[sa[0]] = 0;
44           for(i = 1;i < n;i++)
45               x[sa[i]] = cmp(y,sa[i-1],sa[i],j)?p-1:p++;
46           if(p >= n)break;
47           m = p;//下次基数排序的最大值
48       }
49       int k = 0;
50       n--;
51       for(i = 0;i <= n;i++)rank[sa[i]] = i;
52       for(i = 0;i < n;i++){
53           if(k)k--;
54           j = sa[rank[i]-1];
55           while(str[i+k] == str[j+k])k++;
56           height[rank[i]] = k;
57       }
58   }
59   int rank[MAXN],height[MAXN];
60   int RMQ[MAXN];
61   int mm[MAXN];


62    int best[20][MAXN];
63    void initRMQ(int n){
64        mm[0]=-1;
65        for(int i=1;i<=n;i++)
66            mm[i]=((i&(i-1))==0)?mm[i-1]+1:mm[i-1];
67        for(int i=1;i<=n;i++)best[0][i]=i;
68        for(int i=1;i<=mm[n];i++)
69            for(int j=1;j+(1<<i)-1<=n;j++){
70                 int a=best[i-1][j];
71                 int b=best[i-1][j+(1<<(i-1))];
72                 if(RMQ[a]<RMQ[b])best[i][j]=a;
73                 else best[i][j]=b;
74            }
75    }
76    int askRMQ(int a,int b){
77        int t;
78        t=mm[b-a+1];
79        b-=(1<<t)-1;
80        a=best[t][a];b=best[t][b];
81        return RMQ[a]<RMQ[b]?a:b;
82    }
83    int lcp(int a,int b){
84        a=rank[a];b=rank[b];
85        if(a>b)swap(a,b);
86        return height[askRMQ(a+1,b)];
87    }
88    char str[MAXN];
89    int r[MAXN];
90    int sa[MAXN];
91    int main()
92    {
93        while(scanf("%s",str) == 1){
94            int len = strlen(str);
95            int n = 2*len + 1;
96            for(int i = 0;i < len;i++)r[i] = str[i];
97            for(int i = 0;i < len;i++)r[len + 1 + i] = str[len - 1 - i
                  ];
 98           r[len] = 1;
 99           r[n] = 0;
100           da(r,sa,rank,height,n,128);
101             for(int i=1;i<=n;i++)RMQ[i]=height[i];
102           initRMQ(n);
103           int ans=0,st;
104           int tmp;
105           for(int i=0;i<len;i++){
106                tmp=lcp(i,n-i);//偶对称
107                if(2*tmp>ans){
108                    ans=2*tmp;
109                    st=i-tmp;
110                }
111                tmp=lcp(i,n-i-1);//奇数对称


112                 if(2*tmp-1>ans){
113                     ans=2*tmp-1;
114                     st=i-tmp+1;
115                 }
116             }
117             str[st+ans]=0;
118             printf("%s\n",str+st);
119      }
120      return 0;
121 }
\end{Code}

\subsubsection{DC3}

\begin{Code}
     da[] 和 str[] 数组要开大三倍，相关数组也是三倍
 1   /*
 2    * 后缀数组
 3    * DC3 算法，复杂度 O(n)
 4    * 所有的相关数组都要开三倍
 5    */
 6   const int MAXN = 2010;
 7   #define F(x) ((x)/3+((x)%3==1?0:tb))
 8   #define G(x) ((x)<tb?(x)*3+1:((x)-tb)*3+2)
 9   int wa[MAXN*3],wb[MAXN*3],wv[MAXN*3],wss[MAXN*3];
10   int c0(int *r,int a,int b){
11       return r[a] == r[b] && r[a+1] == r[b+1] && r[a+2] == r[b+2];
12   }
13   int c12(int k,int *r,int a,int b){
14       if(k == 2)
15           return r[a] < r[b] || ( r[a] == r[b] && c12(1,r,a+1,b+1) );
16       else return r[a] < r[b] || ( r[a] == r[b] && wv[a+1] < wv[b+1]
            );
17   }
18   void sort(int *r,int *a,int *b,int n,int m){
19       int i;
20       for(i = 0;i < n;i++)wv[i] = r[a[i]];
21       for(i = 0;i < m;i++)wss[i] = 0;
22       for(i = 0;i < n;i++)wss[wv[i]]++;
23       for(i = 1;i < m;i++)wss[i] += wss[i-1];
24       for(i = n-1;i >= 0;i--)
25           b[--wss[wv[i]]] = a[i];
26   }
27   void dc3(int *r,int *sa,int n,int m){
28       int i, j, *rn = r + n;
29       int *san = sa + n, ta = 0, tb = (n+1)/3, tbc = 0, p;
30       r[n] = r[n+1] = 0;
31       for(i = 0;i < n;i++)if(i %3 != 0)wa[tbc++] = i;
32       sort(r + 2, wa, wb, tbc, m);
33       sort(r + 1, wb, wa, tbc, m);
34       sort(r, wa, wb, tbc, m);
35       for(p = 1, rn[F(wb[0])] = 0, i = 1;i < tbc;i++)
36           rn[F(wb[i])] = c0(r, wb[i-1], wb[i]) ? p - 1 : p++;
37       if(p < tbc)dc3(rn,san,tbc,p);


38         else for(i = 0;i < tbc;i++)san[rn[i]] = i;
39         for(i = 0;i < tbc;i++) if(san[i] < tb)wb[ta++] = san[i] * 3;
40         if(n % 3 == 1)wb[ta++] = n - 1;
41         sort(r, wb, wa, ta, m);
42         for(i = 0;i < tbc;i++)wv[wb[i] = G(san[i])] = i;
43         for(i = 0, j = 0, p = 0;i < ta && j < tbc;p++)
44             sa[p] = c12(wb[j] % 3, r, wa[i], wb[j]) ? wa[i++] : wb[j
                  ++];
45         for(;i < ta;p++)sa[p] = wa[i++];
46         for(;j < tbc;p++)sa[p] = wb[j++];
47   }
48   //str 和 sa 也要三倍
49   void da(int str[],int sa[],int rank[],int height[],int n,int m){
50       for(int i = n;i < n*3;i++)
51           str[i] = 0;
52       dc3(str, sa, n+1, m);
53       int i,j,k = 0;
54       for(i = 0;i <= n;i++)rank[sa[i]] = i;
55       for(i = 0;i < n; i++){
56           if(k) k--;
57           j = sa[rank[i]-1];
58           while(str[i+k] == str[j+k]) k++;
59           height[rank[i]] = k;
60       }
61   }
\end{Code}

\subsection{Tự động cơ hậu tố}

\subsubsection{Các hàm cơ bản}

\begin{Code}
 1   const int CHAR = 26;
 2   const int MAXN = 250010;
 3   struct SAM_Node{
 4       SAM_Node *fa,*next[CHAR];
 5       int len;
 6       long long cnt;
 7       void clear(){
 8           fa = 0;
 9           memset(next,0,sizeof(next));
10           cnt = 0;
11       }
12   }pool[MAXN*2];
13   SAM_Node *root,*tail;
14   SAM_Node* newnode(int len){
15       SAM_Node* cur = tail++;
16       cur->clear();
17       cur->len = len;
18       return cur;
19   }
20   void SAM_init(){
21       tail = pool;


22     root = newnode(0);
23 }
24 SAM_Node* extend(SAM_Node* last,int x){
25     SAM_Node *p = last, *np = newnode(p->len+1);
26     while(p && !p->next[x])
27         p->next[x] = np, p = p->fa;
28     if(!p)np->fa = root;
29     else {
30         SAM_Node* q = p->next[x];
31         if(q->len == p->len+1)np->fa = q;
32         else {
33             SAM_Node* nq = newnode(p->len+1);
34             memcpy(nq->next,q->next,sizeof(q->next));
35             nq->fa = q->fa;
36             q->fa = np->fa = nq;
37             while(p && p->next[x] == q)
38                 p->next[x] = nq, p = p->fa;
39         }
40     }
41     return np;
42 }
\end{Code}

\subsubsection{Ví dụ}

\begin{Code}
     CC TSUBSTR
     给了一个 Trie 树，求 Trie 子树上的第 k 大的子串。

 1   /*
 2    * http://www.codechef.com/problems/TSUBSTR/
 3   Input:
 4   8 4
 5   abcbbaca
 6   1 2
 7   2 3
 8   1 4
 9   4 5
10   4 6
11   4 7
12   1 8
13   abcdefghijklmnopqrstuvwxyz 5
14   abcdefghijklmnopqrstuvwxyz 1
15   bcadefghijklmnopqrstuvwxyz 5
16   abcdefghijklmnopqrstuvwxyz 100
17
18   Output:
19   12
20   aba
21
22   ba
23   -1
24    */
25   const int CHAR = 26;


26   const int MAXN = 250010;
27   struct SAM_Node{
28       SAM_Node *fa,*next[CHAR];
29       int len;
30       long long cnt;
31       void clear(){
32           fa = 0;
33           memset(next,0,sizeof(next));
34           cnt = 0;
35       }
36   }pool[MAXN*2];
37   SAM_Node *root,*tail;
38   SAM_Node* newnode(int len){
39       SAM_Node* cur = tail++;
40       cur->clear();
41       cur->len = len;
42       return cur;
43   }
44   void SAM_init(){
45       tail = pool;
46       root = newnode(0);
47   }
48   SAM_Node* extend(SAM_Node* last,int x){
49       SAM_Node *p = last, *np = newnode(p->len+1);
50       while(p && !p->next[x])
51           p->next[x] = np, p = p->fa;
52       if(!p)np->fa = root;
53       else {
54           SAM_Node* q = p->next[x];
55           if(q->len == p->len+1)np->fa = q;
56           else {
57                SAM_Node* nq = newnode(p->len+1);
58                memcpy(nq->next,q->next,sizeof(q->next));
59                nq->fa = q->fa;
60                q->fa = np->fa = nq;
61                while(p && p->next[x] == q)
62                    p->next[x] = nq, p = p->fa;
63           }
64       }
65       return np;
66   }
67   char str[MAXN];
68   struct Edge
69   {
70       int to,next;
71   }edge[MAXN*2];
72   int head[MAXN],tot;
73   void addedge(int u,int v){
74       edge[tot].to = v;
75       edge[tot].next = head[u];
76       head[u] = tot++;


 77   }
 78
 79   SAM_Node *end[MAXN];
 80   int topcnt[MAXN];// 拓扑排序使用
 81   SAM_Node *topsam[MAXN*2];
 82   char s2[40];
 83   int order[40];
 84
 85   int main()
 86   {
 87       int n,Q;
 88       while(scanf("%d%d",&n,&Q) == 2){
 89           scanf("%s",str+1);
 90           memset(head,-1,sizeof(head));tot = 0;
 91           int u,v;
 92           for(int i = 1;i < n;i++){
 93                scanf("%d%d",&u,&v);
 94                addedge(u,v); addedge(v,u);
 95           }
 96           addedge(0,1);
 97           SAM_init();
 98           memset(end,0,sizeof(end));
 99           end[0] = root;
100           queue<int>q;
101           q.push(0);
102           while(!q.empty()){
103                u = q.front();
104                q.pop();
105                for(int i = head[u];i != -1;i = edge[i].next){
106                    v = edge[i].to;
107                    if(end[v] != 0)continue;
108                    end[v] = extend(end[u],str[v]-'a');
109                    q.push(v);
110                }
111           }
112           memset(topcnt,0,sizeof(topcnt));
113           int num = tail - pool;
114           for(int i = 0;i < num;i++)topcnt[pool[i].len]++;
115           for(int i = 1;i <= n;i++)topcnt[i] += topcnt[i-1];
116           for(int i = 0;i < num;i++)topsam[--topcnt[pool[i].len]] = &
                 pool[i];
117
118              for(int i = num-1;i >= 0;i--){
119                  SAM_Node *p = topsam[i];
120                  p->cnt = 1;
121                  for(int i = 0;i < 26;i++)
122                      if(p->next[i])
123                          p->cnt += p->next[i]->cnt;
124              }
125              printf("%lld\n",root->cnt);
126              long long k;


127             while(Q--){
128                 scanf("%s",s2);
129                 for(int i = 0;i < 26;i++)order[i] = s2[i]-'a';
130                 scanf("%lld",&k);
131                 if(k > root->cnt){
132                     printf("-1\n");
133                     continue;
134                 }
135                 SAM_Node *p = root;
136                 //这里的第 k 个子串是从空串算起的
137                 while( (--k) > 0 ){
138                     for(int i = 0;i < 26;i++)
139                         if(p->next[order[i]]){
140                             if(k <= p->next[order[i]]->cnt){
141                                 printf("%c",'a'+order[i]);
142                                 p = p->next[order[i]];
143                                 break; //这个不要忘记
144                             }
145                             else k -= p->next[order[i]]->cnt;
146                         }
147                 }
148                 printf("\n");
149             }
150      }
151      return 0;
152 }
     CF129 E
     给了 n 个字符串，求每个字符串有多少个至少出现在 k 个字符串中的子串
     fail 树，两遍 dfs, 经典题。

 1   /* http://codeforces.com/contest/204/problem/E
 2   input
 3   3 1
 4   abc
 5   a
 6   ab
 7   output
 8   6 1 3
 9   input
10   7 4
11   rubik
12   furik
13   abab
14   baba
15   aaabbbababa
16   abababababa
17   zero
18   output
19   1 0 9 9 21 30 0
20     */
21   const int CHAR = 26;


22   const int MAXN = 100010;
23   //*************SAM******************
24   struct SAM_Node{
25       SAM_Node *fa,*next[CHAR];
26       int len;
27       void clear(){
28           fa = 0;
29           memset(next,0,sizeof(next));
30       }
31   }pool[MAXN*2];
32   SAM_Node *root,*tail;
33   SAM_Node* newnode(int len){
34       SAM_Node* cur = tail++;
35       cur->clear();
36       cur->len = len;
37       return cur;
38   }
39   void SAM_init(){
40       tail = pool;
41       root = newnode(0);
42   }
43   SAM_Node* extend(SAM_Node* last,int x){
44       SAM_Node *p = last, *np = newnode(p->len+1);
45       while(p && !p->next[x])
46           p->next[x] = np, p = p->fa;
47       if(!p)np->fa = root;
48       else {
49           SAM_Node* q = p->next[x];
50           if(q->len == p->len+1)np->fa = q;
51           else {
52                SAM_Node* nq = newnode(p->len+1);
53                memcpy(nq->next,q->next,sizeof(q->next));
54                nq->fa = q->fa;
55                q->fa = np->fa = nq;
56                while(p && p->next[x] == q)
57                    p->next[x] = nq, p = p->fa;
58           }
59       }
60       return np;
61   }
62   //************Trie************
63   struct Trie_Node{
64       int next[CHAR];
65       vector<int>belongs;
66   }trie[MAXN];
67   int trie_root,trie_tot;
68   int trie_newnode(){
69       memset(trie[trie_tot].next,-1,sizeof(trie[trie_tot].next));
70       trie[trie_tot].belongs.clear();
71       return trie_tot++;
72   }


 73   void Trie_init(){
 74       trie_tot = 0;
 75       trie_root = trie_newnode();
 76   }
 77   void insert(char buf[],int id){
 78       int now = trie_root;
 79       int len = strlen(buf);
 80       for(int i = 0;i < len;i++){
 81           if(trie[now].next[buf[i]-'a'] == -1)
 82                trie[now].next[buf[i]-'a'] = trie_newnode();
 83           now = trie[now].next[buf[i]-'a'];
 84           trie[now].belongs.push_back(id);
 85       }
 86   }
 87   //***** fail 树***************
 88   struct Edge{
 89       int to,next;
 90   }edge[MAXN*2];
 91   int head[MAXN*2],tot;
 92   void addedge(int u,int v){
 93       edge[tot].to = v; edge[tot].next = head[u]; head[u] = tot++;
 94   }
 95   int MtoT[MAXN*2];//SAM 结点映射到 Trie 结点
 96   int cnt[MAXN*2];
 97   int F[MAXN*2];
 98   int find(int x){
 99       if(F[x] == -1)return x;
100       return F[x] = find(F[x]);
101   }
102   void bing(int u,int v)//注意方向性
103   {
104       int t1 = find(u);
105       int t2 = find(v);
106       if(t1 != t2)F[t1] = t2;
107   }
108   int L[MAXN];
109   void Tarjan(int u){
110       for(int i = head[u];i != -1;i = edge[i].next){
111           Tarjan(edge[i].to);
112           bing(edge[i].to,u);
113       }
114       if(MtoT[u]){
115           int tt = MtoT[u];
116           int sz = trie[tt].belongs.size();
117           for(int i = 0;i < sz;i++){
118                int v = trie[tt].belongs[i];
119                cnt[find(L[v])]--;
120                cnt[u]++;
121                L[v] = u;
122           }
123       }


124   }
125   void dfs1(int u){
126       for(int i = head[u];i != -1;i = edge[i].next){
127           dfs1(edge[i].to);
128           cnt[u] += cnt[edge[i].to];
129       }
130   }
131   long long ans[MAXN];
132   void dfs2(int u){
133       for(int i = head[u];i != -1;i = edge[i].next){
134           int v = edge[i].to;
135           cnt[v] += cnt[u];
136           dfs2(v);
137       }
138       if(MtoT[u]){
139           int tt = MtoT[u];
140           int sz = trie[tt].belongs.size();
141           for(int i = 0;i < sz;i++){
142               int v = trie[tt].belongs[i];
143               ans[v] += cnt[u];
144           }
145       }
146   }
147
148   char str[MAXN];
149   SAM_Node *end[MAXN];
150   int main()
151   {
152       int n,k;
153       while(scanf("%d%d",&n,&k) == 2){
154           Trie_init();
155           for(int i = 0;i < n;i++){
156                scanf("%s",str);
157                insert(str,i);
158           }
159           SAM_init();
160           //根据 Trie 建立 SAM
161           memset(end,0,sizeof(end));
162           end[0] = root;
163           memset(MtoT,0,sizeof(MtoT));
164           MtoT[root-pool] = 0;
165           queue<int>q;
166           q.push(trie_root);
167           while(!q.empty()){
168                int u = q.front();
169                q.pop();
170                for(int i = 0;i < 26;i++){
171                    if(trie[u].next[i] == -1)continue;
172                    int v = trie[u].next[i];
173                    end[v] = extend(end[u],i);
174                    MtoT[end[v]-pool] = v;


175                    q.push(v);
176                }
177            }
178            //建立 fail 树
179            int num = tail - pool;
180            memset(head,-1,sizeof(head));
181            tot = 0;
182            for(SAM_Node *p = pool+1;p < tail;p++)
183                addedge(p->fa - pool,p - pool);
184            memset(cnt,0,sizeof(cnt));
185            memset(F,-1,sizeof(F));
186            memset(L,0,sizeof(L));
187            Tarjan(0);
188            dfs1(0);
189            for(int i = 0;i < num;i++){
190                if(cnt[i] >= k)cnt[i] = pool[i].len - pool[i].fa->len;
191                else cnt[i] = 0;
192            }
193            memset(ans,0,sizeof(ans));
194            dfs2(0);
195            for(int i = 0;i < n;i++){
196                printf("%I64d",ans[i]);
197                if(i < n-1)printf(" ");
198                else printf("\n");
199            }
200       }
201       return 0;
202 }
\end{Code}

\subsection{Hash chuỗi}

\begin{Code}
    HDU4622 求区间不相同子串个数
 1 const int HASH = 10007;
 2 const int MAXN = 2010;
 3 struct HASHMAP{
 4     int head[HASH],next[MAXN],size;
 5     unsigned long long state[MAXN];
 6     int f[MAXN];
 7     void init(){
 8         size = 0;
 9         memset(head,-1,sizeof(head));
10     }
11     int insert(unsigned long long val,int _id){
12         int h = val%HASH;
13         for(int i = head[h]; i != -1;i = next[i])
14             if(val == state[i]){
15                  int tmp = f[i];
16                  f[i] = _id;
17                  return tmp;
18             }
19         f[size] = _id;


20              state[size] = val;
21              next[size] = head[h];
22              head[h] = size++;
23              return 0;
24       }
25   }H;
26   const int SEED = 13331;
27   unsigned long long P[MAXN];
28   unsigned long long S[MAXN];
29   char str[MAXN];
30   int ans[MAXN][MAXN];
31   int main(){
32       P[0] = 1;
33       for(int i = 1;i < MAXN;i++)
34           P[i] = P[i-1] * SEED;
35       int T;
36       scanf("%d",&T);
37       while(T--){
38           scanf("%s",str);
39           int n = strlen(str);
40           S[0] = 0;
41           for(int i = 1;i <= n;i++)
42               S[i] = S[i-1]*SEED + str[i-1];
43           memset(ans,0,sizeof(ans));
44           for(int L = 1; L <= n;L++){
45               H.init();
46               for(int i = 1;i + L - 1 <= n;i++){
47                   int l = H.insert(S[i+L-1] - S[i-1]*P[L],i);
48                   ans[i][i+L-1] ++;
49                   ans[l][i+L-1]--;
50               }
51           }
52           for(int i = n;i >= 0;i--)
53               for(int j = i;j <= n;j++)
54                   ans[i][j] += ans[i+1][j] + ans[i][j-1] - ans[i+1][j
                        -1];
55           int m,u,v;
56           scanf("%d",&m);
57           while(m--){
58               scanf("%d%d",&u,&v);
59               printf("%d\n",ans[u][v]);
60           }
61       }
62       return 0;
63   }
\end{Code}

\section{Toán học}

\subsection{Số nguyên tố}

\subsubsection{Sàng số nguyên tố (kiểm tra số nhỏ hơn MAXN có nguyên tố không)}

\begin{Code}
 1   /*
 2    * 素数筛选，判断小于 MAXN 的数是不是素数。
 3    * notprime 是一张表，为 false 表示是素数，true 表示不是素数
 4    */
 5   const int MAXN=1000010;
 6   bool notprime[MAXN];//值为 false 表示素数，值为 true 表示非素数
 7   void init(){
 8       memset(notprime,false,sizeof(notprime));
 9       notprime[0]=notprime[1]=true;
10       for(int i=2;i<MAXN;i++)
11           if(!notprime[i]){
12               if(i>MAXN/i)continue;//防止后面 i*i 溢出 (或者 i,j 用 long
                    long)
13               //直接从 i*i 开始就可以，小于 i 倍的已经筛选过了, 注意是 j+=i
14               for(int j=i*i;j<MAXN;j+=i)
15                   notprime[j]=true;
16           }
17   }
\end{Code}

\subsubsection{Sàng số nguyên tố (liệt kê các số nguyên tố không vượt quá MAXN)}

\begin{Code}
 1   /*
 2    * 素数筛选，存在小于等于 MAXN 的素数
 3    * prime[0] 存的是素数的个数
 4    */
 5   const int MAXN=10000;
 6   int prime[MAXN+1];
 7   void getPrime(){
 8       memset(prime,0,sizeof(prime));
 9       for(int i=2;i<=MAXN;i++){
10           if(!prime[i])prime[++prime[0]]=i;
11           for(int j=1;j<=prime[0]&&prime[j]<=MAXN/i;j++){
12               prime[prime[j]*i]=1;
13               if(i%prime[j]==0) break;
14           }
15       }
16   }
\end{Code}

\subsubsection{Sàng nguyên tố trên khoảng lớn (POJ 2689)}

\begin{Code}
1 /*
2 * POJ 2689 Prime Distance
3 * 给出一个区间 [L,U]，找出区间内容、相邻的距离最近的两个素数和
4 * 距离最远的两个素数。
5 * 1<=L<U<=2,147,483,647 区间长度不超过 1,000,000
6 * 就是要筛选出 [L,U] 之间的素数
7 */


 8 const int MAXN=100010;
 9 int prime[MAXN+1];
10 void getPrime(){
11     memset(prime,0,sizeof(prime));
12     for(int i=2;i<=MAXN;i++){
13         if(!prime[i])prime[++prime[0]]=i;
14         for(int j=1;j<=prime[0]&&prime[j]<=MAXN/i;j++){
15              prime[prime[j]*i]=1;
16              if(i%prime[j]==0)break;
17         }
18     }
19 }
20 bool notprime[1000010];
21 int prime2[1000010];
22 void getPrime2(int L,int R){
23     memset(notprime,false,sizeof(notprime));
24     if(L<2)L=2;
25     for(int i=1;i<=prime[0]&&(long long)prime[i]*prime[i]<=R;i++){
26         int s=L/prime[i]+(L%prime[i]>0);
27         if(s==1)s=2;
28         for(int j=s;(long long)j*prime[i]<=R;j++)
29              if((long long)j*prime[i]>=L)
30                  notprime[j*prime[i]-L]=true;
31     }
32     prime2[0]=0;
33     for(int i=0;i<=R-L;i++)
34         if(!notprime[i])
35              prime2[++prime2[0]]=i+L;
36 }
37 int main(){
38     getPrime();
39     int L,U;
40     while(scanf("%d%d",&L,&U)==2){
41         getPrime2(L,U);
42         if(prime2[0]<2)printf("There are no adjacent primes.\n");
43         else{
44              int x1=0,x2=100000000,y1=0,y2=0;
45              for(int i=1;i<prime2[0];i++){
46                  if(prime2[i+1]-prime2[i]<x2-x1){
47                      x1=prime2[i];
48                      x2=prime2[i+1];
49                  }
50                  if(prime2[i+1]-prime2[i]>y2-y1){
51                      y1=prime2[i];
52                      y2=prime2[i+1];
53                  }
54              }
55              printf("%d,%d are closest, %d,%d are most distant.\n",
                   x1,x2,y1,y2);
56         }
57     }


58 }
\end{Code}

\subsection{Sàng số nguyên tố và phân tích hợp số}

\begin{Code}
 1   //******************************************
 2   //素数筛选和合数分解
 3   const int MAXN=10000;
 4   int prime[MAXN+1];
 5   void getPrime(){
 6       memset(prime,0,sizeof(prime));
 7       for(int i=2;i<=MAXN;i++){
 8           if(!prime[i])prime[++prime[0]]=i;
 9           for(int j=1;j<=prime[0]&&prime[j]<=MAXN/i;j++){
10               prime[prime[j]*i]=1;
11               if(i%prime[j]==0) break;
12           }
13       }
14   }
15   long long factor[100][2];
16   int fatCnt;
17   int getFactors(long long x){
18       fatCnt=0;
19       long long tmp=x;
20       for(int i=1;prime[i]<=tmp/prime[i];i++){
21           factor[fatCnt][1]=0;
22           if(tmp%prime[i]==0){
23               factor[fatCnt][0]=prime[i];
24               while(tmp%prime[i]==0){
25                   factor[fatCnt][1]++;
26                   tmp/=prime[i];
27               }
28               fatCnt++;
29           }
30       }
31       if(tmp!=1){
32           factor[fatCnt][0]=tmp;
33           factor[fatCnt++][1]=1;
34       }
35       return fatCnt;
36   }
37   //******************************************
\end{Code}

\subsection{Thuật toán Euclid mở rộng (tìm nghiệm của ax+by=gcd và nghịch đảo)}

\begin{Code}
1 //******************************
2 //返回 d=gcd(a,b); 和对应于等式 ax+by=d 中的 x,y
3 long long extend_gcd(long long a,long long b,long long &x,long long
      &y){
4     if(a==0&&b==0) return -1;//无最大公约数
5     if(b==0){x=1;y=0;return a;}
6     long long d=extend_gcd(b,a%b,y,x);
7     y-=a/b*x;


 8       return d;
 9   }
10   //********* 求逆元 *******************
11   //ax = 1(mod n)
12   long long mod_reverse(long long a,long long n){
13       long long x,y;
14       long long d=extend_gcd(a,n,x,y);
15       if(d==1) return (x%n+n)%n;
16       else return -1;
17   }
\end{Code}

\subsection{Tìm nghịch đảo}

\subsubsection{Phương pháp Euclid mở rộng}

\begin{Code}
     见上面的写法
\end{Code}

\subsubsection{Cách viết ngắn gọn}

\begin{Code}
     注意：这个只能求 a < m 的情况，而且必须保证 a 和 m 互质
1 //求 ax = 1( mod m) 的 x 值，就是逆元 (0<a<m)
2 long long inv(long long a,long long m){
3     if(a == 1)return 1;
4     return inv(m%a,m)*(m-m/a)%m;
5 }
\end{Code}

\subsubsection{Dùng hàm Euler}

\begin{Code}
     mod 为素数, 而且 a 和 m 互质
1 long long inv(long long a,long long mod)//为素数mod
2 {
3     return pow_m(a,mod-2,mod);
4 }
\end{Code}

\subsection{Hệ phương trình tuyến tính modulo}

\begin{Code}
 1 long long extend_gcd(long long a,long long b,long long &x,long long
        &y){
 2     if(a == 0 && b == 0)return -1;
 3     if(b ==0 ){x = 1; y = 0;return a;}
 4     long long d = extend_gcd(b,a%b,y,x);
 5     y -= a/b*x;
 6     return d;
 7 }
 8 int m[10],a[10];//模数为 m, 余数为 a,X % m = a
 9 bool solve(int &m0,int &a0,int m,int a){
10     long long y,x;
11     int g = extend_gcd(m0,m,x,y);
12     if( abs(a - a0)%g )return false;
13     x *= (a - a0)/g;
14     x %= m/g;


15         a0 = (x*m0 + a0);
16         m0 *= m/g;
17         a0 %= m0;
18         if( a0 < 0 )a0 += m0;
19         return true;
20   }
21   /*
22     * 无解返回 false, 有解返回 true;
23     * 解的形式最后为 a0 + m0 * t (0<=a0<m0)
24     */
25   bool MLES(int &m0 ,int &a0,int n)//解为    X = a0 + m0 * k
26   {
27        bool flag = true;
28        m0 = 1;
29        a0 = 0;
30        for(int i = 0;i < n;i++)
31            if( !solve(m0,a0,m[i],a[i]) )
32        {
33            flag = false;
34            break;
35        }
36        return flag;
37   }
\end{Code}

\subsection{Kiểm tra nguyên tố ngẫu nhiên và phân tích số lớn (POJ 1811)}

\begin{Code}
 1   /* *************************************************
 2    * Miller_Rabin 算法进行素数测试
 3    * 速度快可以判断一个 < 2^63 的数是不是素数
 4    *
 5    **************************************************/
 6   const int S = 8; //随机算法判定次数一般 8∼10 就够了
 7   // 计算 ret = (a*b)%c      a,b,c < 2^63
 8   long long mult_mod(long long a,long long b,long long c){
 9       a %= c;
10       b %= c;
11       long long ret = 0;
12       long long tmp = a;
13       while(b){
14           if(b & 1){
15               ret += tmp;
16               if(ret > c)ret -= c;//直接取模慢很多
17           }
18           tmp <<= 1;
19           if(tmp > c)tmp -= c;
20           b >>= 1;
21       }
22       return ret;
23   }
24   // 计算 ret = (a^n)%mod
25   long long pow_mod(long long a,long long n,long long mod){
26       long long ret = 1;


27       long long temp = a%mod;
28       while(n){
29           if(n & 1)ret = mult_mod(ret,temp,mod);
30           temp = mult_mod(temp,temp,mod);
31           n >>= 1;
32       }
33       return ret;
34   }
35   // 通过 a^(n-1)=1(mod n)来判断 n 是不是素数
36   // n - 1 = x ∗ 2t 中间使用二次判断
37   // 是合数返回 true, 不一定是合数返回 false
38   bool check(long long a,long long n,long long x,long long t){
39       long long ret = pow_mod(a,x,n);
40       long long last = ret;
41       for(int i = 1;i <= t;i++){
42            ret = mult_mod(ret,ret,n);
43            if(ret == 1 && last != 1 && last != n-1)return true;//合数
44            last = ret;
45       }
46       if(ret != 1)return true;
47       else return false;
48   }
49   //**************************************************
50   // Miller_Rabin 算法
51   // 是素数返回 true,(可能是伪素数)
52   // 不是素数返回 false
53   //**************************************************
54   bool Miller_Rabin(long long n){
55       if( n < 2)return false;
56       if( n == 2)return true;
57       if( (n&1) == 0)return false;//偶数
58       long long x = n - 1;
59       long long t = 0;
60       while( (x&1)==0 ){x >>= 1; t++;}
61
62       srand(time(NULL));/* *************** */
63
64       for(int i = 0;i < S;i++){
65           long long a = rand()%(n-1) + 1;
66           if( check(a,n,x,t) )
67               return false;
68       }
69       return true;
70   }
71
72   //**********************************************
73   // pollard_rho 算法进行质因素分解
74   //*********************************************
75   long long factor[100];//质因素分解结果（刚返回时时无序的）
76   int tol;//质因素的个数，编号 0∼tol-1
77


 78   long long gcd(long long a,long long b){
 79       long long t;
 80       while(b){
 81           t = a;
 82           a = b;
 83           b = t%b;
 84       }
 85       if(a >= 0)return a;
 86       else return -a;
 87   }
 88
 89    //找出一个因子
 90   long long pollard_rho(long long x,long long c){
 91        long long i = 1, k = 2;
 92        srand(time(NULL));
 93        long long x0 = rand()%(x-1) + 1;
 94        long long y = x0;
 95        while(1){
 96            i ++;
 97            x0 = (mult_mod(x0,x0,x) + c)%x;
 98            long long d = gcd(y - x0,x);
 99            if( d != 1 && d != x)return d;
100            if(y == x0)return x;
101            if(i == k){y = x0; k += k;}
102        }
103   }
104     //对 n 进行素因子分解，存入 factor. k 设置为 107 左右即可
105   void findfac(long long n,int k){
106        if(n == 1)return;
107        if(Miller_Rabin(n))
108        {
109            factor[tol++] = n;
110            return;
111        }
112        long long p = n;
113        int c = k;
114        while( p >= n)p = pollard_rho(p,c--);//值变化，防止死循环 k
115        findfac(p,k);
116        findfac(n/p,k);
117   }
118     //POJ 1811
119     //给出一个N(2 <= N < 2^54),如果是素数, 输出"Prime", 否则输出最小的素因子
120   int main(){
121        int T;
122        long long n;
123        scanf("%d",&T);
124        while(T--){
125            scanf("%I64d",&n);
126            if(Miller_Rabin(n))printf("Prime\n");
127            else{
128                tol = 0;


129                findfac(n,107);
130                long long ans = factor[0];
131                for(int i = 1;i < tol;i++)
132                ans = min(ans,factor[i]);
133                printf("%I64d\n",ans);
134            }
135        }
136        return 0;
137 }
\end{Code}

\subsection{Hàm Euler}

\subsubsection{Tính hàm Euler bằng phân tích thừa số nguyên tố}

\begin{Code}
 1   getFactors(n);
 2   int ret = n;
 3   for(int i = 0;i < fatCnt;i++){
 4       ret = ret/factor[i][0]*(factor[i][0]-1);
 5   }
\end{Code}

\subsubsection{Tính hàm Euler bằng sàng}

\begin{Code}
 1 int euler[3000001];
 2 void getEuler(){
 3     memset(euler,0,sizeof(euler));
 4     euler[1] = 1;
 5     for(int i = 2;i <= 3000000;i++)
 6         if(!euler[i])
 7            for(int j = i;j <= 3000000; j += i){
 8                if(!euler[j])
 9                     euler[j] = j;
10                euler[j] = euler[j]/i*(i-1);
11            }
12 }
\end{Code}

\subsubsection{Tính hàm Euler cho một số}

\begin{Code}
 1 long long eular(long long n){
 2     long long ans = n;
 3     for(int i = 2;i*i <= n;i++){
 4         if(n % i == 0){
 5             ans -= ans/i;
 6             while(n % i == 0)
 7                 n /= i;
 8         }
 9     }
10     if(n > 1)ans -= ans/n;
11     return ans;
12 }
\end{Code}

\subsubsection{Sàng tuyến tính (đồng thời lấy hàm Euler và bảng nguyên tố)}

\begin{Code}
 1 const int MAXN = 10000000;
 2 bool check[MAXN+10];


 3   int phi[MAXN+10];
 4   int prime[MAXN+10];
 5   int tot;//素数的个数
 6   void phi_and_prime_table(int N){
 7       memset(check,false,sizeof(check));
 8       phi[1] = 1;
 9       tot = 0;
10       for(int i = 2; i <= N; i++){
11           if( !check[i] ){
12                prime[tot++] = i;
13                phi[i] = i-1;
14           }
15           for(int j = 0; j < tot; j++){
16                if(i * prime[j] > N)break;
17                check[i * prime[j]] = true;
18                if( i % prime[j] == 0){
19                    phi[i * prime[j]] = phi[i] * prime[j];
20                    break;
21                }
22                else{
23                    phi[i * prime[j]] = phi[i] * (prime[j] - 1);
24                }
25           }
26       }
27   }
\end{Code}

\subsection{Khử Gauss (số thực)}

\begin{Code}
 1 #define eps 1e-9
 2 const int MAXN=220;
 3 double a[MAXN][MAXN],x[MAXN];//方程的左边的矩阵和等式右边的值, 求解之后 x
      存的就是结果
 4 int equ,var;//方程数和未知数个数
 5 /*
 6 * 返回 0 表示无解,1 表示有解
 7 */
 8 int Gauss(){
 9     int i,j,k,col,max_r;
10     for(k=0,col=0;k<equ&&col<var;k++,col++){
11         max_r=k;
12         for(i=k+1;i<equ;i++)
13           if(fabs(a[i][col])>fabs(a[max_r][col]))
14              max_r=i;
15         if(fabs(a[max_r][col])<eps)return 0;
16         if(k!=max_r){
17              for(j=col;j<var;j++)
18                swap(a[k][j],a[max_r][j]);
19              swap(x[k],x[max_r]);
20         }
21         x[k]/=a[k][col];
22         for(j=col+1;j<var;j++)a[k][j]/=a[k][col];
23         a[k][col]=1;


24              for(i=0;i<equ;i++)
25                if(i!=k){
26                    x[i]-=x[k]*a[i][col];
27                    for(j=col+1;j<var;j++)a[i][j]-=a[k][j]*a[i][col];
28                    a[i][col]=0;
29                }
30       }
31       return 1;
32 }
\end{Code}

\subsection{FFT}

\begin{Code}
 1   //HDU 1402 求高精度乘法
 2   const double PI = acos(-1.0);
 3   //复数结构体
 4   struct Complex{
 5        double x,y;//实部和虚部 x+yi
 6        Complex(double _x = 0.0,double _y = 0.0){
 7            x = _x;
 8            y = _y;
 9        }
10        Complex operator -(const Complex &b)const{
11            return Complex(x-b.x,y-b.y);
12        }
13        Complex operator +(const Complex &b)const{
14            return Complex(x+b.x,y+b.y);
15        }
16        Complex operator *(const Complex &b)const{
17            return Complex(x*b.x-y*b.y,x*b.y+y*b.x);
18        }
19   };
20   /*
21     * 进行 FFT 和 IFFT 前的反转变换。
22     * 位置 i 和（i 二进制反转后位置）互换
23     * len 必须为 2 的幂
24     */
25   void change(Complex y[],int len){
26        int i,j,k;
27        for(i = 1, j = len/2;i <len-1;i++){
28            if(i < j)swap(y[i],y[j]);
29            //交换互为小标反转的元素，i<j 保证交换一次
30            //i 做正常的 +1，j 左反转类型的 +1, 始终保持 i 和 j 是反转的
31            k = len/2;
32            while(j >= k){
33                j -= k;
34                k /= 2;
35            }
36            if(j < k)j += k;
37        }
38   }
39   /*
40     * 做 FFT


41     * len 必须为2^k形式
42     * on==1 时是 DFT，on==-1 时是 IDFT
43     */
44   void fft(Complex y[],int len,int on){
45        change(y,len);
46        for(int h = 2; h <= len; h <<= 1){
47            Complex wn(cos(-on*2*PI/h),sin(-on*2*PI/h));
48            for(int j = 0;j < len;j+=h){
49                Complex w(1,0);
50                for(int k = j;k < j+h/2;k++){
51                    Complex u = y[k];
52                    Complex t = w*y[k+h/2];
53                    y[k] = u+t;
54                    y[k+h/2] = u-t;
55                    w = w*wn;
56                }
57            }
58        }
59        if(on == -1)
60            for(int i = 0;i < len;i++)
61                y[i].x /= len;
62   }
63   const int MAXN = 200010;
64   Complex x1[MAXN],x2[MAXN];
65   char str1[MAXN/2],str2[MAXN/2];
66   int sum[MAXN];
67   int main(){
68        while(scanf("%s%s",str1,str2)==2){
69            int len1 = strlen(str1);
70            int len2 = strlen(str2);
71            int len = 1;
72            while(len < len1*2 || len < len2*2)len<<=1;
73            for(int i = 0;i < len1;i++)
74                x1[i] = Complex(str1[len1-1-i]-'0',0);
75            for(int i = len1;i < len;i++)
76                x1[i] = Complex(0,0);
77            for(int i = 0;i < len2;i++)
78                x2[i] = Complex(str2[len2-1-i]-'0',0);
79            for(int i = len2;i < len;i++)
80                x2[i] = Complex(0,0);
81            //求 DFT
82            fft(x1,len,1);
83            fft(x2,len,1);
84            for(int i = 0;i < len;i++)
85                x1[i] = x1[i]*x2[i];
86            fft(x1,len,-1);
87            for(int i = 0;i < len;i++)
88                sum[i] = (int)(x1[i].x+0.5);
89            for(int i = 0;i < len;i++){
90                sum[i+1]+=sum[i]/10;
91                sum[i]%=10;


 92              }
 93              len = len1+len2-1;
 94              while(sum[len] <= 0 && len > 0)len--;
 95              for(int i = len;i >= 0;i--)
 96                  printf("%c",sum[i]+'0');
 97              printf("\n");
 98       }
 99       return 0;
100   }
101
102   //HDU 4609
103   //给出 n 条线段长度，问任取 3 根，组成三角形的概率。
104   //n<=10^5   用 FFT 求可以组成三角形的取法有几种
105   const int MAXN = 400040;
106   Complex x1[MAXN];
107   int a[MAXN/4];
108   long long num[MAXN];//100000*100000 会超 int
109   long long sum[MAXN];
110   int main(){
111       int T;
112       int n;
113       scanf("%d",&T);
114       while(T--){
115           scanf("%d",&n);
116           memset(num,0,sizeof(num));
117           for(int i = 0;i < n;i++){
118               scanf("%d",&a[i]);
119               num[a[i]]++;
120           }
121           sort(a,a+n);
122           int len1 = a[n-1]+1;
123           int len = 1;
124           while( len < 2*len1 )len <<= 1;
125           for(int i = 0;i < len1;i++)
126               x1[i] = Complex(num[i],0);
127           for(int i = len1;i < len;i++)
128               x1[i] = Complex(0,0);
129           fft(x1,len,1);
130           for(int i = 0;i < len;i++)
131               x1[i] = x1[i]*x1[i];
132           fft(x1,len,-1);
133           for(int i = 0;i < len;i++)
134               num[i] = (long long)(x1[i].x+0.5);
135           len = 2*a[n-1];
136           //减掉取两个相同的组合
137           for(int i = 0;i < n;i++)
138               num[a[i]+a[i]]--;
139           for(int i = 1;i <= len;i++)num[i]/=2;
140           sum[0] = 0;
141           for(int i = 1;i <= len;i++)
142               sum[i] = sum[i-1]+num[i];


143            long long cnt = 0;
144            for(int i = 0;i < n;i++){
145                cnt += sum[len]-sum[a[i]];
146                //减掉一个取大，一个取小的
147                cnt -= (long long)(n-1-i)*i;
148                //减掉一个取本身，另外一个取其它
149                cnt -= (n-1);
150                cnt -= (long long)(n-1-i)*(n-i-2)/2;
151            }
152            long long tot = (long long)n*(n-1)*(n-2)/6;
153            printf("%.7lf\n",(double)cnt/tot);
154     }
155     return 0;
156 }
\end{Code}

\subsection{Dùng khử Gauss để giải hệ phương trình}

\subsubsection{Một dạng bài công tắc, hệ phương trình 0/1 modulo 2}

\begin{Code}
    POJ 1681 需要枚举自由变元，找解中 1 个数最少的
 1 //对 2 取模的 01 方程组
 2 const int MAXN = 300;
 3 //有 equ 个方程，var 个变元。增广矩阵行数为 equ, 列数为 var+1, 分别为 0 到
      var
 4 int equ,var;
 5 int a[MAXN][MAXN]; //增广矩阵
 6 int x[MAXN]; //解集
 7 int free_x[MAXN];//用来存储自由变元（多解枚举自由变元可以使用）
 8 int free_num;//自由变元的个数
 9
10 //返回值为 -1 表示无解，为 0 是唯一解，否则返回自由变元个数
11 int Gauss(){
12     int max_r,col,k;
13     free_num = 0;
14     for(k = 0, col = 0 ; k < equ && col < var ; k++, col++){
15         max_r = k;
16         for(int i = k+1;i < equ;i++){
17              if(abs(a[i][col]) > abs(a[max_r][col]))
18                  max_r = i;
19         }
20         if(a[max_r][col] == 0){
21              k--;
22              free_x[free_num++] = col;//这个是自由变元
23              continue;
24         }
25         if(max_r != k){
26              for(int j = col; j < var+1; j++)
27                  swap(a[k][j],a[max_r][j]);
28         }
29         for(int i = k+1;i < equ;i++){
30              if(a[i][col] != 0){
31                  for(int j = col;j < var+1;j++)


32                       a[i][j] ^= a[k][j];
33                  }
34              }
35       }
36       for(int i = k;i < equ;i++)
37           if(a[i][col] != 0)
38               return -1;//无解
39       if(k < var) return var-k;//自由变元个数
40       //唯一解，回代
41       for(int i = var-1; i >= 0;i--){
42           x[i] = a[i][var];
43           for(int j = i+1;j < var;j++)
44               x[i] ^= (a[i][j] && x[j]);
45       }
46       return 0;
47   }
48   int n;
49   void init(){
50       memset(a,0,sizeof(a));
51       memset(x,0,sizeof(x));
52       equ = n*n;
53       var = n*n;
54       for(int i = 0;i < n;i++)
55           for(int j = 0;j < n;j++){
56               int t = i*n+j;
57               a[t][t] = 1;
58               if(i > 0)a[(i-1)*n+j][t] = 1;
59               if(i < n-1)a[(i+1)*n+j][t] = 1;
60               if(j > 0)a[i*n+j-1][t] = 1;
61               if(j < n-1)a[i*n+j+1][t] = 1;
62           }
63   }
64   void solve(){
65       int t = Gauss();
66       if(t == -1){
67           printf("inf\n");
68           return;
69       }
70       else if(t == 0){
71           int ans = 0;
72           for(int i = 0;i < n*n;i++)
73               ans += x[i];
74           printf("%d\n",ans);
75           return;
76       }
77       else
78       {
79           //枚举自由变元
80           int ans = 0x3f3f3f3f;
81           int tot = (1<<t);
82           for(int i = 0;i < tot;i++){


 83                  int cnt = 0;
 84                  for(int j = 0;j < t;j++){
 85                      if(i&(1<<j)){
 86                          x[free_x[j]] = 1;
 87                          cnt++;
 88                      }
 89                      else x[free_x[j]] = 0;
 90                  }
 91                  for(int j = var-t-1;j >= 0;j--){
 92                      int idx;
 93                      for(idx = j;idx < var;idx++)
 94                          if(a[j][idx])
 95                               break;
 96                      x[idx] = a[j][var];
 97                      for(int l = idx+1;l < var;l++)
 98                          if(a[j][l])
 99                               x[idx] ^= x[l];
100                      cnt += x[idx];
101                  }
102                  ans = min(ans,cnt);
103              }
104              printf("%d\n",ans);
105       }
106   }
107   char str[30][30];
108   int main(){
109       int T;
110       scanf("%d",&T);
111       while(T--){
112           scanf("%d",&n);
113           init();
114           for(int i = 0;i < n;i++){
115               scanf("%s",str[i]);
116               for(int j = 0;j < n;j++){
117                    if(str[i][j] == 'y')
118                        a[i*n+j][n*n] = 0;
119                    else a[i*n+j][n*n] = 1;
120               }
121           }
122           solve();
123       }
124       return 0;
125   }
\end{Code}

\subsubsection{Giải hệ phương trình đồng dư}

\begin{Code}
      POJ 2947 Widget Factory
 1    //求解对 MOD 取模的方程组
 2    const int MOD = 7;
 3    const int MAXN = 400;
 4    int a[MAXN][MAXN];//增广矩阵


 5   int x[MAXN];//最后得到的解集
 6   inline int gcd(int a,int b){
 7       while(b != 0){
 8           int t = b;
 9           b = a%b;
10           a = t;
11       }
12       return a;
13   }
14   inline int lcm(int a,int b){
15       return a/gcd(a,b)*b;
16   }
17   long long inv(long long a,long long m){
18       if(a == 1)return 1;
19       return inv(m%a,m)*(m-m/a)%m;
20   }
21   int Gauss(int equ,int var){
22       int max_r,col,k;
23       for(k = 0, col = 0; k < equ && col < var; k++,col++){
24           max_r = k;
25           for(int i = k+1; i < equ;i++)
26               if(abs(a[i][col]) > abs(a[max_r][col]))
27                    max_r = i;
28           if(a[max_r][col] == 0){
29               k--;
30               continue;
31           }
32           if(max_r != k)
33               for(int j = col; j < var+1;j++)
34                    swap(a[k][j],a[max_r][j]);
35           for(int i = k+1;i < equ;i++){
36               if(a[i][col] != 0){
37                    int LCM = lcm(abs(a[i][col]),abs(a[k][col]));
38                    int ta = LCM/abs(a[i][col]);
39                    int tb = LCM/abs(a[k][col]);
40                    if(a[i][col]*a[k][col] < 0)tb = -tb;
41                    for(int j = col;j < var+1;j++)
42                        a[i][j] = ((a[i][j]*ta - a[k][j]*tb)%MOD + MOD)
                             %MOD;
43               }
44           }
45       }
46       for(int i = k;i < equ;i++)
47           if(a[i][col] != 0)
48               return -1;//无解
49       if(k < var) return var-k;//多解
50       for(int i = var-1;i >= 0;i--){
51           int temp = a[i][var];
52           for(int j = i+1; j < var;j++){
53               if(a[i][j] != 0){
54                    temp -= a[i][j]*x[j];


55                       temp = (temp%MOD + MOD)%MOD;
56                   }
57               }
58               x[i] = (temp*inv(a[i][i],MOD))%MOD;
59        }
60        return 0;
61    }
62    int change(char s[]){
63        if(strcmp(s,"MON") == 0) return 1;
64        else if(strcmp(s,"TUE")==0) return 2;
65        else if(strcmp(s,"WED")==0) return 3;
66        else if(strcmp(s,"THU")==0) return 4;
67        else if(strcmp(s,"FRI")==0) return 5;
68        else if(strcmp(s,"SAT")==0) return 6;
69        else return 7;
70    }
71    int main(){
72        int n,m;
73        while(scanf("%d%d",&n,&m) == 2){
74            if(n == 0 && m == 0)break;
75            memset(a,0,sizeof(a));
76            char str1[10],str2[10];
77            int k;
78            for(int i = 0;i < m;i++){
79                 scanf("%d%s%s",&k,str1,str2);
80                 a[i][n] = ((change(str2) - change(str1) + 1)%MOD + MOD)
                      %MOD;
 81                int t;
 82                while(k--){
 83                    scanf("%d",&t);
 84                    t--;
 85                    a[i][t] ++;
 86                    a[i][t]%=MOD;
 87                }
 88           }
 89           int ans = Gauss(m,n);
 90           if(ans == 0){
 91                for(int i = 0;i < n;i++)
 92                    if(x[i] <= 2)
 93                         x[i] += 7;
 94                for(int i = 0;i < n-1;i++)printf("%d ",x[i]);
 95                printf("%d\n",x[n-1]);
 96           }
 97           else if(ans == -1)printf("Inconsistent data.\n");
 98           else printf("Multiple solutions.\n");
 99       }
100       return 0;
101   }
\end{Code}

\subsection{Phân hoạch số nguyên}

\begin{Code}
1    //HDU 4651
2    //把数 n 拆成几个数（小于等于 n）相加的形式，问有多少种拆法。
3    const int MOD = 1e9+7;
4    int dp[100010];
5    void init(){
6        memset(dp,0,sizeof(dp));
7        dp[0] = 1;
8        for(int i = 1;i <= 100000;i++){
9            for(int j = 1, r = 1; i - (3 * j * j - j) / 2 >= 0; j++, r
                *= -1){
10               dp[i] += dp[i -(3 * j * j - j) / 2] * r;
11               dp[i] %= MOD;
12               dp[i] = (dp[i]+MOD)%MOD;
13               if( i - (3 * j * j + j) / 2 >= 0 ){
14                   dp[i] += dp[i - (3 * j * j + j) / 2] * r;
15                   dp[i] %= MOD;
16                   dp[i] = (dp[i]+MOD)%MOD;
17               }
18           }
19       }
20   }
21   int main(){
22       int T;
23       int n;
24       init();
25       scanf("%d",&T);
26       while(T--){
27           scanf("%d",&n);
28           printf("%d\n",dp[n]);
29       }
30       return 0;
31   }
32
33   //HDU 4658
34   //数 n(<=10^5) 的划分,相同的数重复不能超过 k 个。
35   const int MOD = 1e9+7;
36   int dp[100010];
37   void init(){
38       memset(dp,0,sizeof(dp));
39       dp[0] = 1;
40       for(int i = 1;i <= 100000;i++){
41           for(int j = 1, r = 1; i - (3 * j * j - j) / 2 >= 0; j++, r
                *= -1){
42               dp[i] += dp[i -(3 * j * j - j) / 2] * r;
43               dp[i] %= MOD;
44               dp[i] = (dp[i]+MOD)%MOD;
45               if( i - (3 * j * j + j) / 2 >= 0 ){
46                   dp[i] += dp[i - (3 * j * j + j) / 2] * r;
47                   dp[i] %= MOD;
48                   dp[i] = (dp[i]+MOD)%MOD;
49               }


50         }
51     }
52 }
53 int solve(int n,int k){
54     int ans = dp[n];
55     for(int j = 1, r = -1; n - k*(3 * j * j - j) / 2 >= 0; j++, r
          *= -1){
56         ans += dp[n -k*(3 * j * j - j) / 2] * r;
57         ans %= MOD;
58         ans = (ans+MOD)%MOD;
59         if( n - k*(3 * j * j + j) / 2 >= 0 ){
60              ans += dp[n - k*(3 * j * j + j) / 2] * r;
61              ans %= MOD;
62              ans = (ans+MOD)%MOD;
63         }
64     }
65     return ans;
66 }
67 int main(){
68     init();
69     int T;
70     int n,k;
71     scanf("%d",&T);
72     while(T--){
73         scanf("%d%d",&n,&k);
74         printf("%d\n",solve(n,k));
75     }
76     return 0;
77 }
\end{Code}

\subsection{Tính tổng ước của A\textasciicircum{}B modulo MOD}

\begin{Code}
 1   //参考 POJ 1845
 2   //里面有一种求1+p+p^2+p…^3+p^n的方法。
 3   //需要素数筛选和合数分解的程序，需要先调用 getPrime();
 4   long long pow_m(long long a,long long n){
 5       long long ret = 1;
 6       long long tmp = a%MOD;
 7       while(n){
 8           if(n&1)ret = (ret*tmp)%MOD;
 9           tmp = tmp*tmp%MOD;
10           n >>= 1;
11       }
12       return ret;
13   }
14   //计算1+p+p^2+...+p^n
15   long long sum(long long p,long long n){
16       if(p == 0)return 0;
17       if(n == 0)return 1;
18       if(n & 1){
19           return ((1+pow_m(p,n/2+1))%MOD*sum(p,n/2)%MOD)%MOD;
20       }


21       else return ((1+pow_m(p,n/2+1))%MOD*sum(p,n/2-1)+pow_m(p,n/2)%
            MOD)%MOD;
22   }
23   //返回A^B的约数之和 % MOD
24   long long solve(long long A,long long B){
25       getFactors(A);
26       long long ans = 1;
27       for(int i = 0;i < fatCnt;i++){
28           ans *= sum(factor[i][0],B*factor[i][1])%MOD;
29           ans %= MOD;
30       }
31       return ans;
32   }
\end{Code}

\subsection{Nghịch đảo Mobius}

\subsubsection{Hàm Mobius}

\begin{Code}
 1   const int MAXN = 1000000;
 2   bool check[MAXN+10];
 3   int prime[MAXN+10];
 4   int mu[MAXN+10];
 5   void Moblus(){
 6       memset(check,false,sizeof(check));
 7       mu[1] = 1;
 8       int tot = 0;
 9       for(int i = 2; i <= MAXN; i++){
10           if( !check[i] ){
11               prime[tot++] = i;
12               mu[i] = -1;
13           }
14           for(int j = 0; j < tot; j++){
15               if(i * prime[j] > MAXN) break;
16               check[i * prime[j]] = true;
17               if( i % prime[j] == 0){
18                    mu[i * prime[j]] = 0;
19                    break;
20               }
21               else{
22                    mu[i * prime[j]] = -mu[i];
23               }
24           }
25       }
26   }
\end{Code}

\subsubsection{Ví dụ: BZOJ2301}

\begin{Code}
     对于给出的 n 个询问，每次求有多少个数对 ( x, y)，满足 a <= x <= b, c <= y <= d，且
     gcd( x, y) = k，gcd( x, y) 函数为 x 和 y 的最大公约数。1 <= n <= 50000, 1 <= a <= b <=
     50000, 1 <= c <= d <= 50000, 1 <= k <= 50000
1 const int MAXN = 100000;
2 bool check[MAXN+10];
3 int prime[MAXN+10];


 4   int mu[MAXN+10];
 5   void Moblus(){
 6       memset(check,false,sizeof(check));
 7       mu[1] = 1;
 8       int tot = 0;
 9       for(int i = 2; i <= MAXN; i++){
10           if( !check[i] ){
11               prime[tot++] = i;
12               mu[i] = -1;
13           }
14           for(int j = 0; j < tot; j ++){
15               if( i * prime[j] > MAXN) break;
16               check[i * prime[j]] = true;
17               if( i % prime[j] == 0){
18                    mu[i * prime[j]] = 0;
19                    break;
20               }
21               else{
22                    mu[i * prime[j]] = -mu[i];
23               }
24           }
25       }
26   }
27   int sum[MAXN+10];
28   //找 [1,n],[1,m] 内互质的数的对数
29   long long solve(int n,int m){
30       long long ans = 0;
31       if(n > m)swap(n,m);
32       for(int i = 1, la = 0; i <= n; i = la+1){
33           la = min(n/(n/i),m/(m/i));
34           ans += (long long)(sum[la] - sum[i-1])*(n/i)*(m/i);
35       }
36       return ans;
37   }
38   int main(){
39       Moblus();
40       sum[0] = 0;
41       for(int i = 1;i <= MAXN;i++)
42           sum[i] = sum[i-1] + mu[i];
43       int a,b,c,d,k;
44       int T;
45       scanf("%d",&T);
46       while(T--){
47           scanf("%d%d%d%d%d",&a,&b,&c,&d,&k);
48           long long ans = solve(b/k,d/k) - solve((a-1)/k,d/k) - solve
                (b/k,(c-1)/k) + solve((a-1)/k,(c-1)/k);
49           printf("%lld\n",ans);
50       }
51       return 0;
52   }
\end{Code}

\subsection{Baby-Step Giant-Step}

\begin{Code}
 1   //(POJ 2417,3243)
 2   //baby_step giant_step
 3   // a^x = b (mod n) n 是素数和不是素数都可以
 4   // 求解上式 0<=x < n 的解
 5   #define MOD 76543
 6   int hs[MOD],head[MOD],next[MOD],id[MOD],top;
 7   void insert(int x,int y){
 8       int k = x%MOD;
 9       hs[top] = x, id[top] = y, next[top] = head[k], head[k] = top++;
10   }
11   int find(int x){
12       int k = x%MOD;
13       for(int i = head[k]; i != -1; i = next[i])
14           if(hs[i] == x)
15                return id[i];
16       return -1;
17   }
18   int BSGS(int a,int b,int n){
19       memset(head,-1,sizeof(head));
20       top = 1;
21       if(b == 1)return 0;
22       int m = sqrt(n*1.0), j;
23       long long x = 1, p = 1;
24       for(int i = 0; i < m; ++i, p = p*a%n)insert(p*b%n,i);
25       for(long long i = m; ;i += m){
26           if( (j = find(x = x*p%n)) != -1 )return i-j;
27           if(i > n)break;
28       }
29       return -1;
30   }
\end{Code}

\subsection{Tích phân Simpson thích nghi}

\begin{Code}
 1   double simpson(double a,double b){
 2       double c = a + (b-a)/2;
 3       return (F(a) + 4*F(c) + F(b))*(b-a)/6;
 4   }
 5   double asr(double a,double b,double eps,double A){
 6       double c = a + (b-a)/2;
 7       double L = simpson(a,c), R = simpson(c,b);
 8       if(fabs(L + R - A) <= 15*eps)return L + R + (L + R - A)/15.0;
 9       return asr(a,c,eps/2,L) + asr(c,b,eps/2,R);
10   }
11   double asr(double a,double b,double eps){
12       return asr(a,b,eps,simpson(a,b));
13   }
\end{Code}

\subsection{Chu kỳ modulo của dãy Fibonacci}

\begin{Code}
     必要时要上 unsigned long long
     HDU3977


 1   long long gcd(long long a,long long b){
 2       if(b == 0)return a;
 3       return gcd(b,a%b);
 4   }
 5   long long lcm(long long a,long long b){
 6       return a/gcd(a,b)*b;
 7   }
 8   struct Matrix{
 9       long long mat[2][2];
10   };
11   Matrix mul_M(Matrix a,Matrix b,long long mod){
12       Matrix ret;
13       for(int i = 0;i < 2;i++)
14           for(int j = 0;j < 2;j++){
15               ret.mat[i][j] = 0;
16               for(int k = 0;k < 2;k++){
17                    ret.mat[i][j] += a.mat[i][k]*b.mat[k][j]%mod;
18                    if(ret.mat[i][j] >= mod)ret.mat[i][j] -= mod;
19               }
20           }
21       return ret;
22   }
23   Matrix pow_M(Matrix a,long long n,long long mod){
24       Matrix ret;
25       memset(ret.mat,0,sizeof(ret.mat));
26       for(int i = 0;i < 2;i++)ret.mat[i][i] = 1;
27       Matrix tmp = a;
28       while(n){
29           if(n&1)ret = mul_M(ret,tmp,mod);
30           tmp = mul_M(tmp,tmp,mod);
31           n >>= 1;
32       }
33       return ret;
34   }
35   long long pow_m(long long a,long long n,long long mod)//a^b % mod{
36       long long ret = 1;
37       long long tmp = a%mod;
38       while(n){
39           if(n&1)ret = ret*tmp%mod;
40           tmp = tmp*tmp%mod;
41           n >>= 1;
42       }
43       return ret;
44   }
45   //素数筛选和合数分解
46   const int MAXN = 1000000;
47   int prime[MAXN+1];
48   void getPrime(){
49       memset(prime,0,sizeof(prime));
50       for(int i = 2;i <= MAXN;i++){
51           if(!prime[i])prime[++prime[0]] = i;


 52              for(int j = 1;j <= prime[0] && prime[j] <= MAXN/i;j++){
 53                  prime[prime[j]*i] = 1;
 54                  if(i%prime[j] == 0)break;
 55              }
 56       }
 57   }
 58   long long factor[100][2];
 59   int fatCnt;
 60   int getFactors(long long x){
 61       fatCnt = 0;
 62       long long tmp = x;
 63       for(int i = 1;prime[i] <= tmp/prime[i];i++){
 64           factor[fatCnt][1] = 0;
 65           if(tmp%prime[i] == 0){
 66               factor[fatCnt][0] = prime[i];
 67               while(tmp%prime[i] == 0){
 68                   factor[fatCnt][1]++;
 69                   tmp /= prime[i];
 70               }
 71               fatCnt++;
 72           }
 73       }
 74       if(tmp != 1){
 75           factor[fatCnt][0] = tmp;
 76           factor[fatCnt++][1] = 1;
 77       }
 78       return fatCnt;
 79   }
 80   //勒让德符号
 81   int legendre(long long a,long long p){
 82       if(pow_m(a,(p-1)>>1,p) == 1)return 1;
 83       else return -1;
 84   }
 85   int f0 = 1;
 86   int f1 = 1;
 87   long long getFib(long long n,long long mod){
 88       if(mod == 1)return 0;
 89       Matrix A;
 90       A.mat[0][0] = 0;
 91       A.mat[1][0] = 1;
 92       A.mat[0][1] = 1;
 93       A.mat[1][1] = 1;
 94       Matrix B = pow_M(A,n,mod);
 95       long long ret = f0*B.mat[0][0] + f1*B.mat[1][0];
 96       return ret%mod;
 97   }
 98   long long fac[1000000];
 99   long long G(long long p){
100       long long num;
101       if(legendre(5,p) == 1)num = p-1;
102       else num = 2*(p+1);


103       //找出 num 的所有约数
104       int cnt = 0;
105       for(long long i = 1;i*i <= num;i++)
106           if(num%i == 0){
107               fac[cnt++] = i;
108               if(i*i != num)
109                    fac[cnt++] = num/i;
110           }
111       sort(fac,fac+cnt);
112       long long ans;
113       for(int i = 0;i < cnt;i++){
114           if(getFib(fac[i],p) == f0 && getFib(fac[i]+1,p) == f1){
115               ans = fac[i];
116               break;
117           }
118       }
119       return ans;
120   }
121   long long find_loop(long long n){
122       getFactors(n);
123       long long ans = 1;
124       for(int i = 0;i < fatCnt;i++){
125           long long record = 1;
126           if(factor[i][0] == 2)record = 3;
127           else if(factor[i][0] == 3)record = 8;
128           else if(factor[i][0] == 5)record = 20;
129           else record = G(factor[i][0]);
130           for(int j = 1;j < factor[i][1];j++)
131               record *= factor[i][0];
132           ans = lcm(ans,record);
133       }
134       return ans;
135   }
136   void init(){
137       getPrime();
138   }
139   int main(){
140       init();
141       int T;
142       int iCase = 0;
143       int n;
144       scanf("%d",&T);
145       while(T--){
146           iCase++;
147           scanf("%d",&n);
148           printf("Case #%d: %I64d\n",iCase,find_loop(n));
149       }
150       return 0;
151   }
\end{Code}

\subsection{Căn nguyên thủy}

\begin{Code}
     定义：设 m > 1,gcd( a, m) = 1, 使得 ad ≡ 1(mod m) 成立的最小的正整数 d 为 a 对模 m 的
     阶，记为 δm ( a).
     如果 δm ( a) = φ(m), 则称 a 是模 m 的原根.
     定理：若 m > 1,( a, m) = 1, 正整数 d 满足 ad ≡ 1(mod m), 则 δm ( a) 整除 d.
     定理：模 m 有原根的充要条件是 m = 2, 4, pn , 2pn , 其中 p 是奇质数，n 是任意正整数.
     定理：如果模 m 有原根，那么它一定有 φ( φ(m)) 个原根.
     定理：如果 p 是素数，那么素数 p 一定有原根，并且模 p 的原根的个数为 φ( p - 1).
                                                              a
     求模素数 p 原根的方法：对 p - 1 素因子分解，即 p - 1 = p1a1 p2a2 ...pkk 的标准分解式，若恒
     有
                                 p -1
                                g pi ̸= 1(mod p)
     成立，则 g 就是 p 的原根。（对于合数求原根，只需要把 p - 1 换成 φ( p)）即可。

     求素数的最小原根程序
 1   //******************************************
 2   //素数筛选和合数分解
 3   const int MAXN=100000;
 4   int prime[MAXN+1];
 5   void getPrime(){
 6       memset(prime,0,sizeof(prime));
 7       for(int i=2;i<=MAXN;i++){
 8           if(!prime[i])prime[++prime[0]]=i;
 9           for(int j=1;j<=prime[0]&&prime[j]<=MAXN/i;j++){
10               prime[prime[j]*i]=1;
11               if(i%prime[j]==0) break;
12           }
13       }
14   }
15   long long factor[100][2];
16   int fatCnt;
17   int getFactors(long long x){
18       fatCnt=0;
19       long long tmp=x;
20       for(int i=1;prime[i]<=tmp/prime[i];i++){
21           factor[fatCnt][1]=0;
22           if(tmp%prime[i]==0){
23               factor[fatCnt][0]=prime[i];
24               while(tmp%prime[i]==0){
25                   factor[fatCnt][1]++;
26                   tmp/=prime[i];
27               }
28               fatCnt++;
29           }
30       }
31       if(tmp!=1){
32           factor[fatCnt][0]=tmp;
33           factor[fatCnt++][1]=1;
34       }
35       return fatCnt;
36   }


37   //******************************************
38   long long pow_m(long long a,long long n,long long mod){
39       long long ret = 1;
40       long long tmp = a%mod;
41       while(n){
42           if(n&1)ret = ret*tmp%mod;
43           tmp = tmp*tmp%mod;
44           n >>= 1;
45       }
46       return ret;
47   }
48   //求素数 P 的最小的原根
49   void solve(int P){
50       if(P == 2){
51           printf("1\n");
52           return;
53       }
54       getFactors(P-1);
55       for(int g = 2; g < P;g++){
56           bool flag = true;
57           for(int i = 0;i < fatCnt;i++){
58               int t = (P-1)/factor[i][0];
59               if(pow_m(g,t,P) == 1){
60                    flag = false;
61                    break;
62               }
63           }
64           if(flag){
65               printf("%d\n",g);
66               return;
67           }
68       }
69   }
70   int main(){
71       getPrime();
72       int T;
73       int P;
74       scanf("%d",&T);
75       while(T--){
76           scanf("%d",&P);
77           solve(P);
78       }
79       return 0;
80   }
\end{Code}

\subsection{Biến đổi số học nhanh}

\subsubsection{HDU4656 tích chập modulo}

\begin{Code}
     HDU4656
     xk = b ∗ c(2k) + d, F ( x ) = a0 x0 + a1 x1 + a2 x2 + ... + an-1 xn-1 Given n, b, c, d, a0 , ..., an-1 ,


     calculate F ( x0 ), ..., F ( xn-1 ).

                                   n -1
                        Fxk = ∑ ai (bc2k + d)i
                                   i =0
                            n -1          i
                        = ∑ ai ∑ Ci (bc2k ) j di- j
                                              j

                            i =0      j =0
                            n -1                    n -1
                                          2k j -1
                        = ∑ (bc ) j!                ∑ ai di- j i!(i - j)!-1
                            j =0                    i= j
                            n -1                    n -1- j
                        = ∑ (bc2k ) j j!-1           ∑ an-1-i (n - 1 - i)!dn-1-i- j (n - 1 - i - j)!-1
                            j =0                     i =0
                            n -1
                        = ∑ (bc2k ) j j!-1 p j
                            j =0
                            n -1
                        = ∑ b j j!-1 p j c2jk
                            j =0
                                   n -1
                                   ∑ b j j!-1 p j c j c-(k- j)
                               2                      2          2
                        = ck
                                   j =0
                             k2
                        = c qk

     其中 p j 和 qk 都是卷积，可以使用 NTT 进行快速计算。
 1   //*********************************
 2   //快速数论变换 (NTT)
 3   //求 A 和 B 的卷积，结果对 P 取模
 4   //做长度为 N1 的变换，选取两个质数 P1 和 P2
 5   //P1-1 和 P2-1 必须是 N1 的倍数
 6   //E1 和 E2 分别是 P1,P2 的原根
 7   //F1 是 E1 模 P1 的逆元,F2 是 E2 模 P2 的逆元
 8   //I1 是 N1 对模 P1 的逆元,I2 是 N1 对模 P2 的逆元
 9   //
10   //然后使用中国剩余定理，保证了结果是小于 MM=P1*P2 的
11   //M1 = (P2 对 P1 的逆元)*P2
12   //M2 = (P1 对 P2 的逆元)*P1
13
14   const int P = 1000003;//结果对 P 取模
15   const int N1 = 262144;// 218
16   const int N2 = N1+1;//数组大小
17   const int P1 = 998244353;//P1 = 223 ∗ 7 ∗ 17 + 1
18   const int P2 = 995622913;//P2 = 219 ∗ 3 ∗ 3 ∗ 211 + 1
19   const int E1 = 996173970;
20   const int E2 = 88560779;
21   const int F1 = 121392023;//E1*F1 = 1(mod P1)
22   const int F2 = 840835547;//E2*F2 = 1(mod P2)
23   const int I1 = 998240545;//I1*N1 = 1(mod P1)
24   const int I2 = 995619115;//I2*N1 = 1(mod P2)
25   const long long M1 = 397550359381069386LL;


26   const long long M2 = 596324591238590904LL;
27   const long long MM = 993874950619660289LL;//MM = P1*P2
28   //计算 x*y 对 z 取模
29   long long mul(long long x,long long y,long long z){
30       return (x*y - (long long)(x/(long double)z*y+1e-3)*z+z)%z;
31   }
32   int trf(int x1,int x2){
33       return (mul(M1,x1,MM)+mul(M2,x2,MM))%MM%P;
34   }
35   int A[N2],B[N2],C[N2];
36   int A1[N2],B1[N2],C1[N2];
37   void fft(int *A,int PM,int PW){
38       for(int m = N1,h;h = m/2, m >= 2;PW = (long long)PW*PW%PM,m=h)
39           for(int i = 0,w=1;i < h;i++, w = (long long)w*PW%PM)
40                for(int j = i;j < N1;j += m){
41                    int k = j+h, x = (A[j]-A[k]+PM)%PM;
42                    (A[j]+=A[k])%=PM;
43                    A[k] = (long long)w*x%PM;
44                }
45       for(int i = 0,j = 1;j < N1-1;j++){
46           for(int k = N1/2; k > (i^=k);k /= 2);
47           if(j < i)swap(A[i],A[j]);
48       }
49   }
50   //计算 A 和 B 的卷积，结果保存在 C 中，结果对 P 取模
51   void mul(){
52       memset(C,0,sizeof(C));
53       memcpy(A1,A,sizeof(A));
54       memcpy(B1,B,sizeof(B));
55       fft(A1,P1,E1); fft(B1,P1,E1);
56       for(int i = 0;i < N1;i++)C1[i] = (long long)A1[i]*B1[i]%P1;
57       fft(C1,P1,F1);
58       for(int i = 0;i < N1;i++)C1[i] = (long long)C1[i]*I1%P1;
59       fft(A,P2,E2); fft(B,P2,E2);
60       for(int i = 0;i < N1;i++)C[i] = (long long)A[i]*B[i]%P2;
61       fft(C,P2,F2);
62       for(int i = 0;i < N1;i++)C[i] = (long long)C[i]*I2%P2;
63       for(int i = 0;i < N1;i++)C[i] = trf(C1[i],C[i]);
64   }
65   int INV[P];//逆元
66   const int MAXN = 100010;
67   int F[MAXN];//阶乘
68   int a[MAXN];
69   int pd[MAXN];
70   int pb[MAXN];
71   int pc2[MAXN];
72   int p[MAXN];
73   int main()
74   {
75       //预处理逆元
76       INV[1] = 1;


 77     for(int i = 2;i < P;i++)
 78         INV[i] = (long long)P/i*(P-INV[P%i])%P;
 79     F[0] = 1;
 80     for(int i = 1;i < MAXN;i++)
 81         F[i] = (long long)F[i-1]*i%P;
 82     int n,b,c,d;
 83     while(scanf("%d%d%d%d",&n,&b,&c,&d) == 4){
 84         for(int i = 0;i < n;i++)scanf("%d",&a[i]);
 85         pd[0] = 1;
 86         for(int i = 1;i < n;i++)
 87             pd[i] = (long long)pd[i-1]*d%P;
 88         memset(A,0,sizeof(A));
 89         memset(B,0,sizeof(B));
 90         for(int i = 0;i < n;i++)
 91             A[i] = (long long)a[n-1-i]*F[n-1-i]%P;
 92         for(int i = 0;i < n;i++)
 93             B[i] = (long long)pd[i]*INV[F[i]]%P;
 94         mul();
 95         for(int i = 0;i < n;i++)p[i] = C[i];
 96         reverse(p,p+n);
 97         memset(A,0,sizeof(A));
 98         pb[0] = 1;
 99         for(int i = 1;i < n;i++)
100             pb[i] = (long long)pb[i-1]*b%P;
101         pc2[0] = 1;
102         int c2 = (long long)c*c%P;
103         for(int i = 1, s = c;i < n;i++){
104             pc2[i] = (long long)pc2[i-1]*s%P;
105             s = (long long)s*c2%P;
106         }
107         for(int i = 0;i < n;i++)
108             A[i] = (long long)pb[i]*INV[F[i]]%P*p[i]%P*pc2[i]%P;
109         memset(B,0,sizeof(B));
110         B[0] = 1;
111         for(int i = 1;i < n;i++)
112             B[i] = B[N1-i] = INV[pc2[i]];
113         mul();
114         for(int i = 0;i < n;i++)C[i] = (long long)C[i]*pc2[i]%P;
115         for(int i = 0;i < n;i++)
116             printf("%d\n",C[i]);
117     }
118     return 0;
119 }
\end{Code}

\subsection{Các công thức khác}

\subsubsection{Polya}

\begin{Code}
    设 G 是 p 个对象的一个置换群，用 k 种颜色去染这 p 个对象，若一种染色方案在群 G 的作
    用下变为另一种方案，则这两个方案当作是同一种方案，这样的不同染色方案数为：
    L = |G1 | × Σ(kC( f ) ), f ∈ G


C ( f ) 为循环节，| G | 表示群的置换方法数

对于有 n 个位置的手镯，有 n 种旋转置换和 n 种翻转置换

对于旋转置换：
  C ( f i ) = gcd(n, i ) ,i 表示一次转过 i 颗宝石，i = 0 时 c = n；


对于翻转置换：


    如果 n 为偶数： 则有 n2 个置换 C ( f ) = n2 ，有 n2 个置换 C ( f ) = n2 + 1

    如果 n 为奇数： C ( f ) = n2 + 1
\end{Code}

\section{Cấu trúc dữ liệu}

\subsection{Cây phân hoạch}

\begin{Code}
 1   /*
 2    * 划分树（查询区间第 k 大）
 3    */
 4   const int MAXN = 100010;
 5   int tree[20][MAXN];//表示每层每个位置的值
 6   int sorted[MAXN];//已经排序好的数
 7   int toleft[20][MAXN];//toleft[p][i] 表示第 i 层从 1 到 i 有数分入左边
 8
 9   void build(int l,int r,int dep){
10       if(l == r)return;
11       int mid = (l+r)>>1;
12       int same = mid - l + 1;//表示等于中间值而且被分入左边的个数
13       for(int i = l; i <= r; i++) //注意是 l, 不是 one
14           if(tree[dep][i] < sorted[mid])
15                same--;
16       int lpos = l;
17       int rpos = mid+1;
18       for(int i = l;i <= r;i++){
19           if(tree[dep][i] < sorted[mid])
20                tree[dep+1][lpos++] = tree[dep][i];
21           else if(tree[dep][i] == sorted[mid] && same > 0){
22                tree[dep+1][lpos++] = tree[dep][i];
23                same--;
24           }
25           else
26                tree[dep+1][rpos++] = tree[dep][i];
27           toleft[dep][i] = toleft[dep][l-1] + lpos - l;
28       }
29       build(l,mid,dep+1);
30       build(mid+1,r,dep+1);
31   }
32
33   //查询区间第 k 大的数,[L,R] 是大区间，[l,r] 是要查询的小区间
34   int query(int L,int R,int l,int r,int dep,int k){
35       if(l == r)return tree[dep][l];
36       int mid = (L+R)>>1;
37       int cnt = toleft[dep][r] - toleft[dep][l-1];
38       if(cnt >= k){
39           int newl = L + toleft[dep][l-1] - toleft[dep][L-1];
40           int newr = newl + cnt - 1;
41           return query(L,mid,newl,newr,dep+1,k);
42       }
43       else{
44           int newr = r + toleft[dep][R] - toleft[dep][r];
45           int newl = newr - (r-l-cnt);
46           return query(mid+1,R,newl,newr,dep+1,k-cnt);
47       }


48 }
49 int main(){
50     int n,m;
51     while(scanf("%d%d",&n,&m)==2){
52         memset(tree,0,sizeof(tree));
53         for(int i = 1;i <= n;i++){
54              scanf("%d",&tree[0][i]);
55              sorted[i] = tree[0][i];
56         }
57         sort(sorted+1,sorted+n+1);
58         build(1,n,0);
59         int s,t,k;
60         while(m--){
61              scanf("%d%d%d",&s,&t,&k);
62              printf("%d\n",query(1,n,s,t,0,k));
63         }
64     }
65     return 0;
66 }
\end{Code}

\subsection{RMQ}

\subsubsection{Một chiều}

\begin{Code}
     求最大值，数组下标从 1 开始。
     求最小值，或者最大最小值下标，或者数组从 0 开始对应修改即可。

 1   const int MAXN = 50010;
 2   int dp[MAXN][20];
 3   int mm[MAXN];
 4   //初始化 RMQ, b 数组下标从 1 开始，从 0 开始简单修改
 5   void initRMQ(int n,int b[]){
 6       mm[0] = -1;
 7       for(int i = 1; i <= n;i++){
 8           mm[i] = ((i&(i-1)) == 0)?mm[i-1]+1:mm[i-1];
 9           dp[i][0] = b[i];
10       }
11       for(int j = 1; j <= mm[n];j++)
12           for(int i = 1;i + (1<<j) -1 <= n;i++)
13               dp[i][j] = max(dp[i][j-1],dp[i+(1<<(j-1))][j-1]);
14   }
15   //查询最大值
16   int rmq(int x,int y){
17       int k = mm[y-x+1];
18       return max(dp[x][k],dp[y-(1<<k)+1][k]);
19   }
\end{Code}

\subsubsection{Hai chiều}

\begin{Code}
 1 /*
 2 * 二维 RMQ，预处理复杂度 n*m*log*(n)*log(m)
 3 * 数组下标从 1 开始


 4    */
 5   int val[310][310];
 6   int dp[310][310][9][9];//最大值
 7   int mm[310];//二进制位数减一，使用前初始化
 8   void initRMQ(int n,int m){
 9       for(int i = 1;i <= n;i++)
10           for(int j = 1;j <= m;j++)
11                dp[i][j][0][0] = val[i][j];
12       for(int ii = 0; ii <= mm[n]; ii++)
13           for(int jj = 0; jj <= mm[m]; jj++)
14                if(ii+jj)
15                    for(int i = 1; i + (1<<ii) - 1 <= n;i++)
16                        for(int j = 1; j + (1<<jj) - 1 <= m;j++){
17                            if(ii)dp[i][j][ii][jj] = max(dp[i][j][ii
                                 -1][jj],dp[i+(1<<(ii-1))][j][ii-1][jj]);
18                            else dp[i][j][ii][jj] = max(dp[i][j][ii][jj
                                 -1],dp[i][j+(1<<(jj-1))][ii][jj-1]);
19                        }
20   }
21   //查询矩形内的最大值 (x1<=x2,y1<=y2)
22   int rmq(int x1,int y1,int x2,int y2){
23       int k1 = mm[x2-x1+1];
24       int k2 = mm[y2-y1+1];
25       x2 = x2 - (1<<k1) + 1;
26       y2 = y2 - (1<<k2) + 1;
27       return max(max(dp[x1][y1][k1][k2],dp[x1][y2][k1][k2]),max(dp[x2
            ][y1][k1][k2],dp[x2][y2][k1][k2]));
28   }
29   int main(){
30       //在外面对 mm 数组进行初始化
31       mm[0] = -1;
32       for(int i = 1;i <= 305;i++)
33           mm[i] = ((i&(i-1))==0)?mm[i-1]+1:mm[i-1];
34       int n,m;
35       int Q;
36       int r1,c1,r2,c2;
37       while(scanf("%d%d",&n,&m) == 2){
38           for(int i = 1;i <= n;i++)
39                for(int j = 1;j <= m;j++)
40                    scanf("%d",&val[i][j]);
41           initRMQ(n,m);
42           scanf("%d",&Q);
43           while(Q--){
44                scanf("%d%d%d%d",&r1,&c1,&r2,&c2);
45                if(r1 > r2)swap(r1,r2);
46                if(c1 > c2)swap(c1,c2);
47                int tmp = rmq(r1,c1,r2,c2);
48                printf("%d ",tmp);
49                if(tmp == val[r1][c1] || tmp == val[r1][c2] || tmp ==
                     val[r2][c1] || tmp == val[r2][c2])
50                    printf("yes\n");


51                  else printf("no\n");
52              }
53         }
54         return 0;
55 }
\end{Code}

\subsection{Phân rã cây theo chuỗi nặng}

\subsubsection{Trọng số đỉnh}

\begin{Code}
     基于点权，查询单点值，修改路径的上的点权（HDU 3966 树链剖分 + 树状数组）

 1   const int MAXN = 50010;
 2   struct Edge{
 3       int to,next;
 4   }edge[MAXN*2];
 5   int head[MAXN],tot;
 6   int top[MAXN];//top[v] 表示 v 所在的重链的顶端节点
 7   int fa[MAXN];//父亲节点
 8   int deep[MAXN];//深度
 9   int num[MAXN];//num[v] 表示以 v 为根的子树的节点数
10   int p[MAXN];//p[v] 表示 v 对应的位置
11   int fp[MAXN];//fp 和 p 数组相反
12   int son[MAXN];//重儿子
13   int pos;
14   void init(){
15       tot = 0;
16       memset(head,-1,sizeof(head));
17       pos = 1;//使用树状数组，编号从头 1 开始
18       memset(son,-1,sizeof(son));
19   }
20   void addedge(int u,int v){
21       edge[tot].to = v; edge[tot].next = head[u]; head[u] = tot++;
22   }
23   void dfs1(int u,int pre,int d){
24       deep[u] = d;
25       fa[u] = pre;
26       num[u] = 1;
27       for(int i = head[u];i != -1; i = edge[i].next){
28           int v = edge[i].to;
29           if(v != pre){
30                dfs1(v,u,d+1);
31                num[u] += num[v];
32                if(son[u] == -1 || num[v] > num[son[u]])
33                    son[u] = v;
34           }
35       }
36   }
37   void getpos(int u,int sp){
38       top[u] = sp;
39       p[u] = pos++;
40       fp[p[u]] = u;


41       if(son[u] == -1) return;
42       getpos(son[u],sp);
43       for(int i = head[u];i != -1;i = edge[i].next){
44           int v = edge[i].to;
45           if( v != son[u] && v != fa[u])
46               getpos(v,v);
47       }
48   }
49
50   //树状数组
51   int lowbit(int x){
52       return x&(-x);
53   }
54   int c[MAXN];
55   int n;
56   int sum(int i){
57       int s = 0;
58       while(i > 0)
59       {
60           s += c[i];
61           i -= lowbit(i);
62       }
63       return s;
64   }
65   void add(int i,int val){
66       while(i <= n){
67           c[i] += val;
68           i += lowbit(i);
69       }
70   }
71   //u->v 的路径上点的值改变 val
72   void Change(int u,int v,int val){
73       int f1 = top[u], f2 = top[v];
74       int tmp = 0;
75       while(f1 != f2){
76           if(deep[f1] < deep[f2]){
77                swap(f1,f2);
78                swap(u,v);
79           }
80           add(p[f1],val);
81           add(p[u]+1,-val);
82           u = fa[f1];
83           f1 = top[u];
84       }
85       if(deep[u] > deep[v]) swap(u,v);
86       add(p[u],val);
87       add(p[v]+1,-val);
88   }
89   int a[MAXN];
90   int main(){
91       int M,P;


 92      while(scanf("%d%d%d",&n,&M,&P) == 3){
 93          int u,v;
 94          int C1,C2,K;
 95          char op[10];
 96          init();
 97          for(int i = 1;i <= n;i++){
 98              scanf("%d",&a[i]);
 99          }
100          while(M--){
101              scanf("%d%d",&u,&v);
102              addedge(u,v);
103              addedge(v,u);
104          }
105          dfs1(1,0,0);
106          getpos(1,1);
107          memset(c,0,sizeof(c));
108          for(int i = 1;i <= n;i++){
109              add(p[i],a[i]);
110              add(p[i]+1,-a[i]);
111          }
112          while(P--){
113              scanf("%s",op);
114              if(op[0] == 'Q'){
115                  scanf("%d",&u);
116                  printf("%d\n",sum(p[u]));
117              }
118              else{
119                  scanf("%d%d%d",&C1,&C2,&K);
120                  if(op[0] == 'D')
121                       K = -K;
122                  Change(C1,C2,K);
123              }
124          }
125      }
126      return 0;
127 }
\end{Code}

\subsubsection{Trọng số cạnh}

\begin{Code}
     基于边权，修改单条边权，查询路径边权最大值（SPOJ QTREE 树链剖分 + 线段树）

 1   const int MAXN = 10010;
 2   struct Edge{
 3       int to,next;
 4   }edge[MAXN*2];
 5   int head[MAXN],tot;
 6   int top[MAXN];//top[v] 表示 v 所在的重链的顶端节点
 7   int fa[MAXN]; //父亲节点
 8   int deep[MAXN];//深度
 9   int num[MAXN];//num[v] 表示以 v 为根的子树的节点数
10   int p[MAXN];//p[v] 表示 v 与其父亲节点的连边在线段树中的位置


11   int fp[MAXN];//和 p 数组相反
12   int son[MAXN];//重儿子
13   int pos;
14   void init(){
15       tot = 0;
16       memset(head,-1,sizeof(head));
17       pos = 0;
18       memset(son,-1,sizeof(son));
19   }
20   void addedge(int u,int v){
21       edge[tot].to = v;edge[tot].next = head[u];head[u] = tot++;
22   }
23   //第一遍 dfs 求出 fa,deep,num,son
24   void dfs1(int u,int pre,int d){
25       deep[u] = d;
26       fa[u] = pre;
27       num[u] = 1;
28       for(int i = head[u];i != -1; i = edge[i].next){
29           int v = edge[i].to;
30           if(v != pre){
31                dfs1(v,u,d+1);
32                num[u] += num[v];
33                if(son[u] == -1 || num[v] > num[son[u]])
34                    son[u] = v;
35           }
36       }
37   }
38   //第二遍 dfs 求出 top 和 p
39   void getpos(int u,int sp){
40       top[u] = sp;
41       p[u] = pos++;
42       fp[p[u]] = u;
43       if(son[u] == -1) return;
44       getpos(son[u],sp);
45       for(int i = head[u] ; i != -1; i = edge[i].next){
46           int v = edge[i].to;
47           if(v != son[u] && v != fa[u])
48                getpos(v,v);
49       }
50   }
51
52   //线段树
53   struct Node{
54       int l,r;
55       int Max;
56   }segTree[MAXN*3];
57   void build(int i,int l,int r){
58       segTree[i].l = l;
59       segTree[i].r = r;
60       segTree[i].Max = 0;
61       if(l == r)return;


 62       int mid = (l+r)/2;
 63       build(i<<1,l,mid);
 64       build((i<<1)|1,mid+1,r);
 65   }
 66   void push_up(int i){
 67       segTree[i].Max = max(segTree[i<<1].Max,segTree[(i<<1)|1].Max);
 68   }
 69   // 更新线段树的第 k 个值为 val
 70   void update(int i,int k,int val){
 71       if(segTree[i].l == k && segTree[i].r == k){
 72           segTree[i].Max = val;
 73           return;
 74       }
 75       int mid = (segTree[i].l + segTree[i].r)/2;
 76       if(k <= mid)update(i<<1,k,val);
 77       else update((i<<1)|1,k,val);
 78       push_up(i);
 79   }
 80   //查询线段树中 [l,r] 的最大值
 81   int query(int i,int l,int r){
 82       if(segTree[i].l == l && segTree[i].r == r)
 83           return segTree[i].Max;
 84       int mid = (segTree[i].l + segTree[i].r)/2;
 85       if(r <= mid)return query(i<<1,l,r);
 86       else if(l > mid)return query((i<<1)|1,l,r);
 87       else return max(query(i<<1,l,mid),query((i<<1)|1,mid+1,r));
 88   }
 89   //查询 u->v 边的最大值
 90   int find(int u,int v){
 91       int f1 = top[u], f2 = top[v];
 92       int tmp = 0;
 93       while(f1 != f2){
 94           if(deep[f1] < deep[f2]){
 95               swap(f1,f2);
 96               swap(u,v);
 97           }
 98           tmp = max(tmp,query(1,p[f1],p[u]));
 99           u = fa[f1]; f1 = top[u];
100       }
101       if(u == v)return tmp;
102       if(deep[u] > deep[v]) swap(u,v);
103       return max(tmp,query(1,p[son[u]],p[v]));
104   }
105   int e[MAXN][3];
106   int main(){
107       int T;
108       int n;
109       scanf("%d",&T);
110       while(T--){
111           init();
112           scanf("%d",&n);


113             for(int i = 0;i < n-1;i++){
114                 scanf("%d%d%d",&e[i][0],&e[i][1],&e[i][2]);
115                 addedge(e[i][0],e[i][1]);
116                 addedge(e[i][1],e[i][0]);
117             }
118             dfs1(1,0,0);
119             getpos(1,1);
120             build(1,0,pos-1);
121             for(int i = 0;i < n-1; i++){
122                 if(deep[e[i][0]] > deep[e[i][1]])
123                     swap(e[i][0],e[i][1]);
124                 update(1,p[e[i][1]],e[i][2]);
125             }
126             char op[10];
127             int u,v;
128             while(scanf("%s",op) == 1){
129                 if(op[0] == 'D')break;
130                 scanf("%d%d",&u,&v);
131                 if(op[0] == 'Q')
132                     printf("%d\n",find(u,v));//查询 u->v 路径上边权的最大值
133                 else update(1,p[e[u-1][1]],v);//修改第 u 条边的长度为 v
134             }
135        }
136        return 0;
137 }
\end{Code}

\subsection{Cây splay}

\subsubsection{Ví dụ: HDU1890}

\begin{Code}
 1   const int MAXN = 100010;
 2   struct Node;
 3   Node* null;
 4   struct Node{
 5       Node *ch[2],*fa;
 6       int size;
 7       int rev;
 8       Node(){
 9           ch[0] = ch[1] = fa = null; rev = 0;
10       }
11       inline void push_up(){
12           if(this == null)return;
13           size = ch[0]->size + ch[1]->size + 1;
14       }
15       inline void setc(Node* p,int d){
16           ch[d] = p;
17           p->fa = this;
18       }
19       inline bool d(){
20           return fa->ch[1] == this;
21       }


22       void clear(){
23           size = 1;
24           ch[0] = ch[1] = fa = null;
25           rev = 0;
26       }
27       void Update_Rev(){
28           if(this == null)return;
29           swap(ch[0],ch[1]);
30           rev ^= 1;
31       }
32       inline void push_down(){
33           if(this == null)return;
34           if(rev){
35               ch[0]->Update_Rev();
36               ch[1]->Update_Rev();
37               rev = 0;
38           }
39       }
40       inline bool isroot(){
41           return fa == null || this != fa->ch[0] && this != fa->ch
                [1];
42       }
43   };
44   inline void rotate(Node* x)
45   {
46       Node *f = x->fa, *ff = x->fa->fa;
47       f->push_down();
48       x->push_down();
49       int c = x->d(), cc = f->d();
50       f->setc(x->ch[!c],c);
51       x->setc(f,!c);
52       if(ff->ch[cc] == f)ff->setc(x,cc);
53       else x->fa = ff;
54       f->push_up();
55   }
56   inline void splay(Node* &root,Node* x,Node* goal)
57   {
58       while(x->fa != goal){
59           if(x->fa->fa == goal)rotate(x);
60           else {
61               x->fa->fa->push_down();
62               x->fa->push_down();
63               x->push_down();
64               bool f = x->fa->d();
65               x->d() == f ? rotate(x->fa):rotate(x);
66               rotate(x);
67           }
68       }
69       x->push_up();
70       if(goal == null)root = x;
71   }


 72   Node* get_kth(Node* r,int k)
 73   {
 74       Node* x = r;
 75       x->push_down();
 76       while(x->ch[0]->size+1 != k){
 77           if(k < x->ch[0]->size+1)x = x->ch[0];
 78           else{
 79               k -= x->ch[0]->size+1;
 80               x = x->ch[1];
 81           }
 82           x->push_down();
 83       }
 84       return x;
 85   }
 86   Node* get_next(Node* p){
 87       p->push_down();
 88       p = p->ch[1];
 89       p->push_down();
 90       while(p->ch[0] != null){
 91           p = p->ch[0];
 92           p->push_down();
 93       }
 94       return p;
 95   }
 96   Node pool[MAXN],*tail;
 97   Node *node[MAXN];
 98   Node *root;
 99   void build(Node* &x,int l,int r,Node* fa)
100   {
101       if(l > r)return;
102       int mid = (l+r)/2;
103       x = tail++;
104       x->clear();
105       x->fa = fa;
106       node[mid] = x;
107       build(x->ch[0],l,mid-1,x);
108       build(x->ch[1],mid+1,r,x);
109       x->push_up();
110   }
111   void init(int n)
112   {
113       tail = pool;
114       null = tail++;
115       null->fa = null->ch[0] = null->ch[1] = null;
116       null->size = 0; null->rev = 0;
117       Node *p = tail++;
118       p->clear();
119       root = p;
120       p = tail++;
121       p->clear();
122       root->setc(p,1);


123       build(root->ch[1]->ch[0],1,n,root->ch[1]);
124       root->ch[1]->push_up();
125       root->push_up();
126   }
127   int a[MAXN];
128   int b[MAXN];
129   bool cmp(int i,int j)
130   {
131       if(a[i] != a[j])return a[i] < a[j];
132       else return i < j;
133   }
134   int main()
135   {
136       int n;
137       while(scanf("%d",&n) == 1 && n){
138           for(int i = 1;i <= n;i++){
139                scanf("%d",&a[i]);
140                b[i] = i;
141           }
142           init(n);
143           sort(b+1,b+n+1,cmp);
144           for(int i = 1;i <= n;i++){
145                splay(root,node[b[i]],null);
146                int sz = root->ch[0]->size;
147                printf("%d",root->ch[0]->size);
148                if(i == n)printf("\n");
149                else printf(" ");
150                splay(root,get_kth(root,i),null);
151                splay(root,get_kth(root,sz+2),root);
152                root->ch[1]->ch[0]->Update_Rev();
153           }
154       }
155       return 0;
156   }
\end{Code}

\subsubsection{Ví dụ: HDU3726}

\begin{Code}
 1    const int MAXN = 20010;
 2    struct Node;
 3    Node* null;
 4    struct Node
 5    {
 6        Node *ch[2],*fa;//指向儿子和父亲结点
 7        int size,key;
 8        Node(){
 9            ch[0] = ch[1] = fa = null;
10        }
11        inline void setc(Node* p,int d){
12            ch[d] = p;
13            p->fa = this;
14        }
15        inline bool d(){


16           return fa->ch[1] == this;
17       }
18       void push_up(){
19           size = ch[0]->size + ch[1]->size + 1;
20       }
21       void clear(){
22           size = 1;
23           ch[0] = ch[1] = fa = null;
24       }
25       inline bool isroot()
26       {
27           return fa == null || this != fa->ch[0] && this != fa->ch
                [1];
28       }
29   };
30   inline void rotate(Node* x)
31   {
32       Node *f = x->fa, *ff = x->fa->fa;
33       int c = x->d(), cc = f->d();
34       f->setc(x->ch[!c],c);
35       x->setc(f,!c);
36       if(ff->ch[cc] == f)ff->setc(x,cc);
37       else x->fa = ff;
38       f->push_up();
39   }
40   inline void splay(Node* &root,Node* x,Node* goal)
41   {
42       while(x->fa != goal){
43           if(x->fa->fa == goal)rotate(x);
44           else {
45               bool f = x->fa->d();
46               x->d() == f ? rotate(x->fa) : rotate(x);
47               rotate(x);
48           }
49       }
50       x->push_up();
51       if(goal == null)root = x;
52   }
53   //找到 r 子树里面的第 k 个
54   Node* get_kth(Node* r,int k)
55   {
56       Node* x = r;
57       while(x->ch[0]->size+1 != k){
58           if(k < x->ch[0]->size+1)x = x->ch[0];
59           else {
60               k -= x->ch[0]->size+1;
61               x = x->ch[1];
62           }
63       }
64       return x;
65   }


 66   //在 root 的树中删掉 x
 67   void erase(Node* &root,Node* x)
 68   {
 69       splay(root,x,null);
 70       Node* t = root;
 71       if(t->ch[1] != null){
 72           root = t->ch[1];
 73           splay(root,get_kth(root,1),null);
 74           root->setc(t->ch[0],0);
 75       }
 76       else{
 77           root = root->ch[0];
 78       }
 79       root->fa = null;
 80       if(root != null)root->push_up();
 81   }
 82   void insert(Node* &root,Node* x)
 83   {
 84       if(root == null){
 85           root = x;
 86           return;
 87       }
 88       Node* now = root;
 89       Node* pre = root->fa;
 90       while(now != null){
 91           pre = now;
 92           now = now->ch[x->key >= now->key];
 93       }
 94       x->clear();
 95       pre->setc(x,x->key >= pre->key);
 96       splay(root,x,null);
 97   }
 98   void merge(Node* &A,Node* B)
 99   {
100       if(A->size <= B->size)swap(A,B);
101       queue<Node*>Q;
102       Q.push(B);
103       while(!Q.empty()){
104           Node* fr = Q.front();
105           Q.pop();
106           if(fr->ch[0] != null)Q.push(fr->ch[0]);
107           if(fr->ch[1] != null)Q.push(fr->ch[1]);
108           fr->clear();
109           insert(A,fr);
110       }
111   }
112   Node pool[MAXN],*tail;
113
114   struct Edge
115   {
116       int u,v;


117   }edge[60010];
118   int a[MAXN];
119   bool del[60010];
120   struct QUERY
121   {
122       char op[10];
123       int u,v;
124   }query[500010];
125   int y[500010];
126
127   Node* node[MAXN];
128   Node* root[MAXN];
129   int F[MAXN];
130   int find(int x)
131   {
132       if(F[x] == -1)return x;
133       return F[x] = find(F[x]);
134   }
135   void debug(Node *root)
136   {
137       if(root == null)return;
138       debug(root->ch[0]);
139       printf("size: %d, key = %d\n",root->size,root->key);
140       debug(root->ch[1]);
141   }
142
143   int main()
144   {
145       int n,m;
146       int iCase = 0;
147       while(scanf("%d%d",&n,&m) == 2)
148       {
149           if(n == 0 && m == 0)break;
150           iCase++;
151           memset(F,-1,sizeof(F));
152           tail = pool;
153           null = tail++;
154           null->size = 0; null->ch[0] = null->ch[1] = null->fa = null
                 ;
155           null->key = 0;
156           for(int i = 1;i <= n;i++)scanf("%d",&a[i]);
157           for(int i = 0;i < m;i++)
158           {
159                scanf("%d%d",&edge[i].u,&edge[i].v);
160                del[i] = false;
161           }
162           int Q = 0;
163           while(1)
164           {
165                scanf("%s",&query[Q].op);
166                if(query[Q].op[0] == 'E')break;


167                  if(query[Q].op[0] == 'D'){
168                      scanf("%d",&query[Q].u);
169                      query[Q].u--;
170                      del[query[Q].u] = true;
171                  }
172                  else if(query[Q].op[0] == 'Q'){
173                      scanf("%d%d",&query[Q].u,&query[Q].v);
174                  }
175                  else{
176                      scanf("%d%d",&query[Q].u,&query[Q].v);
177                      y[Q] = a[query[Q].u];
178                      a[query[Q].u] = query[Q].v;
179                  }
180                  Q++;
181              }
182              for(int i = 1;i <= n;i++){
183                  node[i] = tail++;
184                  node[i]->clear();
185                  node[i]->key = a[i];
186                  root[i] = node[i];
187              }
188              for(int i = 0;i < m;i++)
189                  if(!del[i]){
190                      int u = edge[i].u;
191                      int v = edge[i].v;
192                      int t1 = find(u);
193                      int t2 = find(v);
194                      if(t1 == t2)continue;
195                      F[t2] = t1;
196                      merge(root[t1],root[t2]);
197                  }
198              vector<int>ans;
199              for(int i = Q-1;i >= 0;i--){
200                  if(query[i].op[0] == 'D'){
201                      int u = edge[query[i].u].u;
202                      int v = edge[query[i].u].v;
203                      int t1 = find(u);
204                      int t2 = find(v);
205                      if(t1 == t2)continue;
206                      F[t2] = t1;
207                      merge(root[t1],root[t2]);
208                  }
209                  else if(query[i].op[0] == 'Q'){
210                      int u = query[i].u;
211                      int k = query[i].v;
212                      u = find(u);
213                      if(k <= 0 || k > root[u]->size){
214                          ans.push_back(0);
215                      }
216                      else{
217                          k = root[u]->size - k + 1;


218                         Node* p = get_kth(root[u],k);
219                         ans.push_back(p->key);
220                     }
221                 }
222                 else{
223                     int u = query[i].u;
224                     int t1 = find(u);
225                     Node* p = node[u];
226                     erase(root[t1],p);
227                     p->clear();
228                     p->key = y[i];
229                     a[u] = y[i];
230                     insert(root[t1],p);
231                 }
232             }
233             double ret = 0;
234             int sz = ans.size();
235             for(int i = 0;i < sz;i++)ret += ans[i];
236             if(sz)ret /= sz;
237             printf("Case %d: %.6lf\n",iCase,ret);
238        }
239        return 0;
240 }
\end{Code}

\subsection{Cây động}

\subsubsection{SPOJQTREE}

\begin{Code}
     给定一棵 n 个结点的树，树的边上有权。有两种操作：
     1. 修改一条边上的权值。
     2. 查询两个结点 x 和 y 之间的最短路径中经过的最大的边的权值。
     其中n <= 104

 1   // http://www.spoj.com/problems/QTREE/
 2   const int MAXN = 10010;
 3   struct Node *null;
 4   struct Node{
 5       Node *fa,*ch[2];
 6       int Max,key;
 7       inline void push_up(){
 8           if(this == null)return;
 9           Max = max(key,max(ch[0]->Max,ch[1]->Max));
10       }
11       inline void setc(Node *p,int d){
12           ch[d] = p;
13           p->fa = this;
14       }
15       inline bool d(){
16           return fa->ch[1] == this;
17       }
18       inline bool isroot() {


19              return fa == null || fa->ch[0] != this && fa->ch[1] != this
                   ;
20       }
21       inline void rot(){
22           Node *f = fa,*ff = fa->fa;
23           int c = d(), cc = fa->d();
24           f->setc(ch[!c],c);
25           this->setc(f,!c);
26           if(ff->ch[cc] == f)ff->setc(this,cc);
27           else this->fa = ff;
28           f->push_up();
29       }
30       inline Node* splay(){
31           while(!isroot()){
32               if(!fa->isroot())
33                   d()==fa->d() ? fa->rot() : rot();
34               rot();
35           }
36           push_up();
37           return this;
38       }
39       inline Node* access(){
40           for(Node *p = this,*q = null; p != null; q = p, p = p->fa){
41               p->splay()->setc(q,1);
42               p->push_up();
43           }
44           return splay();
45       }
46   };
47   Node pool[MAXN],*tail;
48   Node *node[MAXN];
49   void init(int n){
50       tail = pool;
51       null = tail++;
52       null->fa = null->ch[0] = null->ch[1] = null;
53       null->Max = null->key = 0;
54       for(int i = 1;i <= n;i++){
55           node[i] = tail++;
56           node[i]->fa = node[i]->ch[0] = node[i]->ch[1] = null;
57           node[i]->Max = node[i]->key = 0;
58       }
59   }
60   struct Edge{
61       int to,next;
62       int w,id;
63   }edge[MAXN*2];
64   int head[MAXN],tot;
65   inline int addedge(int u,int v,int w,int id){
66       edge[tot].to = v;
67       edge[tot].w = w;
68       edge[tot].id = id;


 69       edge[tot].next = head[u];
 70       head[u] = tot++;
 71   }
 72   Node *ee[MAXN];
 73   bool vis[MAXN];
 74   void bfs(int n){
 75       for(int i = 1;i <= n;i++)vis[i] = false;
 76       queue<int>q;
 77       q.push(1);
 78       vis[1] = true;
 79       while(!q.empty()){
 80           int u = q.front();
 81           q.pop();
 82           for(int i = head[u];i != -1;i = edge[i].next){
 83               int v = edge[i].to;
 84               if(vis[v])continue;
 85               vis[v] = true;
 86               q.push(v);
 87               ee[edge[i].id] = node[v];
 88               node[v]->key = edge[i].w;
 89               node[v]->push_up();
 90               node[v]->fa = node[u];
 91           }
 92       }
 93   }
 94   inline int ask(Node *x,Node *y){
 95       x->access();
 96       for(x = null; y != null; x = y, y = y->fa){
 97           y->splay();
 98           if(y->fa == null)return max(y->ch[1]->Max,x->Max);
 99           y->setc(x,1);
100           y->push_up();
101       }
102   }
103   int main()
104   {
105       int T;
106       scanf("%d",&T);
107       int n;
108       while(T--){
109           scanf("%d",&n);
110           for(int i = 1;i <= n;i++)head[i] = -1;
111           tot = 0;
112           init(n);
113           int u,v,w;
114           for(int i = 1;i < n;i++){
115               scanf("%d%d%d",&u,&v,&w);
116               addedge(u,v,w,i);
117               addedge(v,u,w,i);
118           }
119           bfs(n);


120             char op[20];
121             int x,y;
122             while(scanf("%s",op) == 1){
123                 if(strcmp(op,"DONE") == 0)break;
124                 scanf("%d%d",&x,&y);
125                 if(op[0] == 'Q'){
126                     printf("%d\n",ask(node[x],node[y]));
127                 }
128                 else {
129                     ee[x]->splay()->key = y;
130                     ee[x]->push_up();
131                 }
132             }
133      }
134      return 0;
135 }
\end{Code}

\subsubsection{SPOJQTREE2}

\begin{Code}
     给定一棵 n 个结点的树，树的边上有权。有两种操作：
     1. 查询两个结点 x 和 y 之间的最短路径长度。
     2. 查询从 x 到 y 的最短路径的第 K 条边的长度。
     其中 n <= 104

 1   // http://www.spoj.com/problems/QTREE2/
 2   const int MAXN = 10010;
 3   struct Node *null;
 4   struct Node{
 5       Node *fa,*ch[2];
 6       int sum,val;
 7       int size;
 8       int id;
 9       void clear(){
10           fa = ch[0] = ch[1] = null;
11           sum = val = 0;
12           size = 1;
13       }
14       inline void push_up(){
15           if(this == null)return;
16           sum = val + ch[0]->sum + ch[1]->sum;
17           size = ch[0]->size + ch[1]->size + 1;
18       }
19       inline void setc(Node *p,int d){
20           ch[d] = p;
21           p->fa = this;
22       }
23       inline bool d(){
24           return fa->ch[1] == this;
25       }
26       inline bool isroot(){


27              return fa == null || fa->ch[0] != this && fa->ch[1] != this
                   ;
28       }
29       inline void rot(){
30           Node *f = fa, *ff = fa->fa;
31           int c = d(), cc = fa->d();
32           f->setc(ch[!c],c);
33           this->setc(f,!c);
34           if(ff->ch[cc] == f)ff->setc(this,cc);
35           else this->fa = ff;
36           f->push_up();
37       }
38       inline Node* splay(){
39           while(!isroot()){
40               if(!fa->isroot())
41                   d()==fa->d() ? fa->rot() : rot();
42               rot();
43           }
44           push_up();
45           return this;
46       }
47       inline Node* access(){
48           for(Node *p = this,*q = null; p != null; q = p, p = p->fa){
49               p->splay()->setc(q,1);
50               p->push_up();
51           }
52           return splay();
53       }
54   };
55   Node pool[MAXN],*tail;
56   Node *node[MAXN];
57   void init(int n){
58       tail = pool;
59       null = tail++;
60       null->fa = null->ch[0] = null->ch[1] = null;
61       null->size = null->sum = null->val = 0;
62       for(int i = 1;i <= n;i++){
63           node[i] = tail++;
64           node[i]->id = i;
65           node[i]->clear();
66       }
67   }
68   struct Edge{
69       int to,next;
70       int w;
71   }edge[MAXN*2];
72   int head[MAXN],tot;
73   inline void addedge(int u,int v,int w){
74       edge[tot].to = v;
75       edge[tot].w = w;
76       edge[tot].next = head[u];


 77       head[u] = tot++;
 78   }
 79   void dfs(int u,int pre){
 80       for(int i = head[u];i != -1;i = edge[i].next){
 81           int v = edge[i].to;
 82           if(v == pre)continue;
 83           dfs(v,u);
 84           node[v]->val = edge[i].w;
 85           node[v]->push_up();
 86           node[v]->fa = node[u];
 87       }
 88   }
 89   //查询 x->y 的距离
 90   inline int query_sum(Node *x,Node *y){
 91       x->access();
 92       for(x = null; y != null; x = y, y = y->fa){
 93           y->splay();
 94           if(y->fa == null)
 95               return y->ch[1]->sum + x->sum;
 96           y->setc(x,1);
 97           y->push_up();
 98       }
 99   }
100   //在 splay 中得到第 k 个点
101   Node* get_kth(Node* r,int k){
102       Node *x = r;
103       while(x->ch[0]->size+1 != k){
104           if(k < x->ch[0]->size+1)x = x->ch[0];
105           else {
106               k -= x->ch[0]->size+1;
107               x = x->ch[1];
108           }
109       }
110       return x;
111   }
112   //查询 x->y 路径上的第 k 个点
113   inline int query_kth(Node *x,Node *y,int k){
114       x->access();
115       for(x = null; y != null; x = y, y = y->fa){
116           y->splay();
117           if(y->fa == null){
118               if(y->ch[1]->size+1 == k)return y->id;
119               else if(y->ch[1]->size+1 > k)
120                    return get_kth(y->ch[1],y->ch[1]->size+1-k)->id;
121               else return get_kth(x,k-(y->ch[1]->size+1))->id;
122           }
123           y->setc(x,1);
124           y->push_up();
125       }
126   }
127   int main()


128 {
129      int T,n;
130      scanf("%d",&T);
131      while(T--){
132          scanf("%d",&n);
133          for(int i = 1;i <= n;i++)head[i] = -1;
134          tot = 0;
135          init(n);
136          int u,v,w;
137          for(int i = 1;i < n;i++){
138               scanf("%d%d%d",&u,&v,&w);
139               addedge(u,v,w);
140               addedge(v,u,w);
141          }
142          dfs(1,1);
143          char op[20];
144          while(scanf("%s",op) == 1){
145               if(strcmp(op,"DONE") == 0)break;
146               if(op[0] == 'D'){
147                   scanf("%d%d",&u,&v);
148                   printf("%d\n",query_sum(node[u],node[v]));
149               }
150               else {
151                   int k; scanf("%d%d%d",&u,&v,&k);
152                   printf("%d\n",query_kth(node[u],node[v],k));
153               }
154          }
155      }
156      return 0;
157 }
\end{Code}

\subsubsection{SPOJQTREE4}

\begin{Code}
     给定一棵 n 个结点的树，树的边上有权，每个结点有黑白两色，初始时所有的结点都是白
     的。有两种操作：
     1. 对一个结点执行反色操作（白变黑，黑变白）
     2. 查询树中距离最远的两个白点的距离。
     其中n <= 105 ，查询数目不超过105 .

 1   //http://www.spoj.com/problems/QTREE4/
 2   const int MAXN = 100010;
 3   const int INF = 0x3f3f3f3f;
 4   struct Node *null;
 5   struct Node{
 6       Node *fa,*ch[2];
 7       multiset<int>st0,st1;//st0 是链，st1 是路径
 8       int dd,d0;//d0 是该点对应边的长度，dd 是重链长度
 9       int w0;//白点值为 0，黑点值为 -INF
10       int ls,rs,ms;
11       inline void clear(){
12           fa = ch[0] = ch[1] = null;


13              st0.clear(); st1.clear();
14              st0.insert(-INF);
15              st0.insert(-INF);
16              st1.insert(-INF);
17              w0 = 0; dd = d0 = 0;
18              ls = rs = ms = -INF;
19       }
20       inline void push_up(){
21           if(this == null)return;
22           dd = d0 + ch[0]->dd + ch[1]->dd;
23           int m0 = max(w0,*st0.rbegin()), ml = max(m0,ch[0]->rs+d0),
                mr = max(m0,ch[1]->ls);
24           ls = max(ch[0]->ls,ch[0]->dd + d0 + mr);
25           rs = max(ch[1]->rs,ch[1]->dd + ml);
26           multiset<int>::reverse_iterator it = st0.rbegin();
27           ++it;
28           int t0 = max((*st0.rbegin()) + (*it) , *st1.rbegin());
29           if(w0 == 0)
30                t0 = max(t0,max(0,*st0.rbegin()));
31           ms = max(max(max(ml+ch[1]->ls,mr+d0+ch[0]->rs),max(ch[0]->
                ms,ch[1]->ms)),t0);
32       }
33       inline void setc(Node *p,int d){
34           ch[d] = p;
35           p->fa = this;
36       }
37       inline bool d(){
38           return fa->ch[1] == this;
39       }
40       inline bool isroot(){
41           return fa == null || fa->ch[0] != this && fa->ch[1] != this
                ;
42       }
43       inline void rot(){
44           Node *f = fa, *ff = fa->fa;
45           int c = d(), cc = fa->d();
46           f->setc(ch[!c],c);
47           this->setc(f,!c);
48           if(ff->ch[cc] == f)ff->setc(this,cc);
49           else this->fa = ff;
50           f->push_up();
51       }
52       inline Node* splay(){
53           while(!isroot()){
54                if(!fa->isroot())
55                    d()==fa->d() ? fa->rot() : rot();
56                rot();
57           }
58           push_up();
59           return this;
60       }


 61       inline Node* access(){
 62           for(Node *p = this,*q = null; p != null; q = p, p = p->fa){
 63               p->splay();
 64               if(p->ch[1] != null){
 65                   p->st0.insert(p->ch[1]->ls);
 66                   p->st1.insert(p->ch[1]->ms);
 67               }
 68               if(q != null){
 69                   p->st0.erase(p->st0.find(q->ls));
 70                   p->st1.erase(p->st1.find(q->ms));
 71               }
 72               p->setc(q,1);
 73               p->push_up();
 74           }
 75           return splay();
 76       }
 77   };
 78   Node pool[MAXN],*tail;
 79   Node *node[MAXN];
 80   inline void init(int n){
 81       tail = pool;
 82       null = tail++;
 83       null->fa = null->ch[0] = null->ch[1] = null;
 84       null->st0.clear(); null->st1.clear();
 85       null->ls = null->rs = null->ms = -INF;
 86       null->w0 = -INF;
 87       null->d0 = null->dd = 0;
 88       for(int i = 1;i <= n;i++){
 89           node[i] = tail++;
 90           node[i]->clear();
 91       }
 92   }
 93   struct Edge{
 94       int to,next,w;
 95   }edge[MAXN*2];
 96   int head[MAXN],tot;
 97   inline void addedge(int u,int v,int w){
 98       edge[tot].to = v;
 99       edge[tot].w = w;
100       edge[tot].next = head[u];
101       head[u] = tot++;
102   }
103   inline void dfs(int u,int pre){
104       for(int i = head[u];i != -1;i = edge[i].next){
105           int v = edge[i].to;
106           if(v == pre)continue;
107           node[v]->fa = node[u];
108           node[v]->d0 = edge[i].w;
109           dfs(v,u);
110           node[u]->st0.insert(node[v]->ls);
111           node[u]->st1.insert(node[v]->ms);


112       }
113       node[u]->push_up();
114   }
115   template <class T>
116   inline bool scan_d(T &ret) {
117      char c; int sgn;
118      if(c=getchar(),c==EOF) return 0;
119      while(c!='-'&&(c<'0'||c>'9')) c=getchar();
120      sgn=(c=='-')?-1:1;
121      ret=(c=='-')?0:(c-'0');
122      while(c=getchar(),c>='0'&&c<='9') ret=ret*10+(c-'0');
123      ret*=sgn;
124      return 1;
125   }
126   int main()
127   {
128       int n;
129       while(scanf("%d",&n) == 1){
130           for(int i = 1;i <= n;i++)head[i] = -1;
131           tot = 0;
132           init(n);
133           int u,v,w;
134           for(int i = 1;i < n;i++){
135                scan_d(u); scan_d(v);scan_d(w);
136                addedge(u,v,w);
137                addedge(v,u,w);
138           }
139           dfs(1,1);
140           int ans = node[1]->ms;
141           int Q;
142           char op[10];
143           scanf("%d",&Q);
144           while(Q--){
145                scanf("%s",op);
146                if(op[0] == 'C'){
147                    scan_d(u);
148                    node[u]->access();
149                    node[u]->splay();
150                    if(node[u]->w0 == 0)node[u]->w0 = -INF;
151                    else node[u]->w0 = 0;
152                    node[u]->push_up();
153                    ans = node[u]->ms;
154                }
155                else{
156                    if(ans < 0)puts("They have disappeared.");
157                    else printf("%d\n",ans);
158                }
159           }
160       }
161       return 0;
162   }
\end{Code}

\subsubsection{SPOJQTREE5}

\begin{Code}
     给定一棵 n 个结点的树，边权均为 1。每个结点有黑白两色，初始时所有结点都是黑的。两
     种查询操作：
     1. 对一个结点执行反色操作（白变黑，黑变白）
     2. 查询距离某个特定结点 i 最远的白点的距离。
     其中<= 105 ，查询数目不超过105 .

 1   //http://www.spoj.com/problems/QTREE5/
 2   const int MAXN = 100010;
 3   const int INF = 0x3f3f3f3f;
 4   struct Node *null;
 5   struct Node{
 6       Node *fa,*ch[2];
 7       multiset<int>st;
 8       int dd,d0;
 9       int w0;
10       int ls,rs;
11       inline void clear(){
12           fa = ch[0] = ch[1] = null;
13           st.clear();
14           st.insert(INF);
15           w0 = INF; dd = d0 = 0;
16           ls = rs = INF;
17       }
18       inline void push_up(){
19           if(this == null)return;
20           dd = d0 + ch[0]->dd + ch[1]->dd;
21           int m0 = min(w0,*st.begin()), ml = min(m0,ch[0]->rs+d0), mr
                 = min(m0,ch[1]->ls);
22           ls = min(ch[0]->ls,ch[0]->dd + d0 + mr);
23           rs = min(ch[1]->rs,ch[1]->dd + ml);
24       }
25       inline void setc(Node *p,int d){
26           ch[d] = p;
27           p->fa = this;
28       }
29       inline bool d(){
30           return fa->ch[1] == this;
31       }
32       inline bool isroot(){
33           return fa == null || fa->ch[0] != this && fa->ch[1] != this
                ;
34       }
35       inline void rot(){
36           Node *f = fa, *ff = fa->fa;
37           int c = d(), cc = fa->d();
38           f->setc(ch[!c],c);
39           this->setc(f,!c);
40           if(ff->ch[cc] == f)ff->setc(this,cc);
41           else this->fa = ff;


42           f->push_up();
43       }
44       inline Node* splay(){
45           while(!isroot()){
46               if(!fa->isroot())
47                   d()==fa->d() ? fa->rot() : rot();
48               rot();
49           }
50           push_up();
51           return this;
52       }
53       inline Node* access(){
54           for(Node *p = this,*q = null; p != null; q = p, p = p->fa){
55               p->splay();
56               if(p->ch[1] != null){
57                   p->st.insert(p->ch[1]->ls);
58               }
59               if(q != null){
60                   p->st.erase(p->st.find(q->ls));
61               }
62               p->setc(q,1);
63               p->push_up();
64           }
65           return splay();
66       }
67   };
68   Node pool[MAXN],*tail;
69   Node *node[MAXN];
70   inline void init(int n){
71       tail = pool;
72       null = tail++;
73       null->fa = null->ch[0] = null->ch[1] = null;
74       null->st.clear();
75       null->ls = null->rs = INF;
76       null->w0 = INF;
77       null->dd = null->d0 = 0;
78       for(int i = 1;i <= n;i++){
79           node[i] = tail++;
80           node[i]->clear();
81       }
82   }
83   struct Edge{
84       int to,next;
85   }edge[MAXN*2];
86   int head[MAXN],tot;
87   inline void addedge(int u,int v){
88       edge[tot].to = v;
89       edge[tot].next = head[u];
90       head[u] = tot++;
91   }
92   inline void dfs(int u,int pre){


 93       for(int i = head[u];i != -1;i = edge[i].next){
 94           int v = edge[i].to;
 95           if(v == pre)continue;
 96           node[v]->fa = node[u];
 97           node[v]->d0 = 1;
 98           dfs(v,u);
 99           node[u]->st.insert(node[v]->ls);
100       }
101       node[u]->push_up();
102   }
103   int main()
104   {
105       int n;
106       while(scanf("%d",&n) == 1){
107           init(n);
108           for(int i = 1;i <= n;i++)head[i] = -1;
109           tot = 0;
110           int u,v;
111           for(int i = 1;i < n;i++){
112               scanf("%d%d",&u,&v);
113               addedge(u,v);
114               addedge(v,u);
115           }
116           dfs(1,1);
117           int Q;
118           scanf("%d",&Q);
119           int op;
120           while(Q--){
121               scanf("%d%d",&op,&v);
122               if(op == 0){
123                    node[v]->access();
124                    node[v]->splay();
125                    if(node[v]->w0 == 0)node[v]->w0 = INF;
126                    else node[v]->w0 = 0;
127                    node[v]->push_up();
128               }
129               else {
130                    node[v]->access();
131                    node[v]->splay();
132                    if(node[v]->rs < INF)printf("%d\n",node[v]->rs);
133                    else printf("-1\n");
134               }
135           }
136       }
137       return 0;
138   }
\end{Code}

\subsubsection{SPOJQTREE6}

\begin{Code}
      给定一棵 n 个结点的树，每个结点有黑白两色，初始时所有结点都是黑的。你被要求支持：
      1. 对一个结点执行反色操作（白变黑，黑变白）


     2. 询问有多少个点与 u 相连。两个结点 u,v 相连当且仅当 u,v 路径上所有点的颜色相同。
     其中n <= 105 ，查询数目不超过105 .

 1   //http://www.spoj.com/problems/QTREE6/
 2   const int MAXN = 100010;
 3   struct Node *null;
 4   struct Node{
 5       Node *fa,*ch[2];
 6       int co;//0 is black, 1 is white
 7       int lco,rco;
 8       int ls,rs;
 9       int s[2];
10       int sum[2];//the sum of black and white
11       inline void clear(){
12           fa = ch[0] = ch[1] = null;
13           co = lco = rco = 0;
14           ls = rs = 1;
15           s[0] = s[1] = 0;
16           sum[0] = 1; sum[1] = 0;
17       }
18       inline void push_up(){
19           if(this == null)return;
20           if(ch[0] != null)lco = ch[0]->lco;
21           else lco = co;
22           if(ch[1] != null)rco = ch[1]->rco;
23           else rco = co;
24           sum[0] = ch[0]->sum[0] + ch[1]->sum[0] + (co == 0);
25           sum[1] = ch[0]->sum[1] + ch[1]->sum[1] + (co == 1);
26           int ml = 1 + s[co] + (co==ch[0]->rco?ch[0]->rs:0);
27           int mr = 1 + s[co] + (co==ch[1]->lco?ch[1]->ls:0);
28           ls = ch[0]->ls;
29           if(lco == co && ch[0]->sum[!co] == 0)ls += mr;
30           rs = ch[1]->rs;
31           if(rco == co && ch[1]->sum[!co] == 0)rs += ml;
32       }
33       inline void setc(Node *p,int d){
34           ch[d] = p;
35           p->fa = this;
36       }
37       inline bool d(){
38           return fa->ch[1] == this;
39       }
40       inline bool isroot(){
41           return fa == null || fa->ch[0] != this && fa->ch[1] != this
                ;
42       }
43       inline void rot(){
44           Node *f = fa, *ff = fa->fa;
45           int c = d(), cc = fa->d();
46           f->setc(ch[!c],c);
47           this->setc(f,!c);


48              if(ff->ch[cc] == f)ff->setc(this,cc);
49              else this->fa = ff;
50              f->push_up();
51       }
52       inline Node* splay(){
53           while(!isroot()){
54               if(!fa->isroot())
55                   d()==fa->d() ? fa->rot() : rot();
56               rot();
57           }
58           push_up();
59           return this;
60       }
61       inline Node* access(){
62           for(Node *p = this,*q = null; p != null; q = p, p = p->fa){
63               p->splay();
64               if(p->ch[1] != null)
65                   p->s[p->ch[1]->lco] += p->ch[1]->ls;
66               if(q != null)
67                   p->s[q->lco] -= q->ls;
68               p->setc(q,1);
69               p->push_up();
70           }
71           return splay();
72       }
73   };
74   Node pool[MAXN],*tail;
75   Node *node[MAXN];
76   void init(int n){
77       tail = pool;
78       null = tail++;
79       null->fa = null->ch[0] = null->ch[1] = null;
80       null->s[0] = null->s[1] = 0;
81       null->ls = null->rs = 0;
82       null->sum[0] = null->sum[1] = 0;
83       null->co = null->lco = null->rco = 0;
84       for(int i = 1;i <= n;i++){
85           node[i] = tail++;
86           node[i]->clear();
87       }
88   }
89   struct Edge{
90       int to,next;
91   }edge[MAXN*2];
92   int head[MAXN],tot;
93   inline void addedge(int u,int v){
94       edge[tot].to = v; edge[tot].next = head[u]; head[u] = tot++;
95   }
96   void dfs(int u,int pre){
97       for(int i = head[u];i != -1;i = edge[i].next){
98           int v = edge[i].to;


 99              if(v == pre)continue;
100              node[v]->fa = node[u];
101              dfs(v,u);
102              node[u]->s[node[v]->lco] += node[v]->ls;
103       }
104       node[u]->push_up();
105   }
106   int main()
107   {
108       int n;
109       while(scanf("%d",&n) == 1){
110           init(n);
111           for(int i = 1;i <= n;i++)head[i] = -1;
112           tot = 0;
113           int u,v;
114           for(int i = 1;i < n;i++){
115               scanf("%d%d",&u,&v);
116               addedge(u,v);
117               addedge(v,u);
118           }
119           dfs(1,1);
120           int Q;
121           int op;
122           scanf("%d",&Q);
123           while(Q--){
124               scanf("%d%d",&op,&u);
125               if(op == 0){
126                    node[u]->access();
127                    node[u]->splay();
128                    printf("%d\n",node[u]->rs);
129               }
130               else{
131                    node[u]->access();
132                    node[u]->splay();
133                    node[u]->co ^= 1;
134                    node[u]->push_up();
135               }
136           }
137           return 0;
138       }
139       return 0;
140   }
\end{Code}

\subsubsection{SPOJQTREE7}

\begin{Code}
      给定一棵 n 个结点的树，每个结点有黑白两色和权值。三种操作：
      1. 对一个结点执行反色操作（白变黑，黑变白）
      2. 询问与 u 相连的点中点权的最大值。两个结点 u,v 相连当且仅当 u,v 路径上所有点的颜色
      相同。
      3. 改变一个点的点权。
      其中n <= 105 ，查询数目不超过105 .


 1   //http://www.spoj.com/problems/QTREE7/
 2   const int MAXN = 100010;
 3   const int INF = 0x3f3f3f3f;
 4   struct Node *null;
 5   struct Node{
 6       Node *fa,*ch[2];
 7       int co;
 8       int lco,rco;
 9       int ls,rs;
10       int w0;
11       multiset<int>st[2];
12       int sum[2];
13       inline void clear(int _co = 0, int _w0 = 0){
14           fa = ch[0] = ch[1] = null;
15           co = lco = rco = _co;
16           w0 = _w0;
17           ls = rs = _w0;
18           st[0].clear(); st[1].clear();
19           st[0].insert(-INF); st[1].insert(-INF);
20           sum[0] = sum[1] = 0; sum[_co]++;
21       }
22       inline void push_up(){
23           if(this == null)return;
24           if(ch[0] != null)lco = ch[0]->lco;
25           else lco = co;
26           if(ch[1] != null)rco = ch[1]->rco;
27           else rco = co;
28           sum[0] = ch[0]->sum[0] + ch[1]->sum[0] + (co == 0);
29           sum[1] = ch[0]->sum[1] + ch[1]->sum[1] + (co == 1);
30           int ml = max(w0,max(*st[co].rbegin(),co==ch[0]->rco?ch[0]->
                rs:-INF));
31           int mr = max(w0,max(*st[co].rbegin(),co==ch[1]->lco?ch[1]->
                ls:-INF));
32           ls = ch[0]->ls;
33           if(lco == co && ch[0]->sum[!co] == 0)ls = max(ls,mr);
34           rs = ch[1]->rs;
35           if(rco == co && ch[1]->sum[!co] == 0)rs = max(rs,ml);
36       }
37       inline void setc(Node *p,int d){
38           ch[d] = p;
39           p->fa = this;
40       }
41       inline bool d(){
42           return fa->ch[1] == this;
43       }
44       inline bool isroot(){
45           return fa == null || fa->ch[0] != this && fa->ch[1] != this
                ;
46       }
47       inline void rot(){


48              Node *f = fa, *ff = fa->fa;
49              int c = d(), cc = fa->d();
50              f->setc(ch[!c],c);
51              this->setc(f,!c);
52              if(ff->ch[cc] == f)ff->setc(this,cc);
53              else this->fa = ff;
54              f->push_up();
55       }
56       inline Node* splay(){
57           while(!isroot()){
58               if(!fa->isroot())
59                   d()==fa->d() ? fa->rot() : rot();
60               rot();
61           }
62           push_up();
63           return this;
64       }
65       inline Node* access(){
66           for(Node *p = this,*q = null; p != null; q = p, p = p->fa){
67               p->splay();
68               if(p->ch[1] != null)
69                   p->st[p->ch[1]->lco].insert(p->ch[1]->ls);
70               if(q != null)
71                   p->st[q->lco].erase(p->st[q->lco].find(q->ls));
72               p->setc(q,1);
73               p->push_up();
74           }
75           return splay();
76       }
77   };
78   Node pool[MAXN],*tail;
79   Node *node[MAXN];
80   int color[MAXN],val[MAXN];
81   void init(int n){
82       tail = pool;
83       null = tail++;
84       null->fa = null->ch[0] = null->ch[1] = null;
85       null->st[0].clear(); null->st[1].clear();
86       null->ls = null->rs = -INF;
87       null->sum[0] = null->sum[1] = 0;
88       null->co = null->lco = null->rco = 0;
89       for(int i = 1;i <= n;i++){
90           node[i] = tail++;
91           node[i]->clear(color[i],val[i]);
92       }
93   }
94   struct Edge{
95       int to,next;
96   }edge[MAXN*2];
97   int head[MAXN],tot;
98   inline void addedge(int u,int v){


 99       edge[tot].to = v; edge[tot].next = head[u]; head[u] = tot++;
100   }
101   void dfs(int u,int pre){
102       for(int i = head[u];i != -1;i = edge[i].next){
103           int v = edge[i].to;
104           if(v == pre)continue;
105           node[v]->fa = node[u];
106           dfs(v,u);
107           node[u]->st[node[v]->lco].insert(node[v]->ls);
108       }
109       node[u]->push_up();
110   }
111   int main()
112   {
113       int n;
114       while(scanf("%d",&n) == 1){
115           for(int i = 1;i <= n;i++)head[i] = -1;
116           tot = 0;
117           int u,v;
118           for(int i = 1;i < n;i++){
119               scanf("%d%d",&u,&v);
120               addedge(u,v);
121               addedge(v,u);
122           }
123           for(int i = 1;i <= n;i++)scanf("%d",&color[i]);
124           for(int i = 1;i <= n;i++)scanf("%d",&val[i]);
125           init(n);
126           dfs(1,1);
127           int Q;
128           int w,op;
129           scanf("%d",&Q);
130           while(Q--){
131               scanf("%d",&op);
132               if(op == 0){
133                    scanf("%d",&u);
134                    node[u]->access(); node[u]->splay();
135                    printf("%d\n",node[u]->rs);
136               }
137               else if(op == 1){
138                    scanf("%d",&u);
139                    node[u]->access(); node[u]->splay();
140                    node[u]->co ^= 1;
141                    node[u]->push_up();
142               }
143               else {
144                    scanf("%d%d",&u,&w);
145                    node[u]->access(); node[u]->splay();
146                    node[u]->w0 = w;
147                    node[u]->push_up();
148               }
149           }


150      }
151      return 0;
152 }
\end{Code}

\subsubsection{HDU4010}

\begin{Code}
     支持：
     1 x y : 如果 x,y 不在同一颗子树中，则通过在 x,y 之间连边的方式，连接这两颗子树
     2 x y : 如果 x,y 在同一颗子树中，且 x!=y, 则将 x 视为这颗子树的根以后，切断 y 与其父亲
     结点的连接
     3 w x y : 如果 x,y 在同一颗子树中，则将 x,y 之间路径上所有点的点权增加 w
     4 x y : 如果 x,y 在同一颗子树中，输出 x,y 之间路径上点权的最大值
     非法操作输出 -1

 1   const int MAXN = 300010;
 2   const int INF = 0x3f3f3f3f;
 3   struct Node *null;
 4   struct Node{
 5       Node *fa,*ch[2];
 6       int Max,val;
 7       int rev;//旋转标记
 8       int add;
 9       inline void clear(int _val){
10           fa = ch[0] = ch[1] = null;
11           val = Max = _val;
12           rev = 0;
13           add = 0;
14       }
15       inline void push_up(){
16           Max = max(val,max(ch[0]->Max,ch[1]->Max));
17       }
18       inline void setc(Node *p,int d){
19           ch[d] = p;
20           p->fa = this;
21       }
22       inline bool d(){
23           return fa->ch[1] == this;
24       }
25       inline bool isroot(){
26           return fa == null || fa->ch[0] != this && fa->ch[1] != this
                ;
27       }
28       //翻转
29       inline void flip(){
30           if(this == null)return;
31           swap(ch[0],ch[1]);
32           rev ^= 1;
33       }
34       inline void update_add(int w){
35           if(this == null)return;
36           val += w;


37              add += w;
38              Max += w;
39       }
40       inline void push_down(){
41           if(rev){
42               ch[0]->flip(); ch[1]->flip(); rev = 0;
43           }
44           if(add){
45               ch[0]->update_add(add); ch[1]->update_add(add);
46               add = 0;
47           }
48       }
49       //直接标记下放
50       inline void go(){
51           if(!isroot())fa->go();
52           push_down();
53       }
54       inline void rot(){
55           Node *f = fa, *ff = fa->fa;
56           int c = d(), cc = fa->d();
57           f->setc(ch[!c],c);
58           this->setc(f,!c);
59           if(ff->ch[cc] == f)ff->setc(this,cc);
60           else this->fa = ff;
61           f->push_up();
62       }
63       inline Node* splay(){
64           go();
65           while(!isroot()){
66               if(!fa->isroot())
67                    d()==fa->d() ? fa->rot() : rot();
68               rot();
69           }
70           push_up();
71           return this;
72       }
73       inline Node* access(){
74           for(Node *p = this,*q = null; p != null; q = p, p = p->fa){
75               p->splay()->setc(q,1);
76               p->push_up();
77           }
78           return splay();
79       }
80       //找该点的根
81       inline Node* find_root(){
82           Node *x;
83           for(x = access(); x->push_down(), x->ch[0] != null; x = x->
                ch[0]);
84           return x;
85       }
86       //变为树根 (换根操作)


 87       void make_root(){
 88           access()->flip();
 89       }
 90       //切断该点和父亲结点的边
 91       void cut(){
 92           access();
 93           ch[0]->fa = null;
 94           ch[0] = null;
 95           push_up();
 96       }
 97       //切断该点以 x 为根时, 该点和父亲结点的根
 98       //要求这个点和 x 在同一颗树而且不能相同
 99       //x 变为所在树的树根
100       void cut(Node* x){
101           if(this == x || find_root() != x->find_root())
102               puts("-1");
103           else {
104               x->make_root();
105               cut();
106           }
107       }
108       //该点连接到 x
109       //假如是有虚边信息的，需要先 x->access() 再连接
110       void link(Node *x){
111           if(find_root() == x->find_root())
112               puts("-1");
113           else {
114               make_root(); fa = x;
115           }
116       }
117   };
118   Node pool[MAXN],*tail;
119   Node *node[MAXN];
120   struct Edge{
121       int to,next;
122   }edge[MAXN*2];
123   int head[MAXN],tot;
124   inline void addedge(int u,int v){
125       edge[tot].to = v;
126       edge[tot].next = head[u];
127       head[u] = tot++;
128   }
129   void dfs(int u,int pre){
130       for(int i = head[u];i != -1;i = edge[i].next){
131           int v = edge[i].to;
132           if(v == pre)continue;
133           node[v]->fa = node[u];
134           dfs(v,u);
135       }
136   }
137   void ADD(Node *x,Node *y,int w){


138       x->access();
139       for(x = null; y != null; x = y, y = y->fa){
140           y->splay();
141           if(y->fa == null){
142               y->ch[1]->update_add(w);
143               x->update_add(w);
144               y->val += w;
145               y->push_up();
146               return;
147           }
148           y->setc(x,1);
149           y->push_up();
150       }
151   }
152   int ask(Node *x,Node *y){
153       x->access();
154       for(x = null; y != null; x = y, y = y->fa){
155           y->splay();
156           if(y->fa == null)
157               return max(y->val,max(y->ch[1]->Max,x->Max));
158           y->setc(x,1);
159           y->push_up();
160       }
161   }
162   int main()
163   {
164       int n;
165       while(scanf("%d",&n) == 1){
166           for(int i = 1;i <= n;i++)head[i] = -1;
167           tot = 0;
168           int u,v;
169           for(int i = 1;i < n;i++){
170               scanf("%d%d",&u,&v);
171               addedge(u,v);
172               addedge(v,u);
173           }
174           tail = pool;
175           null = tail++;
176           null->clear(-INF);
177           for(int i = 1;i <= n;i++){
178               node[i] = tail++;
179               scanf("%d",&v);
180               node[i]->clear(v);
181           }
182           dfs(1,1);
183           int m,op;
184           int x,y,w;
185           scanf("%d",&m);
186           while(m--){
187               scanf("%d",&op);
188               if(op == 1){


189                     scanf("%d%d",&x,&y);
190                     node[x]->link(node[y]);
191                 }
192                 else if(op == 2){
193                     scanf("%d%d",&x,&y);
194                     node[y]->cut(node[x]);
195                 }
196                 else if(op == 3){
197                     scanf("%d%d%d",&w,&x,&y);
198                     if(node[x]->find_root() != node[y]->find_root())
199                         printf("-1\n");
200                     else ADD(node[x],node[y],w);
201                 }
202                 else{
203                     scanf("%d%d",&x,&y);
204                     if(node[x]->find_root() != node[y]->find_root())
205                         printf("-1\n");
206                     else printf("%d\n",ask(node[x],node[y]));
207                 }
208             }
209             printf("\n");
210        }
211        return 0;
212 }
\end{Code}

\subsection{Cây chủ tịch}

\subsubsection{Đếm số giá trị khác nhau trong đoạn}

\begin{Code}
     查询区间有多少个不同的数（SPOJ DQUERY）
 1   /*
 2    * 给出一个序列，查询区间内有多少个不相同的数
 3    */
 4   const int MAXN = 30010;
 5   const int M = MAXN * 100;
 6   int n,q,tot;
 7   int a[MAXN];
 8   int T[MAXN],lson[M],rson[M],c[M];
 9   int build(int l,int r){
10       int root = tot++;
11       c[root] = 0;
12       if(l != r){
13           int mid = (l+r)>>1;
14           lson[root] = build(l,mid);
15           rson[root] = build(mid+1,r);
16       }
17       return root;
18   }
19   int update(int root,int pos,int val){
20       int newroot = tot++, tmp = newroot;
21       c[newroot] = c[root] + val;


22       int l = 1, r = n;
23       while(l < r){
24           int mid = (l+r)>>1;
25           if(pos <= mid){
26               lson[newroot] = tot++; rson[newroot] = rson[root];
27               newroot = lson[newroot]; root = lson[root];
28               r = mid;
29           }
30           else{
31               rson[newroot] = tot++; lson[newroot] = lson[root];
32               newroot = rson[newroot]; root = rson[root];
33               l = mid+1;
34           }
35           c[newroot] = c[root] + val;
36       }
37       return tmp;
38   }
39   int query(int root,int pos){
40       int ret = 0;
41       int l = 1, r = n;
42       while(pos < r){
43           int mid = (l+r)>>1;
44           if(pos <= mid){
45               r = mid;
46               root = lson[root];
47           }
48           else{
49               ret += c[lson[root]];
50               root = rson[root];
51               l = mid+1;
52           }
53       }
54       return ret + c[root];
55   }
56   int main(){
57       while(scanf("%d",&n) == 1){
58           tot = 0;
59           for(int i = 1;i <= n;i++)
60               scanf("%d",&a[i]);
61           T[n+1] = build(1,n);
62           map<int,int>mp;
63           for(int i = n;i>= 1;i--){
64               if(mp.find(a[i]) == mp.end()){
65                    T[i] = update(T[i+1],i,1);
66               }
67               else{
68                    int tmp = update(T[i+1],mp[a[i]],-1);
69                    T[i] = update(tmp,i,1);
70               }
71               mp[a[i]] = i;
72           }


73              scanf("%d",&q);
74              while(q--){
75                  int l,r;
76                  scanf("%d%d",&l,&r);
77                  printf("%d\n",query(T[l],r));
78              }
79       }
80       return 0;
81 }
\end{Code}

\subsubsection{Phần tử lớn thứ k trên đoạn tĩnh}

\begin{Code}
     POJ2104
 1   const int MAXN = 100010;
 2   const int M = MAXN * 30;
 3   int n,q,m,tot;
 4   int a[MAXN], t[MAXN];
 5   int T[MAXN], lson[M], rson[M], c[M];
 6   void Init_hash(){
 7       for(int i = 1; i <= n;i++)
 8           t[i] = a[i];
 9       sort(t+1,t+1+n);
10       m = unique(t+1,t+1+n)-t-1;
11   }
12   int build(int l,int r){
13       int root = tot++;
14       c[root] = 0;
15       if(l != r){
16           int mid = (l+r)>>1;
17           lson[root] = build(l,mid);
18           rson[root] = build(mid+1,r);
19       }
20       return root;
21   }
22   int hash(int x){
23       return lower_bound(t+1,t+1+m,x) - t;
24   }
25   int update(int root,int pos,int val){
26       int newroot = tot++, tmp = newroot;
27       c[newroot] = c[root] + val;
28       int l = 1, r = m;
29       while(l < r){
30           int mid = (l+r)>>1;
31           if(pos <= mid){
32               lson[newroot] = tot++; rson[newroot] = rson[root];
33               newroot = lson[newroot]; root = lson[root];
34               r = mid;
35           }
36           else{
37               rson[newroot] = tot++; lson[newroot] = lson[root];
38               newroot = rson[newroot]; root = rson[root];


39                  l = mid+1;
40              }
41              c[newroot] = c[root] + val;
42       }
43       return tmp;
44   }
45   int query(int left_root,int right_root,int k){
46       int l = 1, r = m;
47       while( l < r){
48           int mid = (l+r)>>1;
49           if(c[lson[left_root]]-c[lson[right_root]] >= k ){
50               r = mid;
51               left_root = lson[left_root];
52               right_root = lson[right_root];
53           }
54           else{
55               l = mid + 1;
56               k -= c[lson[left_root]] - c[lson[right_root]];
57               left_root = rson[left_root];
58               right_root = rson[right_root];
59           }
60       }
61       return l;
62   }
63   int main(){
64       while(scanf("%d%d",&n,&q) == 2){
65           tot = 0;
66           for(int i = 1;i <= n;i++)
67               scanf("%d",&a[i]);
68           Init_hash();
69           T[n+1] = build(1,m);
70           for(int i = n;i ;i--){
71               int pos = hash(a[i]);
72               T[i] = update(T[i+1],pos,1);
73           }
74           while(q--){
75               int l,r,k;
76               scanf("%d%d%d",&l,&r,&k);
77               printf("%d\n",t[query(T[l],T[r+1],k)]);
78           }
79       }
80       return 0;
81   }
\end{Code}

\subsubsection{Trọng số đỉnh lớn thứ k trên đường đi trong cây}

\begin{Code}
     树上路径点权第 k 大（SPOJ COT）
     LCA + 主席树
1 //主席树部分 *****************
2 const int MAXN = 200010;
3 const int M = MAXN * 40;


 4 int n,q,m,TOT;
 5 int a[MAXN], t[MAXN];
 6 int T[MAXN], lson[M], rson[M], c[M];
 7 void Init_hash(){
 8     for(int i = 1; i <= n;i++)
 9         t[i] = a[i];
10     sort(t+1,t+1+n);
11     m = unique(t+1,t+n+1)-t-1;
12 }
13 int build(int l,int r){
14     int root = TOT++;
15     c[root] = 0;
16     if(l != r){
17         int mid = (l+r)>>1;
18         lson[root] = build(l,mid);
19         rson[root] = build(mid+1,r);
20     }
21     return root;
22 }
23 int hash(int x){
24     return lower_bound(t+1,t+1+m,x) - t;
25 }
26 int update(int root,int pos,int val){
27     int newroot = TOT++, tmp = newroot;
28     c[newroot] = c[root] + val;
29     int l = 1, r = m;
30     while( l < r){
31         int mid = (l+r)>>1;
32         if(pos <= mid){
33             lson[newroot] = TOT++; rson[newroot] = rson[root];
34             newroot = lson[newroot]; root = lson[root];
35             r = mid;
36         }
37         else{
38             rson[newroot] = TOT++; lson[newroot] = lson[root];
39             newroot = rson[newroot]; root = rson[root];
40             l = mid+1;
41         }
42         c[newroot] = c[root] + val;
43     }
44     return tmp;
45 }
46 int query(int left_root,int right_root,int LCA,int k){
47     int lca_root = T[LCA];
48     int pos = hash(a[LCA]);
49     int l = 1, r = m;
50     while(l < r){
51         int mid = (l+r)>>1;
52         int tmp = c[lson[left_root]] + c[lson[right_root]] - 2*c[
              lson[lca_root]] + (pos >= l && pos <= mid);
53         if(tmp >= k){


54                  left_root = lson[left_root];
55                  right_root = lson[right_root];
56                  lca_root = lson[lca_root];
57                  r = mid;
58              }
59              else{
60                  k -= tmp;
61                  left_root = rson[left_root];
62                  right_root = rson[right_root];
63                  lca_root = rson[lca_root];
64                  l = mid + 1;
65              }
66       }
67       return l;
68   }
69
70   //LCA 部分
71   int rmq[2*MAXN];//rmq 数组，就是欧拉序列对应的深度序列
72   struct ST{
73       int mm[2*MAXN];
74       int dp[2*MAXN][20];//最小值对应的下标
75       void init(int n){
76           mm[0] = -1;
77           for(int i = 1;i <= n;i++){
78                mm[i] = ((i&(i-1)) == 0)?mm[i-1]+1:mm[i-1];
79                dp[i][0] = i;
80           }
81           for(int j = 1; j <= mm[n];j++)
82                for(int i = 1; i + (1<<j) - 1 <= n; i++)
83                    dp[i][j] = rmq[dp[i][j-1]] < rmq[dp[i+(1<<(j-1))][j
                         -1]]?dp[i][j-1]:dp[i+(1<<(j-1))][j-1];
84       }
85       //查询 [a,b] 之间最小值的下标
86       int query(int a,int b){
87           if(a > b)swap(a,b);
88           int k = mm[b-a+1];
89           return rmq[dp[a][k]] <= rmq[dp[b-(1<<k)+1][k]]?dp[a][k]:dp[
                b-(1<<k)+1][k];
 90      }
 91  };
 92  //边的结构体定义
 93  struct Edge{
 94      int to,next;
 95  };
 96  Edge edge[MAXN*2];
 97  int tot,head[MAXN];
 98
 99 int F[MAXN*2];//欧拉序列，就是 dfs 遍历的顺序，长度为 2*n-1, 下标从 1 开始
100 int P[MAXN];//P[i] 表示点 i 在 F 中第一次出现的位置
101 int cnt;
102


103   ST st;
104   void init(){
105       tot = 0;
106       memset(head,-1,sizeof(head));
107   }
108   //加边，无向边需要加两次
109   void addedge(int u,int v){
110       edge[tot].to = v;
111       edge[tot].next = head[u];
112       head[u] = tot++;
113   }
114   void dfs(int u,int pre,int dep){
115       F[++cnt] = u;
116       rmq[cnt] = dep;
117       P[u] = cnt;
118       for(int i = head[u];i != -1;i = edge[i].next){
119           int v = edge[i].to;
120           if(v == pre)continue;
121           dfs(v,u,dep+1);
122           F[++cnt] = u;
123           rmq[cnt] = dep;
124       }
125   }
126   //查询 LCA 前的初始化
127   void LCA_init(int root,int node_num){
128       cnt = 0;
129       dfs(root,root,0);
130       st.init(2*node_num-1);
131   }
132   //查询 u,v 的 lca 编号
133   int query_lca(int u,int v){
134       return F[st.query(P[u],P[v])];
135   }
136   void dfs_build(int u,int pre){
137       int pos = hash(a[u]);
138       T[u] = update(T[pre],pos,1);
139       for(int i = head[u]; i != -1;i = edge[i].next){
140           int v = edge[i].to;
141           if(v == pre)continue;
142           dfs_build(v,u);
143       }
144   }
145   int main(){
146       while(scanf("%d%d",&n,&q) == 2){
147           for(int i = 1;i <= n;i++)
148                scanf("%d",&a[i]);
149           Init_hash();
150           init();
151           TOT = 0;
152           int u,v;
153           for(int i = 1;i < n;i++){


154                 scanf("%d%d",&u,&v);
155                 addedge(u,v);
156                 addedge(v,u);
157             }
158             LCA_init(1,n);
159             T[n+1] = build(1,m);
160             dfs_build(1,n+1);
161             int k;
162             while(q--){
163                 scanf("%d%d%d",&u,&v,&k);
164                 printf("%d\n",t[query(T[u],T[v],query_lca(u,v),k)]);
165             }
166             return 0;
167      }
168      return 0;
169 }
\end{Code}

\subsubsection{Phần tử lớn thứ k động}

\begin{Code}
     树状数组套主席树 ZOJ 2112
 1   const int MAXN = 60010;
 2   const int M = 2500010;
 3   int n,q,m,tot;
 4   int a[MAXN], t[MAXN];
 5   int T[MAXN], lson[M], rson[M],c[M];
 6   int S[MAXN];
 7   struct Query{
 8       int kind;
 9       int l,r,k;
10   }query[10010];
11
12   void Init_hash(int k){
13       sort(t,t+k);
14       m = unique(t,t+k) - t;
15   }
16   int hash(int x){
17       return lower_bound(t,t+m,x)-t;
18   }
19   int build(int l,int r){
20       int root = tot++;
21       c[root] = 0;
22       if(l != r){
23           int mid = (l+r)/2;
24           lson[root] = build(l,mid);
25           rson[root] = build(mid+1,r);
26       }
27       return root;
28   }
29
30   int Insert(int root,int pos,int val){
31       int newroot = tot++, tmp = newroot;


32       int l = 0, r = m-1;
33       c[newroot] = c[root] + val;
34       while(l < r){
35           int mid = (l+r)>>1;
36           if(pos <= mid){
37               lson[newroot] = tot++; rson[newroot] = rson[root];
38               newroot = lson[newroot]; root = lson[root];
39               r = mid;
40           }
41           else{
42               rson[newroot] = tot++; lson[newroot] = lson[root];
43               newroot = rson[newroot]; root = rson[root];
44               l = mid+1;
45           }
46           c[newroot] = c[root] + val;
47       }
48       return tmp;
49   }
50
51   int lowbit(int x){
52       return x&(-x);
53   }
54   int use[MAXN];
55   void add(int x,int pos,int val){
56       while(x <= n){
57           S[x] = Insert(S[x],pos,val);
58           x += lowbit(x);
59       }
60   }
61   int sum(int x){
62       int ret = 0;
63       while(x > 0){
64           ret += c[lson[use[x]]];
65           x -= lowbit(x);
66       }
67       return ret;
68   }
69   int Query(int left,int right,int k){
70       int left_root = T[left-1];
71       int right_root = T[right];
72       int l = 0, r = m-1;
73       for(int i = left-1;i;i -= lowbit(i)) use[i] = S[i];
74       for(int i = right;i ;i -= lowbit(i)) use[i] = S[i];
75       while(l < r){
76           int mid = (l+r)/2;
77           int tmp = sum(right) - sum(left-1) + c[lson[right_root]] -
                c[lson[left_root]];
78           if(tmp >= k){
79               r = mid;
80               for(int i = left-1; i ;i -= lowbit(i))
81                    use[i] = lson[use[i]];


 82                  for(int i = right; i; i -= lowbit(i))
 83                      use[i] = lson[use[i]];
 84                  left_root = lson[left_root];
 85                  right_root = lson[right_root];
 86              }
 87              else{
 88                  l = mid+1;
 89                  k -= tmp;
 90                  for(int i = left-1; i;i -= lowbit(i))
 91                      use[i] = rson[use[i]];
 92                  for(int i = right;i ;i -= lowbit(i))
 93                      use[i] = rson[use[i]];
 94                  left_root = rson[left_root];
 95                  right_root = rson[right_root];
 96              }
 97       }
 98       return l;
 99   }
100   void Modify(int x,int p,int d){
101       while(x <= n){
102           S[x] = Insert(S[x],p,d);
103           x += lowbit(x);
104       }
105   }
106
107   int main(){
108       int Tcase;
109       scanf("%d",&Tcase);
110       while(Tcase--){
111           scanf("%d%d",&n,&q);
112           tot = 0;
113           m = 0;
114           for(int i = 1;i <= n;i++){
115               scanf("%d",&a[i]);
116               t[m++] = a[i];
117           }
118           char op[10];
119           for(int i = 0;i < q;i++){
120               scanf("%s",op);
121               if(op[0] == 'Q'){
122                    query[i].kind = 0;
123                    scanf("%d%d%d",&query[i].l,&query[i].r,&query[i].k)
                          ;
124               }
125               else{
126                    query[i].kind = 1;
127                    scanf("%d%d",&query[i].l,&query[i].r);
128                    t[m++] = query[i].r;
129               }
130           }
131           Init_hash(m);


132              T[0] = build(0,m-1);
133              for(int i = 1;i <= n;i++)
134                  T[i] = Insert(T[i-1],hash(a[i]),1);
135              for(int i = 1;i <= n;i++)
136                  S[i] = T[0];
137              for(int i = 0;i < q;i++){
138                  if(query[i].kind == 0)
139                      printf("%d\n",t[Query(query[i].l,query[i].r,query[i
                            ].k)]);
140                  else{
141                      Modify(query[i].l,hash(a[query[i].l]),-1);
142                      Modify(query[i].l,hash(query[i].r),1);
143                      a[query[i].l] = query[i].r;
144                  }
145              }
146       }
147       return 0;
148 }
\end{Code}

\subsection{Treap}

\begin{Code}
      ZOJ3765
 1    long long gcd(long long a,long long b){
 2        if(b == 0)return a;
 3        else return gcd(b,a%b);
 4    }
 5    const int MAXN = 300010;
 6    int num[MAXN],st[MAXN];
 7    struct Treap{
 8        int tot1;
 9        int s[MAXN],tot2;//内存池和容量
10        int ch[MAXN][2];
11        int key[MAXN],size[MAXN];
12        int sum0[MAXN],sum1[MAXN];
13        int status[MAXN];
14        void Init(){
15            tot1 = tot2 = 0;
16            size[0] = 0;
17            ch[0][0] = ch[0][1] = 0;
18            sum0[0] = sum1[0] = 0;
19        }
20        bool random(double p){
21            return (double)rand() / RAND_MAX < p;
22        }
23        int newnode(int val,int _status){
24            int r;
25            if(tot2)r = s[tot2--];
26            else r = ++tot1;
27            size[r] = 1;
28            key[r] = val;
29            status[r] = _status;


30              ch[r][0] = ch[r][1] = 0;
31              sum0[r] = sum1[r] = 0;//需要push_up
32              return r;
33       }
34       void del(int r){
35           if(!r)return;
36           s[++tot2] = r;
37           del(ch[r][0]);
38           del(ch[r][1]);
39       }
40       void push_up(int r){
41           int lson = ch[r][0], rson = ch[r][1];
42           size[r] = size[lson] + size[rson] + 1;
43           sum0[r] = gcd(sum0[lson],sum0[rson]);
44           sum1[r] = gcd(sum1[lson],sum1[rson]);
45           if(status[r] == 0)
46               sum0[r] = gcd(sum0[r],key[r]);
47           else sum1[r] = gcd(sum1[r],key[r]);
48       }
49       void merge(int &p,int x,int y){
50           if(!x || !y)
51               p = x|y;
52           else if(random((double)size[x]/(size[x]+size[y]))){
53               merge(ch[x][1],ch[x][1],y);
54               push_up(p=x);
55           }
56           else{
57               merge(ch[y][0],x,ch[y][0]);
58               push_up(p=y);
59           }
60       }
61       void split(int p,int &x,int &y,int k){
62           if(!k){
63               x = 0; y = p;
64               return;
65           }
66           if(size[ch[p][0]] >= k){
67               y = p;
68               split(ch[p][0],x,ch[y][0],k);
69               push_up(y);
70           }
71           else{
72               x = p;
73               split(ch[p][1],ch[x][1],y,k - size[ch[p][0]] - 1);
74               push_up(x);
75           }
76       }
77       void build(int &p,int l,int r){
78           if(l > r)return;
79           int mid = (l + r)/2;
80           p = newnode(num[mid],st[mid]);


81               build(ch[p][0],l,mid-1);
82               build(ch[p][1],mid+1,r);
83               push_up(p);
84        }
85        void debug(int root){
86            if(root == 0)return;
87            printf("%d   左儿子：%d  右儿子: %d  size = %d  key = %d\n",
                 root,ch[root][0],ch[root][1],size[root],key[root]);
 88           debug(ch[root][0]);
 89           debug(ch[root][1]);
 90       }
 91   };
 92   Treap T;
 93   char op[10];
 94   int main(){
 95       int n,q;
 96       while(scanf("%d%d",&n,&q) == 2){
 97            int root = 0;
 98            T.Init();
 99            for(int i = 1;i <= n;i++)
100                scanf("%d%d",&num[i],&st[i]);
101            T.build(root,1,n);
102            while(q--){
103                scanf("%s",op);
104                if(op[0] == 'Q'){
105                    int l,r,s;
106                    scanf("%d%d%d",&l,&r,&s);
107                    int x,y,z;
108                    T.split(root,x,z,r);
109                    T.split(x,x,y,l-1);
110                    if(s == 0)
111                        printf("%d\n",T.sum0[y] == 0? -1:T.sum0[y]);
112                    else
113                        printf("%d\n",T.sum1[y] == 0?-1:T.sum1[y]);
114                    T.merge(x,x,y);
115                    T.merge(root,x,z);
116                }
117                else if(op[0] == 'I'){
118                    int v,s,loc;
119                    scanf("%d%d%d",&loc,&v,&s);
120                    int x,y;
121                    T.split(root,x,y,loc);
122                    T.merge(x,x,T.newnode(v,s));
123                    T.merge(root,x,y);
124                }
125                else if(op[0] == 'D'){
126                    int loc;
127                    scanf("%d",&loc);
128                    int x,y,z;
129                    T.split(root,x,z,loc);
130                    T.split(x,x,y,loc-1);


131                    T.del(y);
132                    T.merge(root,x,z);
133                 }
134                 else if(op[0] == 'R'){
135                     int loc;
136                     scanf("%d",&loc);
137                     int x,y,z;
138                     T.split(root,x,z,loc);
139                     T.split(x,x,y,loc-1);
140                     T.status[y] = 1-T.status[y];
141                     T.push_up(y);
142                     T.merge(x,x,y);
143                     T.merge(root,x,z);
144                 }
145                 else{
146                     int loc,v;
147                     scanf("%d%d",&loc,&v);
148                     int x,y,z;
149                     T.split(root,x,z,loc);
150                     T.split(x,x,y,loc-1);
151                     T.key[y] = v;
152                     T.push_up(y);
153                     T.merge(x,x,y);
154                     T.merge(root,x,z);
155                 }
156             }
157      }
158      return 0;
159 }
\end{Code}

\subsection{Cây KD}

\subsubsection{HDU4347 K láng giềng gần nhất}

\begin{Code}
     模板题，求出最近的 K 个点。
 1   const int MAXN = 50010;
 2   const int DIM = 10;
 3   inline double sqr(double x){return x*x;}
 4   namespace KDTree{
 5       int K;//维数
 6       struct Point{
 7           int x[DIM];
 8           double distance(const Point &b)const{
 9               double ret = 0;
10               for(int i = 0;i < K;i++)
11                   ret += sqr(x[i]-b.x[i]);
12               return ret;
13           }
14           void input(){
15               for(int i = 0;i < K;i++)scanf("%d",&x[i]);
16           }


17              void output(){
18                  for(int i = 0;i < K;i++)
19                      printf("%d%c",x[i],i < K-1?' ':'\n');
20              }
21       };
22       struct qnode{
23           Point p;
24           double dis;
25           qnode(){}
26           qnode(Point _p,double _dis){
27               p = _p; dis = _dis;
28           }
29           bool operator <(const qnode &b)const{
30               return dis < b.dis;
31           }
32       };
33       priority_queue<qnode>q;
34       struct cmpx{
35           int div;
36           cmpx(const int &_div){div = _div;}
37           bool operator()(const Point &a,const Point &b){
38               for(int i = 0;i < K;i++)
39                    if(a.x[(div+i)%K] != b.x[(div+i)%K])
40                        return a.x[(div+i)%K] < b.x[(div+i)%K];
41               return true;
42           }
43       };
44       bool cmp(const Point &a,const Point &b,int div){
45           cmpx cp = cmpx(div);
46           return cp(a,b);
47       }
48       struct Node{
49           Point e;
50           Node *lc,*rc;
51           int div;
52       }pool[MAXN],*tail,*root;
53       void init(){
54           tail = pool;
55       }
56       Node* build(Point *a,int l,int r,int div){
57           if(l >= r)return NULL;
58           Node *p = tail++;
59           p->div = div;
60           int mid = (l+r)/2;
61           nth_element(a+l,a+mid,a+r,cmpx(div));
62           p->e = a[mid];
63           p->lc = build(a,l,mid,(div+1)%K);
64           p->rc = build(a,mid+1,r,(div+1)%K);
65           return p;
66       }
67       void search(Point p,Node *x,int div,int m){


 68              if(!x)return;
 69              if(cmp(p,x->e,div)){
 70                  search(p,x->lc,(div+1)%K,m);
 71                  if(q.size() < m){
 72                      q.push(qnode(x->e,p.distance(x->e)));
 73                      search(p,x->rc,(div+1)%K,m);
 74                  }
 75                  else {
 76                      if(p.distance(x->e) < q.top().dis){
 77                          q.pop();
 78                          q.push(qnode(x->e,p.distance(x->e)));
 79                      }
 80                      if(sqr(x->e.x[div]-p.x[div]) < q.top().dis)
 81                          search(p,x->rc,(div+1)%K,m);
 82                  }
 83              }
 84              else {
 85                  search(p,x->rc,(div+1)%K,m);
 86                  if(q.size() < m){
 87                      q.push(qnode(x->e,p.distance(x->e)));
 88                      search(p,x->lc,(div+1)%K,m);
 89                  }
 90                  else {
 91                      if(p.distance(x->e) < q.top().dis){
 92                          q.pop();
 93                          q.push(qnode(x->e,p.distance(x->e)));
 94                      }
 95                      if(sqr(x->e.x[div]-p.x[div]) < q.top().dis)
 96                          search(p,x->lc,(div+1)%K,m);
 97                  }
 98              }
 99       }
100       void search(Point p,int m){
101           while(!q.empty())q.pop();
102           search(p,root,0,m);
103       }
104   };
105   KDTree::Point p[MAXN];
106   int main()
107   {
108       int n,k;
109       while(scanf("%d%d",&n,&k) == 2){
110           KDTree::K = k;
111           for(int i = 0;i < n;i++)p[i].input();
112           KDTree::init();
113           KDTree::root = KDTree::build(p,0,n,0);
114           int Q;
115           scanf("%d",&Q);
116           KDTree::Point o;
117           while(Q--){
118                o.input();


119                int m;
120                scanf("%d",&m);
121                KDTree::search(o,m);
122                printf("the closest %d points are:\n",m);
123                int cnt = 0;
124                while(!KDTree::q.empty()){
125                    p[cnt++] = KDTree::q.top().p;
126                    KDTree::q.pop();
127                }
128                for(int i = 0;i < m;i++)p[m-1-i].output();
129          }
130      }
131      return 0;
132 }
\end{Code}

\subsubsection{CF44G}

\begin{Code}
     给定若干个靶子 ( xl, xr, yl, yr, z)，z 为该靶子离射击位置的距离，所有靶子都可以看成是二
     维平面上平行于坐标轴的矩形。然后按顺序给定若干个子弹的射击位置 ( x, y)，子弹射到一
     个靶子就会将靶子打碎，并掉落到地上。问每个子弹射到的靶子是谁。保证靶子的 z 值不相
     同。
     找矩形内权值最小的点，支持删除操作
 1 const int MAXN = 100010;
 2 const int INF = 0x3f3f3f3f;
 3 struct Point{
 4     int x,y,id;
 5     bool operator ==(const Point &b)const{
 6         return x == b.x && y == b.y && id == b.id;
 7     }
 8 };
 9 struct Node{
10     Point e;
11     Node *lc,*rc;
12     bool div;
13     int sub,cur;
14     int size;
15     bool exist;
16     void push_up(){
17         size = lc->size + rc->size + exist;
18         sub = min(cur,min(lc->sub,rc->sub));
19     }
20 }pool[MAXN],*root,*tail,*null;
21 inline bool cmpX(const Point &a,const Point &b){return a.x < b.x ||
       (a.x == b.x && a.y < b.y) || (a.x == b.x && a.y == b.y && a.id
      < b.id);}
22 inline bool cmpY(const Point &a,const Point &b){return a.y < b.y ||
       (a.y == b.y && a.x < b.x) || (a.y == b.y && a.x == b.x && a.id
      < b.id);}
23 inline bool cmp(const Point &a,const Point &b,bool div){return div?
      cmpY(a,b):cmpX(a,b);}
24 Node* build(Point *a,int l,int r,bool div){


25       if(l >= r)return null;
26       Node *p = tail++;
27       p->div = div;
28       int mid = (l+r)/2;
29       nth_element(a+l,a+mid,a+r,div?cmpY:cmpX);
30       p->e = a[mid];
31       p->lc = build(a,l,mid,!div);
32       p->rc = build(a,mid+1,r,!div);
33       p->exist = 1;
34       p->cur = p->e.id;
35       p->push_up();
36       return p;
37   }
38   void remove(Node *p,Point o){
39       if(p->e == o){
40            p->exist = 0;
41            p->cur = INF;
42            p->size--;
43       }
44       else {
45            if(cmp(p->e,o,p->div))remove(p->rc,o);
46            else remove(p->lc,o);
47       }
48       p->push_up();
49   }
50   int getMin(Node *p,int xl,int xr,int yl,int yr,int minx,int maxx,
        int miny,int maxy){
51       if(p == null || p->size == 0)return INF;
52       if(xl <= minx && xr >= maxx && yl <= miny && yr >= maxy)return
            p->sub;
53       if(xl > maxx || xr < minx || yl > maxy || yr < miny)return INF;
54       int ret = INF;
55       if(p->e.x >= xl && p->e.x <= xr && p->e.y >= yl && p->e.y <= yr
            )
56            ret = min(ret,p->cur);
57       if(p->div){
58            if(yl <= p->e.y)
59                ret = min(ret,getMin(p->lc,xl,xr,yl,min(yr,p->e.y),minx
                     ,maxx,miny,min(maxy,p->e.y)));
60            if(yr >= p->e.y)
61                ret = min(ret,getMin(p->rc,xl,xr,max(yl,p->e.y),yr,minx
                     ,maxx,max(miny,p->e.y),maxy));
62       }
63       else {
64            if(xl <= p->e.x)
65                ret = min(ret,getMin(p->lc,xl,min(xr,p->e.x),yl,yr,minx
                     ,min(maxx,p->e.x),miny,maxy));
66            if(xr >= p->e.x)
67                ret = min(ret,getMin(p->rc,max(xl,p->e.x),xr,yl,yr,max(
                     minx,p->e.x),maxx,miny,maxy));
68       }


 69       return ret;
 70   }
 71   Point pp[MAXN],pp2[MAXN];
 72   struct REC{
 73       int xl,xr,yl,yr,z;
 74       int id;
 75       void input(){
 76           scanf("%d%d%d%d%d",&xl,&xr,&yl,&yr,&z);
 77       }
 78       bool operator <(const REC &b)const{
 79           return z < b.z;
 80       }
 81   }rec[MAXN];
 82   int ans[MAXN];
 83   int main()
 84   {
 85       int n,m;
 86       while(scanf("%d",&n) == 1){
 87           for(int i = 0;i < n;i++){
 88                rec[i].input();
 89                rec[i].id = i+1;
 90           }
 91           sort(rec,rec+n);
 92           scanf("%d",&m);
 93           for(int i = 0;i < m;i++){
 94                scanf("%d%d",&pp[i].x,&pp[i].y);
 95                pp[i].id = i;
 96                pp2[i] = pp[i];//备份
 97           }
 98           tail = pool;
 99           null = tail++;
100           null->size = 0;
101           null->sub = null->cur = INF;
102           null->lc = null->rc = null;
103           root = build(pp,0,m,0);
104           memset(ans,0,sizeof(ans));
105           for(int i = 0;i < n;i++){
106                int tmp = getMin(root,rec[i].xl,rec[i].xr,rec[i].yl,rec
                      [i].yr,-INF,INF,-INF,INF);
107                if(tmp == INF)continue;
108                ans[tmp] = rec[i].id;
109                remove(root,pp2[tmp]);
110           }
111           for(int i = 0;i < m;i++)printf("%d\n",ans[i]);
112       }
113       return 0;
114   }
\end{Code}

\subsubsection{HDU4742}

\begin{Code}
     三维 LIS。即每个点有个三维坐标，两个点能放在一前一后当且仅当 xi < xj, yi < yj, zi < zj，
     求最长的序列，并该条件下的方案数。
 1   const int MAXN = 100010;
 2   const int MOD = 1<<30;
 3   const int INF = 0x7fffffff;//这个一定要够大
 4   struct Node{
 5       pair<int,int>e,sub,cur;
 6       bool div;
 7       Node *lc,*rc;
 8   };
 9   Node pool[MAXN],*tail;
10   Node *root;
11   bool cmpX(const pair<int,int> &a,const pair<int,int> &b){return a.
        first < b.first || (a.first == b.first && a.second < b.second)
        ;}
12   bool cmpY(const pair<int,int> &a,const pair<int,int> &b){return a.
        second < b.second || (a.second == b.second && a.first < b.first)
        ;}
13   bool cmp(const pair<int,int> &a,const pair<int,int> &b,bool div){
        return div?cmpY(a,b):cmpX(a,b);}
14   Node* build(pair<int,int> *a,int l,int r,bool div){
15       if(l >= r)return NULL;
16       Node *p = tail++;
17       p->div = div;
18       int mid = (l+r)/2;
19       nth_element(a+l,a+mid,a+r,div?cmpY:cmpX);
20       p->e = a[mid];
21       p->cur = p->sub = make_pair(0,0);
22       p->lc = build(a,l,mid,!div);
23       p->rc = build(a,mid+1,r,!div);
24       return p;
25   }
26   inline void update(pair<int,int> &a,pair<int,int> b){
27       if(a.first < b.first)a = b;
28       else if(a.first == b.first){
29           a.second += b.second;
30           if(a.second >= MOD)a.second -= MOD;
31       }
32   }
33   void add(Node *p,pair<int,int> e,pair<int,int> v){
34       update(p->sub,v);
35       if(e == p->e){
36           update(p->cur,v);
37           return;
38       }
39       else {
40           if(cmp(p->e,e,p->div))add(p->rc,e,v);
41           else add(p->lc,e,v);
42       }


43   }
44   pair<int,int>ans;
45   //查询最大值
46   void get(Node *p,pair<int,int>e,int maxx,int maxy){
47       if(!p)return;
48       if(p->sub.first < ans.first)return;
49       if(maxx <= e.first && maxy <= e.second)
50           update(ans,p->sub);
51       else {
52           if(p->e.first <= e.first && p->e.second <= e.second)update(
                ans,p->cur);
53           if(p->div){
54                if(p->e.second <= e.second)get(p->rc,e,maxx,maxy);
55                get(p->lc,e,maxx,min(maxy,p->e.second));
56           }
57           else {
58                if(p->e.first <= e.first)get(p->rc,e,maxx,maxy);
59                get(p->lc,e,min(maxx,p->e.first),maxy);
60           }
61       }
62   }
63   struct TNode{
64       int x,y,z;
65       void input(){
66           scanf("%d%d%d",&x,&y,&z);
67       }
68       bool operator < (const TNode &b)const{
69           if(x != b.x)return x < b.x;
70           else if(y != b.y)return y < b.y;
71           else return z < b.z;
72       }
73   }node[MAXN];
74   pair<int,int>p[MAXN];
75   pair<int,int>dp[MAXN];
76   int main()
77   {
78       int T;
79       int n;
80       scanf("%d",&T);
81       while(T--){
82           scanf("%d",&n);
83           int cnt = 0;
84           for(int i = 0;i < n;i++){
85                node[i].input();
86                p[cnt++] = make_pair(node[i].y,node[i].z);
87           }
88           sort(node,node+n);
89           sort(p,p+cnt);
90           cnt = unique(p,p+cnt)-p;
91           tail = pool;
92           root = build(p,0,cnt,0);


 93            for(int i = 0;i < n;i++)dp[i] = make_pair(1,1);
 94            for(int i = 0;i < n;i++){
 95                ans = make_pair(0,0);
 96                get(root,make_pair(node[i].y,node[i].z),INF,INF);
 97                ans.first++;
 98                update(dp[i],ans);
 99                add(root,make_pair(node[i].y,node[i].z),dp[i]);
100            }
101            printf("%d %d\n",root->sub.first,root->sub.second);
102       }
103       return 0;
104 }
\end{Code}

\subsection{Cây scapegoat}

\subsubsection{CF455D}

\begin{Code}
    http://codeforces.com/contest/455/problem/D
    题意：给了一个序列，1 操作把一个区间的末尾的数插入到头部，2 操作是询问一个区间里
    面等于某个数的个数。
    使用替罪羊树，里面套一个 map 来统计区间的个数。
 1 const int MAXN = 200010;
 2 const double alpha = 0.75;
 3 struct Node{
 4     Node *ch[2];
 5     int size,key,nodeCount;
 6     bool exist;
 7     map<int,int>mp;
 8     bool isBad(){
 9         return ch[0]->nodeCount > alpha*nodeCount+5 || ch[1]->
              nodeCount > alpha*nodeCount + 5;
10     }
11     void push_up(){
12         size = exist + ch[0]->size + ch[1]->size;
13         nodeCount = 1 + ch[0]->nodeCount + ch[1]->nodeCount;
14         mp.clear();
15         if(exist)mp[key]++;
16         for(map<int,int>::iterator it = ch[0]->mp.begin();it != ch
              [0]->mp.end();it++)
17              mp[(*it).first] += (*it).second;
18         for(map<int,int>::iterator it = ch[1]->mp.begin();it != ch
              [1]->mp.end();it++)
19              mp[(*it).first] += (*it).second;
20     }
21 };
22 struct ScapeGoatTree{
23     Node pool[MAXN];
24     Node *tail,*root,*null;
25     Node *bc[MAXN];//内存回收
26     int bc_top;
27     void init(){


28              tail = pool;
29              null = tail++;
30              null->ch[0] = null->ch[1] = null;
31              null->size = null->key = null->nodeCount = 0;
32              null->mp.clear();
33              root = null;
34              bc_top = 0;
35       }
36       inline Node *newNode(int key){
37           Node *p;
38           if(bc_top)p = bc[--bc_top];
39           else p = tail++;
40           p->ch[0] = p->ch[1] = null;
41           p->size = p->nodeCount = 1;
42           p->key = key;
43           p->exist = true;
44           p->mp.clear();
45           p->mp[key] = 1;
46           return p;
47       }
48       Node *buildTree(int *a,int l,int r){
49           if(l >= r)return null;
50           int mid = (l+r)>>1;
51           Node *p = newNode(a[mid]);
52           p->ch[0] = buildTree(a,l,mid);
53           p->ch[1] = buildTree(a,mid+1,r);
54           p->push_up();
55           return p;
56       }
57       inline void Travel(Node *p,vector<Node *>&v){
58           if(p == null)return;
59           Travel(p->ch[0],v);
60           if(p->exist)v.push_back(p);
61           else bc[bc_top++] = p;
62           Travel(p->ch[1],v);
63       }
64       inline Node *divide(vector<Node *>&v,int l,int r){
65           if(l >= r)return null;
66           int mid = (l+r)/2;
67           Node *p = v[mid];
68           p->ch[0] = divide(v,l,mid);
69           p->ch[1] = divide(v,mid+1,r);
70           p->push_up();
71           return p;
72       }
73       //重构，注意 p 要引用
74       inline void rebuild(Node *&p){
75           vector<Node *>v;
76           Travel(p,v);
77           p = divide(v,0,v.size());
78       }


 79       //删除第 id 个元素，返回第 id 个元素的值
 80       inline int erase(Node *p,int id){
 81           if(p->exist && id == p->ch[0]->size + 1){
 82               p->exist = 0;
 83               p->mp[p->key]--;
 84               p->size--;
 85               return p->key;
 86           }
 87           p->size--;
 88           int res;
 89           if(p->ch[0]->size >= id)
 90               res = erase(p->ch[0],id);
 91           else res = erase(p->ch[1],id - p->ch[0]->size - p->exist);
 92           p->mp[res]--;
 93           return res;
 94       }
 95       //删除一定的点以后重构
 96       void check_erase(){
 97           if(root->size < 0.5*root->nodeCount)
 98               rebuild(root);
 99       }
100       Node **insert(Node *&p,int id,int val){
101           if(p == null){
102               p = newNode(val);
103               return &null;
104           }
105           else {
106               p->size++;
107               p->nodeCount++;
108               p->mp[val]++;
109               Node ** res;
110               if(id <= p->ch[0]->size+p->exist)
111                   res = insert(p->ch[0],id,val);
112               else res = insert(p->ch[1],id-p->ch[0]->size-p->exist,
                     val);
113               if(p->isBad())res = &p;
114               return res;
115           }
116       }
117       //在第 id 个位置插入数 val
118       void insert(int id,int val){
119           Node **p = insert(root,id,val);
120           if(*p != null)rebuild(*p);
121       }
122       //查询 [l,r] 之间值为 val 的数的个数
123       int query(Node *p,int l,int r,int val){
124           if(p == null)return 0;
125           if(l <= 1 && p->size <= r)
126               return p->mp.count(val)?p->mp[val]:0;
127           else {
128               int ans = 0;


129                  if(l <= p->ch[0]->size)
130                      ans += query(p->ch[0],l,r,val);
131                  if(r > p->ch[0]->size+p->exist)
132                      ans += query(p->ch[1],l - p->ch[0]->size - p->exist
                            , r - p->ch[0]->size - p->exist,val);
133                  if(p->exist && p->key == val && l <= p->ch[0]->size+1
                        && r >= p->ch[0]->size+1)
134                      ans++;
135                  return ans;
136              }
137       }
138   }tree;
139   int a[MAXN];
140   int main()
141   {
142       int n;
143       while(scanf("%d",&n) == 1){
144           tree.init();
145           for(int i = 0;i < n;i++)scanf("%d",&a[i]);
146           tree.root = tree.buildTree(a,0,n);
147           int m;
148           int op,l,r,k;
149           scanf("%d",&m);
150           int ans = 0;
151           while(m--){
152                scanf("%d",&op);
153                if(op == 1){
154                    scanf("%d%d",&l,&r);
155                    l = ((l+ans-1)%n)+1;
156                    r = ((r+ans-1)%n)+1;
157                    if(l > r)swap(l,r);
158                    int v = tree.erase(tree.root,r);
159                    //tree.check_erase(); //有时候可以加上删除重构
160                    tree.insert(l,v);
161                }
162                else {
163                    scanf("%d%d%d",&l,&r,&k);
164                    l = ((l+ans-1)%n)+1;
165                    r = ((r+ans-1)%n)+1;
166                    k = ((k+ans-1)%n)+1;
167                    if(l > r)swap(l,r);
168                    ans = tree.query(tree.root,l,r,k);
169                    printf("%d\n",ans);
170                }
171           }
172       }
173       return 0;
174   }
\end{Code}

\subsection{Cây KD động}

\begin{Code}
     动态 KD 树就是结合了 KD 树和替罪羊树。支持 KD 树的插入删除操作，用替罪羊树的思想
     来保存平衡。
     UVALive6045
     题意：给了二维平面上的 N 个整点（N ≤ 50000）。每次操作给了点 ( xi , yi ), 需要曼哈顿距
     离小于 E 的点进行一个变换。输出最后的点的坐标，保证变换次数不超过 50000.
 1   const int MAXN = 100010;
 2   const double alpha = 0.75;
 3   struct Point{
 4       int x,y,id;
 5   };
 6   struct Node{
 7       Point e;
 8       int size,nodeCount;
 9       Node *lc,*rc;
10       bool div;
11       bool exist;
12       bool isBad(){
13           return lc->nodeCount > alpha*nodeCount+5 || rc->nodeCount >
                 alpha*nodeCount+5;
14       }
15       inline void push_up(){
16           size = exist + lc->size + rc->size;
17           nodeCount = 1+lc->nodeCount+rc->nodeCount;
18       }
19   };
20   Node pool[MAXN],*tail,*root,*null;
21   Node *bc[MAXN];
22   int bc_top;
23   void init(){
24       tail = pool;
25       null = tail++;
26       null->lc = null->rc = null;
27       null->size = null->nodeCount = 0;
28       root = null;
29       bc_top = 0;
30   }
31   Node *newNode(Point e){
32       Node *p;
33       if(bc_top)p = bc[--bc_top];
34       else p = tail++;
35       p->e = e;
36       p->lc = p->rc = null;
37       p->size = p->nodeCount = 1;
38       p->exist = true;
39       return p;
40   }
41   inline bool cmpX(const Point &a,const Point &b){
42       return a.x < b.x || (a.x == b.x && a.y < b.y) || (a.x == b.x &&
             a.y == b.y && a.id < b.id);


43 }
44 inline bool cmpY(const Point &a,const Point &b){
45     return a.y < b.y || (a.y == b.y && a.x < b.x) || (a.y == b.y &&
           a.x == b.x && a.id < b.id);
46 }
47 inline bool cmp(const Point &a,const Point &b,bool div){
48     return div?cmpY(a,b):cmpX(a,b);
49 }
50 //注意 a 需要备份，否则就乱序
51 Node *build(Point *a,int l,int r,bool div){
52     if(l >= r)return null;
53     int mid = (l+r)/2;
54     nth_element(a+l,a+mid,a+r,div?cmpY:cmpX);
55     Node *p = newNode(a[mid]);
56     p->div = div;
57     p->lc = build(a,l,mid,!div);
58     p->rc = build(a,mid+1,r,!div);
59     p->push_up();
60     return p;
61 }
62 void Travel(Node *p,vector<Point>&v){
63     if(p == null)return;
64     Travel(p->lc,v);
65     if(p->exist)v.push_back(p->e);
66     bc[bc_top++] = p;
67     Travel(p->rc,v);
68 }
69 Node *divide(vector<Point>&v,int l,int r,bool div){
70     if(l >= r)return null;
71     int mid = (l+r)/2;
72     nth_element(v.begin()+l,v.begin()+mid,v.begin()+r,div?cmpY:cmpX
          );
73     Node *p = newNode(v[mid]);
74     p->div = div;
75     p->lc = divide(v,l,mid,!div);
76     p->rc = divide(v,mid+1,r,!div);
77     p->push_up();
78     return p;
79 }
80 inline void rebuild(Node *&p){
81     vector<Point>v;
82     Travel(p,v);
83     p = divide(v,0,v.size(),p->div);
84 }
85 Node **insert(Node *&p,Point a,bool div){
86     if(p == null){
87         p = newNode(a);
88         p->div = div;
89         return &null;
90     }
91     else {


 92              p->nodeCount++;
 93              p->size++;
 94              Node **res;
 95              if(cmp(a,p->e,div))
 96                  res = insert(p->lc,a,!div);
 97              else res = insert(p->rc,a,!div);
 98              if(p->isBad())res = &p;
 99              return res;
100       }
101   }
102   void insert(Point e){
103       Node **p = insert(root,e,0);
104       if(*p != null)rebuild(*p);
105   }
106   vector<int>vec;
107   void getvec(Node *p,int minx,int maxx,int miny,int maxy){
108       if(p->size == 0)return;
109       if(p->exist && minx <= p->e.x && p->e.x <= maxx && miny <= p->e
             .y && p->e.y <= maxy){
110           vec.push_back(p->e.id);
111           p->exist = 0;
112           p->size--;
113       }
114       if(p->div? miny <= p->e.y : minx <= p->e.x)getvec(p->lc,minx,
             maxx,miny,maxy);
115       if(p->div? maxy >= p->e.y : maxx >= p->e.x)getvec(p->rc,minx,
             maxx,miny,maxy);
116       p->push_up();
117   }
118   Point p[MAXN],p2[MAXN];
119   Point p3[MAXN];
120   int main()
121   {
122       int T;
123       scanf("%d",&T);
124       int iCase = 0;
125       int N,Q,W,H;
126       while(T--){
127           iCase++;
128           scanf("%d%d%d%d",&N,&Q,&W,&H);
129           init();
130           for(int i = 0;i < N;i++){
131               scanf("%d%d",&p[i].x,&p[i].y);
132               p[i].id = p2[i].id = i;
133               p2[i].x = p[i].x+p[i].y;
134               p2[i].y = p[i].x-p[i].y;
135               p3[i] = p2[i];
136           }
137           root = build(p3,0,N,0);
138           int X,Y,E,a,b,c,d,e,f;
139           while(Q--){


140                scanf("%d%d%d%d%d%d%d%d%d",&X,&Y,&E,&a,&b,&c,&d,&e,&f);
141                vec.clear();
142                int minx = X+Y-E;
143                int maxx = X+Y+E;
144                int miny = X-Y-E;
145                int maxy = X-Y+E;
146                getvec(root,minx,maxx,miny,maxy);
147                int sz = vec.size();
148                for(int i = 0;i < sz;i++){
149                    int id = vec[i];
150                    long long tx = p[id].x;
151                    long long ty = p[id].y;
152                    p[id].x = (tx*a+ty*b+(long long)(id+1)*c)%W;
153                    p[id].y = (tx*d+ty*e+(long long)(id+1)*f)%H;
154                    p2[id].x = p[id].x+p[id].y;
155                    p2[id].y = p[id].x-p[id].y;
156                    insert(p2[id]);
157                }
158            }
159            printf("Case #%d:\n",iCase);
160            for(int i = 0;i < N;i++)
161                printf("%d %d\n",p[i].x,p[i].y);
162     }
163     return 0;
164 }
\end{Code}

\subsection{Cây lồng cây}

\subsubsection{Cây scapegoat lồng splay}

\begin{Code}
    BZOJ 3065: 带插入区间 K 小值
    带插入、修改的区间 k 小值在线查询。
    1. Q x y k: 询问从左至右第 x 只跳蚤到从左至右第 y 只跳蚤中，弹跳力第 k 小的跳蚤的弹
    跳力是多少。(1 <= x <= y <= m, 1 <= k <= y - x + 1)
    2. M x val: 将从左至右第 x 只跳蚤的弹跳力改为 val。(1 <= x <= m)
    3. I x val: 在从左至右第 x 只跳蚤的前面插入一只弹跳力为 val 的跳蚤。即插入后从左至右
    第 x 只跳蚤是我刚插入的跳蚤。(1 <= x <= m + 1)
 1 const int MAXN = 70010;
 2 namespace Splay{
 3     struct Node *null;
 4     struct Node{
 5         Node *ch[2],*fa;
 6         int size,key,cnt;
 7         inline void setc(Node *p,int d){
 8             ch[d] = p;
 9             p->fa = this;
10         }
11         inline bool d(){
12             return fa->ch[1] == this;
13         }
14         inline void push_up(){


15                  size = ch[0]->size + ch[1]->size + cnt;
16              }
17              void clear(int _key){
18                  size = cnt = 1;
19                  key = _key;
20                  ch[0] = ch[1] = fa = null;
21              }
22              inline bool isroot(){
23                  return fa == null || this != fa->ch[0] && this != fa->
                       ch[1];
24              }
25       };
26       Node pool[MAXN*20],*tail;
27       Node *bc[MAXN*20];
28       int bc_top;//内存回收
29       void init(){
30           tail = pool;
31           bc_top = 0;
32           null = tail++;
33           null->size = null->cnt = 0;
34           null->ch[0] = null->ch[1] = null->fa = null;
35       }
36       inline void rotate(Node *x){
37           Node *f = x->fa, *ff = x->fa->fa;
38           int c = x->d(), cc = f->d();
39           f->setc(x->ch[!c],c);
40           x->setc(f,!c);
41           if(ff->ch[cc] == f)ff->setc(x,cc);
42           else x->fa = ff;
43           f->push_up();
44       }
45       inline void splay(Node* &root,Node* x,Node* goal){
46           while(x->fa != goal){
47               if(x->fa->fa == goal)rotate(x);
48               else {
49                   bool f = x->fa->d();
50                   x->d() == f ? rotate(x->fa) : rotate(x);
51                   rotate(x);
52               }
53           }
54           x->push_up();
55           if(goal == null)root = x;
56       }
57       //找到 r 子树里面的最左边那个
58       Node* get_left(Node* r){
59           Node* x = r;
60           while(x->ch[0] != null)x = x->ch[0];
61           return x;
62       }
63       //在 root 的树中删掉 x
64       void erase(Node* &root,Node* x){


 65              splay(root,x,null);
 66              Node* t = root;
 67              if(t->ch[1] != null){
 68                  root = t->ch[1];
 69                  splay(root,get_left(t->ch[1]),null);
 70                  root->setc(t->ch[0],0);
 71              }
 72              else root = root->ch[0];
 73              bc[bc_top++] = x;
 74              root->fa = null;
 75              if(root != null)root->push_up();
 76       }
 77       Node* newNode(int key){
 78           Node* p;
 79           if(bc_top)p = bc[--bc_top];
 80           else p = tail++;
 81           p->clear(key);
 82           return p;
 83       }
 84       //插入一个值 key
 85       void insert(Node* &root,int key){
 86           if(root == null){
 87               root = newNode(key);
 88               return;
 89           }
 90           Node* now = root;
 91           Node* pre = root->fa;
 92           while(now != null){
 93               if(now->key == key){
 94                    now->cnt++;
 95                    splay(root,now,null);
 96                    return;
 97               }
 98               pre = now;
 99               now = now->ch[key >= now->key];
100           }
101           Node *x = newNode(key);
102           pre->setc(x,key >= pre->key);
103           splay(root,x,null);
104       }
105       //删除一个值 key
106       void erase(Node* &root,int key){
107           Node* now = root;
108           while(now->key != key){
109               now = now->ch[key >= now->key];
110           }
111           now->cnt--;
112           if(now->cnt == 0)erase(root,now);
113           else splay(root,now,null);
114       }
115       void Travel(Node* r){


116         if(r == null)return;
117         Travel(r->ch[0]);
118         bc[bc_top++] = r;
119         Travel(r->ch[1]);
120     }
121     void CLEAR(Node* &root){
122         Travel(root);
123         root = null;
124     }
125     //查询小于等于 val 的个数
126     int query(Node *root,int val){
127         int ans = 0;
128         Node *x = root;
129         while(x != null){
130             if(val < x->key)x = x->ch[0];
131             else{
132                  ans += x->ch[0]->size + x->cnt;
133                  x = x->ch[1];
134             }
135         }
136         return ans;
137     }
138 };
139 namespace ScapeGoatTree{
140     const double alpha = 0.75;
141     struct Node{
142         Node *ch[2];
143         int size,nodeCount,key;
144         Splay::Node *root;
145         bool isBad(){
146             return ch[0]->nodeCount > alpha*nodeCount+5 || ch[1]->
                   nodeCount > alpha*nodeCount+5;
147         }
148         void push_up(){
149             size = 1+ch[0]->size+ch[1]->size;
150             nodeCount = 1+ch[0]->nodeCount+ch[1]->nodeCount;
151         }
152     };
153     Node pool[MAXN];
154     Node *tail,*root,*null;
155     Node *bc[MAXN];
156     int bc_top;
157     void init(){
158         tail = pool;
159         null = tail++;
160         null->ch[0] = null->ch[1] = null;
161         null->size = null->nodeCount = 0;
162         null->root = Splay::null;
163         bc_top = 0;
164     }
165     inline Node* newNode(int key){


166              Node *p;
167              if(bc_top)p = bc[--bc_top];
168              else p = tail++;
169              p->key = key;
170              p->ch[0] = p->ch[1] = null;
171              p->size = p->nodeCount = 1;
172              p->root = Splay::null;
173              return p;
174       }
175       Node *buildTree(int *a,int l,int r){
176           if(l >= r)return null;
177           int mid = (l+r)/2;
178           Node *p = newNode(a[mid]);
179           for(int i = l;i < r;i++)
180               Splay::insert(p->root,a[i]);
181           p->ch[0] = buildTree(a,l,mid);
182           p->ch[1] = buildTree(a,mid+1,r);
183           p->push_up();
184           return p;
185       }
186       void Travel(Node *p,vector<int>&v){
187           if(p == null)return;
188           Travel(p->ch[0],v);
189           v.push_back(p->key);
190           Splay::CLEAR(p->root);
191           bc[bc_top++] = p;
192           Travel(p->ch[1],v);
193       }
194       Node *divide(vector<int>&v,int l,int r){
195           if(l == r)return null;
196           int mid = (l+r)/2;
197           Node *p = newNode(v[mid]);
198           for(int i = l;i < r;i++)
199               Splay::insert(p->root,v[i]);
200           p->ch[0] = divide(v,l,mid);
201           p->ch[1] = divide(v,mid+1,r);
202           p->push_up();
203           return p;
204       }
205       inline void rebuild(Node *&p){
206           vector<int>v;
207           Travel(p,v);
208           p = divide(v,0,v.size());
209       }
210       //将第 id 个值修改为 val
211       int Modify(Node *p,int id,int val){
212           if(id == p->ch[0]->size+1){
213               int v = p->key;
214               Splay::erase(p->root,v);
215               Splay::insert(p->root,val);
216               p->key = val;


217                  return v;
218              }
219              int res;
220              if(p->ch[0]->size >= id)
221                  res = Modify(p->ch[0],id,val);
222              else res = Modify(p->ch[1],id-p->ch[0]->size-1,val);
223              Splay::erase(p->root,res);
224              Splay::insert(p->root,val);
225              return res;
226       }
227       Node **insert(Node *&p,int id,int val){
228           if(p == null){
229               p = newNode(val);
230               Splay::insert(p->root,val);
231               return &null;
232           }
233           else {
234               p->size++;
235               p->nodeCount++;
236               Splay::insert(p->root,val);
237               Node ** res;
238               if(id <= p->ch[0]->size+1)
239                   res = insert(p->ch[0],id,val);
240               else res = insert(p->ch[1],id-p->ch[0]->size-1,val);
241               if(p->isBad())res = &p;
242               return res;
243           }
244       }
245       void insert(int id,int val){
246           Node **p = insert(root,id,val);
247           if(*p != null)rebuild(*p);
248       }
249       //查询在 [l,r] 区间，值小于等于 val 的个数
250       int query(Node *p,int l,int r,int val){
251           if(p == null)return 0;
252           if(l <= 1 && p->size <= r)
253               return Splay::query(p->root,val);
254           else {
255               int ans = 0;
256               if(l <= p->ch[0]->size)
257                   ans += query(p->ch[0],l,r,val);
258               if(r > p->ch[0]->size+1)
259                   ans += query(p->ch[1],l-p->ch[0]->size-1,r-p->ch
                         [0]->size-1,val);
260               if(p->key <= val && l <= p->ch[0]->size+1 && p->ch[0]->
                     size+1 <= r)
261                   ans++;
262               return ans;
263           }
264       }
265       int query(int L,int R,int k){


266              int ans;
267              int l = 0, r = 100000;
268              while(l <= r){
269                  int mid = (l+r)/2;
270                  if(query(root,L,R,mid) >= k){
271                      ans = mid;
272                      r = mid-1;
273                  }
274                  else l = mid+1;
275              }
276              return ans;
277       }
278   };
279   int a[MAXN];
280   int main()
281   {
282       int n;
283       while(scanf("%d",&n) == 1){
284           Splay::init();
285           ScapeGoatTree::init();
286           for(int i = 0;i < n;i++)scanf("%d",&a[i]);
287           ScapeGoatTree::root = ScapeGoatTree::buildTree(a,0,n);
288           int m;
289           char op[10];
290           scanf("%d",&m);
291           int ans = 0;
292           while(m--){
293                scanf("%s",op);
294                if(op[0] == 'Q'){
295                    int x,y,k;
296                    scanf("%d%d%d",&x,&y,&k);
297                    x ^= ans; y ^= ans; k ^= ans;
298                    ans = ScapeGoatTree::query(x,y,k);
299                    printf("%d\n",ans);
300                }
301                else if(op[0] == 'M'){
302                    int x,val;
303                    scanf("%d%d",&x,&val);
304                    x ^= ans; val ^= ans;
305                    ScapeGoatTree::Modify(ScapeGoatTree::root,x,val);
306                }
307                else if(op[0] == 'I'){
308                    int x,val;
309                    scanf("%d%d",&x,&val);
310                    x ^= ans; val ^= ans;
311                    ScapeGoatTree::insert(x,val);
312                }
313           }
314       }
315       return 0;
316   }
\end{Code}

\section{Lý thuyết đồ thị}

\subsection{Đường đi ngắn nhất}

\subsubsection{Dijkstra đường đi ngắn nhất một nguồn}

\begin{Code}
   权值必须是非负
 1 /*
 2 * 单源最短路径，Dijkstra 算法，邻接矩阵形式，复杂度为O(n^2)
 3 * 求出源 beg 到所有点的最短路径，传入图的顶点数，和邻接矩阵 cost[][]
 4 * 返回各点的最短路径 lowcost[], 路径 pre[].pre[i] 记录 beg 到 i 路径上的
       父结点，pre[beg]=-1
 5 * 可更改路径权类型，但是权值必须为非负
 6 */
 7 const int MAXN=1010;
 8 #define typec int
 9 const typec INF=0x3f3f3f3f;//防止后面溢出，这个不能太大
10 bool vis[MAXN];
11 int pre[MAXN];
12 void Dijkstra(typec cost[][MAXN],typec lowcost[],int n,int beg){
13     for(int i=0;i<n;i++){
14         lowcost[i]=INF;vis[i]=false;pre[i]=-1;
15     }
16     lowcost[beg]=0;
17     for(int j=0;j<n;j++){
18         int k=-1;
19         int Min=INF;
20         for(int i=0;i<n;i++)
21             if(!vis[i]&&lowcost[i]<Min){
22                 Min=lowcost[i];
23                 k=i;
24             }
25         if(k==-1)break;
26         vis[k]=true;
27         for(int i=0;i<n;i++)
28             if(!vis[i]&&lowcost[k]+cost[k][i]<lowcost[i]){
29                 lowcost[i]=lowcost[k]+cost[k][i];
30                 pre[i]=k;
31             }
32     }
33 }
\end{Code}

\subsubsection{Dijkstra + tối ưu bằng heap}

\begin{Code}
   使用优先队列优化，复杂度 O (E log E)
1 /*
2 * 使用优先队列优化 Dijkstra 算法
3 * 复杂度 O(ElogE)
4 * 注意对 vector<Edge>E[MAXN] 进行初始化后加边
5 */
6 const int INF=0x3f3f3f3f;


 7   const int MAXN=1000010;
 8   struct qnode{
 9       int v;
10       int c;
11       qnode(int _v=0,int _c=0):v(_v),c(_c){}
12       bool operator <(const qnode &r)const{
13           return c>r.c;
14       }
15   };
16   struct Edge{
17       int v,cost;
18       Edge(int _v=0,int _cost=0):v(_v),cost(_cost){}
19   };
20   vector<Edge>E[MAXN];
21   bool vis[MAXN];
22   int dist[MAXN];
23   //点的编号从 1 开始
24   void Dijkstra(int n,int start){
25       memset(vis,false,sizeof(vis));
26       for(int i=1;i<=n;i++)dist[i]=INF;
27       priority_queue<qnode>que;
28       while(!que.empty())que.pop();
29       dist[start]=0;
30       que.push(qnode(start,0));
31       qnode tmp;
32       while(!que.empty()){
33           tmp=que.top();
34           que.pop();
35           int u=tmp.v;
36           if(vis[u])continue;
37           vis[u]=true;
38           for(int i=0;i<E[u].size();i++){
39                int v=E[tmp.v][i].v;
40                int cost=E[u][i].cost;
41                if(!vis[v]&&dist[v]>dist[u]+cost){
42                    dist[v]=dist[u]+cost;
43                    que.push(qnode(v,dist[v]));
44                }
45           }
46       }
47   }
48   void addedge(int u,int v,int w){
49       E[u].push_back(Edge(v,w));
50   }
\end{Code}

\subsubsection{Bellman-Ford cho đường đi ngắn nhất một nguồn}

\begin{Code}
1 /*
2 * 单源最短路 bellman_ford 算法，复杂度 O(VE)
3 * 可以处理负边权图。
4 * 可以判断是否存在负环回路。返回 true, 当且仅当图中不包含从源点可达的负权回路


 5    * vector<Edge>E; 先 E.clear() 初始化，然后加入所有边
 6    * 点的编号从 1 开始 (从 0 开始简单修改就可以了)
 7    */
 8   const int INF=0x3f3f3f3f;
 9   const int MAXN=550;
10   int dist[MAXN];
11   struct Edge{
12       int u,v;
13       int cost;
14       Edge(int _u=0,int _v=0,int _cost=0):u(_u),v(_v),cost(_cost){}
15   };
16   vector<Edge>E;
17   //点的编号从 1 开始
18   bool bellman_ford(int start,int n){
19       for(int i=1;i<=n;i++)dist[i]=INF;
20       dist[start]=0;
21       //最多做 n-1 次
22       for(int i=1;i<n;i++){
23           bool flag=false;
24           for(int j=0;j<E.size();j++){
25                int u=E[j].u;
26                int v=E[j].v;
27                int cost=E[j].cost;
28                if(dist[v]>dist[u]+cost){
29                    dist[v]=dist[u]+cost;
30                    flag=true;
31                }
32           }
33           if(!flag)return true;//没有负环回路
34       }
35       for(int j=0;j<E.size();j++)
36           if(dist[E[j].v]>dist[E[j].u]+E[j].cost)
37                return false;//有负环回路
38       return true;//没有负环回路
39   }
\end{Code}

\subsubsection{SPFA cho đường đi ngắn nhất một nguồn}

\begin{Code}
 1   /*
 2    * 单源最短路 SPFA
 3    * 时间复杂度 0(kE)
 4    * 这个是队列实现，有时候改成栈实现会更加快，很容易修改
 5    * 这个复杂度是不定的
 6    */
 7   const int MAXN=1010;
 8   const int INF=0x3f3f3f3f;
 9   struct Edge{
10       int v;
11       int cost;
12       Edge(int _v=0,int _cost=0):v(_v),cost(_cost){}
13   };
14   vector<Edge>E[MAXN];


15   void addedge(int u,int v,int w){
16       E[u].push_back(Edge(v,w));
17   }
18   bool vis[MAXN];//在队列标志
19   int cnt[MAXN];//每个点的入队列次数
20   int dist[MAXN];
21   bool SPFA(int start,int n){
22       memset(vis,false,sizeof(vis));
23       for(int i=1;i<=n;i++)dist[i]=INF;
24       vis[start]=true;
25       dist[start]=0;
26       queue<int>que;
27       while(!que.empty())que.pop();
28       que.push(start);
29       memset(cnt,0,sizeof(cnt));
30       cnt[start]=1;
31       while(!que.empty()){
32           int u=que.front();
33           que.pop();
34           vis[u]=false;
35           for(int i=0;i<E[u].size();i++){
36               int v=E[u][i].v;
37               if(dist[v]>dist[u]+E[u][i].cost){
38                   dist[v]=dist[u]+E[u][i].cost;
39                   if(!vis[v]){
40                        vis[v]=true;
41                        que.push(v);
42                        if(++cnt[v]>n)return false;
43                        //cnt[i] 为入队列次数，用来判定是否存在负环回路
44                   }
45               }
46           }
47       }
48       return true;
49   }
\end{Code}

\subsection{Cây khung nhỏ nhất}

\subsubsection{Thuật toán Prim}

\begin{Code}
 1   /*
 2    * Prim 求 MST
 3    * 耗费矩阵 cost[][]，标号从 0 开始，0∼n-1
 4    * 返回最小生成树的权值，返回 -1 表示原图不连通
 5    */
 6   const int INF=0x3f3f3f3f;
 7   const int MAXN=110;
 8   bool vis[MAXN];
 9   int lowc[MAXN];
10   //点是 0 n-1
11   int Prim(int cost[][MAXN],int n){
12       int ans=0;


13       memset(vis,false,sizeof(vis));
14       vis[0]=true;
15       for(int i=1;i<n;i++)lowc[i]=cost[0][i];
16       for(int i=1;i<n;i++){
17           int minc=INF;
18           int p=-1;
19           for(int j=0;j<n;j++)
20               if(!vis[j]&&minc>lowc[j]){
21                   minc=lowc[j];
22                   p=j;
23               }
24           if(minc==INF)return -1;//原图不连通
25           ans+=minc;
26           vis[p]=true;
27           for(int j=0;j<n;j++)
28               if(!vis[j]&&lowc[j]>cost[p][j])
29                   lowc[j]=cost[p][j];
30       }
31       return ans;
32 }
\end{Code}

\subsubsection{Thuật toán Kruskal}

\begin{Code}
 1   /*
 2    * Kruskal 算法求 MST
 3    */
 4   const int MAXN=110;//最大点数
 5   const int MAXM=10000;//最大边数
 6   int F[MAXN];//并查集使用
 7   struct Edge{
 8       int u,v,w;
 9   }edge[MAXM];//存储边的信息，包括起点/终点/权值
10   int tol;//边数，加边前赋值为 0
11   void addedge(int u,int v,int w){
12       edge[tol].u=u;
13       edge[tol].v=v;
14       edge[tol++].w=w;
15   }
16   //排序函数，讲边按照权值从小到大排序
17   bool cmp(Edge a,Edge b){
18       return a.w<b.w;
19   }
20   int find(int x){
21       if(F[x]==-1)return x;
22       else return F[x]=find(F[x]);
23   }
24   //传入点数，返回最小生成树的权值，如果不连通返回 -1
25   int Kruskal(int n){
26       memset(F,-1,sizeof(F));
27       sort(edge,edge+tol,cmp);
28       int cnt=0;//计算加入的边数
29       int ans=0;


30         for(int i=0;i<tol;i++){
31             int u=edge[i].u;
32             int v=edge[i].v;
33             int w=edge[i].w;
34             int t1=find(u);
35             int t2=find(v);
36             if(t1!=t2){
37                 ans+=w;
38                 F[t1]=t2;
39                 cnt++;
40             }
41             if(cnt==n-1)break;
42         }
43         if(cnt<n-1)return -1;//不连通
44         else return ans;
45 }
\end{Code}

\subsection{Cây khung nhỏ thứ hai}

\begin{Code}
 1   /*
 2    * 次小生成树
 3    * 求最小生成树时，用数组 Max[i][j] 来表示 MST 中 i 到 j 最大边权
 4    * 求完后，直接枚举所有不在 MST 中的边，替换掉最大边权的边，更新答案
 5    * 点的编号从 0 开始
 6    */
 7   const int MAXN=110;
 8   const int INF=0x3f3f3f3f;
 9   bool vis[MAXN];
10   int lowc[MAXN];
11   int pre[MAXN];
12   int Max[MAXN][MAXN];//Max[i][j] 表示在最小生成树中从 i 到 j 的路径中的最
        大边权
13   bool used[MAXN][MAXN];
14   int Prim(int cost[][MAXN],int n){
15       int ans=0;
16       memset(vis,false,sizeof(vis));
17       memset(Max,0,sizeof(Max));
18       memset(used,false,sizeof(used));
19       vis[0]=true;
20       pre[0]=-1;
21       for(int i=1;i<n;i++){
22           lowc[i]=cost[0][i];
23           pre[i]=0;
24       }
25       lowc[0]=0;
26       for(int i=1;i<n;i++){
27           int minc=INF;
28           int p=-1;
29           for(int j=0;j<n;j++)
30               if(!vis[j]&&minc>lowc[j]){
31                   minc=lowc[j];
32                   p=j;


33                  }
34              if(minc==INF)return -1;
35              ans+=minc;
36              vis[p]=true;
37              used[p][pre[p]]=used[pre[p]][p]=true;
38              for(int j=0;j<n;j++){
39                  if(vis[j] && j != p)Max[j][p]=Max[p][j]=max(Max[j][pre[
                       p]],lowc[p]);
40                  if(!vis[j]&&lowc[j]>cost[p][j]){
41                      lowc[j]=cost[p][j];
42                      pre[j]=p;
43                  }
44              }
45         }
46         return ans;
47 }
\end{Code}

\subsection{Thành phần liên thông mạnh của đồ thị có hướng}

\subsubsection{Tarjan}

\begin{Code}
 1   /*
 2    * Tarjan 算法
 3    * 复杂度 O(N+M)
 4    */
 5   const int MAXN = 20010;//点数
 6   const int MAXM = 50010;//边数
 7   struct Edge{
 8       int to,next;
 9   }edge[MAXM];
10   int head[MAXN],tot;
11   int Low[MAXN],DFN[MAXN],Stack[MAXN],Belong[MAXN];//Belong 数组的值是
        1 ∼ scc
12   int Index,top;
13   int scc;//强连通分量的个数
14   bool Instack[MAXN];
15   int num[MAXN];//各个强连通分量包含点的个数，数组编号 1 ∼ scc
16   //num 数组不一定需要，结合实际情况
17
18 void addedge(int u,int v){
19     edge[tot].to = v;edge[tot].next = head[u];head[u] = tot++;
20 }
21 void Tarjan(int u){
22     int v;
23     Low[u] = DFN[u] = ++Index;
24     Stack[top++] = u;
25     Instack[u] = true;
26     for(int i = head[u];i != -1;i = edge[i].next){
27         v = edge[i].to;
28         if( !DFN[v] ){
29             Tarjan(v);
30             if( Low[u] > Low[v] )Low[u] = Low[v];


31              }
32              else if(Instack[v] && Low[u] > DFN[v])
33                  Low[u] = DFN[v];
34       }
35       if(Low[u] == DFN[u]){
36           scc++;
37           do{
38               v = Stack[--top];
39               Instack[v] = false;
40               Belong[v] = scc;
41               num[scc]++;
42           }
43           while( v != u);
44       }
45   }
46   void solve(int N){
47       memset(DFN,0,sizeof(DFN));
48       memset(Instack,false,sizeof(Instack));
49       memset(num,0,sizeof(num));
50       Index = scc = top = 0;
51       for(int i = 1;i <= N;i++)
52           if(!DFN[i])
53                Tarjan(i);
54   }
55   void init(){
56       tot = 0;
57       memset(head,-1,sizeof(head));
58   }
\end{Code}

\subsubsection{Kosaraju}

\begin{Code}
 1   /*
 2    * Kosaraju 算法，复杂度 O(N+M)
 3    */
 4   const int MAXN = 20010;
 5   const int MAXM = 50010;
 6   struct Edge{
 7       int to,next;
 8   }edge1[MAXM],edge2[MAXM];
 9   //edge1 是原图 G，edge2 是逆图 GT
10   int head1[MAXN],head2[MAXN];
11   bool mark1[MAXN],mark2[MAXN];
12   int tot1,tot2;
13   int cnt1,cnt2;
14   int st[MAXN];//对原图进行 dfs，点的结束时间从小到大排序
15   int Belong[MAXN];//每个点属于哪个连通分量 (0∼cnt2-1)
16   int num;//中间变量，用来数某个连通分量中点的个数
17   int setNum[MAXN];//强连通分量中点的个数，编号 0∼cnt2-1
18   void addedge(int u,int v){
19       edge1[tot1].to = v;edge1[tot1].next = head1[u];head1[u] = tot1
            ++;


20         edge2[tot2].to = u;edge2[tot2].next = head2[v];head2[v] = tot2
              ++;
21   }
22   void DFS1(int u){
23       mark1[u] = true;
24       for(int i = head1[u];i != -1;i = edge1[i].next)
25           if(!mark1[edge1[i].to])
26               DFS1(edge1[i].to);
27       st[cnt1++] = u;
28   }
29   void DFS2(int u){
30       mark2[u] = true;
31       num++;
32       Belong[u] = cnt2;
33       for(int i = head2[u];i != -1;i = edge2[i].next)
34           if(!mark2[edge2[i].to])
35               DFS2(edge2[i].to);
36   }
37   //点的编号从 1 开始
38   void solve(int n){
39       memset(mark1,false,sizeof(mark1));
40       memset(mark2,false,sizeof(mark2));
41       cnt1 = cnt2 = 0;
42       for(int i = 1;i <= n;i++)
43           if(!mark1[i])
44               DFS1(i);
45       for(int i = cnt1-1;i >= 0; i--)
46           if(!mark2[st[i]]){
47               num = 0;
48               DFS2(st[i]);
49               setNum[cnt2++] = num;
50           }
51   }
\end{Code}

\subsection{Khái niệm cơ bản về điểm cắt, cầu và thành phần song liên thông}

\paragraph{Độ liên thông đỉnh và độ liên thông cạnh.} Trong một đồ thị vô hướng liên thông, nếu xóa một tập đỉnh cùng tất cả các cạnh liên thuộc với các đỉnh đó làm đồ thị ban đầu tách thành nhiều thành phần liên thông, thì tập đỉnh đó được gọi là tập điểm cắt. Độ liên thông đỉnh của đồ thị là số đỉnh trong tập điểm cắt nhỏ nhất. Tương tự, nếu xóa một tập cạnh làm đồ thị tách thành nhiều thành phần liên thông, thì tập cạnh đó được gọi là tập cạnh cắt; độ liên thông cạnh là số cạnh trong tập cạnh cắt nhỏ nhất.

\paragraph{Đồ thị song liên thông, điểm cắt và cầu.} Nếu độ liên thông đỉnh của một đồ thị vô hướng liên thông lớn hơn 1, đồ thị được gọi là song liên thông đỉnh. Một đồ thị có điểm cắt khi và chỉ khi độ liên thông đỉnh bằng 1; phần tử duy nhất của tập điểm cắt được gọi là điểm cắt, hay articulation point. Nếu độ liên thông cạnh lớn hơn 1, đồ thị được gọi là song liên thông cạnh. Một đồ thị có cầu khi và chỉ khi độ liên thông cạnh bằng 1; phần tử duy nhất của tập cạnh cắt được gọi là cầu, hay articulation edge. Trong phần dưới, "song liên thông" có thể chỉ cả hai nghĩa này tùy ngữ cảnh.

\paragraph{Thành phần song liên thông.} Trong mọi đồ thị con G' của G, nếu G' song liên thông thì G' là một đồ thị con song liên thông. Nếu G' không là tập con thực sự của bất kỳ đồ thị con song liên thông nào khác, nó là đồ thị con song liên thông cực đại, còn gọi là thành phần song liên thông. Thành phần song liên thông đỉnh còn được gọi là block.

\paragraph{Tìm điểm cắt và cầu.} Thuật toán do R. Tarjan đề xuất. Khi duyệt DFS, đặt DFS(u) là thứ tự u được thăm trong cây tìm kiếm và Low(u) là đỉnh có thứ tự DFS nhỏ nhất mà u hoặc cây con của u có thể truy ngược tới qua cạnh không phải cha-con. Một đỉnh u là điểm cắt khi: (1) u là gốc và có nhiều hơn một cây con; hoặc (2) u không phải gốc và tồn tại cạnh cây (u,v) sao cho DFS(u) \textless{}= Low(v). Một cạnh vô hướng (u,v) là cầu khi và chỉ khi nó là cạnh cây và DFS(u) \textless{} Low(v).

\paragraph{Tìm thành phần song liên thông.} Với thành phần song liên thông đỉnh, có thể tìm trong quá trình tìm điểm cắt bằng cách dùng một stack lưu các cạnh hiện tại. Mỗi khi gặp cạnh cây hoặc cạnh ngược thì đưa cạnh vào stack. Khi DFS(u) \textless{}= Low(v), u là điểm cắt; lấy các cạnh từ đỉnh stack đến cạnh (u,v), các cạnh này cùng các đỉnh liên thuộc tạo thành một thành phần song liên thông đỉnh. Với thành phần song liên thông cạnh, chỉ cần tìm toàn bộ cầu, xóa chúng, rồi mỗi thành phần liên thông còn lại là một thành phần song liên thông cạnh.

\paragraph{Xây dựng đồ thị song liên thông.} Với một đồ thị liên thông có cầu, trước hết tìm và xóa toàn bộ cầu; mỗi thành phần còn lại là một đồ thị con song liên thông. Co mỗi thành phần thành một đỉnh rồi thêm lại các cầu, ta thu được một cây. Nếu số lá của cây là leaf, cần thêm ít nhất (leaf+1)/2 cạnh để đạt song liên thông cạnh.

\subsection{Điểm cắt và cầu}

\subsubsection{Mẫu}

\begin{Code}
 1   /*
 2   * 求无向图的割点和桥
 3   * 可以找出割点和桥，求删掉每个点后增加的连通块。
 4   * 需要注意重边的处理，可以先用矩阵存，再转邻接表，或者进行判重
 5   */
 6   const int MAXN = 10010;
 7   const int MAXM = 100010;
 8   struct Edge{
 9       int to,next;
10       bool cut;//是否为桥的标记
11   }edge[MAXM];
12   int head[MAXN],tot;
13   int Low[MAXN],DFN[MAXN],Stack[MAXN];
14   int Index,top;
15   bool Instack[MAXN];


16   bool cut[MAXN];
17   int add_block[MAXN];//删除一个点后增加的连通块
18   int bridge;
19   void addedge(int u,int v){
20       edge[tot].to = v;edge[tot].next = head[u];edge[tot].cut = false
            ;
21       head[u] = tot++;
22   }
23   void Tarjan(int u,int pre){
24       int v;
25       Low[u] = DFN[u] = ++Index;
26       Stack[top++] = u;
27       Instack[u] = true;
28       int son = 0;
29       int pre_cnt = 0; //处理重边，如果不需要可以去掉
30       for(int i = head[u];i != -1;i = edge[i].next){
31            v = edge[i].to;
32            if(v == pre && pre_cnt == 0){pre_cnt++;continue;}
33            if( !DFN[v] ){
34                son++;
35                Tarjan(v,u);
36                if(Low[u] > Low[v])Low[u] = Low[v];
37                //桥
38                //一条无向边 (u,v) 是桥，当且仅当 (u,v) 为树枝边，且满足
                     DFS(u)<Low(v)。
39                if(Low[v] > DFN[u]){
40                    bridge++;
41                    edge[i].cut = true;
42                    edge[i^1].cut = true;
43                }
44                //割点
45                //一个顶点 u 是割点，当且仅当满足 (1) 或 (2) (1) u 为树根，且
                     u 有多于一个子树。
46                //(2) u 不为树根，且满足存在 (u,v) 为树枝边 (或称父子边，
47                //即 u 为 v 在搜索树中的父亲)，使得 DFS(u)<=Low(v)
48                if(u != pre && Low[v] >= DFN[u]){//不是树根
49                    cut[u] = true;
50                    add_block[u]++;
51                }
52            }
53            else if( Low[u] > DFN[v])
54                 Low[u] = DFN[v];
55       }
56       //树根，分支数大于 1
57       if(u == pre && son > 1)cut[u] = true;
58       if(u == pre)add_block[u] = son - 1;
59       Instack[u] = false;
60       top--;
61   }
\end{Code}

\subsubsection{Cách gọi}

\begin{Code}
     1）UVA 796 Critical Links 给出一个无向图，按顺序输出桥
 1   void solve(int N){
 2       memset(DFN,0,sizeof(DFN));
 3       memset(Instack,false,sizeof(Instack));
 4       memset(add_block,0,sizeof(add_block));
 5       memset(cut,false,sizeof(cut));
 6       Index = top = 0;
 7       bridge = 0;
 8       for(int i = 1;i <= N;i++)
 9           if( !DFN[i] )
10                Tarjan(i,i);
11       printf("%d critical links\n",bridge);
12       vector<pair<int,int> >ans;
13       for(int u = 1;u <= N;u++)
14           for(int i = head[u];i != -1;i = edge[i].next)
15                if(edge[i].cut && edge[i].to > u)
16                {
17                    ans.push_back(make_pair(u,edge[i].to));
18                }
19       sort(ans.begin(),ans.end());
20       //按顺序输出桥
21       for(int i = 0;i < ans.size();i++)
22           printf("%d - %d\n",ans[i].first-1,ans[i].second-1);
23       printf("\n");
24   }
25   void init(){
26       tot = 0;
27       memset(head,-1,sizeof(head));
28   }
29   //处理重边
30   map<int,int>mapit;
31   inline bool isHash(int u,int v){
32       if(mapit[u*MAXN+v])return true;
33       if(mapit[v*MAXN+u])return true;
34       mapit[u*MAXN+v] = mapit[v*MAXN+u] = 1;
35       return false;
36   }
37   int main(){
38       int n;
39       while(scanf("%d",&n) == 1){
40           init();
41           int u;
42           int k;
43           int v;
44           //mapit.clear();
45           for(int i = 1;i <= n;i++){
46                scanf("%d (%d)",&u,&k);
47                u++;
48                //这样加边，要保证正边和反边是相邻的，建无向图


49                  while(k--){
50                      scanf("%d",&v);
51                      v++;
52                      if(v <= u)continue;
53                      //if(isHash(u,v))continue;
54                      addedge(u,v);
55                      addedge(v,u);
56                  }
57              }
58              solve(n);
59       }
60       return 0;
61 }
     2）POJ 2117 求删除一个点后，图中最多有多少个连通块
 1   void solve(int N){
 2       memset(DFN,0,sizeof(DFN));
 3       memset(Instack,0,sizeof(Instack));
 4       memset(add_block,0,sizeof(add_block));
 5       memset(cut,false,sizeof(cut));
 6       Index = top = 0;
 7       int cnt = 0;//原来的连通块数
 8       for(int i = 1;i <= N;i++)
 9          if( !DFN[i] ){
10              Tarjan(i,i);//找割点调用必须是 Tarjan(i,i)
11              cnt++;
12          }
13       int ans = 0;
14       for(int i = 1;i <= N;i++)
15          ans = max(ans,cnt+add_block[i]);
16       printf("%d\n",ans);
17   }
18   void init(){
19       tot = 0;
20       memset(head,-1,sizeof(head));
21   }
22   int main(){
23       int n,m;
24       int u,v;
25       while(scanf("%d%d",&n,&m)==2){
26           if(n==0 && m == 0)break;
27           init();
28           while(m--){
29                scanf("%d%d",&u,&v);
30                u++;v++;
31                addedge(u,v);
32                addedge(v,u);
33           }
34           solve(n);
35       }
36       return 0;
37   }
\end{Code}

\subsection{Thành phần song liên thông cạnh}

\paragraph{Ghi chú.} Sau khi bỏ các cầu, các thành phần liên thông còn lại chính là thành phần song liên thông cạnh. Với một đồ thị liên thông có cầu, co mỗi đồ thị con song liên thông thành một đỉnh sẽ tạo thành một cây; số cạnh cần thêm là (leaf+1)/2, trong đó leaf là số lá. POJ 3177: cho một đồ thị vô hướng liên thông G, hỏi cần thêm ít nhất bao nhiêu cạnh để nó trở thành song liên thông.

\begin{Code}
1    const int MAXN = 5010;//点数
2    const int MAXM = 20010;//边数，因为是无向图，所以这个值要 *2
3    struct Edge{
4        int to,next;
5        bool cut;//是否是桥标记
6    }edge[MAXM];
7    int head[MAXN],tot;
8    int Low[MAXN],DFN[MAXN],Stack[MAXN],Belong[MAXN];//Belong 数组的值是
        1 ∼ block
 9   int Index,top;
10   int block;//边双连通块数
11   bool Instack[MAXN];
12   int bridge;//桥的数目
13   void addedge(int u,int v){
14       edge[tot].to = v;edge[tot].next = head[u];edge[tot].cut=false;
15       head[u] = tot++;
16   }
17   void Tarjan(int u,int pre){
18       int v;
19       Low[u] = DFN[u] = ++Index;
20       Stack[top++] = u;
21       Instack[u] = true;
22       int pre_cnt = 0;
23       for(int i = head[u];i != -1;i = edge[i].next){
24           v = edge[i].to;
25           if(v == pre && pre_cnt == 0){pre_cnt++; continue;}
26           if( !DFN[v] ){
27                Tarjan(v,u);
28                if( Low[u] > Low[v] )Low[u] = Low[v];
29                if(Low[v] > DFN[u]){
30                    bridge++;
31                    edge[i].cut = true;
32                    edge[i^1].cut = true;
33                }
34           }
35           else if( Instack[v] && Low[u] > DFN[v] )
36                Low[u] = DFN[v];
37       }
38       if(Low[u] == DFN[u]){
39           block++;
40           do{
41                v = Stack[--top];
42                Instack[v] = false;
43                Belong[v] = block;
44           }


45              while( v!=u );
46         }
47   }
48   void init(){
49       tot = 0;
50       memset(head,-1,sizeof(head));
51   }
52   int du[MAXN];//缩点后形成树，每个点的度数
53   void solve(int n){
54       memset(DFN,0,sizeof(DFN));
55       memset(Instack,false,sizeof(Instack));
56       Index = top = block = 0;
57       Tarjan(1,0);
58       int ans = 0;
59       memset(du,0,sizeof(du));
60       for(int i = 1;i <= n;i++)
61          for(int j = head[i];j != -1;j = edge[j].next)
62             if(edge[j].cut)
63                 du[Belong[i]]++;
64       for(int i = 1;i <= block;i++)
65          if(du[i]==1)
66             ans++;
67       //找叶子结点的个数 ans, 构造边双连通图需要加边 (ans+1)/2
68       printf("%d\n",(ans+1)/2);
69   }
70   int main(){
71       int n,m;
72       int u,v;
73       while(scanf("%d%d",&n,&m)==2){
74           init();
75           while(m--){
76                scanf("%d%d",&u,&v);
77                addedge(u,v);
78                addedge(v,u);
79           }
80           solve(n);
81       }
82       return 0;
83   }
\end{Code}

\subsection{Thành phần song liên thông đỉnh}

\paragraph{Ghi chú.} Với thành phần song liên thông đỉnh, thực tế có thể tìm được từng thành phần ngay trong quá trình tìm điểm cắt. Dùng một stack lưu thành phần hiện tại; khi duyệt đồ thị, mỗi cạnh cây hoặc cạnh ngược được đưa vào stack. Khi DFS(u) \textless{}= Low(v), u là một điểm cắt; lấy các cạnh từ đỉnh stack tới cạnh (u,v), các cạnh đó cùng các đỉnh liên thuộc tạo thành một thành phần song liên thông đỉnh. Điểm cắt có thể thuộc nhiều thành phần, còn các đỉnh khác và mỗi cạnh chỉ thuộc đúng một thành phần. POJ 2942 dùng chu trình lẻ, tô màu kiểm tra đồ thị hai phía và thành phần song liên thông đỉnh.

\begin{Code}
1 /*
2 POJ 2942 Knights of the Round Table


 3 亚瑟王要在圆桌上召开骑士会议，为了不引发骑士之间的冲突，
 4 并且能够让会议的议题有令人满意的结果，每次开会前都必须对出席会议的骑士有如下要
      求：
 5 1、相互憎恨的两个骑士不能坐在直接相邻的 2 个位置；
 6 2、出席会议的骑士数必须是奇数，这是为了让投票表决议题时都能有结果。
 7
 8 注意：1、所给出的憎恨关系一定是双向的，不存在单向憎恨关系。
 9 2、由于是圆桌会议，则每个出席的骑士身边必定刚好有 2 个骑士。
10 即每个骑士的座位两边都必定各有一个骑士。
11 3、一个骑士无法开会，就是说至少有 3 个骑士才可能开会。
12
13 首先根据给出的互相憎恨的图中得到补图。
14 然后就相当于找出不能形成奇圈的点。
15 利用下面两个定理：
16 （1）如果一个双连通分量内的某些顶点在一个奇圈中（即双连通分量含有奇圈），
17 那么这个双连通分量的其他顶点也在某个奇圈中；
18 （2）如果一个双连通分量含有奇圈，则他必定不是一个二分图。反过来也成立，这是一个
      充要条件。
19
20 所以本题的做法，就是对补图求点双连通分量。
21 然后对于求得的点双连通分量，使用染色法判断是不是二分图，不是二分图，这个双连通分
      量的点是可以存在的
22 */
23
24 const int MAXN = 1010;
25 const int MAXM = 2000010;
26 struct Edge{
27     int to,next;
28 }edge[MAXM];
29 int head[MAXN],tot;
30 int Low[MAXN],DFN[MAXN],Stack[MAXN],Belong[MAXN];
31 int Index,top;
32 int block;//点双连通分量的个数
33 bool Instack[MAXN];
34 bool can[MAXN];
35 bool ok[MAXN];//标记
36 int tmp[MAXN];//暂时存储双连通分量中的点
37 int cc;//tmp 的计数
38 int color[MAXN];//染色
39 void addedge(int u,int v){
40     edge[tot].to = v;edge[tot].next = head[u];head[u] = tot++;
41 }
42 //染色判断二分图
43 bool dfs(int u,int col){
44     color[u] = col;
45     for(int i = head[u];i != -1;i = edge[i].next){
46         int v = edge[i].to;
47         if( !ok[v] )continue;
48         if(color[v] != -1){
49              if(color[v]==col)return false;
50              continue;


 51              }
 52              if(!dfs(v,!col))return false;
 53       }
 54       return true;
 55   }
 56   void Tarjan(int u,int pre){
 57       int v;
 58       Low[u] = DFN[u] = ++Index;
 59       Stack[top++] = u;
 60       Instack[u] = true;
 61       int pre_cnt = 0;
 62       for(int i = head[u];i != -1;i = edge[i].next){
 63           v = edge[i].to;
 64           if(v == pre && pre_cnt == 0){pre_cnt++; continue;}
 65           if( !DFN[v] ){
 66                Tarjan(v,u);
 67                if(Low[u] > Low[v])Low[u] = Low[v];
 68                if( Low[v] >= DFN[u]){
 69                    block++;
 70                    int vn;
 71                    cc = 0;
 72                    memset(ok,false,sizeof(ok));
 73                    do{
 74                        vn = Stack[--top];
 75                        Belong[vn] = block;
 76                        Instack[vn] = false;
 77                        ok[vn] = true;
 78                        tmp[cc++] = vn;
 79                    }
 80                    while( vn!=v );
 81                    ok[u] = 1;
 82                    memset(color,-1,sizeof(color));
 83                    if( !dfs(u,0) ){
 84                        can[u] = true;
 85                        while(cc--)can[tmp[cc]]=true;
 86                    }
 87                }
 88           }
 89           else if(Instack[v] && Low[u] > DFN[v])
 90               Low[u] = DFN[v];
 91       }
 92   }
 93   void solve(int n){
 94       memset(DFN,0,sizeof(DFN));
 95       memset(Instack,false,sizeof(Instack));
 96       Index = block = top = 0;
 97       memset(can,false,sizeof(can));
 98       for(int i = 1;i <= n;i++)
 99          if(!DFN[i])
100              Tarjan(i,-1);
101       int ans = n;


102         for(int i = 1;i <= n;i++)
103             if(can[i])
104               ans--;
105         printf("%d\n",ans);
106   }
107   void init(){
108       tot = 0;
109       memset(head,-1,sizeof(head));
110   }
111   int g[MAXN][MAXN];
112   int main(){
113       int n,m;
114       int u,v;
115       while(scanf("%d%d",&n,&m)==2){
116           if(n==0 && m==0)break;
117           init();
118           memset(g,0,sizeof(g));
119           while(m--){
120                scanf("%d%d",&u,&v);
121                g[u][v]=g[v][u]=1;
122           }
123           for(int i = 1;i <= n;i++)
124             for(int j = 1;j <= n;j++)
125                 if(i != j && g[i][j]==0)
126                    addedge(i,j);
127           solve(n);
128       }
129       return 0;
130   }
\end{Code}

\subsection{Cây phân nhánh nhỏ nhất}

\begin{Code}
 1    /*
 2     * 最小树形图
 3     * int 型
 4     * 复杂度 O(NM)
 5     * 点从 0 开始
 6     */
 7    const int INF = 0x3f3f3f3f;
 8    const int MAXN = 1010;
 9    const int MAXM = 40010;
10    struct Edge{
11        int u,v,cost;
12    };
13    Edge edge[MAXM];
14    int pre[MAXN],id[MAXN],visit[MAXN],in[MAXN];
15    int zhuliu(int root,int n,int m,Edge edge[]){
16        int res = 0,u,v;
17        while(1){
18            for(int i = 0;i < n;i++)
19                 in[i] = INF;


20              for(int i = 0;i < m;i++)
21                  if(edge[i].u != edge[i].v && edge[i].cost < in[edge[i].
                       v]){
22                      pre[edge[i].v] = edge[i].u;
23                      in[edge[i].v] = edge[i].cost;
24                  }
25              for(int i = 0;i < n;i++)
26                  if(i != root && in[i] == INF)
27                      return -1;//不存在最小树形图
28              int tn = 0;
29              memset(id,-1,sizeof(id));
30              memset(visit,-1,sizeof(visit));
31              in[root] = 0;
32              for(int i = 0;i < n;i++){
33                  res += in[i];
34                  v = i;
35                  while( visit[v] != i && id[v] == -1 && v != root){
36                      visit[v] = i;
37                      v = pre[v];
38                  }
39                  if( v != root && id[v] == -1 ){
40                      for(int u = pre[v]; u != v ;u = pre[u])
41                          id[u] = tn;
42                      id[v] = tn++;
43                  }
44              }
45              if(tn == 0)break;//没有有向环
46              for(int i = 0;i < n;i++)
47                  if(id[i] == -1)
48                      id[i] = tn++;
49              for(int i = 0;i < m;){
50                  v = edge[i].v;
51                  edge[i].u = id[edge[i].u];
52                  edge[i].v = id[edge[i].v];
53                  if(edge[i].u != edge[i].v)
54                      edge[i++].cost -= in[v];
55                  else
56                      swap(edge[i],edge[--m]);
57              }
58              n = tn;
59              root = id[root];
60     }
61     return res;
62 }
63 int g[MAXN][MAXN];
64 int main(){
65     int n,m;
66     int iCase = 0;
67     int T;
68     scanf("%d",&T);
69     while( T-- ){


70              iCase ++;
71              scanf("%d%d",&n,&m);
72              for(int i = 0;i < n;i++)
73                  for(int j = 0;j < n;j++)
74                      g[i][j] = INF;
75              int u,v,cost;
76              while(m--){
77                  scanf("%d%d%d",&u,&v,&cost);
78                  if(u == v)continue;
79                  g[u][v] = min(g[u][v],cost);
80              }
81              int L = 0;
82              for(int i = 0;i < n;i++)
83                  for(int j = 0;j < n;j++)
84                      if(g[i][j] < INF){
85                          edge[L].u = i;
86                          edge[L].v = j;
87                          edge[L++].cost = g[i][j];
88                      }
89              int ans = zhuliu(0,n,L,edge);
90              printf("Case #%d: ",iCase);
91              if(ans == -1)printf("Possums!\n");
92              else printf("%d\n",ans);
93       }
94       return 0;
95 }
\end{Code}

\subsection{Ghép cặp đồ thị hai phía}

\subsubsection{Ma trận kề (thuật toán Hungary)}

\begin{Code}
 1   /* ***********************************************************
 2   //二分图匹配（匈牙利算法的 DFS 实现）(邻接矩阵形式)
 3   //初始化：g[][] 两边顶点的划分情况
 4   //建立 g[i][j] 表示 i->j 的有向边就可以了，是左边向右边的匹配
 5   //g 没有边相连则初始化为 0
 6   //uN 是匹配左边的顶点数，vN 是匹配右边的顶点数
 7   //调用：res=hungary(); 输出最大匹配数
 8   //优点：适用于稠密图，DFS 找增广路，实现简洁易于理解
 9   //时间复杂度:O(VE)
10   //*************************************************************/
11   //顶点编号从 0 开始的
12   const int MAXN = 510;
13   int uN,vN;//u,v 的数目，使用前面必须赋值
14   int g[MAXN][MAXN];//邻接矩阵
15   int linker[MAXN];
16   bool used[MAXN];
17   bool dfs(int u){
18       for(int v = 0; v < vN;v++)
19           if(g[u][v] && !used[v]){
20               used[v] = true;
21               if(linker[v] == -1 || dfs(linker[v])){


22                  linker[v] = u;
23                  return true;
24             }
25         }
26     return false;
27 }
28 int hungary(){
29     int res = 0;
30     memset(linker,-1,sizeof(linker));
31     for(int u = 0;u < uN;u++){
32         memset(used,false,sizeof(used));
33         if(dfs(u))res++;
34     }
35     return res;
36 }
\end{Code}

\subsubsection{Danh sách kề (thuật toán Hungary)}

\begin{Code}
 1   /*
 2    * 匈牙利算法邻接表形式
 3    * 使用前用 init() 进行初始化，给 uN 赋值
 4    * 加边使用函数 addedge(u,v)
 5    *
 6    */
 7   const int MAXN = 5010;//点数的最大值
 8   const int MAXM = 50010;//边数的最大值
 9   struct Edge{
10       int to,next;
11   }edge[MAXM];
12   int head[MAXN],tot;
13   void init(){
14       tot = 0;
15       memset(head,-1,sizeof(head));
16   }
17   void addedge(int u,int v){
18       edge[tot].to = v; edge[tot].next = head[u];
19       head[u] = tot++;
20   }
21   int linker[MAXN];
22   bool used[MAXN];
23   int uN;
24   bool dfs(int u){
25       for(int i = head[u]; i != -1 ;i = edge[i].next){
26           int v = edge[i].to;
27           if(!used[v]){
28                used[v] = true;
29                if(linker[v] == -1 || dfs(linker[v])){
30                    linker[v] = u;
31                    return true;
32                }
33           }
34       }


35     return false;
36 }
37 int hungary(){
38     int res = 0;
39     memset(linker,-1,sizeof(linker));
40     //点的编号 0∼uN-1
41     for(int u = 0; u < uN;u++){
42         memset(used,false,sizeof(used));
43         if(dfs(u))res++;
44     }
45     return res;
46 }
\end{Code}

\subsubsection{Thuật toán Hopcroft-Karp}

\begin{Code}
 1   /* *******************************
 2    * 二分图匹配（Hopcroft-Karp 算法）
 3    * 复杂度 O(sqrt(n)*E)
 4    * 邻接表存图，vector 实现
 5    * vector 先初始化，然后假如边
 6    * uN 为左端的顶点数，使用前赋值 (点编号 0 开始)
 7    */
 8   const int MAXN = 3000;
 9   const int INF = 0x3f3f3f3f;
10   vector<int>G[MAXN];
11   int uN;
12   int Mx[MAXN],My[MAXN];
13   int dx[MAXN],dy[MAXN];
14   int dis;
15   bool used[MAXN];
16   bool SearchP(){
17       queue<int>Q;
18       dis = INF;
19       memset(dx,-1,sizeof(dx));
20       memset(dy,-1,sizeof(dy));
21       for(int i = 0 ; i < uN; i++)
22           if(Mx[i] == -1){
23               Q.push(i);
24               dx[i] = 0;
25           }
26       while(!Q.empty()){
27           int u = Q.front();
28           Q.pop();
29           if(dx[u] > dis)break;
30           int sz = G[u].size();
31           for(int i = 0;i < sz;i++){
32               int v = G[u][i];
33               if(dy[v] == -1){
34                    dy[v] = dx[u] + 1;
35                    if(My[v] == -1)dis = dy[v];
36                    else{
37                        dx[My[v]] = dy[v] + 1;


38                       Q.push(My[v]);
39                   }
40               }
41           }
42       }
43       return dis != INF;
44   }
45   bool DFS(int u){
46       int sz = G[u].size();
47       for(int i = 0;i < sz;i++){
48           int v = G[u][i];
49           if(!used[v] && dy[v] == dx[u] + 1){
50               used[v] = true;
51               if(My[v] != -1 && dy[v] == dis)continue;
52               if(My[v] == -1 || DFS(My[v])){
53                    My[v] = u;
54                    Mx[u] = v;
55                    return true;
56               }
57           }
58       }
59       return false;
60   }
61   int MaxMatch(){
62       int res = 0;
63       memset(Mx,-1,sizeof(Mx));
64       memset(My,-1,sizeof(My));
65       while(SearchP()){
66           memset(used,false,sizeof(used));
67           for(int i = 0;i < uN;i++)
68               if(Mx[i] == -1 && DFS(i))
69                    res++;
70       }
71       return res;
72   }
\end{Code}

\subsection{Đa ghép cặp trên đồ thị hai phía}

\begin{Code}
 1   const int MAXN = 1010;
 2   const int MAXM = 510;
 3   int uN,vN;
 4   int g[MAXN][MAXM];
 5   int linker[MAXM][MAXN];
 6   bool used[MAXM];
 7   int num[MAXM];//右边最大的匹配数
 8   bool dfs(int u){
 9       for(int v = 0; v < vN;v++)
10           if(g[u][v] && !used[v]){
11               used[v] = true;
12               if(linker[v][0] < num[v]){
13                    linker[v][++linker[v][0]] = u;
14                    return true;


15             }
16             for(int i = 1;i <= num[v];i++)
17                  if(dfs(linker[v][i])){
18                      linker[v][i] = u;
19                      return true;
20                  }
21         }
22     return false;
23 }
24 int hungary(){
25     int res = 0;
26     for(int i = 0;i < vN;i++)
27         linker[i][0] = 0;
28     for(int u = 0; u < uN; u++){
29         memset(used,false,sizeof(used));
30         if(dfs(u))res++;
31     }
32     return res;
33 }
\end{Code}

\subsection{Ghép cặp trọng số lớn nhất trên đồ thị hai phía (thuật toán KM)}

\begin{Code}
 1   /* KM 算法
 2    *    复杂度 O(nx*nx*ny)
 3    * 求最大权匹配
 4    *    若求最小权匹配，可将权值取相反数，结果取相反数
 5    * 点的编号从 0 开始
 6    */
 7   const int N = 310;
 8   const int INF = 0x3f3f3f3f;
 9   int nx,ny;//两边的点数
10   int g[N][N];//二分图描述
11   int linker[N],lx[N],ly[N];//y 中各点匹配状态，x,y 中的点标号
12   int slack[N];
13   bool visx[N],visy[N];
14   bool DFS(int x){
15       visx[x] = true;
16       for(int y = 0; y < ny; y++){
17           if(visy[y])continue;
18           int tmp = lx[x] + ly[y] - g[x][y];
19           if(tmp == 0){
20               visy[y] = true;
21               if(linker[y] == -1 || DFS(linker[y])){
22                    linker[y] = x;
23                    return true;
24               }
25           }
26           else if(slack[y] > tmp)
27               slack[y] = tmp;
28       }
29       return false;
30   }


31   int KM(){
32       memset(linker,-1,sizeof(linker));
33       memset(ly,0,sizeof(ly));
34       for(int i = 0;i < nx;i++){
35           lx[i] = -INF;
36           for(int j = 0;j < ny;j++)
37               if(g[i][j] > lx[i])
38                    lx[i] = g[i][j];
39       }
40       for(int x = 0;x < nx;x++){
41           for(int i = 0;i < ny;i++)
42               slack[i] = INF;
43           while(true){
44               memset(visx,false,sizeof(visx));
45               memset(visy,false,sizeof(visy));
46               if(DFS(x))break;
47               int d = INF;
48               for(int i = 0;i < ny;i++)
49                    if(!visy[i] && d > slack[i])
50                        d = slack[i];
51               for(int i = 0;i < nx;i++)
52                    if(visx[i])
53                        lx[i] -= d;
54               for(int i = 0;i < ny;i++){
55                    if(visy[i])ly[i] += d;
56                    else slack[i] -= d;
57               }
58           }
59       }
60       int res = 0;
61       for(int i = 0;i < ny;i++)
62           if(linker[i] != -1)
63               res += g[linker[i]][i];
64       return res;
65   }
66   //HDU 2255
67   int main(){
68       int n;
69       while(scanf("%d",&n) == 1){
70           for(int i = 0;i < n;i++)
71               for(int j = 0;j < n;j++)
72                    scanf("%d",&g[i][j]);
73           nx = ny = n;
74           printf("%d\n",KM());
75       }
76       return 0;
77   }
\end{Code}

\subsection{Ghép cặp đồ thị tổng quát bằng cây hoa}

\begin{Code}
     URAL 1099
1 const int MAXN = 250;


 2   int N; //点的个数，点的编号从 1 到 N
 3   bool Graph[MAXN][MAXN];
 4   int Match[MAXN];
 5   bool InQueue[MAXN],InPath[MAXN],InBlossom[MAXN];
 6   int Head,Tail;
 7   int Queue[MAXN];
 8   int Start,Finish;
 9   int NewBase;
10   int Father[MAXN],Base[MAXN];
11   int Count;//匹配数，匹配对数是 Count/2
12   void CreateGraph(){
13       int u,v;
14       memset(Graph,false,sizeof(Graph));
15       scanf("%d",&N);
16       while(scanf("%d%d",&u,&v) == 2){
17           Graph[u][v] = Graph[v][u] = true;
18       }
19   }
20   void Push(int u){
21       Queue[Tail] = u;
22       Tail++;
23       InQueue[u] = true;
24   }
25   int Pop(){
26       int res = Queue[Head];
27       Head++;
28       return res;
29   }
30   int FindCommonAncestor(int u,int v){
31       memset(InPath,false,sizeof(InPath));
32       while(true){
33           u = Base[u];
34           InPath[u] = true;
35           if(u == Start) break;
36           u = Father[Match[u]];
37       }
38       while(true){
39           v = Base[v];
40           if(InPath[v])break;
41           v = Father[Match[v]];
42       }
43       return v;
44   }
45   void ResetTrace(int u){
46       int v;
47       while(Base[u] != NewBase){
48           v = Match[u];
49           InBlossom[Base[u]] = InBlossom[Base[v]] = true;
50           u = Father[v];
51           if(Base[u] != NewBase) Father[u] = v;
52       }


53    }
54    void BloosomContract(int u,int v){
55        NewBase = FindCommonAncestor(u,v);
56        memset(InBlossom,false,sizeof(InBlossom));
57        ResetTrace(u);
58        ResetTrace(v);
59        if(Base[u] != NewBase) Father[u] = v;
60        if(Base[v] != NewBase) Father[v] = u;
61        for(int tu = 1; tu <= N; tu++)
62            if(InBlossom[Base[tu]]){
63                Base[tu] = NewBase;
64                if(!InQueue[tu]) Push(tu);
65            }
66    }
67    void FindAugmentingPath(){
68        memset(InQueue,false,sizeof(InQueue));
69        memset(Father,0,sizeof(Father));
70        for(int i = 1;i <= N;i++)
71            Base[i] = i;
72        Head = Tail = 1;
73        Push(Start);
74        Finish = 0;
75        while(Head < Tail){
76            int u = Pop();
77            for(int v = 1; v <= N; v++)
78                if(Graph[u][v] && (Base[u] != Base[v]) && (Match[u] !=
                     v)){
79                     if((v == Start) || ((Match[v] > 0) && Father[Match[
                          v]] > 0))
 80                        BloosomContract(u,v);
 81                    else if(Father[v] == 0){
 82                        Father[v] = u;
 83                        if(Match[v] > 0)
 84                            Push(Match[v]);
 85                        else{
 86                            Finish = v;
 87                            return;
 88                        }
 89                    }
 90               }
 91       }
 92   }
 93   void AugmentPath(){
 94       int u,v,w;
 95       u = Finish;
 96       while(u > 0){
 97           v = Father[u];
 98           w = Match[v];
 99           Match[v] = u;
100           Match[u] = v;
101           u = w;


102       }
103   }
104   void Edmonds(){
105       memset(Match,0,sizeof(Match));
106       for(int u = 1; u <= N; u++)
107           if(Match[u] == 0){
108               Start = u;
109               FindAugmentingPath();
110               if(Finish > 0)AugmentPath();
111           }
112   }
113   void PrintMatch(){
114       Count = 0;
115       for(int u = 1; u <= N;u++)
116           if(Match[u] > 0)
117               Count++;
118       printf("%d\n",Count);
119       for(int u = 1; u <= N; u++)
120           if(u < Match[u])
121               printf("%d %d\n",u,Match[u]);
122   }
123   int main(){
124       CreateGraph();//建图
125       Edmonds();//进行匹配
126       PrintMatch();//输出匹配数和匹配
127       return 0;
128   }
\end{Code}

\subsection{Ghép cặp có trọng số lớn nhất trên đồ thị tổng quát}

\begin{Code}
 1    //一般图的最大加权匹配模板
 2    //注意 G 的初始化, 需要偶数个点，刚好可以形成 n/2 个匹配
 3    //如果要求最小权匹配，可以取相反数，或者稍加修改就可以了
 4    //点的编号从 0 开始的
 5    const int MAXN = 110;
 6    const int INF = 0x3f3f3f3f;
 7    int G[MAXN][MAXN];
 8    int cnt_node;//点的个数
 9    int dis[MAXN];
10    int match[MAXN];
11    bool vis[MAXN];
12    int sta[MAXN],top;//堆栈
13    bool dfs(int u){
14        sta[top++] = u;
15        if(vis[u])return true;
16        vis[u] = true;
17        for(int i = 0;i < cnt_node;i++)
18            if(i != u && i != match[u] && !vis[i]){
19                int t = match[i];
20                if(dis[t] < dis[u] + G[u][i] - G[i][t]){
21                     dis[t] = dis[u] + G[u][i] - G[i][t];


22                   if(dfs(t))return true;
23               }
24           }
25       top--;
26       vis[u] = false;
27       return false;
28   }
29   int P[MAXN];
30   //返回最大匹配权值
31   int get_Match(int N){
32       cnt_node = N;
33       for(int i = 0;i < cnt_node;i++)P[i] = i;
34       for(int i = 0;i < cnt_node;i += 2){
35           match[i] = i+1;
36           match[i+1] = i;
37       }
38       int cnt = 0;
39       while(1){
40           memset(dis,0,sizeof(dis));
41           memset(vis,false,sizeof(vis));
42           top = 0;
43           bool update = false;
44           for(int i = 0;i < cnt_node;i++)
45                if(dfs(P[i])){
46                    update = true;
47                    int tmp = match[sta[top-1]];
48                    int j = top-2;
49                    while(sta[j] != sta[top-1]){
50                        match[tmp] = sta[j];
51                        swap(tmp,match[sta[j]]);
52                        j--;
53                    }
54                    match[tmp] = sta[j];
55                    match[sta[j]] = tmp;
56                    break;
57                }
58           if(!update){
59                cnt++;
60                if(cnt >= 3)break;
61                random_shuffle(P,P+cnt_node);
62           }
63       }
64       int ans = 0;
65       for(int i = 0;i < cnt_node;i++){
66           int v = match[i];
67           if(i < v)ans += G[i][v];
68       }
69       return ans;
70   }
\end{Code}

\subsection{Đếm cây khung}

\begin{Code}
     Matrix-Tree 定理 (Kirchhoff 矩阵 -树定理)
     1、G 的度数矩阵 D[G] 是一个 n*n 的矩阵，并且满足：当 i�j 时,dij=0；当 i=j 时，dij 等于
     vi 的度数。
     2、G 的邻接矩阵 A[G] 也是一个 n*n 的矩阵，并且满足：如果 vi、vj 之间有边直接相连，
     则 aij=1，否则为 0。
     我们定义 G 的 Kirchhoff 矩阵 (也称为拉普拉斯算子)C[G] 为 C[G]=D[G]-A[G]，则 Matrix-
     Tree 定理可以描述为：
     G 的所有不同的生成树的个数等于其 Kirchhoff 矩阵 C[G] 任何一个 n-1 阶主子式的行列式的
     绝对值。
     所谓 n-1 阶主子式，就是对于 r(1�r�n)，将 C[G] 的第 r 行、第 r 列同时去掉后得到的新矩阵，
     用 Cr[G] 表示。

     //HDU 4305
     //求生成树计数部分代码，计数对 10007 取模
 1 const int MOD = 10007;
 2 int INV[MOD];
 3 //求 ax = 1( mod m) 的 x 值，就是逆元 (0<a<m)
 4 long long inv(long long a,long long m){
 5     if(a == 1)return 1;
 6     return inv(m%a,m)*(m-m/a)%m;
 7 }
 8 struct Matrix{
 9     int mat[330][330];
10     void init(){
11         memset(mat,0,sizeof(mat));
12     }
13     //求行列式的值模上，需要使用逆元MOD
14     int det(int n){
15         for(int i = 0;i < n;i++)
16             for(int j = 0;j < n;j++)
17                 mat[i][j] = (mat[i][j]%MOD+MOD)%MOD;
18         int res = 1;
19         for(int i = 0;i < n;i++){
20             for(int j = i;j < n;j++)
21                 if(mat[j][i]!=0){
22                      for(int k = i;k < n;k++)
23                          swap(mat[i][k],mat[j][k]);
24                      if(i != j)
25                          res = (-res+MOD)%MOD;
26                      break;
27                 }
28             if(mat[i][i] == 0){
29                 res = -1;//不存在也就是行列式值为(0)
30                 break;
31             }
32             for(int j = i+1;j < n;j++){
33                 //int mut = (mat[j][i]*INV[mat[i][i]])%MOD打表逆
                      元;//
34                 int mut = (mat[j][i]*inv(mat[i][i],MOD))%MOD;


35                     for(int k = i;k < n;k++)
36                         mat[j][k] = (mat[j][k]-(mat[i][k]*mut)%MOD+MOD)
                              %MOD;
37                  }
38                  res = (res * mat[i][i])%MOD;
39              }
40              return res;
41       }
42 };
43
44              Matrix ret;
45              ret.init();
46              for(int i = 0;i < n;i++)
47                  for(int j = 0;j < n;j++)
48                      if(i != j && g[i][j]){
49                          ret.mat[i][j] = -1;
50                          ret.mat[i][i]++;
51                      }
52              printf("%d\n",ret.det(n-1));
     计算生成树个数，不取模，SPOJ 104
 1   const double eps = 1e-8;
 2   const int MAXN = 110;
 3   int sgn(double x){
 4       if(fabs(x) < eps)return 0;
 5       if(x < 0)return -1;
 6       else return 1;
 7   }
 8   double b[MAXN][MAXN];
 9   double det(double a[][MAXN],int n){
10       int i, j, k, sign = 0;
11       double ret = 1;
12       for(i = 0;i < n;i++)
13           for(j = 0;j < n;j++)
14               b[i][j] = a[i][j];
15       for(i = 0;i < n;i++){
16           if(sgn(b[i][i]) == 0){
17               for(j = i + 1; j < n;j++)
18                   if(sgn(b[j][i]) != 0)
19                       break;
20               if(j == n)return 0;
21               for(k = i;k < n;k++)
22                   swap(b[i][k],b[j][k]);
23               sign++;
24           }
25           ret *= b[i][i];
26           for(k = i + 1;k < n;k++)
27               b[i][k]/=b[i][i];
28           for(j = i+1;j < n;j++)
29               for(k = i+1;k < n;k++)
30                   b[j][k] -= b[j][i]*b[i][k];
31       }


32       if(sign & 1)ret = -ret;
33       return ret;
34   }
35   double a[MAXN][MAXN];
36   int g[MAXN][MAXN];
37   int main(){
38       int T;
39       int n,m;
40       int u,v;
41       scanf("%d",&T);
42       while(T--){
43           scanf("%d%d",&n,&m);
44           memset(g,0,sizeof(g));
45           while(m--){
46                scanf("%d%d",&u,&v);
47                u--;v--;
48                g[u][v] = g[v][u] = 1;
49           }
50           memset(a,0,sizeof(a));
51           for(int i = 0;i < n;i++)
52                for(int j = 0;j < n;j++)
53                    if(i != j && g[i][j]){
54                        a[i][i]++;
55                        a[i][j] = -1;
56                    }
57           double ans = det(a,n-1);
58           printf("%.0lf\n",ans);
59       }
60       return 0;
61   }
\end{Code}

\subsection{Luồng cực đại}

\subsubsection{SAP dạng ma trận kề}

\begin{Code}
 1   /*
 2    * SAP 算法（矩阵形式）
 3    * 结点编号从 0 开始
 4    */
 5   const int MAXN=1100;
 6   int maze[MAXN][MAXN];
 7   int gap[MAXN],dis[MAXN],pre[MAXN],cur[MAXN];
 8   int sap(int start,int end,int nodenum){
 9       memset(cur,0,sizeof(cur));
10       memset(dis,0,sizeof(dis));
11       memset(gap,0,sizeof(gap));
12       int u=pre[start]=start,maxflow=0,aug=-1;
13       gap[0]=nodenum;
14       while(dis[start]<nodenum){
15           loop:
16             for(int v=cur[u];v<nodenum;v++)


17               if(maze[u][v] && dis[u]==dis[v]+1){
18                   if(aug==-1 || aug>maze[u][v])aug=maze[u][v];
19                   pre[v]=u;
20                   u=cur[u]=v;
21                   if(v==end){
22                       maxflow+=aug;
23                       for(u=pre[u];v!=start;v=u,u=pre[u]){
24                           maze[u][v]-=aug;
25                           maze[v][u]+=aug;
26                       }
27                       aug=-1;
28                   }
29                   goto loop;
30               }
31               int mindis=nodenum-1;
32               for(int v=0;v<nodenum;v++)
33                  if(maze[u][v]&&mindis>dis[v]){
34                      cur[u]=v;
35                      mindis=dis[v];
36                  }
37               if((--gap[dis[u]])==0)break;
38               gap[dis[u]=mindis+1]++;
39               u=pre[u];
40       }
41       return maxflow;
42 }
\end{Code}

\subsubsection{SAP dạng ma trận kề 2}

\begin{Code}
     保留原矩阵，可用于多次使用最大流
 1   /*
 2    * SAP 邻接矩阵形式
 3    * 点的编号从 0 开始
 4    * 增加个 flow 数组，保留原矩阵 maze, 可用于多次使用最大流
 5    */
 6   const int MAXN=1100;
 7   int maze[MAXN][MAXN];
 8   int gap[MAXN],dis[MAXN],pre[MAXN],cur[MAXN];
 9   int flow[MAXN][MAXN];//存最大流的容量
10   int sap(int start,int end,int nodenum){
11       memset(cur,0,sizeof(cur));
12       memset(dis,0,sizeof(dis));
13       memset(gap,0,sizeof(gap));
14       memset(flow,0,sizeof(flow));
15       int u=pre[start]=start,maxflow=0,aug=-1;
16       gap[0]=nodenum;
17       while(dis[start]<nodenum){
18           loop:
19             for(int v=cur[u];v<nodenum;v++)
20               if(maze[u][v]-flow[u][v] && dis[u]==dis[v]+1){
21                   if(aug==-1 || aug>maze[u][v]-flow[u][v])aug=maze[u
                        ][v]-flow[u][v];


22                   pre[v]=u;
23                   u=cur[u]=v;
24                   if(v==end){
25                       maxflow+=aug;
26                       for(u=pre[u];v!=start;v=u,u=pre[u]){
27                           flow[u][v]+=aug;
28                           flow[v][u]-=aug;
29                       }
30                       aug=-1;
31                   }
32                   goto loop;
33               }
34               int mindis=nodenum-1;
35               for(int v=0;v<nodenum;v++)
36                  if(maze[u][v]-flow[u][v]&&mindis>dis[v]){
37                      cur[u]=v;
38                      mindis=dis[v];
39                  }
40               if((--gap[dis[u]])==0)break;
41               gap[dis[u]=mindis+1]++;
42               u=pre[u];
43       }
44       return maxflow;
45 }
\end{Code}

\subsubsection{ISAP dạng danh sách kề}

\begin{Code}
 1   const int MAXN = 100010;//点数的最大值
 2   const int MAXM = 400010;//边数的最大值
 3   const int INF = 0x3f3f3f3f;
 4   struct Edge{
 5       int to,next,cap,flow;
 6   }edge[MAXM];//注意是 MAXM
 7   int tol;
 8   int head[MAXN];
 9   int gap[MAXN],dep[MAXN],pre[MAXN],cur[MAXN];
10   void init(){
11       tol = 0;
12       memset(head,-1,sizeof(head));
13   }
14   //加边，单向图三个参数，双向图四个参数
15   void addedge(int u,int v,int w,int rw=0){
16       edge[tol].to = v;edge[tol].cap = w;edge[tol].next = head[u];
17       edge[tol].flow = 0;head[u] = tol++;
18       edge[tol].to = u;edge[tol].cap = rw;edge[tol].next = head[v];
19       edge[tol].flow = 0;head[v]=tol++;
20   }
21   //输入参数：起点、终点、点的总数
22   //点的编号没有影响，只要输入点的总数
23   int sap(int start,int end,int N){
24       memset(gap,0,sizeof(gap));


25       memset(dep,0,sizeof(dep));
26       memcpy(cur,head,sizeof(head));
27       int u = start;
28       pre[u] = -1;
29       gap[0] = N;
30       int ans = 0;
31       while(dep[start] < N){
32           if(u == end){
33               int Min = INF;
34               for(int i = pre[u];i != -1; i = pre[edge[i^1].to])
35                    if(Min > edge[i].cap - edge[i].flow)
36                            Min = edge[i].cap - edge[i].flow;
37               for(int i = pre[u];i != -1; i = pre[edge[i^1].to]){
38                    edge[i].flow += Min;
39                    edge[i^1].flow -= Min;
40               }
41               u = start;
42               ans += Min;
43               continue;
44           }
45           bool flag = false;
46           int v;
47           for(int i = cur[u]; i != -1;i = edge[i].next){
48               v = edge[i].to;
49               if(edge[i].cap - edge[i].flow && dep[v]+1 == dep[u])
50               {
51                    flag = true;
52                    cur[u] = pre[v] = i;
53                    break;
54               }
55           }
56           if(flag){
57               u = v;
58               continue;
59           }
60           int Min = N;
61           for(int i = head[u]; i != -1;i = edge[i].next)
62               if(edge[i].cap - edge[i].flow && dep[edge[i].to] < Min)
                    {
63                    Min = dep[edge[i].to];
64                    cur[u] = i;
65               }
66           gap[dep[u]]--;
67           if(!gap[dep[u]])return ans;
68           dep[u] = Min+1;
69           gap[dep[u]]++;
70           if(u != start) u = edge[pre[u]^1].to;
71       }
72       return ans;
73 }
\end{Code}

\subsubsection{ISAP + khởi tạo bfs + tối ưu bằng stack}

\begin{Code}
 1   const int MAXN = 100010;//点数的最大值
 2   const int MAXM = 400010;//边数的最大值
 3   const int INF = 0x3f3f3f3f;
 4   struct Edge{
 5       int to,next,cap,flow;
 6   }edge[MAXM];//注意是 MAXM
 7   int tol;
 8   int head[MAXN];
 9   int gap[MAXN],dep[MAXN],cur[MAXN];
10   void init(){
11       tol = 0;
12       memset(head,-1,sizeof(head));
13   }
14   void addedge(int u,int v,int w,int rw = 0){
15       edge[tol].to = v; edge[tol].cap = w; edge[tol].flow = 0;
16       edge[tol].next = head[u]; head[u] = tol++;
17       edge[tol].to = u; edge[tol].cap = rw; edge[tol].flow = 0;
18       edge[tol].next = head[v]; head[v] = tol++;
19   }
20   int Q[MAXN];
21   void BFS(int start,int end){
22       memset(dep,-1,sizeof(dep));
23       memset(gap,0,sizeof(gap));
24       gap[0] = 1;
25       int front = 0, rear = 0;
26       dep[end] = 0;
27       Q[rear++] = end;
28       while(front != rear){
29           int u = Q[front++];
30           for(int i = head[u]; i != -1; i = edge[i].next){
31                int v = edge[i].to;
32                if(dep[v] != -1)continue;
33                Q[rear++] = v;
34                dep[v] = dep[u] + 1;
35                gap[dep[v]]++;
36           }
37       }
38   }
39   int S[MAXN];
40   int sap(int start,int end,int N){
41       BFS(start,end);
42       memcpy(cur,head,sizeof(head));
43       int top = 0;
44       int u = start;
45       int ans = 0;
46       while(dep[start] < N){
47           if(u == end){
48                int Min = INF;
49                int inser;


50                  for(int i = 0;i < top;i++)
51                      if(Min > edge[S[i]].cap - edge[S[i]].flow){
52                          Min = edge[S[i]].cap - edge[S[i]].flow;
53                          inser = i;
54                      }
55                  for(int i = 0;i < top;i++){
56                      edge[S[i]].flow += Min;
57                      edge[S[i]^1].flow -= Min;
58                  }
59                  ans += Min;
60                  top = inser;
61                  u = edge[S[top]^1].to;
62                  continue;
63              }
64              bool flag = false;
65              int v;
66              for(int i = cur[u]; i != -1; i = edge[i].next){
67                  v = edge[i].to;
68                  if(edge[i].cap - edge[i].flow && dep[v]+1 == dep[u]){
69                       flag = true;
70                       cur[u] = i;
71                       break;
72                  }
73              }
74              if(flag){
75                  S[top++] = cur[u];
76                  u = v;
77                  continue;
78              }
79              int Min = N;
80              for(int i = head[u]; i != -1; i = edge[i].next)
81                  if(edge[i].cap - edge[i].flow && dep[edge[i].to] < Min)
                       {
82                       Min = dep[edge[i].to];
83                       cur[u] = i;
84                  }
85              gap[dep[u]]--;
86              if(!gap[dep[u]])return ans;
87              dep[u] = Min + 1;
88              gap[dep[u]]++;
89              if(u != start)u = edge[S[--top]^1].to;
90       }
91       return ans;
92 }
\end{Code}

\subsubsection{Dinic}

\begin{Code}
1    const int MAXN = 2010;//点数的最大值
2    const int MAXM = 1200010;//边数的最大值
3    const int INF = 0x3f3f3f3f;
4    struct Edge{
5        int to,next,cap,flow;


 6   }edge[MAXM];//注意是 MAXM
 7   int tol;
 8   int head[MAXN];
 9   void init(){
10       tol = 2;
11       memset(head,-1,sizeof(head));
12   }
13   void addedge(int u,int v,int w,int rw = 0){
14       edge[tol].to = v; edge[tol].cap = w; edge[tol].flow = 0;
15       edge[tol].next = head[u]; head[u] = tol++;
16       edge[tol].to = u; edge[tol].cap = rw; edge[tol].flow = 0;
17       edge[tol].next = head[v]; head[v] = tol++;
18   }
19   int Q[MAXN];
20   int dep[MAXN],cur[MAXN],sta[MAXN];
21   bool bfs(int s,int t,int n){
22       int front = 0,tail = 0;
23       memset(dep,-1,sizeof(dep[0])*(n+1));
24       dep[s] = 0;
25       Q[tail++] = s;
26       while(front < tail){
27           int u = Q[front++];
28           for(int i = head[u]; i != -1;i = edge[i].next){
29                int v = edge[i].to;
30                if(edge[i].cap > edge[i].flow && dep[v] == -1){
31                    dep[v] = dep[u] + 1;
32                    if(v == t)return true;
33                    Q[tail++] = v;
34                }
35           }
36       }
37       return false;
38   }
39   int dinic(int s,int t,int n){
40       int maxflow = 0;
41       while(bfs(s,t,n)){
42           for(int i = 0;i < n;i++)cur[i] = head[i];
43           int u = s, tail = 0;
44           while(cur[s] != -1){
45                if(u == t){
46                    int tp = INF;
47                    for(int i = tail-1;i >= 0;i--)
48                        tp = min(tp,edge[sta[i]].cap-edge[sta[i]].flow)
                             ;
49                    maxflow += tp;
50                    for(int i = tail-1;i >= 0;i--){
51                        edge[sta[i]].flow += tp;
52                        edge[sta[i]^1].flow -= tp;
53                        if(edge[sta[i]].cap-edge[sta[i]].flow == 0)
54                             tail = i;
55                    }


56                      u = edge[sta[tail]^1].to;
57                  }
58                  else if(cur[u] != -1 && edge[cur[u]].cap > edge[cur[u
                       ]].flow && dep[u] + 1 == dep[edge[cur[u]].to]){
59                      sta[tail++] = cur[u];
60                      u = edge[cur[u]].to;
61                  }
62                  else {
63                      while(u != s && cur[u] == -1)
64                          u = edge[sta[--tail]^1].to;
65                      cur[u] = edge[cur[u]].next;
66                  }
67              }
68       }
69       return maxflow;
70 }
\end{Code}

\subsubsection{Kiểm tra nhiều nghiệm cho luồng cực đại}

\begin{Code}
 1   //判断最大流多解就是在残留网络中找正环
 2   bool vis[MAXN],no[MAXN];
 3   int Stack[MAXN],top;
 4   bool dfs(int u,int pre,bool flag){
 5       vis[u] = 1;
 6       Stack[top++] = u;
 7       for(int i = head[u];i != -1;i = edge[i].next)
 8       {
 9           int v = edge[i].to;
10           if(edge[i].cap <= edge[i].flow)continue;
11           if(v == pre)continue;
12           if(!vis[v]){
13               if(dfs(v,u,edge[i^1].flow < edge[i^1].cap))return true;
14           }
15           else if(!no[v])return true;
16       }
17       if(!flag){
18           while(1){
19               int v = Stack[--top];
20               no[v] = true;
21               if(v == u)break;
22           }
23       }
24       return false;
25   }
26   //执行完最大流后可进行调用
27           memset(vis,false,sizeof(vis));
28           memset(no,false,sizeof(no));
29           top = 0;
30           bool flag = dfs(end,end,0);
\end{Code}

\subsection{Luồng cực đại chi phí nhỏ nhất}

\subsubsection{Luồng chi phí bản SPFA}

\begin{Code}
     最小费用最大流，求最大费用只需要取相反数，结果取相反数即可。
     点的总数为 N，点的编号 0 ∼ N-1
 1   const int MAXN = 10000;
 2   const int MAXM = 100000;
 3   const int INF = 0x3f3f3f3f;
 4   struct Edge{
 5       int to,next,cap,flow,cost;
 6   }edge[MAXM];
 7   int head[MAXN],tol;
 8   int pre[MAXN],dis[MAXN];
 9   bool vis[MAXN];
10   int N;//节点总个数，节点编号从 0∼N-1
11   void init(int n){
12       N = n;
13       tol = 0;
14       memset(head,-1,sizeof(head));
15   }
16   void addedge(int u,int v,int cap,int cost){
17       edge[tol].to = v;
18       edge[tol].cap = cap;
19       edge[tol].cost = cost;
20       edge[tol].flow = 0;
21       edge[tol].next = head[u];
22       head[u] = tol++;
23       edge[tol].to = u;
24       edge[tol].cap = 0;
25       edge[tol].cost = -cost;
26       edge[tol].flow = 0;
27       edge[tol].next = head[v];
28       head[v] = tol++;
29   }
30   bool spfa(int s,int t){
31       queue<int>q;
32       for(int i = 0;i < N;i++){
33           dis[i] = INF;
34           vis[i] = false;
35           pre[i] = -1;
36       }
37       dis[s] = 0;
38       vis[s] = true;
39       q.push(s);
40       while(!q.empty()){
41           int u = q.front();
42           q.pop();
43           vis[u] = false;
44           for(int i = head[u]; i != -1;i = edge[i].next){
45                int v = edge[i].to;


46               if(edge[i].cap > edge[i].flow && dis[v] > dis[u] + edge
                    [i].cost )
47               {
48                   dis[v] = dis[u] + edge[i].cost;
49                   pre[v] = i;
50                   if(!vis[v]){
51                       vis[v] = true;
52                       q.push(v);
53                   }
54               }
55           }
56       }
57       if(pre[t] == -1)return false;
58       else return true;
59   }
60   //返回的是最大流，cost 存的是最小费用
61   int minCostMaxflow(int s,int t,int &cost){
62       int flow = 0;
63       cost = 0;
64       while(spfa(s,t)){
65           int Min = INF;
66           for(int i = pre[t];i != -1;i = pre[edge[i^1].to]){
67               if(Min > edge[i].cap - edge[i].flow)
68                    Min = edge[i].cap - edge[i].flow;
69           }
70           for(int i = pre[t];i != -1;i = pre[edge[i^1].to]){
71               edge[i].flow += Min;
72               edge[i^1].flow -= Min;
73               cost += edge[i].cost * Min;
74           }
75           flow += Min;
76       }
77       return flow;
78   }
\end{Code}

\subsubsection{Luồng chi phí zkw}

\begin{Code}
     对于二分图类型的比较高效
1    const int MAXN = 100;
2    const int MAXM = 20000;
3    const int INF = 0x3f3f3f3f;
4    struct Edge{
5        int to,next,cap,flow,cost;
6        Edge(int _to = 0,int _next = 0,int _cap = 0,int _flow = 0,int
            _cost = 0):
 7           to(_to),next(_next),cap(_cap),flow(_flow),cost(_cost){}
 8   }edge[MAXM];
 9   struct ZKW_MinCostMaxFlow{
10       int head[MAXN],tot;
11       int cur[MAXN];
12       int dis[MAXN];


13       bool vis[MAXN];
14       int ss,tt,N;//源点、汇点和点的总个数（编号是 0∼N-1）, 不需要额外赋值，
            调用会直接赋值
15       int min_cost, max_flow;
16       void init(){
17           tot = 0;
18           memset(head,-1,sizeof(head));
19       }
20       void addedge(int u,int v,int cap,int cost){
21           edge[tot] = Edge(v,head[u],cap,0,cost);
22           head[u] = tot++;
23           edge[tot] = Edge(u,head[v],0,0,-cost);
24           head[v] = tot++;
25       }
26       int aug(int u,int flow){
27           if(u == tt)return flow;
28           vis[u] = true;
29           for(int i = cur[u];i != -1;i = edge[i].next){
30               int v = edge[i].to;
31               if(edge[i].cap > edge[i].flow && !vis[v] && dis[u] ==
                    dis[v] + edge[i].cost){
32                    int tmp = aug(v,min(flow,edge[i].cap-edge[i].flow))
                         ;
33                    edge[i].flow += tmp;
34                    edge[i^1].flow -= tmp;
35                    cur[u] = i;
36                    if(tmp)return tmp;
37               }
38           }
39           return 0;
40       }
41       bool modify_label(){
42           int d = INF;
43           for(int u = 0;u < N;u++)
44               if(vis[u])
45                    for(int i = head[u];i != -1;i = edge[i].next){
46                         int v = edge[i].to;
47                         if(edge[i].cap>edge[i].flow && !vis[v])
48                             d = min(d,dis[v]+edge[i].cost-dis[u]);
49                    }
50           if(d == INF)return false;
51           for(int i = 0;i < N;i++)
52               if(vis[i]){
53                    vis[i] = false;
54                    dis[i] += d;
55
56                  }
57              return true;
58       }
59       /*
60         * 直接调用获取最小费用和最大流


61       * 输入: start-源点，end-汇点，n-点的总个数（编号从 0 开始）
62       * 返回值: pair<int,int> 第一个是最小费用，第二个是最大流
63       */
64     pair<int,int> mincostmaxflow(int start,int end,int n){
65          ss = start, tt = end, N = n;
66          min_cost = max_flow = 0;
67          for(int i = 0;i < n;i++)dis[i] = 0;
68          while(1){
69              for(int i = 0;i < n;i++)cur[i] = head[i];
70              while(1){
71                  for(int i = 0;i < n;i++)vis[i] = false;
72                  int tmp = aug(ss,INF);
73                  if(tmp == 0)break;
74                  max_flow += tmp;
75                  min_cost += tmp*dis[ss];
76              }
77              if(!modify_label())break;
78          }
79          return make_pair(min_cost,max_flow);
80     }
81 }solve;
\end{Code}

\subsection{2-SAT}

\subsubsection{Phương pháp tô màu (có thể nhận nghiệm nhỏ nhất theo thứ tự từ điển)}

\begin{Code}
     HDU 1814
 1   const int MAXN = 20020;
 2   const int MAXM = 100010;
 3   struct Edge
 4   {
 5       int to,next;
 6   }edge[MAXM];
 7   int head[MAXN],tot;
 8   void init()
 9   {
10       tot = 0;
11       memset(head,-1,sizeof(head));
12   }
13   void addedge(int u,int v)
14   {
15       edge[tot].to = v;edge[tot].next = head[u];head[u] = tot++;
16   }
17   bool vis[MAXN];//染色标记，为 true 表示选择
18   int S[MAXN],top;//栈
19   bool dfs(int u)
20   {
21       if(vis[u^1])return false;
22       if(vis[u])return true;
23       vis[u] = true;
24       S[top++] = u;


25       for(int i = head[u];i != -1;i = edge[i].next)
26           if(!dfs(edge[i].to))
27               return false;
28       return true;
29   }
30   bool Twosat(int n)
31   {
32       memset(vis,false,sizeof(vis));
33       for(int i = 0;i < n;i += 2)
34       {
35           if(vis[i] || vis[i^1])continue;
36           top = 0;
37           if(!dfs(i))
38           {
39                while(top)vis[S[--top]] = false;
40                if(!dfs(i^1)) return false;
41           }
42       }
43       return true;
44   }
45   int main()
46   {
47       int n,m;
48       int u,v;
49       while(scanf("%d%d",&n,&m) == 2)
50       {
51           init();
52           while(m--)
53           {
54                scanf("%d%d",&u,&v);
55                u--;v--;
56                addedge(u,v^1);
57                addedge(v,u^1);
58           }
59           if(Twosat(2*n))
60           {
61                for(int i = 0;i < 2*n;i++)
62                    if(vis[i])
63                       printf("%d\n",i+1);
64           }
65           else printf("NIE\n");
66       }
67       return 0;
68   }
\end{Code}

\subsubsection{Co thành phần liên thông mạnh (sắp xếp topo chỉ cho một nghiệm bất kỳ)}

\begin{Code}
     POJ 3648 Wedding
1 //******************************************
2 //2-SAT 强连通缩点
3 const int MAXN = 1010;


 4   const int MAXM = 100010;
 5   struct Edge
 6   {
 7       int to,next;
 8   }edge[MAXM];
 9   int head[MAXN],tot;
10   void init()
11   {
12       tot = 0;
13       memset(head,-1,sizeof(head));
14   }
15   void addedge(int u,int v)
16   {
17       edge[tot].to = v; edge[tot].next = head[u]; head[u] = tot++;
18   }
19   int Low[MAXN],DFN[MAXN],Stack[MAXN],Belong[MAXN];//Belong 数组的值 1
        ∼ scc
20   int Index,top;
21   int scc;
22   bool Instack[MAXN];
23   int num[MAXN];
24   void Tarjan(int u)
25   {
26       int v;
27       Low[u] = DFN[u] = ++Index;
28       Stack[top++] = u;
29       Instack[u] = true;
30       for(int i = head[u];i != -1;i = edge[i].next)
31       {
32           v = edge[i].to;
33           if( !DFN[v] )
34           {
35                Tarjan(v);
36                if(Low[u] > Low[v])Low[u] = Low[v];
37           }
38           else if(Instack[v] && Low[u] > DFN[v])
39                Low[u] = DFN[v];
40       }
41       if(Low[u] == DFN[u])
42       {
43           scc++;
44           do
45           {
46                v = Stack[--top];
47                Instack[v] = false;
48                Belong[v] = scc;
49                num[scc]++;
50           }
51           while(v != u);
52       }
53   }


 54   bool solvable(int n)//n 是总个数, 需要选择一半
 55   {
 56       memset(DFN,0,sizeof(DFN));
 57       memset(Instack,false,sizeof(Instack));
 58       memset(num,0,sizeof(num));
 59       Index = scc = top = 0;
 60       for(int i = 0;i < n;i++)
 61           if(!DFN[i])
 62               Tarjan(i);
 63       for(int i = 0;i < n;i += 2)
 64       {
 65           if(Belong[i] == Belong[i^1])
 66               return false;
 67       }
 68       return true;
 69   }
 70   //*************************************************
 71
 72   //拓扑排序求任意一组解部分
 73   queue<int>q1,q2;
 74   vector<vector<int> > dag;//缩点后的逆向 DAG 图
 75   char color[MAXN];//染色，为'R' 是选择的
 76   int indeg[MAXN];//入度
 77   int cf[MAXN];
 78   void solve(int n)
 79   {
 80       dag.assign(scc+1,vector<int>());
 81       memset(indeg,0,sizeof(indeg));
 82       memset(color,0,sizeof(color));
 83       for(int u = 0;u < n;u++)
 84           for(int i = head[u];i != -1;i = edge[i].next)
 85           {
 86               int v = edge[i].to;
 87               if(Belong[u] != Belong[v])
 88               {
 89                   dag[Belong[v]].push_back(Belong[u]);
 90                   indeg[Belong[u]]++;
 91               }
 92           }
 93       for(int i = 0;i < n;i += 2)
 94       {
 95           cf[Belong[i]] = Belong[i^1];
 96           cf[Belong[i^1]] = Belong[i];
 97       }
 98       while(!q1.empty())q1.pop();
 99       while(!q2.empty())q2.pop();
100       for(int i = 1;i <= scc;i++)
101           if(indeg[i] == 0)
102               q1.push(i);
103       while(!q1.empty())
104       {


105              int u = q1.front();
106              q1.pop();
107              if(color[u] == 0)
108              {
109                  color[u] = 'R';
110                  color[cf[u]] = 'B';
111              }
112              int sz = dag[u].size();
113              for(int i = 0;i < sz;i++)
114              {
115                  indeg[dag[u][i]]--;
116                  if(indeg[dag[u][i]] == 0)
117                      q1.push(dag[u][i]);
118              }
119       }
120   }
121
122   int change(char s[])
123   {
124       int ret = 0;
125       int i = 0;
126       while(s[i]>='0' && s[i]<='9')
127       {
128           ret *= 10;
129           ret += s[i]-'0';
130           i++;
131       }
132       if(s[i] == 'w')return 2*ret;
133       else return 2*ret+1;
134   }
135   int main()
136   {
137       int n,m;
138       char s1[10],s2[10];
139       while(scanf("%d%d",&n,&m) == 2)
140       {
141           if(n == 0 && m == 0)break;
142           init();
143           while(m--)
144           {
145                scanf("%s%s",s1,s2);
146                int u = change(s1);
147                int v = change(s2);
148                addedge(u^1,v);
149                addedge(v^1,u);
150           }
151           addedge(1,0);
152           if(solvable(2*n))
153           {
154                solve(2*n);
155                for(int i = 1;i < n;i++)


156                 {
157                     //注意这一定是判断 color[Belong]
158                     if(color[Belong[2*i]] == 'R')printf("%dw",i);
159                     else printf("%dh",i);
160                     if(i < n-1)printf(" ");
161                     else printf("\n");
162                 }
163             }
164             else printf("bad luck\n");
165      }
166      return 0;
167 }
\end{Code}

\subsection{Cây khung nhỏ nhất Manhattan}

\begin{Code}
     POJ 3241 求曼哈顿最小生成树上第 k 大的边
 1   const int MAXN = 100010;
 2   const int INF = 0x3f3f3f3f;
 3   struct Point{
 4       int x,y,id;
 5   }p[MAXN];
 6   bool cmp(Point a,Point b){
 7       if(a.x != b.x) return a.x < b.x;
 8       else return a.y < b.y;
 9   }
10   //树状数组，找 y-x 大于当前的，但是 y+x 最小的
11   struct BIT{
12       int min_val,pos;
13       void init()
14       {
15           min_val = INF;
16           pos = -1;
17       }
18   }bit[MAXN];
19   //所有有效边
20   struct Edge{
21       int u,v,d;
22   }edge[MAXN<<2];
23   bool cmpedge(Edge a,Edge b){
24       return a.d < b.d;
25   }
26   int tot;
27   int n;
28   int F[MAXN];
29   int find(int x){
30       if(F[x] == -1) return x;
31       else return F[x] = find(F[x]);
32   }
33   void addedge(int u,int v,int d){
34       edge[tot].u = u;
35       edge[tot].v = v;


36       edge[tot++].d = d;
37   }
38   int lowbit(int x){
39       return x&(-x);
40   }
41   void update(int i,int val,int pos){
42       while(i > 0){
43           if(val < bit[i].min_val){
44                bit[i].min_val = val;
45                bit[i].pos = pos;
46           }
47           i -= lowbit(i);
48       }
49   }
50   //查询 [i,m] 的最小值位置
51   int ask(int i,int m){
52       int min_val = INF,pos = -1;
53       while(i <= m){
54           if(bit[i].min_val < min_val){
55                min_val = bit[i].min_val;
56                pos = bit[i].pos;
57           }
58           i += lowbit(i);
59       }
60       return pos;
61   }
62   int dist(Point a,Point b){
63       return abs(a.x - b.x) + abs(a.y - b.y);
64   }
65   void Manhattan_minimum_spanning_tree(int n,Point p[]){
66       int a[MAXN],b[MAXN];
67       tot = 0;
68       for(int dir = 0; dir < 4;dir++){
69           //4 种坐标变换
70           if(dir == 1 || dir == 3){
71                for(int i = 0;i < n;i++)
72                    swap(p[i].x,p[i].y);
73           }
74           else if(dir == 2){
75                for(int i = 0;i < n;i++)
76                    p[i].x = -p[i].x;
77           }
78           sort(p,p+n,cmp);
79           for(int i = 0;i < n;i++)
80                a[i] = b[i] = p[i].y - p[i].x;
81           sort(b,b+n);
82           int m = unique(b,b+n) - b;
83           for(int i = 1;i <= m;i++)
84                bit[i].init();
85           for(int i = n-1 ;i >= 0;i--){
86                int pos = lower_bound(b,b+m,a[i]) - b + 1;


 87                  int ans = ask(pos,m);
 88                  if(ans != -1)
 89                      addedge(p[i].id,p[ans].id,dist(p[i],p[ans]));
 90                  update(pos,p[i].x+p[i].y,i);
 91              }
 92       }
 93   }
 94   int solve(int k){
 95       Manhattan_minimum_spanning_tree(n,p);
 96       memset(F,-1,sizeof(F));
 97       sort(edge,edge+tot,cmpedge);
 98       for(int i = 0;i < tot;i++){
 99           int u = edge[i].u;
100           int v = edge[i].v;
101           int t1 = find(u), t2 = find(v);
102           if(t1 != t2){
103               F[t1] = t2;
104               k--;
105               if(k == 0)return edge[i].d;
106           }
107       }
108   }
109   int main()
110   {
111       //freopen("in.txt","r",stdin);
112       //freopen("out.txt","w",stdout);
113       int k;
114       while(scanf("%d%d",&n,&k)==2 && n){
115           for(int i = 0;i < n;i++){
116               scanf("%d%d",&p[i].x,&p[i].y);
117               p[i].id = i;
118           }
119           printf("%d\n",solve(n-k));
120       }
121       return 0;
122   }
\end{Code}

\subsection{LCA}

\subsubsection{Thuật toán trực tuyến dfs+ST}

\begin{Code}
 1    /*
 2     * LCA (POJ 1330)
 3     * 在线算法 DFS + ST
 4     */
 5    const int MAXN = 10010;
 6    int rmq[2*MAXN];//rmq 数组，就是欧拉序列对应的深度序列
 7    struct ST{
 8        int mm[2*MAXN];
 9        int dp[2*MAXN][20];//最小值对应的下标
10        void init(int n){


11              mm[0] = -1;
12              for(int i = 1;i <= n;i++){
13                  mm[i] = ((i&(i-1)) == 0)?mm[i-1]+1:mm[i-1];
14                  dp[i][0] = i;
15              }
16              for(int j = 1; j <= mm[n];j++)
17                  for(int i = 1; i + (1<<j) - 1 <= n; i++)
18                      dp[i][j] = rmq[dp[i][j-1]] < rmq[dp[i+(1<<(j-1))][j
                           -1]]?dp[i][j-1]:dp[i+(1<<(j-1))][j-1];
19       }
20       //查询 [a,b] 之间最小值的下标
21       int query(int a,int b)
22       {
23           if(a > b)swap(a,b);
24           int k = mm[b-a+1];
25           return rmq[dp[a][k]] <= rmq[dp[b-(1<<k)+1][k]]?dp[a][k]:dp[
                b-(1<<k)+1][k];
26       }
27   };
28   //边的结构体定义
29   struct Edge{
30       int to,next;
31   };
32   Edge edge[MAXN*2];
33   int tot,head[MAXN];
34
35   int F[MAXN*2];//欧拉序列，就是 dfs 遍历的顺序，长度为 2*n-1, 下标从 1 开始
36   int P[MAXN];//P[i] 表示点 i 在 F 中第一次出现的位置
37   int cnt;
38   ST st;
39   void init(){
40       tot = 0;
41       memset(head,-1,sizeof(head));
42   }
43   //加边，无向边需要加两次
44   void addedge(int u,int v){
45       edge[tot].to = v;
46       edge[tot].next = head[u];
47       head[u] = tot++;
48   }
49   void dfs(int u,int pre,int dep){
50       F[++cnt] = u;
51       rmq[cnt] = dep;
52       P[u] = cnt;
53       for(int i = head[u];i != -1;i = edge[i].next){
54           int v = edge[i].to;
55           if(v == pre)continue;
56           dfs(v,u,dep+1);
57           F[++cnt] = u;
58           rmq[cnt] = dep;
59       }


60   }
61   //查询 LCA 前的初始化
62   void LCA_init(int root,int node_num){
63       cnt = 0;
64       dfs(root,root,0);
65       st.init(2*node_num-1);
66   }
67   //查询 u,v 的 lca 编号
68   int query_lca(int u,int v){
69       return F[st.query(P[u],P[v])];
70   }
71   bool flag[MAXN];
72   int main(){
73       int T;
74       int N;
75       int u,v;
76       scanf("%d",&T);
77       while(T--){
78           scanf("%d",&N);
79           init();
80           memset(flag,false,sizeof(flag));
81           for(int i = 1; i < N;i++){
82                scanf("%d%d",&u,&v);
83                addedge(u,v);
84                addedge(v,u);
85                flag[v] = true;
86           }
87           int root;
88           for(int i = 1; i <= N;i++)
89                if(!flag[i]){
90                    root = i;
91                    break;
92                }
93           LCA_init(root,N);
94           scanf("%d%d",&u,&v);
95           printf("%d\n",query_lca(u,v));
96       }
97       return 0;
98   }
\end{Code}

\subsubsection{Thuật toán Tarjan offline}

\begin{Code}
 1 /*
 2 * POJ 1470
 3 * 给出一颗有向树，Q 个查询
 4 * 输出查询结果中每个点出现次数
 5 */
 6 /*
 7 * 离线算法，LCATarjan
 8 * 复杂度O(n+Q);
 9 */
10 const int MAXN = 1010;


11   const int MAXQ = 500010;//查询数的最大值
12
13   //并查集部分
14   int F[MAXN];//需要初始化为 -1
15   int find(int x){
16       if(F[x] == -1)return x;
17       return F[x] = find(F[x]);
18   }
19   void bing(int u,int v){
20       int t1 = find(u);
21       int t2 = find(v);
22       if(t1 != t2)
23           F[t1] = t2;
24   }
25   //************************
26   bool vis[MAXN];//访问标记
27   int ancestor[MAXN];//祖先
28   struct Edge{
29       int to,next;
30   }edge[MAXN*2];
31   int head[MAXN],tot;
32   void addedge(int u,int v){
33       edge[tot].to = v;
34       edge[tot].next = head[u];
35       head[u] = tot++;
36   }
37
38   struct Query{
39       int q,next;
40       int index;//查询编号
41   }query[MAXQ*2];
42   int answer[MAXQ];//存储最后的查询结果，下标 0 Q-1
43   int h[MAXQ];
44   int tt;
45   int Q;
46
47   void add_query(int u,int v,int index){
48       query[tt].q = v;
49       query[tt].next = h[u];
50       query[tt].index = index;
51       h[u] = tt++;
52       query[tt].q = u;
53       query[tt].next = h[v];
54       query[tt].index = index;
55       h[v] = tt++;
56   }
57
58   void init(){
59       tot = 0;
60       memset(head,-1,sizeof(head));
61       tt = 0;


 62       memset(h,-1,sizeof(h));
 63       memset(vis,false,sizeof(vis));
 64       memset(F,-1,sizeof(F));
 65       memset(ancestor,0,sizeof(ancestor));
 66   }
 67   void LCA(int u){
 68       ancestor[u] = u;
 69       vis[u] = true;
 70       for(int i = head[u];i != -1;i = edge[i].next){
 71           int v = edge[i].to;
 72           if(vis[v])continue;
 73           LCA(v);
 74           bing(u,v);
 75           ancestor[find(u)] = u;
 76       }
 77       for(int i = h[u];i != -1;i = query[i].next){
 78           int v = query[i].q;
 79           if(vis[v]){
 80               answer[query[i].index] = ancestor[find(v)];
 81           }
 82       }
 83   }
 84   bool flag[MAXN];
 85   int Count_num[MAXN];
 86   int main(){
 87       int n;
 88       int u,v,k;
 89       while(scanf("%d",&n) == 1){
 90           init();
 91           memset(flag,false,sizeof(flag));
 92           for(int i = 1;i <= n;i++){
 93               scanf("%d:(%d)",&u,&k);
 94               while(k--){
 95                    scanf("%d",&v);
 96                    flag[v] = true;
 97                    addedge(u,v);
 98                    addedge(v,u);
 99               }
100           }
101           scanf("%d",&Q);
102           for(int i = 0;i < Q;i++){
103               char ch;
104               cin>>ch;
105               scanf("%d %d)",&u,&v);
106               add_query(u,v,i);
107           }
108           int root;
109           for(int i = 1;i <= n;i++)
110               if(!flag[i]){
111                    root = i;
112                    break;


113                 }
114             LCA(root);
115             memset(Count_num,0,sizeof(Count_num));
116             for(int i = 0;i < Q;i++)
117                 Count_num[answer[i]]++;
118             for(int i = 1;i <= n;i++)
119                 if(Count_num[i] > 0)
120                     printf("%d:%d\n",i,Count_num[i]);
121      }
122      return 0;
123 }
\end{Code}

\subsubsection{Phương pháp nhân đôi cho LCA}

\begin{Code}
 1   /*
 2    * POJ 1330
 3    * LCA 在线算法
 4    */
 5   const int MAXN = 10010;
 6   const int DEG = 20;
 7
 8   struct Edge{
 9       int to,next;
10   }edge[MAXN*2];
11   int head[MAXN],tot;
12   void addedge(int u,int v){
13       edge[tot].to = v;
14       edge[tot].next = head[u];
15       head[u] = tot++;
16   }
17   void init(){
18       tot = 0;
19       memset(head,-1,sizeof(head));
20   }
21   int fa[MAXN][DEG];//fa[i][j] 表示结点 i 的第2^j个祖先
22   int deg[MAXN];//深度数组
23
24   void BFS(int root){
25       queue<int>que;
26       deg[root] = 0;
27       fa[root][0] = root;
28       que.push(root);
29       while(!que.empty()){
30           int tmp = que.front();
31           que.pop();
32           for(int i = 1;i < DEG;i++)
33               fa[tmp][i] = fa[fa[tmp][i-1]][i-1];
34           for(int i = head[tmp]; i != -1;i = edge[i].next){
35               int v = edge[i].to;
36               if(v == fa[tmp][0])continue;
37               deg[v] = deg[tmp] + 1;
38               fa[v][0] = tmp;


39                  que.push(v);
40              }
41
42       }
43   }
44   int LCA(int u,int v){
45       if(deg[u] > deg[v])swap(u,v);
46       int hu = deg[u], hv = deg[v];
47       int tu = u, tv = v;
48       for(int det = hv-hu, i = 0; det ;det>>=1, i++)
49           if(det&1)
50                tv = fa[tv][i];
51       if(tu == tv)return tu;
52       for(int i = DEG-1; i >= 0; i--){
53           if(fa[tu][i] == fa[tv][i])
54                continue;
55           tu = fa[tu][i];
56           tv = fa[tv][i];
57       }
58       return fa[tu][0];
59   }
60   bool flag[MAXN];
61   int main(){
62       int T;
63       int n;
64       int u,v;
65       scanf("%d",&T);
66       while(T--){
67           scanf("%d",&n);
68           init();
69           memset(flag,false,sizeof(flag));
70           for(int i = 1;i < n;i++){
71                scanf("%d%d",&u,&v);
72                addedge(u,v);
73                addedge(v,u);
74                flag[v] = true;
75           }
76           int root;
77           for(int i = 1;i <= n;i++)
78                if(!flag[i]){
79                    root = i;
80                    break;
81                }
82           BFS(root);
83           scanf("%d%d",&u,&v);
84           printf("%d\n",LCA(u,v));
85       }
86       return 0;
87   }
\end{Code}

\subsection{Đường đi Euler}

\begin{Code}
     欧拉回路：每条边只经过一次，而且回到起点
     欧拉路径：每条边只经过一次，不要求回到起点

     欧拉回路判断：
     无向图：连通（不考虑度为 0 的点），每个顶点度数都为偶数。
     有向图：基图连通（把边当成无向边，同样不考虑度为 0 的点），每个顶点出度等于入度。
     混合图（有无向边和有向边）：首先是基图连通（不考虑度为 0 的点），然后需要借助网络流
     判定。
     首先给原图中的每条无向边随便指定一个方向（称为初始定向），将原图改为有向图 G’，然
     后的任务就是改变 G’ 中某些边的方向（当然是无向边转化来的，原混合图中的有向边不能
     动）使其满足每个点的入度等于出度。
     设 D[i] 为 G’ 中 (点 i 的出度 - 点 i 的入度）。可以发现，在改变 G’ 中边的方向的过程中，任
     何点的 D 值的奇偶性都不会发生改变（设将边 <i, j> 改为 <j, i>，则 i 入度加 1 出度减 1，
     j 入度减 1 出度加 1，两者之差加 2 或减 2，奇偶性不变）！而最终要求的是每个点的入度等
     于出度，即每个点的 D 值都为 0，是偶数，故可得：若初始定向得到的 G’ 中任意一个点的
     D 值是奇数，那么原图中一定不存在欧拉环！

     若初始 D 值都是偶数，则将 G’ 改装成网络：设立源点 S 和汇点 T，对于每个 D[i]>0 的点
     i，连边 <S, i>，容量为 D[i]/2；对于每个 D[j]<0 的点 j，连边 <j, T>，容量为 -D[j]/2；G’
     中的每条边在网络中仍保留，容量为 1（表示该边最多只能被改变方向一次）。求这个网络的
     最大流，若 S 引出的所有边均满流，则原混合图是欧拉图，将网络中所有流量为 1 的中间边
     （就是不与 S 或 T 关联的边）在 G’ 中改变方向，形成的新图 G” 一定是有向欧拉图；若 S
     引出的边中有的没有满流，则原混合图不是欧拉图。

     欧拉路径的判断：
     无向图：连通（不考虑度为 0 的点），每个顶点度数都为偶数或者仅有两个点的度数为偶数。
     有向图：基图连通（把边当成无向边，同样不考虑度为 0 的点），每个顶点出度等于入度或
     者有且仅有一个点的出度比入度多 1，有且仅有一个点的出度比入度少 1，其余出度等于入
     度。
     混合图：如果存在欧拉回路，一点存在欧拉路径了。否则如果有且仅有两个点的（出度 -入
     度）是奇数，那么给这个两个点加边，判断是否存在欧拉回路。
\end{Code}

\subsubsection{Đồ thị có hướng}

\begin{Code}
     POJ 2337
     给出 n 个小写字母组成的单词，要求将 n 个单词连接起来，使得前一个单词的最后一个字母
     和后一个单词的第一个字母相同。输出字典序最小的解。
 1   struct Edge{
 2       int to,next;
 3       int index;
 4       bool flag;
 5   }edge[2010];
 6   int head[30],tot;
 7   void init(){
 8       tot = 0;
 9       memset(head,-1,sizeof(head));
10   }
11   void addedge(int u,int v,int index){
12       edge[tot].to = v;


13       edge[tot].next = head[u];
14       edge[tot].index = index;
15       edge[tot].flag = false;
16       head[u] = tot++;
17   }
18   string str[1010];
19   int in[30],out[30];
20   int cnt;
21   int ans[1010];
22   void dfs(int u){
23       for(int i = head[u] ;i != -1;i = edge[i].next)
24           if(!edge[i].flag)
25           {
26                edge[i].flag = true;
27                dfs(edge[i].to);
28                ans[cnt++] = edge[i].index;
29           }
30   }
31   int main(){
32       int T,n;
33       scanf("%d",&T);
34       while(T--){
35           scanf("%d",&n);
36           for(int i = 0;i < n;i++)
37                cin>>str[i];
38           sort(str,str+n);//要输出字典序最小的解，先按照字典序排序
39           init();
40           memset(in,0,sizeof(in));
41           memset(out,0,sizeof(out));
42           int start = 100;
43           for(int i = n-1;i >= 0;i--)//字典序大的先加入
44           {
45                int u = str[i][0] - 'a';
46                int v = str[i][str[i].length() - 1] - 'a';
47                addedge(u,v,i);
48                out[u]++;
49                in[v]++;
50                if(u < start)start = u;
51                if(v < start)start = v;
52           }
53           int cc1 = 0, cc2 = 0;
54           for(int i = 0;i < 26;i++){
55                if(out[i] - in[i] == 1){
56                    cc1++;
57                    start = i;//如果有一个出度比入度大 1 的点，就从这个点出发，
                         否则从最小的点出发
58                }
59                else if(out[i] - in[i] == -1)
60                    cc2++;
61                else if(out[i] - in[i] != 0)
62                    cc1 = 3;


63              }
64              if(! ( (cc1 == 0 && cc2 == 0) || (cc1 == 1 && cc2 == 1) )){
65                  printf("***\n");
66                  continue;
67              }
68              cnt = 0;
69              dfs(start);
70              if(cnt != n)//判断是否连通
71              {
72                  printf("***\n");
73                  continue;
74              }
75              for(int i = cnt-1; i >= 0;i--){
76                  cout<<str[ans[i]];
77                  if(i > 0)printf(".");
78                  else printf("\n");
79              }
80       }
81       return 0;
82 }
\end{Code}

\subsubsection{Đồ thị vô hướng}

\begin{Code}
     SGU 101
 1   struct Edge{
 2       int to,next;
 3       int index;
 4       int dir;
 5       bool flag;
 6   }edge[220];
 7   int head[10],tot;
 8   void init(){
 9       memset(head,-1,sizeof(head));
10       tot = 0;
11   }
12   void addedge(int u,int v,int index){
13       edge[tot].to = v;
14       edge[tot].next = head[u];
15       edge[tot].index = index;
16       edge[tot].dir = 0;
17       edge[tot].flag = false;
18       head[u] = tot++;
19       edge[tot].to = u;
20       edge[tot].next = head[v];
21       edge[tot].index = index;
22       edge[tot].dir = 1;
23       edge[tot].flag = false;
24       head[v] = tot++;
25   }
26   int du[10];
27   vector<int>ans;


28   void dfs(int u){
29       for(int i = head[u]; i != -1;i = edge[i].next)
30           if(!edge[i].flag ){
31               edge[i].flag = true;
32               edge[i^1].flag = true;
33               dfs(edge[i].to);
34               ans.push_back(i);
35           }
36   }
37   int main(){
38       int n;
39       while(scanf("%d",&n) == 1){
40           init();
41           int u,v;
42           memset(du,0,sizeof(du));
43           for(int i = 1;i <= n;i++){
44               scanf("%d%d",&u,&v);
45               addedge(u,v,i);
46               du[u]++;
47               du[v]++;
48           }
49           int s = -1;
50           int cnt = 0;
51           for(int i = 0;i <= 6;i++){
52               if(du[i]&1) {cnt++; s = i;}
53               if(du[i] > 0 && s == -1)
54                    s = i;
55           }
56           bool ff = true;
57           if(cnt != 0 && cnt != 2){
58               printf("No solution\n");
59               continue;
60           }
61           ans.clear();
62           dfs(s);
63           if(ans.size() != n){
64               printf("No solution\n");
65               continue;
66           }
67           for(int i = 0;i < ans.size();i++){
68               printf("%d ",edge[ans[i]].index);
69               if(edge[ans[i]].dir == 0)printf("-\n");
70               else printf("+\n");
71           }
72       }
73       return 0;
74   }
\end{Code}

\subsubsection{Đồ thị hỗn hợp}

\begin{Code}
     POJ 1637 （本题保证了连通，故不需要判断连通，否则要判断连通）


 1   const int MAXN = 210;
 2   //最大流部分ISAP
 3   const int MAXM = 20100;
 4   const int INF = 0x3f3f3f3f;
 5   struct Edge
 6   {
 7       int to,next,cap,flow;
 8   }edge[MAXM];
 9   int tol;
10   int head[MAXN];
11   int gap[MAXN],dep[MAXN],pre[MAXN],cur[MAXN];
12   void init()
13   {
14       tol = 0;
15       memset(head,-1,sizeof(head));
16   }
17   void addedge(int u,int v,int w,int rw = 0)
18   {
19       edge[tol].to = v;
20       edge[tol].cap = w;
21       edge[tol].next = head[u];
22       edge[tol].flow = 0;
23       head[u] = tol++;
24       edge[tol].to = u;
25       edge[tol].cap = rw;
26       edge[tol].next = head[v];
27       edge[tol].flow = 0;
28       head[v] = tol++;
29   }
30   int sap(int start,int end,int N)
31   {
32       memset(gap,0,sizeof(gap));
33       memset(dep,0,sizeof(dep));
34       memcpy(cur,head,sizeof(head));
35       int u = start;
36       pre[u] = -1;
37       gap[0] = N;
38       int ans = 0;
39       while(dep[start] < N)
40       {
41           if(u == end)
42           {
43                int Min = INF;
44                for(int i = pre[u]; i != -1;i = pre[edge[i^1].to])
45                    if(Min > edge[i].cap - edge[i].flow)
46                        Min = edge[i].cap - edge[i].flow;
47                for(int i = pre[u];i != -1;i = pre[edge[i^1].to])
48                {
49                    edge[i].flow += Min;
50                    edge[i^1].flow -= Min;
51                }


 52                  u = start;
 53                  ans += Min;
 54                  continue;
 55              }
 56              bool flag = false;
 57              int v;
 58              for(int i = cur[u];i != -1;i = edge[i].next)
 59              {
 60                  v = edge[i].to;
 61                  if(edge[i].cap - edge[i].flow && dep[v] + 1 == dep[u])
 62                  {
 63                       flag = true;
 64                       cur[u] = pre[v] = i;
 65                       break;
 66                  }
 67              }
 68              if(flag)
 69              {
 70                  u = v;
 71                  continue;
 72              }
 73              int Min = N;
 74              for(int i = head[u];i != -1;i = edge[i].next)
 75                  if(edge[i].cap - edge[i].flow && dep[edge[i].to] < Min)
 76                  {
 77                       Min = dep[edge[i].to];
 78                       cur[u] = i;
 79                  }
 80              gap[dep[u]]--;
 81              if(!gap[dep[u]])return ans;
 82              dep[u] = Min+1;
 83              gap[dep[u]]++;
 84              if(u != start)u = edge[pre[u]^1].to;
 85       }
 86       return ans;
 87   }
 88   //the end of 最大流部分
 89
 90   int in[MAXN],out[MAXN];//每个点的出度和入度
 91
 92   int main()
 93   {
 94       //freopen("in.txt","r",stdin);
 95       //freopen("out.txt","w",stdout);
 96       int T;
 97       int n,m;
 98       scanf("%d",&T);
 99       while(T--)
100       {
101           scanf("%d%d",&n,&m);
102           init();


103             int u,v,w;
104             memset(in,0,sizeof(in));
105             memset(out,0,sizeof(out));
106             while(m--)
107             {
108                 scanf("%d%d%d",&u,&v,&w);
109                 out[u]++; in[v]++;
110                 if(w == 0)//双向
111                     addedge(u,v,1);
112             }
113             bool flag = true;
114             for(int i = 1;i <= n;i++)
115             {
116                 if(out[i] - in[i] > 0)
117                     addedge(0,i,(out[i] - in[i])/2);
118                 else if(in[i] - out[i] > 0)
119                     addedge(i,n+1,(in[i] - out[i])/2);
120                 if((out[i] - in[i]) & 1)
121                     flag = false;
122             }
123             if(!flag)
124             {
125                 printf("impossible\n");
126                 continue;
127             }
128             sap(0,n+1,n+2);
129             for(int i = head[0]; i != -1;i = edge[i].next)
130                 if(edge[i].cap > 0 && edge[i].cap > edge[i].flow)
131                 {
132                     flag = false;
133                     break;
134                 }
135             if(flag)printf("possible\n");
136             else printf("impossible\n");
137      }
138      return 0;
139 }
\end{Code}

\subsection{Chia để trị trên cây}

\subsubsection{Chia để trị theo điểm - HDU5016}

\begin{Code}
     HDU 5016 给定一个边权树，初始时，一些结点上已经建立了市场。每个结点会被距离自己
     最近的市场所支配（距离相同时，被标号最小的市场支配）可以新建一个市场，问新建的市
     场最多可以支配多少点。
 1   const int MAXN = 100010;
 2   const int INF = 0x3f3f3f3f;
 3   struct Edge{
 4       int to,next,w;
 5   }edge[MAXN*2];
 6   int head[MAXN],tot;


 7   void init(){
 8       tot = 0;
 9       memset(head,-1,sizeof(head));
10   }
11   inline void addedge(int u,int v,int w){
12       edge[tot].to = v;
13       edge[tot].w = w;
14       edge[tot].next = head[u];
15       head[u] = tot++;
16   }
17   int size[MAXN],vis[MAXN],fa[MAXN],que[MAXN];
18   int TT;//时间戳
19   //找重心
20   inline int getroot(int u){
21       int Min = MAXN, root = 0;
22       int l,r;
23       que[l = r = 1] = u;
24       fa[u] = 0;
25       for(;l <= r;l++)
26           for(int i = head[que[l]]; i != -1;i = edge[i].next){
27                int v = edge[i].to;
28                if(v == fa[que[l]] || vis[v] == TT)continue;
29                que[++r] = v;
30                fa[v] = que[l];
31           }
32       for(l--;l;l--){
33           int x = que[l], Max = 0;
34           size[x] = 1;
35           for(int i = head[x];i != -1;i = edge[i].next){
36                int v = edge[i].to;
37                if(v == fa[x] || vis[v] == TT)continue;
38                Max = max(Max,size[v]);
39                size[x] += size[v];
40           }
41           Max = max(Max,r - size[x]);
42           if(Max < Min){
43                Min = Max; root = x;
44           }
45       }
46       return root;
47   }
48
49   int ans[MAXN];
50   pair<int,int>pp[MAXN];
51   pair<int,int>np[MAXN];
52   int dis[MAXN];
53   int type[MAXN];
54   inline void go(int u,int pre,int w,int tt){
55       int l,r;
56       que[l = r = 1] = u;
57       fa[u] = pre; dis[u] = w;


 58       for(;l <= r;l++)
 59           for(int i = head[que[l]];i != -1;i = edge[i].next){
 60               int v = edge[i].to;
 61               if(v == fa[que[l]] || vis[v] == TT)continue;
 62               que[++r] = v;
 63               fa[v] = que[l];
 64               dis[v] = dis[que[l]] + edge[i].w;
 65           }
 66       int cnt = 0;
 67       for(int i = 1;i <= r;i++){
 68           int x = que[i];
 69           pp[cnt++] = make_pair(np[x].first-dis[x],np[x].second);
 70       }
 71       sort(pp,pp+cnt);
 72       for(int i = 1;i <= r;i++){
 73           int x = que[i];
 74           if(type[x])continue;
 75           int id = lower_bound(pp,pp+cnt,make_pair(dis[x],x)) - pp;
 76           ans[x] += (cnt-id)*tt;
 77       }
 78   }
 79   void solve(int u){
 80       int root = getroot(u);
 81       vis[root] = TT;
 82       go(root,0,0,1);
 83       for(int i = head[root];i != -1;i = edge[i].next){
 84           int v = edge[i].to;
 85           if(vis[v] == TT)continue;
 86           go(v,root,edge[i].w,-1);
 87       }
 88       for(int i = head[root];i != -1;i = edge[i].next){
 89           int v = edge[i].to;
 90           if(vis[v] == TT)continue;
 91           solve(v);
 92       }
 93   }
 94   bool ff[MAXN];
 95   int main()
 96   {
 97       int n;
 98       memset(vis,0,sizeof(vis));
 99       TT = 0;
100       while(scanf("%d",&n) == 1){
101           init();
102           int u,v,w;
103           for(int i = 1;i < n;i++){
104               scanf("%d%d%d",&u,&v,&w);
105               addedge(u,v,w);
106               addedge(v,u,w);
107           }
108           for(int i = 1;i <= n;i++)scanf("%d",&type[i]);


109              queue<int>q;
110              for(int i = 1;i <= n;i++){
111                  if(type[i]){
112                      np[i] = make_pair(0,i);
113                      ff[i] = true;
114                      q.push(i);
115                  }
116                  else{
117                      np[i] = make_pair(INF,0);
118                      ff[i] = false;
119                  }
120              }
121              while(!q.empty()){
122                  u = q.front();
123                  q.pop();
124                  ff[u] = false;
125                  for(int i = head[u];i != -1;i = edge[i].next){
126                      v = edge[i].to;
127                      pair<int,int>tmp = make_pair(np[u].first+edge[i].w,
                            np[u].second);
128                      if(tmp < np[v]){
129                           np[v] = tmp;
130                           if(!ff[v]){
131                               ff[v] = true;
132                               q.push(v);
133                           }
134                      }
135                  }
136              }
137              TT++;
138              for(int i = 1;i <= n;i++)ans[i] = 0;
139              solve(1);
140              int ret = 0;
141              for(int i = 1;i <= n;i++)ret = max(ret,ans[i]);
142              printf("%d\n",ret);
143       }
144       return 0;
145 }
\end{Code}

\subsubsection{* Chia để trị theo điểm - HDU4918}

\begin{Code}
      HDU 4918
      题意：给出一颗 n 个点的树，每个点有一个权值，有两种操作，一种是将某个点的权值修改
      为 v，另一种是查询距离点 u 不超过 d 的点的权值和。
 1    const int MAXN = 100010;
 2    const int MAXD = 40;
 3    int cc[MAXN*MAXD];
 4    int *cc_tail; //记得初始化 cc_tail = cc
 5    //0-based BinaryIndexTree
 6    struct BIT{
 7        int *c;


 8       int n;
 9       void init(int _n){
10           n = _n;
11           c = cc_tail;
12           cc_tail = cc_tail + n;
13           memset(c,0,sizeof(int)*n);
14       }
15       void add(int i,int val){
16           while(i < n){
17                c[i] += val;
18                i += ~i & i+1;
19           }
20       }
21       int sum(int i){
22           i = min(i,n-1);
23           int s = 0;
24           while(i >= 0){
25                s += c[i];
26                i -= ~i & i+1;
27           }
28           return s;
29       }
30   }bits[MAXN<<1];
31   struct Edge{
32       int to,next;
33   }edge[MAXN*2];
34   int head[MAXN],tot;
35   void init(){
36       tot = 0;
37       memset(head,-1,sizeof(head));
38   }
39   inline void addedge(int u,int v){
40       edge[tot].to = v;
41       edge[tot].next = head[u];
42       head[u] = tot++;
43   }
44   int size[MAXN],vis[MAXN],fa[MAXN],que[MAXN];
45   int TT;
46   inline int getroot(int u,int &tot){
47       int Min = MAXN, root = 0;
48       int l,r;
49       que[l = r = 1] = u;
50       fa[u] = 0;
51       for(;l <= r;l++)
52           for(int i = head[que[l]];i != -1;i = edge[i].next){
53                int v = edge[i].to;
54                if(v == fa[que[l]] || vis[v] == TT)continue;
55                que[++r] = v;
56                fa[v] = que[l];
57           }
58       tot = r;


 59       for(l--;l;l--){
 60           int x = que[l], Max = 0;
 61           size[x] = 1;
 62           for(int i = head[x];i != -1;i = edge[i].next){
 63               int v = edge[i].to;
 64               if(v == fa[x] || vis[v] == TT)continue;
 65               Max = max(Max,size[v]);
 66               size[x] += size[v];
 67           }
 68           Max = max(Max,r - size[x]);
 69           if(Max < Min){
 70               Min = Max, root = x;
 71           }
 72       }
 73       return root;
 74   }
 75   struct Node{
 76       int root,subtree,dis;
 77       Node(int _root = 0, int _subtree = 0,int _dis = 0){
 78           root = _root;
 79           subtree = _subtree;
 80           dis = _dis;
 81       }
 82   };
 83   vector<Node>vec[MAXN];
 84   int id[MAXN];
 85   int dist[MAXN];
 86   int val[MAXN];
 87   int cnt;
 88   inline void go(int u,int root,int subtree){
 89       int l,r;
 90       que[l = r = 1] = u;
 91       fa[u] = 0; dist[u] = 1;
 92       for(; l <= r;l++){
 93           u = que[l];
 94           vec[u].push_back(Node(id[root],subtree,dist[u]));
 95           for(int i = head[u];i != -1;i = edge[i].next){
 96                int v = edge[i].to;
 97                if(v == fa[u] || vis[v] == TT)continue;
 98                que[++r] = v;
 99                fa[v] = u;
100                dist[v] = dist[u] + 1;
101           }
102       }
103       bits[subtree].init(r+1);
104       for(int i = 1;i <= r;i++){
105           u = que[i];
106           bits[id[root]].add(dist[u],val[u]);
107           bits[subtree].add(dist[u],val[u]);
108       }
109   }


110 void solve(int u){
111     int tot;
112     int root = getroot(u,tot);
113     vis[root] = TT;
114     id[root] = cnt++;
115     vec[root].push_back(Node(id[root],-1,0));
116     bits[id[root]].init(tot);
117     bits[id[root]].add(0,val[root]);
118     for(int i = head[root];i != -1;i = edge[i].next){
119         int v = edge[i].to;
120         if(vis[v] == TT)continue;
121         go(v,root,cnt);
122         cnt++;
123     }
124     for(int i = head[root];i != -1;i = edge[i].next){
125         int v = edge[i].to;
126         if(vis[v] == TT)continue;
127         solve(v);
128     }
129 }
130 int main(){
131     int n,m;
132     memset(vis,0,sizeof(vis));
133     TT = 0;
134     while(scanf("%d%d",&n,&m) == 2){
135         init();
136         TT++;
137         cc_tail = cc;
138         cnt = 0;
139         for(int i = 1;i <= n;i++)vec[i].clear();
140         for(int i = 1;i <= n;i++)scanf("%d",&val[i]);
141         int u,v;
142         for(int i = 1;i < n;i++){
143              scanf("%d%d",&u,&v);
144              addedge(u,v);
145              addedge(v,u);
146         }
147         solve(1);
148         char op[10];
149         int d;
150         while(m--){
151              scanf("%s%d%d",op,&u,&d);
152              if(op[0] == '!'){
153                  int dv = d - val[u];
154                  int sz = vec[u].size();
155                  for(int i = 0;i < sz;i++){
156                      Node tmp = vec[u][i];
157                      bits[tmp.root].add(tmp.dis,dv);
158                      if(tmp.subtree != -1)
159                          bits[tmp.subtree].add(tmp.dis,dv);
160                  }


161                  val[u] += dv;
162              }
163              else {
164                  int ans = 0;
165                  int sz = vec[u].size();
166                  for(int i = 0;i < sz;i++){
167                      Node tmp = vec[u][i];
168                      ans += bits[tmp.root].sum(d-tmp.dis);
169                      if(tmp.subtree != -1)
170                          ans -= bits[tmp.subtree].sum(d-tmp.dis);
171                  }
172                  printf("%d\n",ans);
173              }
174          }
175      }
176      return 0;
177 }
\end{Code}

\subsubsection{Chia để trị theo chuỗi - HDU5039}

\begin{Code}
     HDU 5039 一颗树，每条边的属性为 0 或 1，求有多少条路径经过奇数条属性为 1 的边。
     一种是查询操作，一种是修改边的属性。
     虽然这题有更加简单的方法，但是用来练习链分治还是不错的。
 1   const int MAXN = 30010;
 2   const int INF = 0x3f3f3f3f;
 3   struct Edge{
 4       int to,next;
 5       int f;
 6   }edge[MAXN*2];
 7   int head[MAXN],tot;
 8   void init(){
 9       tot = 0;
10       memset(head,-1,sizeof(head));
11   }
12   void addedge(int u,int v,int f){
13       edge[tot].to = v;
14       edge[tot].next = head[u];
15       edge[tot].f = f;
16       head[u] = tot++;
17   }
18   long long ans;
19   int num0[MAXN],num1[MAXN];
20   long long tnum[MAXN];
21   struct Node{
22       int l0,l1;
23       int r0,r1;
24       int cc;
25       long long sum;
26       Node gao(int u){
27           l0 = r0 = num0[u];
28           l1 = r1 = num1[u];


29              sum = tnum[u];
30              cc = 0;
31              return *this;
32       }
33   };
34   int pos[MAXN];
35   int val[MAXN];
36   int fa[MAXN];
37   int cnt[MAXN];
38   int col[MAXN];
39   int link[MAXN];
40   int CHANGEU;
41   struct chain{
42       vector<int>uu;
43       vector<Node>nde;
44       int n;
45       void init(){
46           n = uu.size();
47           nde.resize(n << 2);
48           for(int i = 0;i < n;i++)pos[uu[i]] = i;
49           build(0,n-1,1);
50       }
51       void up(int l,int r,int p){
52           int mid = (l+r)>>1;
53           nde[p].cc = nde[p<<1].cc ^ nde[(p<<1)|1].cc ^ val[uu[mid]];
54           nde[p].l0 = nde[p<<1].l0;
55           nde[p].l1 = nde[p<<1].l1;
56           if(nde[p<<1].cc^val[uu[mid]]){
57                nde[p].l0 += nde[(p<<1)|1].l1;
58                nde[p].l1 += nde[(p<<1)|1].l0;
59           }
60           else {
61                nde[p].l0 += nde[(p<<1)|1].l0;
62                nde[p].l1 += nde[(p<<1)|1].l1;
63           }
64           nde[p].r0 = nde[(p<<1)|1].r0;
65           nde[p].r1 = nde[(p<<1)|1].r1;
66           if(nde[(p<<1)|1].cc^val[uu[mid]]){
67                nde[p].r0 += nde[p<<1].r1;
68                nde[p].r1 += nde[p<<1].r0;
69           }
70           else {
71                nde[p].r0 += nde[p<<1].r0;
72                nde[p].r1 += nde[p<<1].r1;
73           }
74           if(val[uu[mid]] == 0){
75                nde[p].sum = nde[p<<1].sum + nde[(p<<1)|1].sum +
76                    (long long)nde[p<<1].r0*nde[(p<<1)|1].l1 +
77                    (long long)nde[p<<1].r1*nde[(p<<1)|1].l0;
78           }
79           else {


 80               nde[p].sum = nde[p<<1].sum + nde[(p<<1)|1].sum +
 81                   (long long)nde[p<<1].r0*nde[(p<<1)|1].l0 +
 82                   (long long)nde[p<<1].r1*nde[(p<<1)|1].l1;
 83           }
 84       }
 85       void build(int l,int r,int p){
 86           if(l == r){
 87               nde[p].gao(uu[l]);
 88               return;
 89           }
 90           int mid = (l+r)/2;
 91           build(l,mid,p<<1);
 92           build(mid+1,r,(p<<1)|1);
 93           up(l,r,p);
 94       }
 95       void update(int k,int l,int r,int p){
 96           if(l == r){
 97               nde[p].gao(uu[k]);
 98               return;
 99           }
100           int mid = (l+r)/2;
101           if(k <= mid)update(k,l,mid,p<<1);
102           else update(k,mid+1,r,(p<<1)|1);
103           up(l,r,p);
104       }
105       int change(int y){
106           int x = uu.back();
107           int p = fa[x];
108           if(p){
109               if(x == CHANGEU)val[x] ^= 1;
110               if(val[x]){
111                   tnum[p] -= (long long)nde[1].r0*(num0[p]-nde[1].r1)
                         ;
112                   tnum[p] -= (long long)nde[1].r1*(num1[p]-nde[1].r0)
                         ;
113                   num0[p] -= nde[1].r1;
114                   num1[p] -= nde[1].r0;
115               }
116               else {
117                   tnum[p] -= (long long)nde[1].r1*(num0[p]-nde[1].r0)
                         ;
118                   tnum[p] -= (long long)nde[1].r0*(num1[p]-nde[1].r1)
                         ;
119                   num0[p] -= nde[1].r0;
120                   num1[p] -= nde[1].r1;
121               }
122               if(x == CHANGEU)val[x] ^= 1;
123           }
124           ans -= nde[1].sum;
125           update(pos[y],0,n-1,1);
126           if(p){


127                  if(val[x]){
128                      tnum[p] += (long long)nde[1].r0*num0[p];
129                      tnum[p] += (long long)nde[1].r1*num1[p];
130                      num0[p] += nde[1].r1;
131                      num1[p] += nde[1].r0;
132                  }
133                  else {
134                      tnum[p] += (long long)nde[1].r0*num1[p];
135                      tnum[p] += (long long)nde[1].r1*num0[p];
136                      num0[p] += nde[1].r0;
137                      num1[p] += nde[1].r1;
138                  }
139              }
140              ans += nde[1].sum;
141              return p;
142       }
143   }ch[MAXN];
144   void dfs1(int u,int pre){
145       chain &c = ch[u];
146       c.uu.clear();
147       int v, x = 0;
148       cnt[u] = 1;
149       for(int i = head[u];i != -1;i = edge[i].next){
150           v = edge[i].to;
151           if(v == pre)continue;
152           dfs1(v,u);
153           link[i/2] = v;
154           val[v] = edge[i].f;
155           cnt[u] += cnt[v];
156           fa[v] = u;
157           if(cnt[v] > cnt[x]) x = v;
158       }
159       if(!x)col[u] = u;
160       else col[u] = col[x];
161       ch[col[u]].uu.push_back(u);
162       num0[u] = 1;
163       num1[u] = 0;
164       tnum[u] = 0;
165
166   }
167   void dfs2(int x){
168       x = col[x];
169       chain &c = ch[x];
170       int n = c.uu.size();
171       int u,v;
172       for(int i = 1;i < n;i++){
173           u = c.uu[i];
174           for(int j = head[u];j != -1;j = edge[j].next){
175                v = edge[j].to;
176                if(v == c.uu[i-1] || fa[u] == v)continue;
177                dfs2(v);


178                  if(val[v]){
179                      tnum[u] += (long long)num0[u]*ch[col[v]].nde[1].r0
                            + (long long)num1[u]*ch[col[v]].nde[1].r1;
180                      num0[u] += ch[col[v]].nde[1].r1;
181                      num1[u] += ch[col[v]].nde[1].r0;
182                  }
183                  else {
184                      tnum[u] += (long long)num1[u]*ch[col[v]].nde[1].r0
                            + (long long)num0[u]*ch[col[v]].nde[1].r1;
185                      num0[u] += ch[col[v]].nde[1].r0;
186                      num1[u] += ch[col[v]].nde[1].r1;
187                  }
188              }
189       }
190       c.init();
191       ans += c.nde[1].sum;
192   }
193   char str[100];
194   char str1[100],str2[100];
195   int main()
196   {
197       int T;
198       int iCase = 0;
199       scanf("%d",&T);
200       int n;
201       while(T--){
202           ans = 0;
203           iCase++;
204           scanf("%d",&n);
205           map<string,int>mp;
206           init();
207           for(int i = 1;i <= n;i++){
208               scanf("%s",str);
209               mp[str] = i;
210           }
211           int u,v,f;
212           for(int i = 1;i < n;i++){
213               scanf("%s%s%d",str1,str2,&f);
214               addedge(mp[str1],mp[str2],f);
215               addedge(mp[str2],mp[str1],f);
216           }
217           int Q;
218           char op[10];
219           scanf("%d",&Q);
220           printf("Case #%d:\n",iCase);
221           val[1] = 0;
222           fa[1] = 0;
223           dfs1(1,1);
224           dfs2(1);
225           while(Q--){
226               scanf("%s",op);


227             if(op[0] == 'Q'){
228                 printf("%I64d\n",ans*2);
229             }
230             else {
231                 int id ;
232                 scanf("%d",&id);
233                 id--;
234                 u = link[id];
235                 val[u] ^= 1;
236                 CHANGEU = u;
237                 while(u)
238                     u = ch[col[u]].change(u);
239             }
240         }
241     }
242     return 0;
243 }
\end{Code}

\section{Tìm kiếm}

\subsection{Dancing Links}

\subsubsection{Phủ chính xác}

\begin{Code}
 1   /*
 2    * POJ3074
 3    */
 4   const int N = 9; //3*3 数独
 5   const int MaxN = N*N*N + 10;
 6   const int MaxM = N*N*4 + 10;
 7   const int maxnode = MaxN*4 + MaxM + 10;
 8   char g[MaxN];
 9   struct DLX{
10       int n,m,size;
11       int U[maxnode],D[maxnode],R[maxnode],L[maxnode],Row[maxnode],
            Col[maxnode];
12       int H[MaxN],S[MaxM];
13       int ansd,ans[MaxN];
14       void init(int _n,int _m){
15           n = _n;
16           m = _m;
17           for(int i = 0;i <= m;i++){
18               S[i] = 0;
19               U[i] = D[i] = i;
20               L[i] = i-1;
21               R[i] = i+1;
22           }
23           R[m] = 0; L[0] = m;
24           size = m;
25           for(int i = 1;i <= n;i++)H[i] = -1;
26       }
27       void Link(int r,int c){
28           ++S[Col[++size]=c];
29           Row[size] = r;
30           D[size] = D[c];
31           U[D[c]] = size;
32           U[size] = c;
33           D[c] = size;
34           if(H[r] < 0)H[r] = L[size] = R[size] = size;
35           else{
36               R[size] = R[H[r]];
37               L[R[H[r]]] = size;
38               L[size] = H[r];
39               R[H[r]] = size;
40           }
41       }
42       void remove(int c){
43           L[R[c]] = L[c]; R[L[c]] = R[c];
44           for(int i = D[c];i != c;i = D[i])
45               for(int j = R[i];j != i;j = R[j]){


46                  U[D[j]] = U[j];
47                  D[U[j]] = D[j];
48                  --S[Col[j]];
49               }
50       }
51       void resume(int c){
52           for(int i = U[c];i != c;i = U[i])
53               for(int j = L[i];j != i;j = L[j])
54                   ++S[Col[U[D[j]]=D[U[j]]=j]];
55           L[R[c]] = R[L[c]] = c;
56       }
57       bool Dance(int d){
58           if(R[0] == 0){
59               for(int i = 0;i < d;i++)g[(ans[i]-1)/9] = (ans[i]-1)%9
                    + '1';
60               for(int i = 0;i < N*N;i++)printf("%c",g[i]);
61               printf("\n");
62               return true;
63           }
64           int c = R[0];
65           for(int i = R[0];i != 0;i = R[i])
66               if(S[i] < S[c])
67                   c = i;
68           remove(c);
69           for(int i = D[c];i != c;i = D[i]){
70               ans[d] = Row[i];
71               for(int j = R[i];j != i;j = R[j])remove(Col[j]);
72               if(Dance(d+1))return true;
73               for(int j = L[i];j != i;j = L[j])resume(Col[j]);
74           }
75           resume(c);
76           return false;
77       }
78 };
79 void place(int &r,int &c1,int &c2,int &c3,int &c4,int i,int j,int k
      ){
80     r = (i*N+j)*N + k; c1 = i*N+j+1; c2 = N*N+i*N+k;
81     c3 = N*N*2+j*N+k; c4 = N*N*3+((i/3)*3+(j/3))*N+k;
82 }
83 DLX dlx;
84 int main(){
85     while(scanf("%s",g) == 1){
86          if(strcmp(g,"end") == 0)break;
87          dlx.init(N*N*N,N*N*4);
88          int r,c1,c2,c3,c4;
89          for(int i = 0;i < N;i++)
90              for(int j = 0;j < N;j++)
91                  for(int k = 1;k <= N;k++)
92                      if(g[i*N+j] == '.' || g[i*N+j] == '0'+k){
93                          place(r,c1,c2,c3,c4,i,j,k);
94                          dlx.Link(r,c1);


 95                             dlx.Link(r,c2);
 96                             dlx.Link(r,c3);
 97                             dlx.Link(r,c4);
 98                         }
 99             dlx.Dance(0);
100      }
101      return 0;
102 }
\end{Code}

\subsubsection{Phủ lặp}

\begin{Code}
 1   /*
 2    * FZU1686
 3    */
 4   const int MaxM = 15*15+10;
 5   const int MaxN = 15*15+10;
 6   const int maxnode = MaxN * MaxM;
 7   const int INF = 0x3f3f3f3f;
 8   struct DLX{
 9       int n,m,size;
10       int U[maxnode],D[maxnode],R[maxnode],L[maxnode],Row[maxnode],
            Col[maxnode];
11       int H[MaxN],S[MaxM];
12       int ansd;
13       void init(int _n,int _m){
14           n = _n;
15           m = _m;
16           for(int i = 0;i <= m;i++){
17               S[i] = 0;
18               U[i] = D[i] = i;
19               L[i] = i-1;
20               R[i] = i+1;
21           }
22           R[m] = 0; L[0] = m;
23           size = m;
24           for(int i = 1;i <= n;i++)H[i] = -1;
25       }
26       void Link(int r,int c){
27           ++S[Col[++size]=c];
28           Row[size] = r;
29           D[size] = D[c];
30           U[D[c]] = size;
31           U[size] = c;
32           D[c] = size;
33           if(H[r] < 0)H[r] = L[size] = R[size] = size;
34           else{
35               R[size] = R[H[r]];
36               L[R[H[r]]] = size;
37               L[size] = H[r];
38               R[H[r]] = size;
39           }
40       }


41       void remove(int c){
42           for(int i = D[c];i != c;i = D[i])
43                L[R[i]] = L[i], R[L[i]] = R[i];
44       }
45       void resume(int c){
46           for(int i = U[c];i != c;i = U[i])
47                L[R[i]] = R[L[i]] = i;
48       }
49       bool v[MaxM];
50       int f(){
51           int ret = 0;
52           for(int c = R[0]; c != 0;c = R[c])v[c] = true;
53           for(int c = R[0]; c != 0;c = R[c])
54                if(v[c])
55                {
56                    ret++;
57                    v[c] = false;
58                    for(int i = D[c];i != c;i = D[i])
59                         for(int j = R[i];j != i;j = R[j])
60                             v[Col[j]] = false;
61                }
62           return ret;
63       }
64       void Dance(int d){
65           if(d + f() >= ansd)return;
66           if(R[0] == 0){
67                if(d < ansd)ansd = d;
68                return;
69           }
70           int c = R[0];
71           for(int i = R[0];i != 0;i = R[i])
72                if(S[i] < S[c])
73                    c = i;
74           for(int i = D[c];i != c;i = D[i]){
75                remove(i);
76                for(int j = R[i];j != i;j = R[j])remove(j);
77                Dance(d+1);
78                for(int j = L[i];j != i;j = L[j])resume(j);
79                resume(i);
80           }
81       }
82   };
83   DLX g;
84   int a[20][20];
85   int id[20][20];
86   int main(){
87       int n,m;
88       while(scanf("%d%d",&n,&m) == 2){
89           int sz = 0;
90           memset(id,0,sizeof(id));
91           for(int i = 0;i < n;i++)


 92                 for(int j = 0;j < m;j++){
 93                     scanf("%d",&a[i][j]);
 94                     if(a[i][j] == 1)id[i][j] = (++sz);
 95                 }
 96             g.init(n*m,sz);
 97             sz = 1;
 98             int n1,m1;
 99             scanf("%d%d",&n1,&m1);
100             for(int i = 0;i < n;i++)
101                 for(int j = 0;j < m;j++){
102                     for(int x = 0;x < n1 && i + x < n;x++)
103                         for(int y = 0;y < m1 && j + y < m;y++)
104                             if(id[i+x][j+y])
105                                 g.Link(sz,id[i+x][j+y]);
106                     sz++;
107                 }
108             g.ansd = INF;
109             g.Dance(0);
110             printf("%d\n",g.ansd);
111        }
112        return 0;
113 }
\end{Code}

\subsection{Bài toán tám ô số}

\subsubsection{HDU1043 tìm kiếm ngược}

\begin{Code}
 1   /*
 2   HDU 1043 Eight
 3   八数码，输出路径
 4   思路：反向搜索，从目标状态找回状态对应的路径
 5   用康托展开判重
 6   */
 7   const int MAXN=1000000;//最多是 9!/2
 8   int fac[]={1,1,2,6,24,120,720,5040,40320,362880};//康拖展开判重
 9   //         0!1!2!3! 4! 5! 6! 7!     8!    9!
10   bool vis[MAXN];//标记
11   string path[MAXN];//记录路径
12   //康拖展开求该序列的 hash 值
13   int cantor(int s[]){
14       int sum=0;
15       for(int i=0;i<9;i++){
16           int num=0;
17           for(int j=i+1;j<9;j++)
18             if(s[j]<s[i])num++;
19           sum+=(num*fac[9-i-1]);
20       }
21       return sum+1;
22   }
23   struct Node{
24       int s[9];
25       int loc;//'0' 的位置


26       int status;//康拖展开的 hash 值
27       string path;//路径
28   };
29   int move[4][2]={{-1,0},{1,0},{0,-1},{0,1}};//u,d,l,r
30   char indexs[5]="durl";//和上面的要相反，因为是反向搜索
31   int aim=46234;//123456780 对应的康拖展开的 hash 值
32   void bfs(){
33       memset(vis,false,sizeof(vis));
34       Node cur,next;
35       for(int i=0;i<8;i++)cur.s[i]=i+1;
36       cur.s[8]=0;
37       cur.loc=8;
38       cur.status=aim;
39       cur.path="";
40       queue<Node>q;
41       q.push(cur);
42       path[aim]="";
43       while(!q.empty()){
44           cur=q.front();
45           q.pop();
46           int x=cur.loc/3;
47           int y=cur.loc%3;
48           for(int i=0;i<4;i++){
49               int tx=x+move[i][0];
50               int ty=y+move[i][1];
51               if(tx<0||tx>2||ty<0||ty>2)continue;
52               next=cur;
53               next.loc=tx*3+ty;
54               next.s[cur.loc]=next.s[next.loc];
55               next.s[next.loc]=0;
56               next.status=cantor(next.s);
57               if(!vis[next.status]){
58                    vis[next.status]=true;
59                    next.path=indexs[i]+next.path;
60                    q.push(next);
61                    path[next.status]=next.path;
62               }
63           }
64       }
65
66   }
67   int main(){
68       char ch;
69       Node cur;
70       bfs();
71       while(cin>>ch){
72           if(ch=='x') {cur.s[0]=0;cur.loc=0;}
73           else cur.s[0]=ch-'0';
74           for(int i=1;i<9;i++){
75                cin>>ch;
76                if(ch=='x'){


77                    cur.s[i]=0;
78                    cur.loc=i;
79                }
80                else cur.s[i]=ch-'0';
81            }
82            cur.status=cantor(cur.s);
83            if(vis[cur.status]){
84                cout<<path[cur.status]<<endl;
85            }
86            else cout<<"unsolvable"<<endl;
87     }
88     return 0;
89 }
\end{Code}

\section{Quy hoạch động}

\subsection{Dãy con tăng dài nhất O(nlogn)}

\begin{Code}
 1 const int MAXN=500010;
 2 int a[MAXN],b[MAXN];
 3
 4 //用二分查找的方法找到一个位置，使得 num>b[i-1] 并且 num<b[i], 并用 num 代
      替 b[i]
 5 int Search(int num,int low,int high){
 6     int mid;
 7     while(low<=high){
 8         mid=(low+high)/2;
 9         if(num>=b[mid]) low=mid+1;
10         else    high=mid-1;
11     }
12     return low;
13 }
14 int DP(int n){
15     int i,len,pos;
16     b[1]=a[1];
17     len=1;
18     for(i=2;i<=n;i++){
19         if(a[i]>=b[len])//如果 a[i] 比 b[] 数组中最大还大直接插入到后面即
              可
20         {
21              len=len+1;
22              b[len]=a[i];
23         }
24         else//用二分的方法在 b[] 数组中找出第一个比 a[i] 大的位置并且让
              a[i] 替代这个位置
25         {
26              pos=Search(a[i],1,len);
27              b[pos]=a[i];
28         }
29     }
30     return len;
31 }
\end{Code}

\subsection{Balo}

\begin{Code}
 1   int nValue,nKind;
 2
 3   //0-1 背包，代价为 cost, 获得的价值为 weight
 4   void ZeroOnePack(int cost,int weight){
 5       for(int i=nValue;i>=cost;i--)
 6         dp[i]=max(dp[i],dp[i-cost]+weight);
 7   }
 8
 9   //完全背包，代价为 cost, 获得的价值为 weight
10   void CompletePack(int cost,int weight){
11       for(int i=cost;i<=nValue;i++)


12       dp[i]=max(dp[i],dp[i-cost]+weight);
13 }
14
15 //多重背包
16 void MultiplePack(int cost,int weight,int amount){
17     if(cost*amount>=nValue) CompletePack(cost,weight);
18     else{
19         int k=1;
20         while(k<amount){
21             ZeroOnePack(k*cost,k*weight);
22             amount-=k;
23             k<<=1;
24         }
25         ZeroOnePack(amount*cost,amount*weight);//这个不要忘记了，经常
              掉了
26     }
27 }
     分组背包：
     for k = 1 to K
     for v = V to 0
     for item i in group k
     F[v] = maxF[v],F[v-Ci ]+Wi
\end{Code}

\subsection{Plug DP}

\subsubsection{HDU 4285}

\paragraph{Ghi chú.} Đếm số cách tạo đúng K chu trình và không cho phép chu trình lồng nhau. Thêm một cờ để ghi số chu trình đã tạo; khi tạo chu trình mới, cần bảo đảm số plug ở hai phía là chẵn để tránh trường hợp vòng lồng vòng.

\begin{Code}
 1   /*
 2   HDU 4285
 3   要形成刚好 K 条回路的方法数
 4   要避免环套环的情况。
 5   所以形成回路时，要保证两边的插头数是偶数
 6   G++ 11265ms 11820K
 7   C++ 10656ms 11764K
 8   */
 9   const int MAXD=15;
10   const int STATE=1000010;
11   const int HASH=300007;//这个大一点可以防止 TLE, 但是容易 MLE
12   const int MOD=1000000007;
13   int N,M,K;
14   int maze[MAXD][MAXD];
15   int code[MAXD];
16   int ch[MAXD];
17   int num;//圈的个数
18   struct HASHMAP
19   {
20       int head[HASH],next[STATE],size;
21       long long state[STATE];


22       int f[STATE];
23       void init()
24       {
25           size=0;
26           memset(head,-1,sizeof(head));
27       }
28       void push(long long st,int ans)
29       {
30           int i;
31           int h=st%HASH;
32           for(i=head[h];i!=-1;i=next[i])
33             if(state[i]==st)
34             {
35                  f[i]+=ans;
36                  f[i]%=MOD;
37                  return;
38             }
39           state[size]=st;
40           f[size]=ans;
41           next[size]=head[h];
42           head[h]=size++;
43       }
44   }hm[2];
45   void decode(int *code,int m,long long st)
46   {
47       num=st&63;
48       st>>=6;
49       for(int i=m;i>=0;i--)
50       {
51           code[i]=st&7;
52           st>>=3;
53       }
54   }
55   long long encode(int *code,int m)//最小表示法
56   {
57       int cnt=1;
58       memset(ch,-1,sizeof(ch));
59       ch[0]=0;
60       long long st=0;
61       for(int i=0;i<=m;i++)
62       {
63           if(ch[code[i]]==-1)ch[code[i]]=cnt++;
64           code[i]=ch[code[i]];
65           st<<=3;
66           st|=code[i];
67       }
68       st<<=6;
69       st|=num;
70       return st;
71   }
72   void shift(int *code,int m)


73    {
74        for(int i=m;i>0;i--)code[i]=code[i-1];
75        code[0]=0;
76    }
77    void dpblank(int i,int j,int cur)
78    {
79        int k,left,up;
80        for(k=0;k<hm[cur].size;k++)
81        {
82            decode(code,M,hm[cur].state[k]);
83            left=code[j-1];
84            up=code[j];
85            if(left&&up)
86            {
87                if(left==up)
88                {
89                    if(num>=K)continue;
90                    int t=0;
91                    //要避免环套环的情况，需要两边插头数为偶数
92                    for(int p=0;p<j-1;p++)
93                       if(code[p])t++;
94                    if(t&1)continue;
95                    if(num<K)
96                    {
97                         num++;
98                         code[j-1]=code[j]=0;
99                         hm[cur^1].push(encode(code,j==M?M-1:M),hm[cur].
                              f[k]);
100                   }
101               }
102               else
103               {
104                   code[j-1]=code[j]=0;
105                   for(int t=0;t<=M;t++)
106                      if(code[t]==up)
107                        code[t]=left;
108                   hm[cur^1].push(encode(code,j==M?M-1:M),hm[cur].f[k
                          ]);
109               }
110           }
111           else if(left||up)
112           {
113               int t;
114               if(left)t=left;
115               else t=up;
116               if(maze[i][j+1])
117               {
118                   code[j-1]=0;
119                   code[j]=t;
120                   hm[cur^1].push(encode(code,M),hm[cur].f[k]);
121               }


122                     if(maze[i+1][j])
123                     {
124                         code[j]=0;
125                         code[j-1]=t;
126                         hm[cur^1].push(encode(code,j==M?M-1:M),hm[cur].f[k
                               ]);
127                     }
128              }
129              else
130              {
131                     if(maze[i][j+1]&&maze[i+1][j])
132                     {
133                         code[j-1]=code[j]=13;
134                         hm[cur^1].push(encode(code,j==M?M-1:M),hm[cur].f[k
                               ]);
135                     }
136              }
137       }
138   }
139   void dpblock(int i,int j,int cur)
140   {
141       int k;
142       for(k=0;k<hm[cur].size;k++)
143       {
144           decode(code,M,hm[cur].state[k]);
145           code[j-1]=code[j]=0;
146           hm[cur^1].push(encode(code,j==M?M-1:M),hm[cur].f[k]);
147       }
148   }
149   char str[20];
150   void init()
151   {
152       scanf("%d%d%d",&N,&M,&K);
153       memset(maze,0,sizeof(maze));
154       for(int i=1;i<=N;i++)
155       {
156           scanf("%s",&str);
157           for(int j=1;j<=M;j++)
158              if(str[j-1]=='.')
159                maze[i][j]=1;
160       }
161   }
162   void solve()
163   {
164       int i,j,cur=0;
165       hm[cur].init();
166       hm[cur].push(0,1);
167       for(i=1;i<=N;i++)
168         for(j=1;j<=M;j++)
169         {
170              hm[cur^1].init();


171              if(maze[i][j])dpblank(i,j,cur);
172              else dpblock(i,j,cur);
173              cur^=1;
174          }
175        int ans=0;
176        for(i=0;i<hm[cur].size;i++)
177          if(hm[cur].state[i]==K)
178          {
179              ans+=hm[cur].f[i];
180              ans%=MOD;
181          }
182        printf("%d\n",ans);
183
184   }
185   int main()
186   {
187       int T;
188       scanf("%d",&T);
189       while(T--)
190       {
191           init();
192           solve();
193       }
194       return 0;
195   }
196
197   /*
198   Sample Input
199   4 4 1
200   **..
201   ....
202   ....
203   ....
204   4 1
205   ....
206   ....
207   ....
208   ....
209
210   Sample Output
211   6
212
213   */
\end{Code}

\section{Hình học tính toán}

\subsection{Hình học hai chiều}

\begin{Code}
 1   // 计算几何模板
 2   const double eps = 1e-8;
 3   const double inf = 1e20;
 4   const double pi = acos(-1.0);
 5   const int maxp = 1010;
 6   //Compares a double to zero
 7   int sgn(double x){
 8        if(fabs(x) < eps)return 0;
 9        if(x < 0)return -1;
10        else return 1;
11   }
12   //square of a double
13   inline double sqr(double x){return x*x;}
14   /*
15     * Point
16     * Point()               - Empty constructor
17     * Point(double _x,double _y) - constructor
18     * input()             - double input
19     * output()            - %.2f output
20     * operator ==         - compares x and y
21     * operator <          - compares first by x, then by y
22     * operator -          - return new Point after subtracting
          curresponging x and y
23     * operator ^          - cross product of 2d points
24     * operator *          - dot product
25     * len()               - gives length from origin
26     * len2()              - gives square of length from origin
27     * distance(Point p)   - gives distance from p
28     * operator + Point b - returns new Point after adding
          curresponging x and y
29     * operator * double k - returns new Point after multiplieing x and
           y by k
30     * operator / double k - returns new Point after divideing x and y
          by k
31     * rad(Point a,Point b)- returns the angle of Point a and Point b
          from this Point
32     * trunc(double r)     - return Point that if truncated the
          distance from center to r
33     * rotleft()           - returns 90 degree ccw rotated point
34     * rotright()          - returns 90 degree cw rotated point
35     * rotate(Point p,double angle) - returns Point after rotateing the
           Point centering at p by angle radian ccw
36     */
37   struct Point{
38        double x,y;
39        Point(){}
40        Point(double _x,double _y){


41              x = _x;
42              y = _y;
43       }
44       void input(){
45           scanf("%lf%lf",&x,&y);
46       }
47       void output(){
48           printf("%.2f %.2f\n",x,y);
49       }
50       bool operator == (Point b)const{
51           return sgn(x-b.x) == 0 && sgn(y-b.y) == 0;
52       }
53       bool operator < (Point b)const{
54           return sgn(x-b.x)== 0?sgn(y-b.y)<0:x<b.x;
55       }
56       Point operator -(const Point &b)const{
57           return Point(x-b.x,y-b.y);
58       }
59       //叉积
60       double operator ^(const Point &b)const{
61           return x*b.y - y*b.x;
62       }
63       //点积
64       double operator *(const Point &b)const{
65           return x*b.x + y*b.y;
66       }
67       //返回长度
68       double len(){
69           return hypot(x,y);//库函数
70       }
71       //返回长度的平方
72       double len2(){
73           return x*x + y*y;
74       }
75       //返回两点的距离
76       double distance(Point p){
77           return hypot(x-p.x,y-p.y);
78       }
79       Point operator +(const Point &b)const{
80           return Point(x+b.x,y+b.y);
81       }
82       Point operator *(const double &k)const{
83           return Point(x*k,y*k);
84       }
85       Point operator /(const double &k)const{
86           return Point(x/k,y/k);
87       }
88       //计算 pa 和 pb 的夹角
89       //就是求这个点看 a,b 所成的夹角
90       //测试 LightOJ1203
91       double rad(Point a,Point b){


 92        Point p = *this;
 93        return fabs(atan2( fabs((a-p)^(b-p)),(a-p)*(b-p) ));
 94    }
 95    //化为长度为 r 的向量
 96    Point trunc(double r){
 97        double l = len();
 98        if(!sgn(l))return *this;
 99        r /= l;
100        return Point(x*r,y*r);
101    }
102    //逆时针旋转 90 度
103    Point rotleft(){
104        return Point(-y,x);
105    }
106    //顺时针旋转 90 度
107    Point rotright(){
108        return Point(y,-x);
109    }
110    //绕着 p 点逆时针旋转 angle
111    Point rotate(Point p,double angle){
112        Point v = (*this) - p;
113        double c = cos(angle), s = sin(angle);
114        return Point(p.x + v.x*c - v.y*s,p.y + v.x*s + v.y*c);
115    }
116 };
117 /*
118 * Stores two points
119 * Line()                         - Empty constructor
120 * Line(Point _s,Point _e)        - Line through _s and _e
121 * operator ==                    - checks if two points are same
122 * Line(Point p,double angle)     - one end p , another end at
       angle degree
123 * Line(double a,double b,double c) - Line of equation ax + by + c
       = 0
124 * input()                        - inputs s and e
125 * adjust()                       - orders in such a way that s < e
126 * length()                       - distance of se
127 * angle()                        - return 0 <= angle < pi
128 * relation(Point p)              - 3 if point is on line
129 *                                  1 if point on the left of line
130 *                                  2 if point on the right of line
131 * pointonseg(double p)           - return true if point on segment
132 * parallel(Line v)               - return true if they are
       parallel
133 * segcrossseg(Line v)            - returns 0 if does not intersect
134 *                                  returns 1 if non-standard
       intersection
135 *                                  returns 2 if intersects
136 * linecrossseg(Line v)           - line and seg
137 * linecrossline(Line v)          - 0 if parallel
138 *                                  1 if coincides


139    *                                   2 if intersects
140    * crosspoint(Line v)              - returns intersection point
141    * dispointtoline(Point p)         - distance from point p to the
          line
142    * dispointtoseg(Point p)          - distance from p to the segment
143    * dissegtoseg(Line v)             - distance of two segment
144    * lineprog(Point p)               - returns projected point p on se
           line
145    * symmetrypoint(Point p)          - returns reflection point of p
          over se
146    *
147    */
148   struct Line{
149       Point s,e;
150       Line(){}
151       Line(Point _s,Point _e){
152            s = _s;
153            e = _e;
154       }
155       bool operator ==(Line v){
156            return (s == v.s)&&(e == v.e);
157       }
158       //根据一个点和倾斜角 angle 确定直线,0<=angle<pi
159       Line(Point p,double angle){
160            s = p;
161            if(sgn(angle-pi/2) == 0){
162                e = (s + Point(0,1));
163            }
164            else{
165                e = (s + Point(1,tan(angle)));
166            }
167       }
168       //ax+by+c=0
169       Line(double a,double b,double c){
170            if(sgn(a) == 0){
171                s = Point(0,-c/b);
172                e = Point(1,-c/b);
173            }
174            else if(sgn(b) == 0){
175                s = Point(-c/a,0);
176                e = Point(-c/a,1);
177            }
178            else{
179                s = Point(0,-c/b);
180                e = Point(1,(-c-a)/b);
181            }
182       }
183       void input(){
184            s.input();
185            e.input();
186       }


187       void adjust(){
188           if(e < s)swap(s,e);
189       }
190       //求线段长度
191       double length(){
192           return s.distance(e);
193       }
194       //返回直线倾斜角 0<=angle<pi
195       double angle(){
196           double k = atan2(e.y-s.y,e.x-s.x);
197           if(sgn(k) < 0)k += pi;
198           if(sgn(k-pi) == 0)k -= pi;
199           return k;
200       }
201       //点和直线关系
202       //1 在左侧
203       //2 在右侧
204       //3 在直线上
205       int relation(Point p){
206           int c = sgn((p-s)^(e-s));
207           if(c < 0)return 1;
208           else if(c > 0)return 2;
209           else return 3;
210       }
211       // 点在线段上的判断
212       bool pointonseg(Point p){
213           return sgn((p-s)^(e-s)) == 0 && sgn((p-s)*(p-e)) <= 0;
214       }
215       //两向量平行 (对应直线平行或重合)
216       bool parallel(Line v){
217           return sgn((e-s)^(v.e-v.s)) == 0;
218       }
219       //两线段相交判断
220       //2 规范相交
221       //1 非规范相交
222       //0 不相交
223       int segcrossseg(Line v){
224           int d1 = sgn((e-s)^(v.s-s));
225           int d2 = sgn((e-s)^(v.e-s));
226           int d3 = sgn((v.e-v.s)^(s-v.s));
227           int d4 = sgn((v.e-v.s)^(e-v.s));
228           if( (d1^d2)==-2 && (d3^d4)==-2 )return 2;
229           return (d1==0 && sgn((v.s-s)*(v.s-e))<=0) ||
230               (d2==0 && sgn((v.e-s)*(v.e-e))<=0) ||
231               (d3==0 && sgn((s-v.s)*(s-v.e))<=0) ||
232               (d4==0 && sgn((e-v.s)*(e-v.e))<=0);
233       }
234       //直线和线段相交判断
235       //-*this line -v seg
236       //2 规范相交
237       //1 非规范相交


238       //0 不相交
239       int linecrossseg(Line v){
240           int d1 = sgn((e-s)^(v.s-s));
241           int d2 = sgn((e-s)^(v.e-s));
242           if((d1^d2)==-2) return 2;
243           return (d1==0||d2==0);
244       }
245       //两直线关系
246       //0 平行
247       //1 重合
248       //2 相交
249       int linecrossline(Line v){
250           if((*this).parallel(v))
251               return v.relation(s)==3;
252           return 2;
253       }
254       //求两直线的交点
255       //要保证两直线不平行或重合
256       Point crosspoint(Line v){
257           double a1 = (v.e-v.s)^(s-v.s);
258           double a2 = (v.e-v.s)^(e-v.s);
259           return Point((s.x*a2-e.x*a1)/(a2-a1),(s.y*a2-e.y*a1)/(a2-a1
                 ));
260       }
261       //点到直线的距离
262       double dispointtoline(Point p){
263           return fabs((p-s)^(e-s))/length();
264       }
265       //点到线段的距离
266       double dispointtoseg(Point p){
267           if(sgn((p-s)*(e-s))<0 || sgn((p-e)*(s-e))<0)
268               return min(p.distance(s),p.distance(e));
269           return dispointtoline(p);
270       }
271       //返回线段到线段的距离
272       //前提是两线段不相交，相交距离就是 0 了
273       double dissegtoseg(Line v){
274           return min(min(dispointtoseg(v.s),dispointtoseg(v.e)),min(v
                 .dispointtoseg(s),v.dispointtoseg(e)));
275       }
276       //返回点 p 在直线上的投影
277       Point lineprog(Point p){
278           return s + ( ((e-s)*((e-s)*(p-s)))/((e-s).len2()) );
279       }
280       //返回点 p 关于直线的对称点
281       Point symmetrypoint(Point p){
282           Point q = lineprog(p);
283           return Point(2*q.x-p.x,2*q.y-p.y);
284       }
285 };
286 //圆


287 struct circle{
288     Point p;//圆心
289     double r;//半径
290     circle(){}
291     circle(Point _p,double _r){
292         p = _p;
293         r = _r;
294     }
295     circle(double x,double y,double _r){
296         p = Point(x,y);
297         r = _r;
298     }
299     //三角形的外接圆
300     //需要 Point 的 + / rotate() 以及 Line 的 crosspoint()
301     //利用两条边的中垂线得到圆心
302     //测试：UVA12304
303     circle(Point a,Point b,Point c){
304         Line u = Line((a+b)/2,((a+b)/2)+((b-a).rotleft()));
305         Line v = Line((b+c)/2,((b+c)/2)+((c-b).rotleft()));
306         p = u.crosspoint(v);
307         r = p.distance(a);
308     }
309     //三角形的内切圆
310     //参数 bool t 没有作用，只是为了和上面外接圆函数区别
311     //测试：UVA12304
312     circle(Point a,Point b,Point c,bool t){
313         Line u,v;
314         double m = atan2(b.y-a.y,b.x-a.x), n = atan2(c.y-a.y,c.x-a.
               x);
315         u.s = a;
316         u.e = u.s + Point(cos((n+m)/2),sin((n+m)/2));
317         v.s = b;
318         m = atan2(a.y-b.y,a.x-b.x) , n = atan2(c.y-b.y,c.x-b.x);
319         v.e = v.s + Point(cos((n+m)/2),sin((n+m)/2));
320         p = u.crosspoint(v);
321         r = Line(a,b).dispointtoseg(p);
322     }
323     //输入
324     void input(){
325         p.input();
326         scanf("%lf",&r);
327     }
328     //输出
329     void output(){
330         printf("%.2lf %.2lf %.2lf\n",p.x,p.y,r);
331     }
332     bool operator == (circle v){
333         return (p==v.p) && sgn(r-v.r)==0;
334     }
335     bool operator < (circle v)const{
336         return ((p<v.p)||((p==v.p)&&sgn(r-v.r)<0));


337       }
338       //面积
339       double area(){
340           return pi*r*r;
341       }
342       //周长
343       double circumference(){
344           return 2*pi*r;
345       }
346       //点和圆的关系
347       //0 圆外
348       //1 圆上
349       //2 圆内
350       int relation(Point b){
351           double dst = b.distance(p);
352           if(sgn(dst-r) < 0)return 2;
353           else if(sgn(dst-r)==0)return 1;
354           return 0;
355       }
356       //线段和圆的关系
357       //比较的是圆心到线段的距离和半径的关系
358       int relationseg(Line v){
359           double dst = v.dispointtoseg(p);
360           if(sgn(dst-r) < 0)return 2;
361           else if(sgn(dst-r) == 0)return 1;
362           return 0;
363       }
364       //直线和圆的关系
365       //比较的是圆心到直线的距离和半径的关系
366       int relationline(Line v){
367           double dst = v.dispointtoline(p);
368           if(sgn(dst-r) < 0)return 2;
369           else if(sgn(dst-r) == 0)return 1;
370           return 0;
371       }
372       //两圆的关系
373       //5 相离
374       //4 外切
375       //3 相交
376       //2 内切
377       //1 内含
378       //需要 Point 的 distance
379       //测试：UVA12304
380       int relationcircle(circle v){
381           double d = p.distance(v.p);
382           if(sgn(d-r-v.r) > 0)return 5;
383           if(sgn(d-r-v.r) == 0)return 4;
384           double l = fabs(r-v.r);
385           if(sgn(d-r-v.r)<0 && sgn(d-l)>0)return 3;
386           if(sgn(d-l)==0)return 2;
387           if(sgn(d-l)<0)return 1;


388       }
389       //求两个圆的交点，返回 0 表示没有交点，返回 1 是一个交点，2 是两个交点
390       //需要 relationcircle
391       //测试：UVA12304
392       int pointcrosscircle(circle v,Point &p1,Point &p2){
393           int rel = relationcircle(v);
394           if(rel == 1 || rel == 5)return 0;
395           double d = p.distance(v.p);
396           double l = (d*d+r*r-v.r*v.r)/(2*d);
397           double h = sqrt(r*r-l*l);
398           Point tmp = p + (v.p-p).trunc(l);
399           p1 = tmp + ((v.p-p).rotleft().trunc(h));
400           p2 = tmp + ((v.p-p).rotright().trunc(h));
401           if(rel == 2 || rel == 4)
402               return 1;
403           return 2;
404       }
405       //求直线和圆的交点，返回交点个数
406       int pointcrossline(Line v,Point &p1,Point &p2){
407           if(!(*this).relationline(v))return 0;
408           Point a = v.lineprog(p);
409           double d = v.dispointtoline(p);
410           d = sqrt(r*r-d*d);
411           if(sgn(d) == 0){
412               p1 = a;
413               p2 = a;
414               return 1;
415           }
416           p1 = a + (v.e-v.s).trunc(d);
417           p2 = a - (v.e-v.s).trunc(d);
418           return 2;
419       }
420       //得到过 a,b 两点，半径为 r1 的两个圆
421       int gercircle(Point a,Point b,double r1,circle &c1,circle &c2){
422           circle x(a,r1),y(b,r1);
423           int t = x.pointcrosscircle(y,c1.p,c2.p);
424           if(!t)return 0;
425           c1.r = c2.r = r;
426           return t;
427       }
428       //得到与直线 u 相切，过点 q, 半径为 r1 的圆
429       //测试：UVA12304
430       int getcircle(Line u,Point q,double r1,circle &c1,circle &c2){
431           double dis = u.dispointtoline(q);
432           if(sgn(dis-r1*2)>0)return 0;
433           if(sgn(dis) == 0){
434               c1.p = q + ((u.e-u.s).rotleft().trunc(r1));
435               c2.p = q + ((u.e-u.s).rotright().trunc(r1));
436               c1.r = c2.r = r1;
437               return 2;
438           }


439              Line u1 = Line((u.s + (u.e-u.s).rotleft().trunc(r1)),(u.e +
                     (u.e-u.s).rotleft().trunc(r1)));
440              Line u2 = Line((u.s + (u.e-u.s).rotright().trunc(r1)),(u.e
                    + (u.e-u.s).rotright().trunc(r1)));
441              circle cc = circle(q,r1);
442              Point p1,p2;
443              if(!cc.pointcrossline(u1,p1,p2))cc.pointcrossline(u2,p1,p2)
                    ;
444              c1 = circle(p1,r1);
445              if(p1 == p2){
446                   c2 = c1;
447                   return 1;
448              }
449              c2 = circle(p2,r1);
450              return 2;
451       }
452       //同时与直线 u,v 相切，半径为 r1 的圆
453       //测试：UVA12304
454       int getcircle(Line u,Line v,double r1,circle &c1,circle &c2,
             circle &c3,circle &c4){
455           if(u.parallel(v))return 0;//两直线平行
456           Line u1 = Line(u.s + (u.e-u.s).rotleft().trunc(r1),u.e + (u
                 .e-u.s).rotleft().trunc(r1));
457           Line u2 = Line(u.s + (u.e-u.s).rotright().trunc(r1),u.e + (
                 u.e-u.s).rotright().trunc(r1));
458           Line v1 = Line(v.s + (v.e-v.s).rotleft().trunc(r1),v.e + (v
                 .e-v.s).rotleft().trunc(r1));
459           Line v2 = Line(v.s + (v.e-v.s).rotright().trunc(r1),v.e + (
                 v.e-v.s).rotright().trunc(r1));
460           c1.r = c2.r = c3.r = c4.r = r1;
461           c1.p = u1.crosspoint(v1);
462           c2.p = u1.crosspoint(v2);
463           c3.p = u2.crosspoint(v1);
464           c4.p = u2.crosspoint(v2);
465           return 4;
466       }
467       //同时与不相交圆 cx,cy 相切，半径为 r1 的圆
468       //测试：UVA12304
469       int getcircle(circle cx,circle cy,double r1,circle &c1,circle &
             c2){
470           circle x(cx.p,r1+cx.r),y(cy.p,r1+cy.r);
471           int t = x.pointcrosscircle(y,c1.p,c2.p);
472           if(!t)return 0;
473           c1.r = c2.r = r1;
474           return t;
475       }
476
477       //过一点作圆的切线 (先判断点和圆的关系)
478       //测试：UVA12304
479       int tangentline(Point q,Line &u,Line &v){
480           int x = relation(q);


481              if(x == 2)return 0;
482              if(x == 1){
483                   u = Line(q,q + (q-p).rotleft());
484                   v = u;
485                   return 1;
486              }
487              double d = p.distance(q);
488              double l = r*r/d;
489              double h = sqrt(r*r-l*l);
490              u = Line(q,p + ((q-p).trunc(l) + (q-p).rotleft().trunc(h)))
                    ;
491              v = Line(q,p + ((q-p).trunc(l) + (q-p).rotright().trunc(h))
                    );
492              return 2;
493       }
494       //求两圆相交的面积
495       double areacircle(circle v){
496           int rel = relationcircle(v);
497           if(rel >= 4)return 0.0;
498           if(rel <= 2)return min(area(),v.area());
499           double d = p.distance(v.p);
500           double hf = (r+v.r+d)/2.0;
501           double ss = 2*sqrt(hf*(hf-r)*(hf-v.r)*(hf-d));
502           double a1 = acos((r*r+d*d-v.r*v.r)/(2.0*r*d));
503           a1 = a1*r*r;
504           double a2 = acos((v.r*v.r+d*d-r*r)/(2.0*v.r*d));
505           a2 = a2*v.r*v.r;
506           return a1+a2-ss;
507       }
508       //求圆和三角形 pab 的相交面积
509       //测试：POJ3675 HDU3982 HDU2892
510       double areatriangle(Point a,Point b){
511           if(sgn((p-a)^(p-b)) == 0)return 0.0;
512           Point q[5];
513           int len = 0;
514           q[len++] = a;
515           Line l(a,b);
516           Point p1,p2;
517           if(pointcrossline(l,q[1],q[2])==2){
518               if(sgn((a-q[1])*(b-q[1]))<0)q[len++] = q[1];
519               if(sgn((a-q[2])*(b-q[2]))<0)q[len++] = q[2];
520           }
521           q[len++] = b;
522           if(len == 4 && sgn((q[0]-q[1])*(q[2]-q[1]))>0)swap(q[1],q
                 [2]);
523           double res = 0;
524           for(int i = 0;i < len-1;i++){
525               if(relation(q[i])==0||relation(q[i+1])==0){
526                    double arg = p.rad(q[i],q[i+1]);
527                    res += r*r*arg/2.0;
528               }


529              else{
530                  res += fabs((q[i]-p)^(q[i+1]-p))/2.0;
531              }
532         }
533         return res;
534     }
535 };
536
537 /*
538 * n,p Line l for each side
539 * input(int _n)                          - inputs _n size polygon
540 * add(Point q)                           - adds a point at end of
        the list
541 * getline()                              - populates line array
542 * cmp                                    - comparision in
        convex_hull order
543 * norm()                                 - sorting in convex_hull
        order
544 * getconvex(polygon &convex)             - returns convex hull in
        convex
545 * Graham(polygon &convex)                - returns convex hull in
        convex
546 * isconvex()                             - checks if convex
547 * relationpoint(Point q)                 - returns 3 if q is a
        vertex
548 *                                                  2 if on a side
549 *                                                  1 if inside
550 *                                                  0 if outside
551 * convexcut(Line u,polygon &po)          - left side of u in po
552 * gercircumference()                     - returns side length
553 * getarea()                              - returns area
554 * getdir()                               - returns 0 for cw, 1 for
        ccw
555 * getbarycentre()                        - returns barycenter
556 *
557 */
558 struct polygon{
559     int n;
560     Point p[maxp];
561     Line l[maxp];
562     void input(int _n){
563         n = _n;
564         for(int i = 0;i < n;i++)
565              p[i].input();
566     }
567     void add(Point q){
568         p[n++] = q;
569     }
570     void getline(){
571         for(int i = 0;i < n;i++){
572              l[i] = Line(p[i],p[(i+1)%n]);


573           }
574       }
575       struct cmp{
576           Point p;
577           cmp(const Point &p0){p = p0;}
578           bool operator()(const Point &aa,const Point &bb){
579               Point a = aa, b = bb;
580               int d = sgn((a-p)^(b-p));
581               if(d == 0){
582                    return sgn(a.distance(p)-b.distance(p)) < 0;
583               }
584               return d > 0;
585           }
586       };
587       //进行极角排序
588       //首先需要找到最左下角的点
589       //需要重载号好 Point 的 < 操作符 (min 函数要用)
590       void norm(){
591           Point mi = p[0];
592           for(int i = 1;i < n;i++)mi = min(mi,p[i]);
593           sort(p,p+n,cmp(mi));
594       }
595       //得到凸包
596       //得到的凸包里面的点编号是 0∼n-1 的
597       //两种凸包的方法
598       //注意如果有影响，要特判下所有点共点，或者共线的特殊情况
599       //测试 LightOJ1203 LightOJ1239
600       void getconvex(polygon &convex){
601           sort(p,p+n);
602           convex.n = n;
603           for(int i = 0;i < min(n,2);i++){
604               convex.p[i] = p[i];
605           }
606           if(convex.n == 2 && (convex.p[0] == convex.p[1]))convex.n
                 --;//特
                 判
607           if(n <= 2)return;
608           int &top = convex.n;
609           top = 1;
610           for(int i = 2;i < n;i++){
611               while(top && sgn((convex.p[top]-p[i])^(convex.p[top-1]-
                     p[i])) <= 0)
612                    top--;
613               convex.p[++top] = p[i];
614           }
615           int temp = top;
616           convex.p[++top] = p[n-2];
617           for(int i = n-3;i >= 0;i--){
618               while(top != temp && sgn((convex.p[top]-p[i])^(convex.p
                     [top-1]-p[i])) <= 0)
619                    top--;


620                  convex.p[++top] = p[i];
621              }
622              if(convex.n == 2 && (convex.p[0] == convex.p[1]))convex.n
                    --;//特
                    判
623              convex.norm();//原来得到的是顺时针的点，排序后逆时针
624       }
625       //得到凸包的另外一种方法
626       //测试 LightOJ1203 LightOJ1239
627       void Graham(polygon &convex){
628           norm();
629           int &top = convex.n;
630           top = 0;
631           if(n == 1){
632               top = 1;
633               convex.p[0] = p[0];
634               return;
635           }
636           if(n == 2){
637               top = 2;
638               convex.p[0] = p[0];
639               convex.p[1] = p[1];
640               if(convex.p[0] == convex.p[1])top--;
641               return;
642           }
643           convex.p[0] = p[0];
644           convex.p[1] = p[1];
645           top = 2;
646           for(int i = 2;i < n;i++){
647               while( top > 1 && sgn((convex.p[top-1]-convex.p[top-2])
                     ^(p[i]-convex.p[top-2])) <= 0 )
648                    top--;
649               convex.p[top++] = p[i];
650           }
651           if(convex.n == 2 && (convex.p[0] == convex.p[1]))convex.n
                 --;//特
                 判
652       }
653       //判断是不是凸的
654       bool isconvex(){
655           bool s[2];
656           memset(s,false,sizeof(s));
657           for(int i = 0;i < n;i++){
658               int j = (i+1)%n;
659               int k = (j+1)%n;
660               s[sgn((p[j]-p[i])^(p[k]-p[i]))+1] = true;
661               if(s[0] && s[2])return false;
662           }
663           return true;
664       }
665       //判断点和任意多边形的关系


666       // 3 点上
667       // 2 边上
668       // 1 内部
669       // 0 外部
670       int relationpoint(Point q){
671           for(int i = 0;i < n;i++){
672               if(p[i] == q)return 3;
673           }
674           getline();
675           for(int i = 0;i < n;i++){
676               if(l[i].pointonseg(q))return 2;
677           }
678           int cnt = 0;
679           for(int i = 0;i < n;i++){
680               int j = (i+1)%n;
681               int k = sgn((q-p[j])^(p[i]-p[j]));
682               int u = sgn(p[i].y-q.y);
683               int v = sgn(p[j].y-q.y);
684               if(k > 0 && u < 0 && v >= 0)cnt++;
685               if(k < 0 && v < 0 && u >= 0)cnt--;
686           }
687           return cnt != 0;
688       }
689       //直线 u 切割凸多边形左侧
690       //注意直线方向
691       //测试：HDU3982
692       void convexcut(Line u,polygon &po){
693           int &top = po.n;//注意引用
694           top = 0;
695           for(int i = 0;i < n;i++){
696               int d1 = sgn((u.e-u.s)^(p[i]-u.s));
697               int d2 = sgn((u.e-u.s)^(p[(i+1)%n]-u.s));
698               if(d1 >= 0)po.p[top++] = p[i];
699               if(d1*d2 < 0)po.p[top++] = u.crosspoint(Line(p[i],p[(i
                     +1)%n]));
700           }
701       }
702       //得到周长
703       //测试 LightOJ1239
704       double getcircumference(){
705           double sum = 0;
706           for(int i = 0;i < n;i++){
707               sum += p[i].distance(p[(i+1)%n]);
708           }
709           return sum;
710       }
711       //得到面积
712       double getarea(){
713           double sum = 0;
714           for(int i = 0;i < n;i++){
715               sum += (p[i]^p[(i+1)%n]);


716              }
717              return fabs(sum)/2;
718       }
719       //得到方向
720       // 1 表示逆时针，0 表示顺时针
721       bool getdir(){
722           double sum = 0;
723           for(int i = 0;i < n;i++)
724               sum += (p[i]^p[(i+1)%n]);
725           if(sgn(sum) > 0)return 1;
726           return 0;
727       }
728       //得到重心
729       Point getbarycentre(){
730           Point ret(0,0);
731           double area = 0;
732           for(int i = 1;i < n-1;i++){
733               double tmp = (p[i]-p[0])^(p[i+1]-p[0]);
734               if(sgn(tmp) == 0)continue;
735               area += tmp;
736               ret.x += (p[0].x+p[i].x+p[i+1].x)/3*tmp;
737               ret.y += (p[0].y+p[i].y+p[i+1].y)/3*tmp;
738           }
739           if(sgn(area)) ret = ret/area;
740           return ret;
741       }
742       //多边形和圆交的面积
743       //测试：POJ3675 HDU3982 HDU2892
744       double areacircle(circle c){
745           double ans = 0;
746           for(int i = 0;i < n;i++){
747               int j = (i+1)%n;
748               if(sgn( (p[j]-c.p)^(p[i]-c.p) ) >= 0)
749                   ans += c.areatriangle(p[i],p[j]);
750               else ans -= c.areatriangle(p[i],p[j]);
751           }
752           return fabs(ans);
753       }
754       //多边形和圆关系
755       // 2 圆完全在多边形内
756       // 1 圆在多边形里面，碰到了多边形边界
757       // 0 其它
758       int relationcircle(circle c){
759           getline();
760           int x = 2;
761           if(relationpoint(c.p) != 1)return 0;//圆心不在内部
762           for(int i = 0;i < n;i++){
763               if(c.relationseg(l[i])==2)return 0;
764               if(c.relationseg(l[i])==1)x = 1;
765           }
766           return x;


767       }
768   };
769   //AB X AC
770   double cross(Point A,Point B,Point C){
771       return (B-A)^(C-A);
772   }
773   //AB*AC
774   double dot(Point A,Point B,Point C){
775       return (B-A)*(C-A);
776   }
777   //最小矩形面积覆盖
778   // A 必须是凸包 (而且是逆时针顺序)
779   // 测试 UVA 10173
780   double minRectangleCover(polygon A){
781       //要特判 A.n < 3 的情况
782       if(A.n < 3)return 0.0;
783       A.p[A.n] = A.p[0];
784       double ans = -1;
785       int r = 1, p = 1, q;
786       for(int i = 0;i < A.n;i++){
787           //卡出离边 A.p[i] - A.p[i+1] 最远的点
788           while( sgn( cross(A.p[i],A.p[i+1],A.p[r+1]) - cross(A.p[i],
                 A.p[i+1],A.p[r]) ) >= 0 )
789               r = (r+1)%A.n;
790           //卡出 A.p[i] - A.p[i+1] 方向上正向 n 最远的点
791           while(sgn( dot(A.p[i],A.p[i+1],A.p[p+1]) - dot(A.p[i],A.p[i
                 +1],A.p[p]) ) >= 0 )
792               p = (p+1)%A.n;
793           if(i == 0)q = p;
794           //卡出 A.p[i] - A.p[i+1] 方向上负向最远的点
795           while(sgn(dot(A.p[i],A.p[i+1],A.p[q+1]) - dot(A.p[i],A.p[i
                 +1],A.p[q])) <= 0)
796               q = (q+1)%A.n;
797           double d = (A.p[i] - A.p[i+1]).len2();
798           double tmp = cross(A.p[i],A.p[i+1],A.p[r]) *
799               (dot(A.p[i],A.p[i+1],A.p[p]) - dot(A.p[i],A.p[i+1],A.p[
                     q]))/d;
800           if(ans < 0 || ans > tmp)ans = tmp;
801       }
802       return ans;
803   }
804
805   //直线切凸多边形
806   //多边形是逆时针的，在 q1q2 的左侧
807   //测试:HDU3982
808   vector<Point> convexCut(const vector<Point> &ps,Point q1,Point q2){
809       vector<Point>qs;
810       int n = ps.size();
811       for(int i = 0;i < n;i++){
812           Point p1 = ps[i], p2 = ps[(i+1)%n];
813           int d1 = sgn((q2-q1)^(p1-q1)), d2 = sgn((q2-q1)^(p2-q1));


814              if(d1 >= 0)
815                  qs.push_back(p1);
816              if(d1 * d2 < 0)
817                  qs.push_back(Line(p1,p2).crosspoint(Line(q1,q2)));
818       }
819       return qs;
820   }
821   //半平面交
822   //测试 POJ3335 POJ1474 POJ1279
823   //***************************
824   struct halfplane:public Line{
825       double angle;
826       halfplane(){}
827       //表示向量 s->e 逆时针 (左侧) 的半平面
828       halfplane(Point _s,Point _e){
829           s = _s;
830           e = _e;
831       }
832       halfplane(Line v){
833           s = v.s;
834           e = v.e;
835       }
836       void calcangle(){
837           angle = atan2(e.y-s.y,e.x-s.x);
838       }
839       bool operator <(const halfplane &b)const{
840           return angle < b.angle;
841       }
842   };
843   struct halfplanes{
844       int n;
845       halfplane hp[maxp];
846       Point p[maxp];
847       int que[maxp];
848       int st,ed;
849       void push(halfplane tmp){
850           hp[n++] = tmp;
851       }
852       //去重
853       void unique(){
854           int m = 1;
855           for(int i = 1;i < n;i++){
856               if(sgn(hp[i].angle-hp[i-1].angle) != 0)
857                    hp[m++] = hp[i];
858               else if(sgn( (hp[m-1].e-hp[m-1].s)^(hp[i].s-hp[m-1].s)
                     ) > 0)
859                    hp[m-1] = hp[i];
860           }
861           n = m;
862       }
863       bool halfplaneinsert(){


864              for(int i = 0;i < n;i++)hp[i].calcangle();
865              sort(hp,hp+n);
866              unique();
867              que[st=0] = 0;
868              que[ed=1] = 1;
869              p[1] = hp[0].crosspoint(hp[1]);
870              for(int i = 2;i < n;i++){
871                  while(st<ed && sgn((hp[i].e-hp[i].s)^(p[ed]-hp[i].s))
                        <0)ed--;
872                  while(st<ed && sgn((hp[i].e-hp[i].s)^(p[st+1]-hp[i].s))
                        <0)st++;
873                  que[++ed] = i;
874                  if(hp[i].parallel(hp[que[ed-1]]))return false;
875                  p[ed]=hp[i].crosspoint(hp[que[ed-1]]);
876              }
877              while(st<ed && sgn((hp[que[st]].e-hp[que[st]].s)^(p[ed]-hp[
                    que[st]].s))<0)ed--;
878              while(st<ed && sgn((hp[que[ed]].e-hp[que[ed]].s)^(p[st+1]-
                    hp[que[ed]].s))<0)st++;
879              if(st+1>=ed)return false;
880              return true;
881       }
882       //得到最后半平面交得到的凸多边形
883       //需要先调用 halfplaneinsert() 且返回 true
884       void getconvex(polygon &con){
885           p[st] = hp[que[st]].crosspoint(hp[que[ed]]);
886           con.n = ed-st+1;
887           for(int j = st,i = 0;j <= ed;i++,j++)
888               con.p[i] = p[j];
889       }
890   };
891   //***************************
892
893   const int maxn = 1010;
894   struct circles{
895       circle c[maxn];
896       double ans[maxn];//ans[i] 表示被覆盖了 i 次的面积
897       double pre[maxn];
898       int n;
899       circles(){}
900       void add(circle cc){
901           c[n++] = cc;
902       }
903       //x 包含在 y 中
904       bool inner(circle x,circle y){
905           if(x.relationcircle(y) != 1)return 0;
906           return sgn(x.r-y.r)<=0?1:0;
907       }
908       //圆的面积并去掉内含的圆
909       void init_or(){
910           bool mark[maxn] = {0};


911              int i,j,k=0;
912              for(i = 0;i < n;i++){
913                  for(j = 0;j < n;j++)
914                      if(i != j && !mark[j]){
915                          if( (c[i]==c[j])||inner(c[i],c[j]) )break;
916                      }
917                  if(j < n)mark[i] = 1;
918              }
919              for(i = 0;i < n;i++)
920                  if(!mark[i])
921                      c[k++] = c[i];
922              n = k;
923       }
924       //圆的面积交去掉内含的圆
925       void init_add(){
926           int i,j,k;
927           bool mark[maxn] = {0};
928           for(i = 0;i < n;i++){
929               for(j = 0;j < n;j++)
930                   if(i != j && !mark[j]){
931                       if( (c[i]==c[j])||inner(c[j],c[i]) )break;
932                   }
933               if(j < n)mark[i] = 1;
934           }
935           for(i = 0;i < n;i++)
936               if(!mark[i])
937                   c[k++] = c[i];
938           n = k;
939       }
940       //半径为 r 的圆，弧度为 th 对应的弓形的面积
941       double areaarc(double th,double r){
942           return 0.5*r*r*(th-sin(th));
943       }
944       //测试 SPOJVCIRCLES SPOJCIRUT
945       //SPOJVCIRCLES 求 n 个圆并的面积，需要加上 init_or() 去掉重复圆（否则
             WA）
946       //SPOJCIRUT 是求被覆盖 k 次的面积，不能加 init_or()
947       //对于求覆盖多少次面积的问题，不能解决相同圆，而且不能 init_or()
948       //求多圆面积并，需要 init_or, 其中一个目的就是去掉相同圆
949       void getarea(){
950           memset(ans,0,sizeof(ans));
951           vector<pair<double,int> >v;
952           for(int i = 0;i < n;i++){
953               v.clear();
954               v.push_back(make_pair(-pi,1));
955               v.push_back(make_pair(pi,-1));
956               for(int j = 0;j < n;j++)
957                   if(i != j){
958                       Point q = (c[j].p - c[i].p);
959                       double ab = q.len(),ac = c[i].r, bc = c[j].r;
960                       if(sgn(ab+ac-bc)<=0){


961                              v.push_back(make_pair(-pi,1));
962                              v.push_back(make_pair(pi,-1));
963                              continue;
964                          }
965                          if(sgn(ab+bc-ac)<=0)continue;
966                          if(sgn(ab-ac-bc)>0)continue;
967                          double th = atan2(q.y,q.x), fai = acos((ac*ac+
                                ab*ab-bc*bc)/(2.0*ac*ab));
968                          double a0 = th-fai;
969                          if(sgn(a0+pi)<0)a0+=2*pi;
970                          double a1 = th+fai;
971                          if(sgn(a1-pi)>0)a1-=2*pi;
972                          if(sgn(a0-a1)>0){
973                              v.push_back(make_pair(a0,1));
974                              v.push_back(make_pair(pi,-1));
975                              v.push_back(make_pair(-pi,1));
976                              v.push_back(make_pair(a1,-1));
977                          }
978                          else{
979                              v.push_back(make_pair(a0,1));
980                              v.push_back(make_pair(a1,-1));
981                          }
982                      }
983                  sort(v.begin(),v.end());
984                  int cur = 0;
985                  for(int j = 0;j < v.size();j++){
986                      if(cur && sgn(v[j].first-pre[cur])){
987                          ans[cur] += areaarc(v[j].first-pre[cur],c[i].r)
                                ;
988                          ans[cur] += 0.5*(Point(c[i].p.x+c[i].r*cos(pre[
                                cur]),c[i].p.y+c[i].r*sin(pre[cur]))^Point(c
                                [i].p.x+c[i].r*cos(v[j].first),c[i].p.y+c[i
                                ].r*sin(v[j].first)));
989                      }
990                      cur += v[j].second;
991                      pre[cur] = v[j].first;
992                  }
993              }
994              for(int i = 1;i < n;i++)
995                  ans[i] -= ans[i+1];
996         }
997 };
\end{Code}

\subsection{Hình học ba chiều}

\begin{Code}
 1    const double eps = 1e-8;
 2    int sgn(double x){
 3        if(fabs(x) < eps)return 0;
 4        if(x < 0)return -1;
 5        else return 1;
 6    }
 7    struct Point3{


 8       double x,y,z;
 9       Point3(double _x = 0,double _y = 0,double _z = 0){
10           x = _x;
11           y = _y;
12           z = _z;
13       }
14       void input(){
15           scanf("%lf%lf%lf",&x,&y,&z);
16       }
17       void output(){
18           scanf("%.2lf %.2lf %.2lf\n",x,y,z);
19       }
20       bool operator ==(const Point3 &b)const{
21           return sgn(x-b.x) == 0 && sgn(y-b.y) == 0 && sgn(z-b.z) ==
                0;
22       }
23       bool operator <(const Point3 &b)const{
24           return sgn(x-b.x)==0?(sgn(y-b.y)==0?sgn(z-b.z)<0:y<b.y):x<b
                .x;
25       }
26       double len(){
27           return sqrt(x*x+y*y+z*z);
28       }
29       double len2(){
30           return x*x+y*y+z*z;
31       }
32       double distance(const Point3 &b)const{
33           return sqrt((x-b.x)*(x-b.x)+(y-b.y)*(y-b.y)+(z-b.z)*(z-b.z)
                );
34       }
35       Point3 operator -(const Point3 &b)const{
36           return Point3(x-b.x,y-b.y,z-b.z);
37       }
38       Point3 operator +(const Point3 &b)const{
39           return Point3(x+b.x,y+b.y,z+b.z);
40       }
41       Point3 operator *(const double &k)const{
42           return Point3(x*k,y*k,z*k);
43       }
44       Point3 operator /(const double &k)const{
45           return Point3(x/k,y/k,z/k);
46       }
47       //点乘
48       double operator *(const Point3 &b)const{
49           return x*b.x+y*b.y+z*b.z;
50       }
51       //叉乘
52       Point3 operator ^(const Point3 &b)const{
53           return Point3(y*b.z-z*b.y,z*b.x-x*b.z,x*b.y-y*b.x);
54       }
55       double rad(Point3 a,Point3 b){


56              Point3 p = (*this);
57              return acos( ( (a-p)*(b-p) )/ (a.distance(p)*b.distance(p))
                    );
 58     }
 59     //变换长度
 60     Point3 trunc(double r){
 61         double l = len();
 62         if(!sgn(l))return *this;
 63         r /= l;
 64         return Point3(x*r,y*r,z*r);
 65     }
 66 };
 67 struct Line3
 68 {
 69     Point3 s,e;
 70     Line3(){}
 71     Line3(Point3 _s,Point3 _e)
 72     {
 73         s = _s;
 74         e = _e;
 75     }
 76     bool operator ==(const Line3 v)
 77     {
 78         return (s==v.s)&&(e==v.e);
 79     }
 80     void input()
 81     {
 82         s.input();
 83         e.input();
 84     }
 85     double length()
 86     {
 87         return s.distance(e);
 88     }
 89     //点到直线距离
 90     double dispointtoline(Point3 p)
 91     {
 92         return ((e-s)^(p-s)).len()/s.distance(e);
 93     }
 94     //点到线段距离
 95     double dispointtoseg(Point3 p)
 96     {
 97         if(sgn((p-s)*(e-s)) < 0 || sgn((p-e)*(s-e)) < 0)
 98              return min(p.distance(s),e.distance(p));
 99         return dispointtoline(p);
100     }
101     //返回点 p 在直线上的投影
102     Point3 lineprog(Point3 p)
103     {
104         return s + ( ((e-s)*((e-s)*(p-s)))/((e-s).len2()) );
105     }


106       //p 绕此向量逆时针 arg 角度
107       Point3 rotate(Point3 p,double ang)
108       {
109           if(sgn(((s-p)^(e-p)).len()) == 0)return p;
110           Point3 f1 = (e-s)^(p-s);
111           Point3 f2 = (e-s)^(f1);
112           double len = ((s-p)^(e-p)).len()/s.distance(e);
113           f1 = f1.trunc(len); f2 = f2.trunc(len);
114           Point3 h = p+f2;
115           Point3 pp = h+f1;
116           return h + ((p-h)*cos(ang)) + ((pp-h)*sin(ang));
117       }
118       //点在直线上
119       bool pointonseg(Point3 p)
120       {
121           return sgn( ((s-p)^(e-p)).len() ) == 0 && sgn((s-p)*(e-p))
                 == 0;
122       }
123 };
124 struct Plane
125 {
126     Point3 a,b,c,o;//平面上的三个点，以及法向量
127     Plane(){}
128     Plane(Point3 _a,Point3 _b,Point3 _c)
129     {
130         a = _a;
131         b = _b;
132         c = _c;
133         o = pvec();
134     }
135     Point3 pvec()
136     {
137         return (b-a)^(c-a);
138     }
139     //ax+by+cz+d = 0
140     Plane(double _a,double _b,double _c,double _d)
141     {
142         o = Point3(_a,_b,_c);
143         if(sgn(_a) != 0)
144              a = Point3((-_d-_c-_b)/_a,1,1);
145         else if(sgn(_b) != 0)
146              a = Point3(1,(-_d-_c-_a)/_b,1);
147         else if(sgn(_c) != 0)
148              a = Point3(1,1,(-_d-_a-_b)/_c);
149     }
150     //点在平面上的判断
151     bool pointonplane(Point3 p)
152     {
153         return sgn((p-a)*o) == 0;
154     }
155     //两平面夹角


156        double angleplane(Plane f)
157        {
158            return acos(o*f.o)/(o.len()*f.o.len());
159        }
160        //平面和直线的交点，返回值是交点个数
161        int crossline(Line3 u,Point3 &p)
162        {
163            double x = o*(u.e-a);
164            double y = o*(u.s-a);
165            double d = x-y;
166            if(sgn(d) == 0)return 0;
167            p = ((u.s*x)-(u.e*y))/d;
168            return 1;
169        }
170        //点到平面最近点 (也就是投影)
171        Point3 pointtoplane(Point3 p)
172        {
173            Line3 u = Line3(p,p+o);
174            crossline(u,p);
175            return p;
176        }
177        //平面和平面的交线
178        int crossplane(Plane f,Line3 &u)
179        {
180            Point3 oo = o^f.o;
181            Point3 v = o^oo;
182            double d = fabs(f.o*v);
183            if(sgn(d) == 0)return 0;
184            Point3 q = a + (v*(f.o*(f.a-a))/d);
185            u = Line3(q,q+oo);
186            return 1;
187        }
188 };
\end{Code}

\subsection{Cặp điểm gần nhất trên mặt phẳng}

\begin{Code}
     HDU1007/ZOJ2107
 1   const int MAXN = 100010;
 2   const double eps = 1e-8;
 3   const double INF = 1e20;
 4   struct Point{
 5       double x,y;
 6       void input(){
 7           scanf("%lf%lf",&x,&y);
 8       }
 9   };
10   double dist(Point a,Point b){
11       return sqrt((a.x-b.x)*(a.x-b.x) + (a.y-b.y)*(a.y-b.y));
12   }
13   Point p[MAXN];
14   Point tmpt[MAXN];
15   bool cmpx(Point a,Point b){


16       return a.x < b.x || (a.x == b.x && a.y < b.y);
17   }
18   bool cmpy(Point a,Point b){
19       return a.y < b.y || (a.y == b.y && a.x < b.x);
20   }
21   double Closest_Pair(int left,int right){
22       double d = INF;
23       if(left == right)return d;
24       if(left+1 == right)return dist(p[left],p[right]);
25       int mid = (left+right)/2;
26       double d1 = Closest_Pair(left,mid);
27       double d2 = Closest_Pair(mid+1,right);
28       d = min(d1,d2);
29       int cnt = 0;
30       for(int i = left;i <= right;i++){
31           if(fabs(p[mid].x - p[i].x) <= d)
32               tmpt[cnt++] = p[i];
33       }
34       sort(tmpt,tmpt+cnt,cmpy);
35       for(int i = 0;i < cnt;i++){
36           for(int j = i+1;j < cnt && tmpt[j].y - tmpt[i].y < d;j++)
37               d = min(d,dist(tmpt[i],tmpt[j]));
38       }
39       return d;
40   }
41   int main(){
42       int n;
43       while(scanf("%d",&n) == 1 && n){
44           for(int i = 0;i < n;i++)p[i].input();
45           sort(p,p+n,cmpx);
46           printf("%.2lf\n",Closest_Pair(0,n-1));
47       }
48       return 0;
49   }
\end{Code}

\subsection{Bao lồi ba chiều}

\subsubsection{HDU4273}

\begin{Code}
     HDU 4273 给一个三维凸包，求重心到表面的最短距离。
 1   const double eps = 1e-8;
 2   const int MAXN = 550;
 3   int sgn(double x){
 4       if(fabs(x) < eps)return 0;
 5       if(x < 0)return -1;
 6       else return 1;
 7   }
 8   struct Point3{
 9       double x,y,z;
10       Point3(double _x = 0, double _y = 0, double _z = 0){
11           x = _x;


12              y = _y;
13              z = _z;
14       }
15       void input(){
16           scanf("%lf%lf%lf",&x,&y,&z);
17       }
18       bool operator ==(const Point3 &b)const{
19           return sgn(x-b.x) == 0 && sgn(y-b.y) == 0 && sgn(z-b.z) ==
                0;
20       }
21       double len(){
22           return sqrt(x*x+y*y+z*z);
23       }
24       double len2(){
25           return x*x+y*y+z*z;
26       }
27       double distance(const Point3 &b)const{
28           return sqrt((x-b.x)*(x-b.x)+(y-b.y)*(y-b.y)+(z-b.z)*(z-b.z)
                );
29       }
30       Point3 operator -(const Point3 &b)const{
31           return Point3(x-b.x,y-b.y,z-b.z);
32       }
33       Point3 operator +(const Point3 &b)const{
34           return Point3(x+b.x,y+b.y,z+b.z);
35       }
36       Point3 operator *(const double &k)const{
37           return Point3(x*k,y*k,z*k);
38       }
39       Point3 operator /(const double &k)const{
40           return Point3(x/k,y/k,z/k);
41       }
42       //点乘
43       double operator *(const Point3 &b)const{
44           return x*b.x + y*b.y + z*b.z;
45       }
46       //叉乘
47       Point3 operator ^(const Point3 &b)const{
48           return Point3(y*b.z-z*b.y,z*b.x-x*b.z,x*b.y-y*b.x);
49       }
50 };
51 struct CH3D{
52     struct face{
53         //表示凸包一个面上的三个点的编号
54         int a,b,c;
55         //表示该面是否属于最终的凸包上的面
56         bool ok;
57     };
58     //初始顶点数
59     int n;
60     Point3 P[MAXN];


 61       //凸包表面的三角形数
 62       int num;
 63       //凸包表面的三角形
 64       face F[8*MAXN];
 65       int g[MAXN][MAXN];
 66       //叉乘
 67       Point3 cross(const Point3 &a,const Point3 &b,const Point3 &c){
 68           return (b-a)^(c-a);
 69       }
 70       //三角形面积 *2
 71       double area(Point3 a,Point3 b,Point3 c){
 72           return ((b-a)^(c-a)).len();
 73       }
 74       //四面体有向面积 *6
 75       double volume(Point3 a,Point3 b,Point3 c,Point3 d){
 76           return ((b-a)^(c-a))*(d-a);
 77       }
 78       //正：点在面同向
 79       double dblcmp(Point3 &p,face &f){
 80           Point3 p1 = P[f.b] - P[f.a];
 81           Point3 p2 = P[f.c] - P[f.a];
 82           Point3 p3 = p - P[f.a];
 83           return (p1^p2)*p3;
 84       }
 85       void deal(int p,int a,int b){
 86           int f = g[a][b];
 87           face add;
 88           if(F[f].ok){
 89                if(dblcmp(P[p],F[f]) > eps)
 90                    dfs(p,f);
 91                else {
 92                    add.a = b;
 93                    add.b = a;
 94                    add.c = p;
 95                    add.ok = true;
 96                    g[p][b] = g[a][p] = g[b][a] = num;
 97                    F[num++] = add;
 98                }
 99           }
100       }
101       //递归搜索所有应该从凸包内删除的面
102       void dfs(int p,int now){
103           F[now].ok = false;
104           deal(p,F[now].b,F[now].a);
105           deal(p,F[now].c,F[now].b);
106           deal(p,F[now].a,F[now].c);
107       }
108       bool same(int s,int t){
109           Point3 &a = P[F[s].a];
110           Point3 &b = P[F[s].b];
111           Point3 &c = P[F[s].c];


112              return fabs(volume(a,b,c,P[F[t].a])) < eps &&
113                  fabs(volume(a,b,c,P[F[t].b])) < eps &&
114                  fabs(volume(a,b,c,P[F[t].c])) < eps;
115       }
116       //构建三维凸包
117       void create(){
118           num = 0;
119           face add;
120
121              //***********************************
122              //此段是为了保证前四个点不共面
123              bool flag = true;
124              for(int i = 1;i < n;i++){
125                  if(!(P[0] == P[i])){
126                       swap(P[1],P[i]);
127                       flag = false;
128                       break;
129                  }
130              }
131              if(flag)return;
132              flag = true;
133              for(int i = 2;i < n;i++){
134                  if( ((P[1]-P[0])^(P[i]-P[0])).len() > eps ){
135                       swap(P[2],P[i]);
136                       flag = false;
137                       break;
138                  }
139              }
140              if(flag)return;
141              flag = true;
142              for(int i = 3;i < n;i++){
143                  if(fabs( ((P[1]-P[0])^(P[2]-P[0]))*(P[i]-P[0]) ) > eps)
                        {
144                       swap(P[3],P[i]);
145                       flag = false;
146                       break;
147                  }
148              }
149              if(flag)return;
150              //**********************************
151
152              for(int i = 0;i < 4;i++){
153                  add.a = (i+1)%4;
154                  add.b = (i+2)%4;
155                  add.c = (i+3)%4;
156                  add.ok = true;
157                  if(dblcmp(P[i],add) > 0)swap(add.b,add.c);
158                  g[add.a][add.b] = g[add.b][add.c] = g[add.c][add.a] =
                        num;
159                  F[num++] = add;
160              }


161              for(int i = 4;i < n;i++)
162                  for(int j = 0;j < num;j++)
163                       if(F[j].ok && dblcmp(P[i],F[j]) > eps){
164                           dfs(i,j);
165                           break;
166                       }
167              int tmp = num;
168              num = 0;
169              for(int i = 0;i < tmp;i++)
170                  if(F[i].ok)
171                       F[num++] = F[i];
172       }
173       //表面积
174       //测试：HDU3528
175       double area(){
176           double res = 0;
177           if(n == 3){
178               Point3 p = cross(P[0],P[1],P[2]);
179               return p.len()/2;
180           }
181           for(int i = 0;i < num;i++)
182               res += area(P[F[i].a],P[F[i].b],P[F[i].c]);
183           return res/2.0;
184       }
185       double volume(){
186           double res = 0;
187           Point3 tmp = Point3(0,0,0);
188           for(int i = 0;i < num;i++)
189               res += volume(tmp,P[F[i].a],P[F[i].b],P[F[i].c]);
190           return fabs(res/6);
191       }
192       //表面三角形个数
193       int triangle(){
194           return num;
195       }
196       //表面多边形个数
197       //测试：HDU3662
198       int polygon(){
199           int res = 0;
200           for(int i = 0;i < num;i++){
201               bool flag = true;
202               for(int j = 0;j < i;j++)
203                   if(same(i,j)){
204                        flag = 0;
205                        break;
206                   }
207               res += flag;
208           }
209           return res;
210       }
211       //重心


212       //测试：HDU4273
213       Point3 barycenter(){
214           Point3 ans = Point3(0,0,0);
215           Point3 o = Point3(0,0,0);
216           double all = 0;
217           for(int i = 0;i < num;i++){
218                double vol = volume(o,P[F[i].a],P[F[i].b],P[F[i].c]);
219                ans = ans + (((o+P[F[i].a]+P[F[i].b]+P[F[i].c])/4.0)*
                      vol);
220                all += vol;
221           }
222           ans = ans/all;
223           return ans;
224       }
225       //点到面的距离
226       //测试：HDU4273
227       double ptoface(Point3 p,int i){
228           double tmp1 = fabs(volume(P[F[i].a],P[F[i].b],P[F[i].c],p))
                 ;
229           double tmp2 = ((P[F[i].b]-P[F[i].a])^(P[F[i].c]-P[F[i].a]))
                 .len();
230           return tmp1/tmp2;
231       }
232   };
233   CH3D hull;
234   int main()
235   {
236       while(scanf("%d",&hull.n) == 1){
237           for(int i = 0;i < hull.n;i++)hull.P[i].input();
238           hull.create();
239           Point3 p = hull.barycenter();
240           double ans = 1e20;
241           for(int i = 0;i < hull.num;i++)
242               ans = min(ans,hull.ptoface(p,i));
243           printf("%.3lf\n",ans);
244       }
245       return 0;
246   }
\end{Code}

\section{Khác}

\subsection{Số lớn}

\begin{Code}
   高精度，支持乘法和加法
 1 /*
 2 * 高精度，支持乘法和加法
 3 */
 4 struct BigInt{
 5     const static int mod = 10000;
 6     const static int DLEN = 4;
 7     int a[600],len;
 8     BigInt(){
 9         memset(a,0,sizeof(a));
10         len = 1;
11     }
12     BigInt(int v){
13         memset(a,0,sizeof(a));
14         len = 0;
15         do{
16             a[len++] = v%mod;
17             v /= mod;
18         }while(v);
19     }
20     BigInt(const char s[]){
21         memset(a,0,sizeof(a));
22         int L = strlen(s);
23         len = L/DLEN;
24         if(L%DLEN)len++;
25         int index = 0;
26         for(int i = L-1;i >= 0;i -= DLEN){
27             int t = 0;
28             int k = i - DLEN + 1;
29             if(k < 0)k = 0;
30             for(int j = k;j <= i;j++)
31                  t = t*10 + s[j] - '0';
32             a[index++] = t;
33         }
34     }
35     BigInt operator +(const BigInt &b)const{
36         BigInt res;
37         res.len = max(len,b.len);
38         for(int i = 0;i <= res.len;i++)
39             res.a[i] = 0;
40         for(int i = 0;i < res.len;i++){
41             res.a[i] += ((i < len)?a[i]:0)+((i < b.len)?b.a[i]:0);
42             res.a[i+1] += res.a[i]/mod;
43             res.a[i] %= mod;
44         }
45         if(res.a[res.len] > 0)res.len++;
46         return res;


47       }
48       BigInt operator *(const BigInt &b)const{
49           BigInt res;
50           for(int i = 0; i < len;i++){
51               int up = 0;
52               for(int j = 0;j < b.len;j++){
53                   int temp = a[i]*b.a[j] + res.a[i+j] + up;
54                   res.a[i+j] = temp%mod;
55                   up = temp/mod;
56               }
57               if(up != 0)
58                   res.a[i + b.len] = up;
59           }
60           res.len = len + b.len;
61           while(res.a[res.len - 1] == 0 &&res.len > 1)res.len--;
62           return res;
63       }
64       void output(){
65           printf("%d",a[len-1]);
66           for(int i = len-2;i >=0 ;i--)
67               printf("%04d",a[i]);
68           printf("\n");
69       }
70 };
\end{Code}

\subsection{Số lớn đầy đủ}

\begin{Code}
   HDU 1134 求卡特兰数
 1 /*
 2 * 完全大数模板
 3 * 输入 cin>>a
 4 * 输出 a.print();
 5 * 注意这个输入不能自动去掉前导 0 的，可以先读入到 char 数组，去掉前导 0，再用
       构造函数。
 6 */
 7 #define MAXN 9999
 8 #define MAXSIZE 1010
 9 #define DLEN 4
10 class BigNum{
11 private:
12     int a[500]; //可以控制大数的位数
13     int len;
14 public:
15     BigNum(){len=1;memset(a,0,sizeof(a));} //构造函数
16     BigNum(const int);     //将一个 int 类型的变量转化成大数
17     BigNum(const char*);   //将一个字符串类型的变量转化为大数
18     BigNum(const BigNum &); //拷贝构造函数
19     BigNum &operator=(const BigNum &); //重载赋值运算符，大数之间进行赋
          值运算
20     friend istream& operator>>(istream&,BigNum&); //重载输入运算符
21     friend ostream& operator<<(ostream&,BigNum&); //重载输出运算符


22
23       BigNum operator+(const BigNum &)const;    //重载加法运算符，两个大数
            之间的相加运算
24       BigNum operator-(const BigNum &)const;    //重载减法运算符，两个大数
            之间的相减运算
25       BigNum operator*(const BigNum &)const;    //重载乘法运算符，两个大数
            之间的相乘运算
26       BigNum operator/(const int &)const;       //重载除法运算符，大数对一
            个整数进行相除运算
27
28       BigNum operator^(const int &)const;       //大数的 n 次方运算
29       int operator%(const int &)const;          //大数对一个类型的变量进行
            取模运算int
30       bool operator>(const BigNum &T)const;     //大数和另一个大数的大小比
            较
31       bool operator>(const int &t)const;        //大数和一个 int 类型的变
            量的大小比较
32
33       void print();         //输出大数
34   };
35   //将一个 int 类型的变量转化为大数
36   BigNum::BigNum(const int b){
37       int c,d=b;
38       len=0;
39       memset(a,0,sizeof(a));
40       while(d>MAXN){
41           c=d-(d/(MAXN+1))*(MAXN+1);
42           d=d/(MAXN+1);
43           a[len++]=c;
44       }
45       a[len++]=d;
46   }
47   //将一个字符串类型的变量转化为大数
48   BigNum::BigNum(const char *s){
49       int t,k,index,L,i;
50       memset(a,0,sizeof(a));
51       L=strlen(s);
52       len=L/DLEN;
53       if(L%DLEN)len++;
54       index=0;
55       for(i=L-1;i>=0;i-=DLEN){
56           t=0;
57           k=i-DLEN+1;
58           if(k<0)k=0;
59           for(int j=k;j<=i;j++)
60                t=t*10+s[j]-'0';
61           a[index++]=t;
62       }
63   }
64   //拷贝构造函数
65   BigNum::BigNum(const BigNum &T):len(T.len){


 66       int i;
 67       memset(a,0,sizeof(a));
 68       for(i=0;i<len;i++)
 69           a[i]=T.a[i];
 70   }
 71   //重载赋值运算符，大数之间赋值运算
 72   BigNum & BigNum::operator=(const BigNum &n){
 73       int i;
 74       len=n.len;
 75       memset(a,0,sizeof(a));
 76       for(i=0;i<len;i++)
 77           a[i]=n.a[i];
 78       return *this;
 79   }
 80   istream& operator>>(istream &in,BigNum &b){
 81       char ch[MAXSIZE*4];
 82       int i=-1;
 83       in>>ch;
 84       int L=strlen(ch);
 85       int count=0,sum=0;
 86       for(i=L-1;i>=0;){
 87           sum=0;
 88           int t=1;
 89           for(int j=0;j<4&&i>=0;j++,i--,t*=10){
 90               sum+=(ch[i]-'0')*t;
 91           }
 92           b.a[count]=sum;
 93           count++;
 94       }
 95       b.len=count++;
 96       return in;
 97   }
 98   //重载输出运算符
 99   ostream& operator<<(ostream& out,BigNum& b){
100       int i;
101       cout<<b.a[b.len-1];
102       for(i=b.len-2;i>=0;i--){
103           printf("%04d",b.a[i]);
104       }
105       return out;
106   }
107   //两个大数之间的相加运算
108   BigNum BigNum::operator+(const BigNum &T)const{
109       BigNum t(*this);
110       int i,big;
111       big=T.len>len?T.len:len;
112       for(i=0;i<big;i++){
113           t.a[i]+=T.a[i];
114           if(t.a[i]>MAXN){
115               t.a[i+1]++;
116               t.a[i]-=MAXN+1;


117           }
118       }
119       if(t.a[big]!=0)
120          t.len=big+1;
121       else t.len=big;
122       return t;
123   }
124   //两个大数之间的相减运算
125   BigNum BigNum::operator-(const BigNum &T)const{
126       int i,j,big;
127       bool flag;
128       BigNum t1,t2;
129       if(*this>T){
130           t1=*this;
131           t2=T;
132           flag=0;
133       }
134       else{
135           t1=T;
136           t2=*this;
137           flag=1;
138       }
139       big=t1.len;
140       for(i=0;i<big;i++){
141           if(t1.a[i]<t2.a[i]){
142                j=i+1;
143                while(t1.a[j]==0)
144                    j++;
145                t1.a[j--]--;
146                while(j>i)
147                    t1.a[j--]+=MAXN;
148                t1.a[i]+=MAXN+1-t2.a[i];
149           }
150           else t1.a[i]-=t2.a[i];
151       }
152       t1.len=big;
153       while(t1.a[t1.len-1]==0 && t1.len>1){
154           t1.len--;
155           big--;
156       }
157       if(flag)
158           t1.a[big-1]=0-t1.a[big-1];
159       return t1;
160   }
161   //两个大数之间的相乘
162   BigNum BigNum::operator*(const BigNum &T)const{
163       BigNum ret;
164       int i,j,up;
165       int temp,temp1;
166       for(i=0;i<len;i++){
167           up=0;


168              for(j=0;j<T.len;j++){
169                  temp=a[i]*T.a[j]+ret.a[i+j]+up;
170                  if(temp>MAXN){
171                      temp1=temp-temp/(MAXN+1)*(MAXN+1);
172                      up=temp/(MAXN+1);
173                      ret.a[i+j]=temp1;
174                  }
175                  else{
176                      up=0;
177                      ret.a[i+j]=temp;
178                  }
179              }
180              if(up!=0)
181                 ret.a[i+j]=up;
182       }
183       ret.len=i+j;
184       while(ret.a[ret.len-1]==0 && ret.len>1)ret.len--;
185       return ret;
186   }
187   //大数对一个整数进行相除运算
188   BigNum BigNum::operator/(const int &b)const{
189       BigNum ret;
190       int i,down=0;
191       for(i=len-1;i>=0;i--){
192           ret.a[i]=(a[i]+down*(MAXN+1))/b;
193           down=a[i]+down*(MAXN+1)-ret.a[i]*b;
194       }
195       ret.len=len;
196       while(ret.a[ret.len-1]==0 && ret.len>1)
197           ret.len--;
198       return ret;
199   }
200   //大数对一个 int 类型的变量进行取模
201   int BigNum::operator%(const int &b)const{
202       int i,d=0;
203       for(i=len-1;i>=0;i--)
204           d=((d*(MAXN+1))%b+a[i])%b;
205       return d;
206   }
207   //大数的 n 次方运算
208   BigNum BigNum::operator^(const int &n)const{
209       BigNum t,ret(1);
210       int i;
211       if(n<0)exit(-1);
212       if(n==0)return 1;
213       if(n==1)return *this;
214       int m=n;
215       while(m>1){
216           t=*this;
217           for(i=1;(i<<1)<=m;i<<=1)
218              t=t*t;


219              m-=i;
220              ret=ret*t;
221              if(m==1)ret=ret*(*this);
222       }
223       return ret;
224   }
225   //大数和另一个大数的大小比较
226   bool BigNum::operator>(const BigNum &T)const{
227       int ln;
228       if(len>T.len)return true;
229       else if(len==T.len){
230           ln=len-1;
231           while(a[ln]==T.a[ln]&&ln>=0)
232              ln--;
233           if(ln>=0 && a[ln]>T.a[ln])
234               return true;
235           else
236               return false;
237       }
238       else
239          return false;
240   }
241   //大数和一个 int 类型的变量的大小比较
242   bool BigNum::operator>(const int &t)const{
243       BigNum b(t);
244       return *this>b;
245   }
246   //输出大数
247   void BigNum::print(){
248       int i;
249       printf("%d",a[len-1]);
250       for(i=len-2;i>=0;i--)
251         printf("%04d",a[i]);
252       printf("\n");
253   }
254   BigNum f[110];//卡特兰数
255
256   int main(){
257       f[0]=1;
258       for(int i=1;i<=100;i++)
259           f[i]=f[i-1]*(4*i-2)/(i+1);//卡特兰数递推式
260       int n;
261       while(scanf("%d",&n)==1){
262           if(n==-1)break;
263           f[n].print();
264       }
265       return 0;
266   }
\end{Code}

\subsection{Kết hợp strtok và sscanf để nhập dữ liệu}

\paragraph{Ghi chú.} Dùng dấu cách làm ký tự phân tách để đọc các số nguyên trên một dòng.

\begin{Code}
1               gets(buf);
2               int v;
3               char *p = strtok(buf," ");
4               while(p)
5               {
6                   sscanf(p,"%d",&v);
7                   p = strtok(NULL," ");
8               }
\end{Code}

\subsection{Tránh tràn stack bằng cách tăng stack thủ công}

\paragraph{Ghi chú.} Cách tốt nhất để tránh tràn stack là đổi cách viết sang BFS hoặc tự mô phỏng stack. Tăng stack thủ công chỉ là biện pháp phụ thuộc môi trường, cần dùng thận trọng.

\begin{Code}
1 #pragma comment(linker, "/STACK:1024000000,1024000000")
     G++
     放在主函数里面 (汇编开栈不一定适用，和系统有关。需谨慎!)
1 int __size__ = 256<<20;
2 char *__p__ = (char *)malloc(__size__)+__size__;
3 __asm__("movl %0,%%esp\n"::"r"(__p__));
\end{Code}

\subsection{STL}

\subsubsection{Hàng đợi ưu tiên priority\_queue}

\paragraph{Ghi chú.} empty() trả true nếu hàng đợi rỗng; pop() xóa phần tử đầu; push() thêm phần tử; size() trả số phần tử; top() trả phần tử đầu. Với priority\_queue mặc định, phần tử có độ ưu tiên cao hơn ra trước; với int mặc định, số lớn hơn ra trước.

\begin{Code}
1 priority_queue<int>q1;//大的先出对
2 priority_queue<int,vector<int>,greater<int> >q2; //小的先出队
     自定义比较函数：

 1   struct cmp
 2   {
 3       bool operator ()(int x, int y)
 4       {
 5           return x > y; // x 小的优先级高
 6         //也可以写成其他方式，如：return p[x] > p[y]; 表示 p[i] 小的优先级高
 7   }
 8   };
 9   priority_queue<int, vector<int>, cmp>q;//定义方法
10   //其中，第二个参数为容器类型。第三个参数为比较函数。


 1   struct node
 2   {
 3       int x, y;
 4       friend bool operator < (node a, node b)
 5       {
 6           return a.x > b.x; //结构体中，x 小的优先级高
 7       }
 8   };
 9   priority_queue<node>q;//定义方法
10   //在该结构中，y 为值, x 为优先级。
11   //通过自定义 operator< 操作符来比较元素中的优先级。
12   //在重载”<”时，最好不要重载”>”，可能会发生编译错误
\end{Code}

\subsubsection{set và multiset}

\paragraph{Ghi chú.} set và multiset dùng gần như giống nhau, nhưng multiset cho phép phần tử lặp. Phần tử được tự động sắp xếp theo quy tắc so sánh, mặc định là less\textless{}\textgreater{}. Không thể sửa trực tiếp giá trị phần tử trong container, chỉ có thể chèn và xóa.

\begin{Code}
1 struct classcomp {
2   bool operator() (const int& lhs, const int& rhs) const
3   {return lhs>rhs;}
4 };//这里有个逗号的，注意
5 multiset<int,classcomp> fifth;                // class as Compare
     上面这样就定义成了从大到小排列了。
     结构体自定义排序函数：
     （定义 set 或者 multiset 的时候定义了排序函数，定义迭代器时一样带上排序函数）

 1   struct Node
 2   {
 3       int x,y;
 4   };
 5   struct classcomp//先按照 x 从小到大排序，x 相同则按照 y 从大到小排序
 6   {
 7       bool operator()(const Node &a,const Node &b)const
 8       {
 9           if(a.x!=b.x)return a.x<b.x;
10           else return a.y>b.y;
11       }
12   }; //注意这里有个逗号
13   multiset<Node,classcomp>mt;
14   multiset<Node,classcomp>::iterator it;

1
2    主要函数:
3    begin() 返回指向第一个元素的迭代器
4    clear() 清除所有元素
5    count() 返回某个值元素的个数
6    empty() 如果集合为空，返回 true


 7   end() 返回指向最后一个元素的迭代器
 8   erase() 删除集合中的元素 (参数是一个元素值，或者迭代器)
 9   find() 返回一个指向被查找到元素的迭代器
10   insert() 在集合中插入元素
11   size() 集合中元素的数目
12   lower_bound() 返回指向大于（或等于）某值的第一个元素的迭代器
13   upper_bound() 返回大于某个值元素的迭代器
14   equal_range() 返回集合中与给定值相等的上下限的两个迭代器
15   (注意对于 multiset 删除操作之间删除值会把所以这个值的都删掉，删除一个要用迭代
        器)
\end{Code}

\subsection{Bộ nhập xuất nhanh}

\begin{Code}
 1
 2   //适用于正负整数
 3   template <class T>
 4   inline bool scan_d(T &ret) {
 5      char c; int sgn;
 6      if(c=getchar(),c==EOF) return 0; //EOF
 7      while(c!='-'&&(c<'0'||c>'9')) c=getchar();
 8      sgn=(c=='-')?-1:1;
 9      ret=(c=='-')?0:(c-'0');
10      while(c=getchar(),c>='0'&&c<='9') ret=ret*10+(c-'0');
11      ret*=sgn;
12      return 1;
13   }
14
15   inline void out(int x) {
16      if(x>9) out(x/10);
17      putchar(x%10+'0');
18   }
\end{Code}

\subsection{Thuật toán Mo}

\begin{Code}
     莫队算法，可以解决一类静态，离线区间查询问题。
     BZOJ 2038: [2009 国家集训队] 小 Z 的袜子 (hose)
     Description
     作为一个生活散漫的人，小 Z 每天早上都要耗费很久从一堆五颜六色的袜子中找出一双来
     穿。终于有一天，小 Z 再也无法忍受这恼人的找袜子过程，于是他决定听天由命……具体来
     说，小 Z 把这 N 只袜子从 1 到 N 编号，然后从编号 L 到 R(L
     Input
     输入文件第一行包含两个正整数 N 和 M。N 为袜子的数量，M 为小 Z 所提的询问的数量。
     接下来一行包含 N 个正整数 Ci，其中 Ci 表示第 i 只袜子的颜色，相同的颜色用相同的数字
     表示。再接下来 M 行，每行两个正整数 L，R 表示一个询问。
     Output
     包含 M 行，对于每个询问在一行中输出分数 A/B 表示从该询问的区间 [L,R] 中随机抽出两
     只袜子颜色相同的概率。若该概率为 0 则输出 0/1，否则输出的 A/B 必须为最简分数。（详
     见样例）
     Sample Input
     64
     123332


     26
     13
     35
     16
     Sample Output
     2/5
     0/1
     1/1
     4/15

     题解：
     只需要统计区间内各个数出现次数的平方和

     莫队算法，两种方法，一种是直接分成 sqrt(n) 块，分块排序。
     另外一种是求得曼哈顿距离最小生成树，根据 manhattan MST 的 dfs 序求解。
\end{Code}

\subsubsection{Chia khối}

\begin{Code}
 1   const int MAXN = 50010;
 2   const int MAXM = 50010;
 3   struct Query
 4   {
 5       int L,R,id;
 6   }node[MAXM];
 7   long long gcd(long long a,long long b){
 8       if(b == 0)return a;
 9       return gcd(b,a%b);
10   }
11   struct Ans{
12       long long a,b;//分数 a/b
13       void reduce()//分数化简
14       {
15           long long d = gcd(a,b);
16           a /= d; b /= d;
17       }
18   }ans[MAXM];
19   int a[MAXN];
20   int num[MAXN];
21   int n,m,unit;
22   bool cmp(Query a,Query b){
23       if(a.L/unit != b.L/unit)return a.L/unit < b.L/unit;
24       else return a.R < b.R;
25   }
26   void work(){
27       long long temp = 0;
28       memset(num,0,sizeof(num));
29       int L = 1;
30       int R = 0;
31       for(int i = 0;i < m;i++){
32           while(R < node[i].R){


33             R++;
34             temp -= (long long)num[a[R]]*num[a[R]];
35             num[a[R]]++;
36             temp += (long long)num[a[R]]*num[a[R]];
37         }
38         while(R > node[i].R){
39             temp -= (long long)num[a[R]]*num[a[R]];
40             num[a[R]]--;
41             temp += (long long)num[a[R]]*num[a[R]];
42             R--;
43         }
44         while(L < node[i].L){
45             temp -= (long long)num[a[L]]*num[a[L]];
46             num[a[L]]--;
47             temp += (long long)num[a[L]]*num[a[L]];
48             L++;
49         }
50         while(L > node[i].L){
51             L--;
52             temp -= (long long)num[a[L]]*num[a[L]];
53             num[a[L]]++;
54             temp += (long long)num[a[L]]*num[a[L]];
55         }
56         ans[node[i].id].a = temp - (R-L+1);
57         ans[node[i].id].b = (long long)(R-L+1)*(R-L);
58         ans[node[i].id].reduce();
59     }
60 }
61 int main(){
62     while(scanf("%d%d",&n,&m) == 2){
63         for(int i = 1;i <= n;i++)
64             scanf("%d",&a[i]);
65         for(int i = 0;i < m;i++){
66             node[i].id = i;
67             scanf("%d%d",&node[i].L,&node[i].R);
68         }
69         unit = (int)sqrt(n);
70         sort(node,node+m,cmp);
71         work();
72         for(int i = 0;i < m;i++)
73             printf("%lld/%lld\n",ans[i].a,ans[i].b);
74     }
75     return 0;
76 }
\end{Code}

\subsubsection{Giải bằng thứ tự dfs của Manhattan MST}

\begin{Code}
1   const int MAXN = 50010;
2   const int MAXM = 50010;
3   const int INF = 0x3f3f3f3f;
4   struct Point{
5       int x,y,id;


 6   }p[MAXN],pp[MAXN];
 7   bool cmp(Point a,Point b)
 8   {
 9       if(a.x != b.x) return a.x < b.x;
10       else return a.y < b.y;
11   }
12   //树状数组，找 y-x 大于当前的，但是 y+x 最小的
13   struct BIT{
14       int min_val,pos;
15       void init()
16       {
17           min_val = INF;
18           pos = -1;
19       }
20   }bit[MAXN];
21   struct Edge{
22       int u,v,d;
23   }edge[MAXN<<2];
24   bool cmpedge(Edge a,Edge b){
25       return a.d < b.d;
26   }
27   int tot;
28   int n;
29   int F[MAXN];
30   int find(int x){
31       if(F[x] == -1) return x;
32       else return F[x] = find(F[x]);
33   }
34   void addedge(int u,int v,int d){
35       edge[tot].u = u;
36       edge[tot].v = v;
37       edge[tot++].d = d;
38   }
39   struct Graph{
40       int to,next;
41   }e[MAXN<<1];
42   int total,head[MAXN];
43   void _addedge(int u,int v){
44       e[total].to = v;
45       e[total].next = head[u];
46       head[u] = total++;
47   }
48   int lowbit(int x){
49       return x&(-x);
50   }
51   void update(int i,int val,int pos){
52       while(i > 0){
53           if(val < bit[i].min_val){
54                bit[i].min_val = val;
55                bit[i].pos = pos;
56           }


 57              i -= lowbit(i);
 58       }
 59   }
 60   int ask(int i,int m){
 61       int min_val = INF,pos = -1;
 62       while(i <= m){
 63           if(bit[i].min_val < min_val){
 64                min_val = bit[i].min_val;
 65                pos = bit[i].pos;
 66           }
 67           i += lowbit(i);
 68       }
 69       return pos;
 70   }
 71   int dist(Point a,Point b){
 72       return abs(a.x - b.x) + abs(a.y - b.y);
 73   }
 74   void Manhattan_minimum_spanning_tree(int n,Point p[]){
 75       int a[MAXN],b[MAXN];
 76       tot = 0;
 77       for(int dir = 0;dir < 4;dir++){
 78           if(dir == 1 || dir == 3){
 79                for(int i = 0;i < n;i++)
 80                    swap(p[i].x,p[i].y);
 81           }
 82           else if(dir == 2){
 83                for(int i = 0;i < n;i++)
 84                    p[i].x = -p[i].x;
 85           }
 86           sort(p,p+n,cmp);
 87           for(int i = 0;i < n;i++)
 88                a[i] = b[i] = p[i].y - p[i].x;
 89           sort(b,b+n);
 90           int m = unique(b,b+n) - b;
 91           for(int i = 1;i <= m;i++)
 92                bit[i].init();
 93           for(int i = n-1;i >= 0;i--){
 94                int pos = lower_bound(b,b+m,a[i]) - b + 1;
 95                int ans = ask(pos,m);
 96                if(ans != -1)
 97                    addedge(p[i].id,p[ans].id,dist(p[i],p[ans]));
 98                update(pos,p[i].x+p[i].y,i);
 99           }
100       }
101       memset(F,-1,sizeof(F));
102       sort(edge,edge+tot,cmpedge);
103       total = 0;
104       memset(head,-1,sizeof(head));
105       for(int i = 0;i < tot;i++){
106           int u = edge[i].u, v = edge[i].v;
107           int t1 = find(u), t2 = find(v);


108              if(t1 != t2){
109                  F[t1] = t2;
110                  _addedge(u,v);
111                  _addedge(v,u);
112              }
113       }
114   }
115   int m;
116   int a[MAXN];
117   struct Ans{
118       long long a,b;
119   }ans[MAXM];
120   long long temp ;
121   int num[MAXN];
122   void add(int l,int r){
123       for(int i = l;i <= r;i++){
124           temp -= (long long)num[a[i]]*num[a[i]];
125           num[a[i]]++;
126           temp += (long long)num[a[i]]*num[a[i]];
127       }
128   }
129   void del(int l,int r){
130       for(int i = l;i <= r;i++){
131           temp -= (long long)num[a[i]]*num[a[i]];
132           num[a[i]]--;
133           temp += (long long)num[a[i]]*num[a[i]];
134       }
135   }
136   void dfs(int l1,int r1,int l2,int r2,int idx,int pre){
137       if(l2 < l1) add(l2,l1-1);
138       if(r2 > r1) add(r1+1,r2);
139       if(l2 > l1) del(l1,l2-1);
140       if(r2 < r1) del(r2+1,r1);
141       ans[pp[idx].id].a = temp - (r2-l2+1);
142       ans[pp[idx].id].b = (long long)(r2-l2+1)*(r2-l2);
143       for(int i = head[idx];i != -1;i = e[i].next){
144           int v = e[i].to;
145           if(v == pre) continue;
146           dfs(l2,r2,pp[v].x,pp[v].y,v,idx);
147       }
148       if(l2 < l1)del(l2,l1-1);
149       if(r2 > r1)del(r1+1,r2);
150       if(l2 > l1)add(l1,l2-1);
151       if(r2 < r1)add(r2+1,r1);
152   }
153   long long gcd(long long a,long long b){
154       if(b == 0) return a;
155       else return gcd(b,a%b);
156   }
157   int main(){
158       while(scanf("%d%d",&n,&m) == 2){


159             for(int i = 1;i <= n;i++)
160                 scanf("%d",&a[i]);
161             for(int i = 0;i < m;i++){
162                 scanf("%d%d",&p[i].x,&p[i].y);
163                 p[i].id = i;
164                 pp[i] = p[i];
165             }
166             Manhattan_minimum_spanning_tree(m,p);
167             memset(num,0,sizeof(num));
168             temp = 0;
169             dfs(1,0,pp[0].x,pp[0].y,0,-1);
170             for(int i = 0;i < m;i++){
171                 long long d = gcd(ans[i].a,ans[i].b);
172                 printf("%lld/%lld\n",ans[i].a/d,ans[i].b/d);
173             }
174      }
175      return 0;
176 }
\end{Code}

\subsection{Cấu hình VIM}

\begin{Code}
 1   set nu
 2   set history=1000000
 3
 4   set tabstop=4
 5   set shiftwidth=4
 6   set smarttab
 7
 8   set cindent
 9
10   colo evening
11
12   set nobackup
13   set noswapfile
14
15   set mouse=a
16   map <F6> :call CR()<CR>
17   func! CR()
18   exec "w"
19   exec "!g++ % -o %<"
20   exec "! ./%<"
21   endfunc
22
23   imap <c-]> {<cr>}<c-o>O<left><right>
24   map <F2> :call SetTitle()<CR>
25   func SetTitle()
26   let l = 0
27   let l = l + 1 | call setline(l,'#include <stdio.h>')
28   let l = l + 1 | call setline(l,'#include <string.h>')
29   let l = l + 1 | call setline(l,'#include <iostream>')
30   let l = l + 1 | call setline(l,'#include <algorithm>')
31   let l = l + 1 | call setline(l,'#include <vector>')


32   let l = l + 1 | call setline(l,'#include <queue>')
33   let l = l + 1 | call setline(l,'#include <set>')
34   let l = l + 1 | call setline(l,'#include <map>')
35   let l = l + 1 | call setline(l,'#include <string>')
36   let l = l + 1 | call setline(l,'#include <math.h>')
37   let l = l + 1 | call setline(l,'#include <stdlib.h>')
38   let l = l + 1 | call setline(l,'#include <time.h>')
39   let l = l + 1 | call setline(l,'using namespace std;')
40   let l = l + 1 | call setline(l,'')
41   let l = l + 1 | call setline(l,'int main()')
42   let l = l + 1 | call setline(l,'{')
43   let l = l + 1 | call setline(l,'    //freopen("in.txt","r",stdin);'
        )
44   let l = l + 1 | call setline(l,'    //freopen("out.txt","w",stdout)
        ;')
45   let l = l + 1 | call setline(l,'    ')
46   let l = l + 1 | call setline(l,'    return 0;')
47   let l = l + 1 | call setline(l,'}')
48   endfunc
     现场赛配置：
1    syntax on
2    set nu
3    set tabstop=4
4    set shiftwidth=4
5    set cin
6    colo evening
7    set mouse=a
\end{Code}

\end{document}
